% ============ 第 I 部分: \index{拓扑群}拓扑群的\index{表示}表示 (Representations of Topological Groups) ============
% 原书 I.1 - I.5 (原书 pp.7-118)
% 编号: 无子节的节用 \section 后定理编号 I.1.x; 有子节的节用 \subsection 后编号 I.3.1.x

% 本部分文件结构 (与原书一致):
%   I.1 Topological groups (pp.7-10)
%   I.2 Haar measures (pp.11-18)
%   I.3 Representations (pp.19-29): I.3.1-I.3.5
%   I.4 The convolution product and the group algebra (pp.30-44): I.4.1-I.4.3
%   I.5 Exercises (pp.44-118): I.5.1-I.5.6

\section{拓扑群}
% 原书 I.1 (pp.7-10)

\begin{definition}[拓扑群 (topological group)]
\label{def:I.1.1}
一个拓扑群 (topological group) 是一个带群结构的拓扑空间 $G$, 使得乘法映射
$G \times G \to G$, $(x,y) \longmapsto xy$ 与取逆映射 $G \to G$, $x \longmapsto x^{-1}$ 都是连续的.
\end{definition}

我们通常用 $1$ 或 $e$ 表示 $G$ 的单位元.

\begin{example}[拓扑群的例子]
\label{exa:I.1.2}
\begin{itemize}
\item 任何赋予离散拓扑的群都是拓扑群. 常用的例子包括有限群、自由群 (自由交换群与自由非交换群) 以及``算术''矩阵群, 如 $\GL_n(\Z)$ 与 $\SL_n(\Z)$.
\item $\R$ 和 $\CC$ 的加法群是拓扑群.
\item 群 $\GL_n(\CC)$ 赋以 $\CC$-向量空间 $M_n(\CC)$ 上任意范数给出的拓扑, 是一个拓扑群 (见习题 \ref{ex:I.5.1.1}), 因此其一切子群在诱导拓扑下也是拓扑群. 例如 $S^1 := \{z \in \CC \mid |z| = 1\}$, $\GL_n(\R)$, $\SU(n)$, $\SO(n)$ 等等.
\item $\Q_p$ 的加法群与群 $\GL_n(\Q_p)$ 是拓扑群. (见习题 \ref{ex:I.5.1.4}.)
\end{itemize}
\end{example}

\begin{definition}[\index{局部紧}局部紧 (locally compact)]
\label{def:I.1.3}
设 $X$ 是一个拓扑空间. 若 $X$ 的每一点都有一个紧邻域, 则称 $X$ 是局部紧 (locally compact) 的.
\end{definition}

\begin{remark}[局部紧与紧邻域基]
\label{rem:I.1.4}
若 $X$ 是 Hausdorff 的, 则这等价于 $X$ 的每一点都有一个由紧邻域组成的邻域基.\footnote{参考文献 ?}

注意我们并不假定拓扑空间中点的邻域是开的.
\end{remark}

\begin{notation}[子集的运算]
\label{not:I.1.5}
设 $G$ 是一个群, $A$, $B \subset G$, $x \in G$, $n \geq 1$. 我们使用以下记号:
\[
xA = \{xy \mid y \in A\} \quad \text{与} \quad Ax = \{yx \mid y \in A\},
\]
\[
AB = \{yz \mid y \in A,\, z \in B\},
\]
\[
A^n = AA \cdots A \quad (\text{$n$ 个因子}),
\]
\[
A^{-1} = \{y^{-1} \mid y \in A\}.
\]
\end{notation}

\begin{definition}[\index{对称子集}对称子集 (symmetric subset)]
\label{def:I.1.6}
设 $G$ 是一个群. 若 $G$ 的子集 $A$ 满足 $A = A^{-1}$, 则称 $A$ 是对称 (symmetric) 的.
\end{definition}

\begin{proposition}[拓扑群的基本性质]
\label{prop:I.1.7}
设 $G$ 是一个拓扑群.
\begin{enumerate}
\item[(1)] 若 $U$ 是 $G$ 的开子集而 $A$ 是 $G$ 的任意子集, 则集合 $UA$, $AU$ 与 $U^{-1}$ 都是开的.
\item[(2)] 若 $U$ 是 $1$ 在 $G$ 中的一个邻域, 则存在 $1$ 的开对称邻域 $V$ 使得 $V^2 \subset U$.
\item[(3)] 若 $H$ 是 $G$ 的子群, 则其闭包 $\overline{H}$ 也是 $G$ 的子群.
\item[(4)] 若 $H$ 是 $G$ 的开子群, 则它也是闭的.
\item[(5)] 若 $A$ 和 $B$ 是 $G$ 的紧子集, 则集合 $AB$ 也是紧的.
\item[(6)] 设 $H$ 是 $G$ 的子群. 则商 $G/H$ (赋予商拓扑) 具有以下性质:
  \begin{enumerate}
  \item[(a)] 若 $H$ 是闭的, 则 $G/H$ 是 Hausdorff 的;
  \item[(b)] 若 $G$ 是局部紧的, 则 $G/H$ 是局部紧的;
  \item[(c)] 若 $H$ 是正规的, 则 $G/H$ 是一个拓扑群.
  \end{enumerate}
\end{enumerate}
\end{proposition}

\begin{proof}
(1). 对 $x \in G$, 记 $l_x \colon G \to G$ (相应地 $r_x \colon G \to G$) 为左 (相应地右) 乘以 $x$. 又记 $\iota \colon G \to G$ 为映射 $x \longmapsto x^{-1}$. 由拓扑群的公理, 所有这些映射都是连续的.

注意到 $U^{-1} = \iota^{-1}(U)$, $AU = \bigcup_{x \in A} l_{x^{-1}}^{-1}(U)$, $UA = \bigcup_{x \in A} r_{x^{-1}}^{-1}(U)$. 因此 $U^{-1}$, $AU$ 与 $UA$ 都是开的.

(2). 我们可以假定 $U$ 是开的. 设 $m \colon G \times G \to G$, $(x,y) \longmapsto xy$. 则 $m$ 连续, 所以 $W := m^{-1}(U)$ 是开的. 因为 $1^2 = 1 \in U$, 我们有 $(1,1) \in W$. 由 $G \times G$ 上乘积拓扑的定义, 存在 $G$ 的开子集 $\Omega \ni 1$ 使得 $\Omega \times \Omega \subset W$. 由 $W$ 的定义有 $\Omega^2 \subset U$. 令 $V = \Omega \cap \Omega^{-1}$. 由 (1) 知 $\Omega^{-1}$ 是开的, 所以 $V$ 是开的, 且由定义知它是对称的. 显然 $1 \in V$ 且 $V^2 \subset \Omega^2 \subset U$.

(3). 考虑映射 $u \colon G \times G \to G$, $(x,y) \longmapsto xy^{-1}$; 则 $G$ 的非空子集 $A$ 是子群当且仅当 $u(A \times A) \subset A$. 又由拓扑群的公理, 映射 $u$ 连续. 因此对每个 $Z \subset G \times G$, 有 $u(\overline{Z}) \subset \overline{u(Z)}$. 将其应用于 $H \times H$ (其闭包为 $\overline{H} \times \overline{H}$), 我们得到 $\overline{H}$ 是 $G$ 的子群.

(4). 我们有 $G = H \sqcup ((G - H)H)$. 若 $H$ 是开的, 则由 (1) 知 $(G - H)H$ 也是开的, 因此 $H$ 是闭的.

(5). 乘法映射 $m \colon G \times G \to G$ 由假设连续. 由于 $AB = m(A \times B)$ 且 $A \times B$ 紧, 集合 $AB$ 也是紧的.

(6). (a) 设 $x, y \in G$ 满足 $xH \neq yH$. 由 (1), $x(G - H)y^{-1}$ 是开的, 所以其余集 $xHy^{-1}$ 是闭的. 又由假设 $xH \neq yH$, 单位元 $1$ 不在 $xHy^{-1}$ 中. 由 (2), 存在对称开集 $1 \in U$ 使得 $U^2 \subset G - xHy^{-1}$. 我们来证明 $UxH \cap UyH = \emptyset$, 这将证明结论, 因为 $UxH$ (相应地 $UyH$) 是 $xH$ (相应地 $yH$) 在 $G/H$ 中的开邻域. 若 $UxH \cap UyH \neq \emptyset$, 则存在 $u_1, u_2 \in U$ 与 $h_1, h_2 \in H$ 使得 $u_1xh_1 = u_2yh_2$. 但此时 $xh_1h_2^{-1}y^{-1} = u_1^{-1}u_2 \in xHy^{-1} \cap U^2$, 这是不可能的.

(b) 设 $xH \in G/H$. 若 $K$ 是 $x$ 在 $G$ 中的紧邻域, 则其在 $G/H$ 中的像是 $xH$ 在 $G/H$ 中的紧邻域.

(c) 若 $H$ 正规, 则 $G/H$ 是一个群. 我们来证明其乘法连续. 设 $x, y \in G$. $xyH$ 在 $G/H$ 中的任何开邻域都具有 $UxyH$ 的形式, 其中 $U$ 是 $xy$ 在 $G$ 中的开邻域. 由 $G$ 上乘法的连续性, 存在 $x$ 与 $y$ 在 $G$ 中的开邻域 $V$ 与 $W$ 使得 $VW \subset U$. 则 $VH$ 与 $WH$ 是 $xH$ 与 $yH$ 在 $G/H$ 中的开邻域, 且有 $(VH)(WH) \subset UH$. (记住, 由于 $H$ 正规, 对 $G$ 的任何子集 $A$ 有 $AH = HA$.) 我们来证明 $G/H$ 上的取逆连续. 设 $x \in G$. $x^{-1}H$ 在 $G/H$ 中的任何开邻域都具有 $UH$ 的形式, 其中 $U$ 是 $x^{-1}$ 在 $G$ 中的开邻域. 由 (1), 集合 $U^{-1}$ 是开的, 所以 $U^{-1}H$ 是 $xH$ 在 $G/H$ 中的开邻域, 且有 $(U^{-1}H)^{-1} = HU = UH$.
\end{proof}

\begin{remark}[Hausdorff 化]
\label{rem:I.1.8}
特别地, 若 $G$ 是拓扑群, 则 $G/\overline{\{1\}}$ 是 Hausdorff 拓扑群. 我们感兴趣的是 $G$ 在向量空间上的连续群作用, 因此可以把 $G$ 换成 $G/\overline{\{1\}}$ 来研究它们. 因此在以下的叙述中, 我们只考虑 Hausdorff 拓扑群 (除非另有说明).
\end{remark}

\begin{definition}[\index{紧群}紧群 (compact group)]
\label{def:I.1.9}
紧群 (compact group) (相应地局部紧群 (locally compact group)) 是 Hausdorff 且紧 (相应地局部紧) 的拓扑群.
\end{definition}

\begin{example}[紧群与局部紧群的例子]
\label{exa:I.1.10}
在例 \ref{exa:I.1.2} 的群中, 有限离散群与群 $S^1$, $\SU(n)$, $\SO(n)$ 是紧的. 其余所有群都是局部紧的. 若考虑一个无穷维巴拿赫空间的可逆有界线性自同态群, 则得到一个非局部紧的群 (见习题 \ref{ex:I.5.1.1}).
\end{example}

\textbf{平移算子 (Translation operators).} 设 $G$ 是一个群, $x \in G$, $f \colon G \to \CC$ 是一个函数. 我们定义两个函数 $L_xf, R_xf \colon G \to \CC$ 如下:
\[
L_xf(y) = f(x^{-1}y) \quad \text{与} \quad R_xf(y) = f(yx).
\]
我们这样选择约定是为了 $L_{xy} = L_x \circ L_y$ 且 $R_{xy} = R_x \circ R_y$. 注意, 若 $G$ 是拓扑群且 $f$ 连续, 则 $L_xf$ 与 $R_xf$ 也连续.

\textbf{函数空间 (Function spaces).} 设 $X$ 是一个拓扑空间. 若 $f \colon X \to \CC$ 是一个函数, 我们记
\[
\|f\|_{\infty} = \sup_{x \in X} |f(x)| \in [0, +\infty].
\]
我们还使用以下记号:
\begin{itemize}
\item $\mathcal{C}(X)$: 从 $X$ 到 $\CC$ 的连续函数 $f \colon X \to \CC$ 的集合;
\item $\mathcal{C}_b(X)$: 有界连续函数 $f \colon X \to \CC$ 的集合 (即满足 $\|f\|_{\infty} < +\infty$ 的 $\mathcal{C}(X)$ 的元素);
\item $\mathcal{C}_0(X)$: 在无穷远处消失 (vanish at infinity) 的连续函数 $X \to \CC$ 的集合 (即对每个 $\varepsilon > 0$, 存在 $X$ 的紧子集 $K$ 使得对所有 $x \notin K$ 有 $|f(x)| < \varepsilon$);
\item $\mathcal{C}_c(X)$: 具有紧支撑 (compact support) 的连续函数 $X \to \CC$ 的集合.
\end{itemize}
注意我们有 $\mathcal{C}(X) \supset \mathcal{C}_b(X) \supset \mathcal{C}_0(X) \supset \mathcal{C}_c(X)$, 且若 $X$ 紧, 则这些包含都是等式. 函数 $\|\cdot\|_{\infty}$ 是 $\mathcal{C}_b(X)$ 及其子空间上的范数, 且 $\mathcal{C}_b(X)$ 与 $\mathcal{C}_0(X)$ 对此范数完备 (但 $\mathcal{C}_c(X)$ 不完备, 除非 $X$ 紧).\footnote{参考文献 ?}

\begin{definition}[\index{一致连续}一致连续 (uniformly continuous)]
\label{def:I.1.11}
设 $G$ 是一个拓扑群. 若当 $x \to 1$ 时 $\|L_xf - f\|_{\infty} \to 0$ (相应地当 $x \to 1$ 时 $\|R_xf - f\|_{\infty} \to 0$), 则称函数 $f \colon G \to \CC$ 是左 (相应地右) 一致连续 (left (resp.\ right) uniformly continuous) 的.
\end{definition}

\begin{proposition}[紧支撑函数的均匀连续性]
\label{prop:I.1.12}
若 $f \in \mathcal{C}_c(G)$, 则 $f$ 既是左一致连续又是右一致连续的.
\end{proposition}

\begin{proof}
我们证明 $f$ 是右一致连续的 (左一致连续的证明类似). 设 $K$ 是 $f$ 的支撑. 设 $\varepsilon > 0$. 对每个 $x \in K$, 我们选取 $1$ 的一个邻域 $U_x$ 使得对所有 $y \in U_x$ 有 $|f(xy) - f(x)| < \frac{\varepsilon}{2}$; 由命题 \ref{prop:I.1.7}, 可以找到 $1$ 的对称开邻域 $V_x$ 使得 $V_x^2 \subset U_x$. 我们有 $K \subset \bigcup_{x \in K} xV_x$. 由于 $K$ 紧, 可以找到 $x_1, \ldots, x_n \in K$ 使得 $K \subset \bigcup_{i=1}^n x_iV_{x_i}$. 令 $V = \bigcup_{i=1}^n V_{x_i}$, 这是 $1$ 的一个对称开邻域.

我们断言: 若 $y \in V$, 则 $\|R_yf - f\|_{\infty} < \varepsilon$. 事实上, 设 $y \in V$, $x \in G$. 先假定 $x \in K$. 则存在 $i \in \{1, \ldots, n\}$ 使得 $x \in x_iV_{x_i}$. 于是 $xy \in x_iV_{x_i}V_{x_i} \subset x_iU_{x_i}$, 从而
\[
|f(xy) - f(x)| \leq |f(xy) - f(x_i)| + |f(x_i) - f(x)| < \frac{\varepsilon}{2} + \frac{\varepsilon}{2} = \varepsilon.
\]
再假定 $xy \in K$. 则存在 $i \in \{1, \ldots, n\}$ 使得 $xy \in x_iV_{x_i}$, 且有 $x = xyy^{-1} \in x_iV_{x_i}V_{x_i} \subset x_iU_{x_i}$. 因此
\[
|f(xy) - f(x)| \leq |f(xy) - f(x_i)| + |f(x_i) - f(x)| < \frac{\varepsilon}{2} + \frac{\varepsilon}{2} = \varepsilon.
\]
最后, 若 $x, xy \notin K$, 则 $f(x) = f(xy) = 0$, 当然有 $|f(xy) - f(x)| < \varepsilon$.
\end{proof}

\begin{remark}[一致连续与连续性的等价刻画]
\label{rem:I.1.13}
我们在 $\mathcal{C}_b(G)$ 上赋予由 $\|\cdot\|_{\infty}$ 给出的拓扑. 则函数 $f \in \mathcal{C}_b(G)$ 是左 (相应地右) 一致连续的, 当且仅当映射 $G \to \mathcal{C}_b(G)$, $x \longmapsto L_xf$ (相应地 $x \longmapsto R_xf$) 在 $G$ 的单位元处连续.

利用 $L_{xy} = L_x \circ L_y$ 与 $R_{xy} = R_x \circ R_y$, 以及算子 $L_x$ 与 $R_x$ 保持 $\mathcal{C}_c(G)$, 我们看到上述命题蕴含: 若 $f \in \mathcal{C}_c(G)$, 则把 $x \in G$ 分别映到 $L_xf$ 与 $R_xf$ 的两个映射 $G \to \mathcal{C}_c(G)$ 都是连续的.
\end{remark}

% ============ I.2 \index{Haar 测度}Haar 测度 (原书 pp.11-18) ============
\section{Haar 测度}

\begin{definition}[\index{Borel 测度}Borel 测度与\index{正则 Borel 测度}正则 Borel 测度]
\label{def:I.2.1}
设 $X$ 是一个拓扑空间.
\begin{enumerate}
\item[(1)] $X$ 上 Borel 集的 $\sigma$-代数 (the $\sigma$-algebra of Borel sets) 是由 $X$ 的开子集生成的 $\sigma$-代数. $X$ 上的 Borel 测度 (Borel measure) 是该 $\sigma$-代数上的测度.
\item[(2)] $X$ 上的正则 Borel 测度 (regular Borel measure) 是 $X$ 的 Borel 集的 $\sigma$-代数上满足下列性质的测度 $\mu$:
  \begin{enumerate}
  \item[(a)] 对 $X$ 的每个紧子集 $K$, $\mu(K) < +\infty$;
  \item[(b)] $\mu$ 是外正则 (outer regular) 的: 对 $X$ 的每个 Borel 子集 $E$, 有 $\mu(E) = \inf\{\mu(U) \mid U \supset E \text{ 开}\}$;
  \item[(c)] $\mu$ 是内正则 (inner regular) 的: 对每个 $E \subset X$, 只要 $E$ 是有限测度的 Borel 集或是开集, 就有 $\mu(E) = \sup\{\mu(K) \mid K \subset E \text{ 紧}\}$.
  \end{enumerate}
\end{enumerate}
\end{definition}

\begin{notation}[非负紧支撑函数]
\label{not:I.2.2}
我们记 $\mathcal{C}_c^+(X)$ 为所有非零的 $f \in \mathcal{C}_c(X)$ 且 $f(X) \subset \R_{\geq 0}$ 组成的集合.
\end{notation}

\begin{theorem}[Riesz 表示定理 (Riesz representation theorem)]
\label{thm:I.2.3}
设 $X$ 是局部紧 Hausdorff 空间, 设 $\Lambda \colon \mathcal{C}_c(X) \to \CC$ 是满足对每个 $f \in \mathcal{C}_c^+(X)$ 都有 $\Lambda(f) \geq 0$ 的线性泛函.\footnote{这样的线性泛函称为正的 (positive).} 则存在 $X$ 上唯一的正则 Borel 测度 $\mu$ 使得对每个 $f \in \mathcal{C}_c(X)$ 有
\[
\Lambda(f) = \int_X f \,\dd\mu.
\]
\end{theorem}

\begin{definition}[左/右 Haar 测度]
\label{def:I.2.4}
设 $G$ 是局部紧群. $G$ 上的左 (相应地右) Haar 测度 (left (resp.\ right) Haar measure) 是非零正则 Borel 测度 $\mu$, 使得对 $G$ 的每个 Borel 集 $E$ 和每个 $x \in G$, 有 $\mu(xE) = \mu(E)$ (相应地 $\mu(Ex) = \mu(E)$).
\end{definition}

\begin{example}[Haar 测度的例子]
\label{exa:I.2.5}
\begin{enumerate}
\item[(1)] 若 $G$ 是离散群, 则计数测度 (counting measure) 是 $G$ 上的左和右 Haar 测度.
\item[(2)] Lebesgue 测度是 $\R$ 的加法群上的左和右 Haar 测度.
\end{enumerate}
\end{example}

\begin{proposition}[Haar 测度的刻画]
\label{prop:I.2.6}
设 $G$ 是局部紧群, $\mu$ 是 $G$ 上的正则 Borel 测度.
\begin{enumerate}
\item[(1)] 设 $\tilde{\mu}$ 是由 $\tilde{\mu}(E) = \mu(E^{-1})$ 定义的 Borel 测度. 则 $\mu$ 是左 Haar 测度当且仅当 $\tilde{\mu}$ 是右 Haar 测度.
\item[(2)] $\mu$ 是 $G$ 上的左 Haar 测度当且仅当: 对每个 $f \in \mathcal{C}_c(G)$, 对每个 $y \in G$, 有 $\int_G L_yf \,\dd\mu = \int_G f \,\dd\mu$.
\item[(3)] 若 $\mu$ 是 $G$ 上的左 Haar 测度, 则对 $G$ 的每个非空开子集有 $\mu(U) > 0$, 且对每个 $f \in \mathcal{C}_c^+(G)$ 有 $\int_G f \,\dd\mu > 0$.
\end{enumerate}
\end{proposition}

\begin{proof}
(1). 首先注意 $\tilde{\mu}$ 是 $G$ 上的正则 Borel 测度, 因为 $x \longmapsto x^{-1}$ 是 $G$ 到自身的同胚.

若 $E \subset G$ 是 Borel 集且 $x \in G$, 则 $\tilde{\mu}(Ex) = \mu(x^{-1}E^{-1})$. 这就蕴含了结论.

(2). 设 $x \in G$, 令 $\mu_x$ 是由 $\mu_x(E) = \mu(xE)$ 定义的 Borel 测度. (这确实是 $G$ 上的正则 Borel 测度, 因为 $y \longmapsto xy$ 是 $G$ 到自身的同胚.) 则对每个可测函数 $f \colon G \to \CC$, 有 $\int_G f \,\dd\mu_x = \int_G L_xf \,\dd\mu$. (这对 Borel 子集的特征函数是显然的, 一般情形通过对 $f$ 用特征函数的线性组合逼近得到.)

一方面, 测度 $\mu$ 是左 Haar 测度当且仅当对每个 $x \in G$ 有 $\mu = \mu_x$. 另一方面, 由 Riesz 表示定理中的唯一性 (以及上一段), 对 $x \in G$, $\mu = \mu_x$ 当且仅当对每个 $f \in \mathcal{C}_c(G)$ 有 $\int_G f \,\dd\mu = \int_G L_xf \,\dd\mu$. 结论由此得出.

(3). 假设存在 $G$ 的非空开子集 $U$ 使得 $\mu(U) = 0$. 则对每个 $x \in G$ 有 $\mu(xU) = 0$, 所以我们可以假设 $1 \in U$. 设 $K$ 是 $G$ 的紧子集. 则 $K \subset \bigcup_{x \in K} xU$, 所以存在 $x_1, \ldots, x_n \in K$ 使得 $K \subset \bigcup_{i=1}^n x_iU$. 由于对每个 $i$ 有 $\mu(x_iU) = 0$, 这蕴含 $\mu(K) = 0$. 但由 $\mu$ 的内正则性, 我们得到 $\mu(G) = 0$, 这与 $\mu$ 非零矛盾.

设 $f \in \mathcal{C}_c^+(G)$. 则 $U := \{x \in G \mid f(x) > \frac{1}{2}\|f\|_{\infty}\}$ 是 $G$ 的非空开子集, 所以 $\mu(U) > 0$. 但 $f \geq \frac{1}{2}\|f\|_{\infty} \mathbb{1}_U$, 因此 $\int_G f \,\dd\mu \geq \frac{1}{2}\|f\|_{\infty}\mu(U) > 0$.
\end{proof}

\begin{theorem}[左 Haar 测度的存在性与唯一性]
\label{thm:I.2.7}
设 $G$ 是局部紧群. 则:
\begin{enumerate}
\item[(1)] $G$ 上存在左 Haar 测度.
\item[(2)] 若 $\mu_1$ 和 $\mu_2$ 是 $G$ 上的两个左 Haar 测度, 则存在 $c \in \R_{>0}$ 使得 $\mu_2 = c\mu_1$.
\end{enumerate}
\end{theorem}

由命题 \ref{prop:I.2.6}, 此定理蕴含右 Haar 测度的类似结论.

\begin{proof}
我们先证明存在性. 其想法与 $\R$ 上 Lebesgue 测度的构造非常相似. 假设 $c > 0$, 且 $\phi \in \mathcal{C}_c^+(\R)$ 被 $1$ 控制且非常接近区间 $[0,c]$ 的特征函数. 若 $f \in \mathcal{C}_c(\R)$ 在长度为 $c$ 的区间上变化不太快, 则我们可以用 $\phi$ 的左平移的线性组合逼近 $f$: $f \simeq \sum c_iL_{x_i}\phi$, 然后 $\int f \,\dd\mu \simeq \sum c_i \int L_{x_i}\phi \,\dd\mu$. 当 $c \to 0$ 时, 我们将能逼近每个 $f \in \mathcal{C}_c(\R)$ (因为我们知道这些函数是一致连续的), 并且能通过取极限来定义 $\int f \,\dd\mu$. 在一般的局部紧群上, 我们把区间换成 $1$ 的越来越小的紧邻域.

下面是严格的证明. 设 $f, \phi \in \mathcal{C}_c^+(G)$. 则 $U := \{x \in G \mid \phi(x) > \frac{1}{2}\|\phi\|_{\infty}\}$ 是 $G$ 的非空开子集, 且 $\phi \geq \frac{1}{2}\|\phi\|_{\infty}\mathbb{1}_U$. 由于 $f$ 的支撑紧, 它可被 $U$ 的有限多个平移覆盖, 所以存在 $x_1, \ldots, x_n \in G$ 和 $c_1, \ldots, c_n \in \R_{\geq 0}$ 使得 $f \leq \sum_{i=1}^n c_iL_{x_i}\phi$. 因此, 若定义 $(f : \phi)$ 为所有有限和 $\sum_{i=1}^n c_i$ (其中 $c_1, \ldots, c_n \in \R_{\geq 0}$ 且存在 $x_1, \ldots, x_n \in G$ 使得 $f \leq \sum_{i=1}^n c_iL_{x_i}\phi$) 的下确界, 则 $(f : \phi) < +\infty$. 我们断言:
\begin{align}
\label{eq:I.2.1}
(f : \phi) &= (L_xf : \phi) \quad \forall x \in G, \\
\label{eq:I.2.2}
(f_1 + f_2 : \phi) &\leq (f_1 : \phi) + (f_2 : \phi), \\
\label{eq:I.2.3}
(cf : \phi) &= c(f : \phi) \quad \forall c \geq 0, \\
\label{eq:I.2.4}
(f_1 : \phi) &\leq (f_2 : \phi) \quad \text{若 } f_1 \leq f_2, \\
\label{eq:I.2.5}
(f : \phi) &\geq \frac{\|f\|_{\infty}}{\|\phi\|_{\infty}}, \\
\label{eq:I.2.6}
(f : \phi) &\leq (f : \psi)(\psi : \phi) \quad \forall \psi \in \mathcal{C}_c^+(G) - \{0\}.
\end{align}

前四条性质是容易的. 对第五条性质, 注意若 $f \leq \sum_{i=1}^n c_iL_{x_i}\phi$, 则
\[
\|f\|_{\infty} \leq \sum_{i=1}^n c_i\|L_{x_i}\phi\|_{\infty} = \sum_{i=1}^n c_i\|\phi\|_{\infty}.
\]

最后, 最后一条性质是以下事实的推论: 设 $\psi \in \mathcal{C}_c^+(G)$. 若 $f \leq \sum_{i=1}^n c_iL_{x_i}\psi$ 且 $\psi \leq \sum_{j=1}^m d_jL_{y_j}\phi$, 则 $f \leq \sum_{i=1}^n \sum_{j=1}^m c_id_jL_{x_iy_j}\phi$.

现在我们固定 $f_0 \in \mathcal{C}_c^+(G)$. 由 \eqref{eq:I.2.5}, 我们知道 $(f_0 : \phi) > 0$. 我们定义 $I_{\phi} \colon \mathcal{C}_c^+(G) \to \R_{\geq 0}$ 为
\[
I_{\phi}(f) = \frac{(f : \phi)}{(f_0 : \phi)}.
\]
由 \eqref{eq:I.2.1}-\eqref{eq:I.2.4}, 我们有
\begin{align*}
I_{\phi}(f) &= I_{\phi}(L_xf) \quad \forall x \in G, \\
I_{\phi}(f_1 + f_2) &\leq I_{\phi}(f_1) + I_{\phi}(f_2), \\
I_{\phi}(cf) &= cI_{\phi}(f) \quad \forall c \geq 0, \\
I_{\phi}(f_1) &\leq I_{\phi}(f_2) \quad \text{若 } f_1 \leq f_2.
\end{align*}

若第二个不等式是等式 (即若 $I_{\phi}$ 可加), 我们就能把 $I_{\phi}$ 延拓成 $\mathcal{C}_c(G)$ 上的\index{正线性泛函}正线性泛函并应用 Riesz 表示定理. 这并不完全成立, 但我们有以下结果:

\textbf{断言 (Claim):} 对所有 $f_1, f_2 \in \mathcal{C}_c^+(G)$ 和 $\varepsilon > 0$, 存在 $1$ 在 $G$ 中的邻域 $V$ 使得只要 $\operatorname{supp}(\phi) \subset V$, 就有 $I_{\phi}(f_1) + I_{\phi}(f_2) \leq I_{\phi}(f_1 + f_2) + \varepsilon$.

我们首先证明这个断言. 选取函数 $g \in \mathcal{C}_c^+(G)$ 使得对每个 $x \in \operatorname{supp}(f_1 + f_2)$ 有 $g(x) = 1$, 并设 $\delta$ 是一个正实数. 令 $h = f_1 + f_2 + \delta g$. 我们定义函数 $h_1, h_2 \colon G \to \R_{\geq 0}$ 为
\[
h_i(x) = \begin{cases} \frac{f_i(x)}{h(x)} & \text{若 } f_i(x) \neq 0, \\ 0 & \text{若 } f_i(x) = 0. \end{cases}
\]
注意 $h_i$ 等于 $\frac{f_i}{h}$, 因此在开子集 $\{x \in G \mid h(x) \neq 0\}$ 上连续. 由于 $G$ 是这个开子集与开子集 $G - \operatorname{supp}(f_i)$ (在其上 $h_i$ 也连续) 的并, 这证明 $h_i$ 连续, 因此 $h_i \in \mathcal{C}_c^+(G)$. 又注意 $f_i = h_ih$.

由命题 \ref{prop:I.1.12}, 存在 $1$ 的邻域 $V$ 使得对 $i \in \{1, 2\}$ 及满足 $y^{-1}x \in V$ 的 $x, y \in G$, 有 $|h_i(x) - h_i(y)| < \delta$. 设 $\phi \in \mathcal{C}_c^+(G)$ 满足 $\operatorname{supp}(\phi) \subset V$. 若 $c_1, \ldots, c_n \in \R_{\geq 0}$ 与 $x_1, \ldots, x_n \in G$ 满足 $h \leq \sum_{j=1}^n c_jL_{x_j}\phi$, 则对每个 $x \in G$ 和 $i \in \{1, 2\}$,
\[
f_i(x) = h(x)h_i(x) \leq \sum_{j=1}^n c_j\phi(x_j^{-1}x)h_i(x) \leq \sum_{j=1}^n c_j\phi(x_j^{-1}x)(h_i(x_j) + \delta),
\]
因为 $\phi(x_j^{-1}x) = 0$ 除非 $x_j^{-1}x \in V$. 因此
\[
(f_1 : \phi) + (f_2 : \phi) \leq \sum_{j=1}^n c_j(h_1(x_j) + h_2(x_j) + 2\delta).
\]
由于 $h_1 + h_2 \leq 1$, 我们得到
\[
(f_1 : \phi) + (f_2 : \phi) \leq (1 + 2\delta)\sum_{j=1}^n c_j,
\]
从而, 对族 $(c_1, \ldots, c_n)$ 取下确界并除以 $(f_0 : \phi)$, 得到
\[
I_{\phi}(f_1) + I_{\phi}(f_2) \leq (1 + 2\delta)I_{\phi}(h) \leq (1 + 2\delta)(I_{\phi}(f_1 + f_2) + \delta I_{\phi}(g)).
\]
当 $\delta$ 趋于 $0$ 时, 右端趋于 $I_{\phi}(f_1 + f_2)$, 所以取足够小的 $\delta$ 就得到所要的不等式. 这就完成了断言的证明.

我们回到左 Haar 测度的构造. 对每个 $f \in \mathcal{C}_c^+(G)$, 设 $X_f = [(f_0 : f)^{-1}, (f : f_0)] \subset \R$. 令 $X = \prod_{f \in \mathcal{C}_c^+(G)} X_f$, 赋予乘积拓扑. 由 Tychonoff 定理,\footnote{参考文献 ?} $X$ 是紧 Hausdorff 空间. 它正是函数 $I \colon \mathcal{C}_c^+(G) \to \R$ (满足对每个 $f$ 有 $I(f) \in X_f$) 的空间 (赋予逐点收敛拓扑). 又由 \eqref{eq:I.2.6}, 对每个 $\phi \in \mathcal{C}_c^+(G)$ 有 $I_{\phi} \in X$. 对 $1$ 在 $G$ 中的每个邻域 $V$, 设 $K(V)$ 为 $\{I_{\phi} \mid \operatorname{supp}(\phi) \subset V\}$ 在 $X$ 中的闭包. 对每个 $V$ 有 $K(V) \neq \emptyset$, 所以对 $1$ 在 $G$ 中的任何有限族 $V_1, \ldots, V_n$ 个邻域, $K(V_1) \cap \cdots \cap K(V_n) \supset K(\bigcap_{i=1}^n V_i) \neq \emptyset$. 由于 $X$ 紧, 这蕴含所有集合 $K(V)$ 的交非空. 我们选取这个交中的一个元素 $I$.

我们来证明 $I$ 在左平移下不变、可加且是一次齐次的. (也就是说, 它和 $I_{\phi}$ 有相同的性质, 但它是可加的面不是次可加的.) 设 $f_1, f_2 \in \mathcal{C}_c^+(G)$, $c \in \R_{\geq 0}$, $x \in G$, $\varepsilon > 0$. 选取 $1$ 在 $G$ 中的邻域 $V$ 使得只要 $\operatorname{supp}(\phi) \subset V$ 就有 $I_{\phi}(f_1) + I_{\phi}(f_2) \leq I_{\phi}(f_1 + f_2) + \varepsilon$; 这样的 $V$ 由断言保证存在. 由 $I$ 的定义, 它属于 $\{I_{\phi} \mid \operatorname{supp}(\phi) \subset V\}$ 的闭包, 这意味着存在 $\phi \in \mathcal{C}_c^+(G)$ 使得 $\operatorname{supp}(\phi) \subset V$ 且对 $g \in \{f_1, f_2, f_1 + f_2\}$, $y \in \{1, x\}$, $a \in \{1, c\}$ 有 $|I(aL_yg) - I_{\phi}(aL_yg)| < \varepsilon$. 于是我们得到:
\begin{align*}
|I(L_xf_1) - I(f_1)| &\leq |I(L_xf_1) - I_{\phi}(L_xf_1)| + |I_{\phi}(L_xf_1) - I_{\phi}(f_1)| + |I_{\phi}(f_1) - I(f_1)| < 2\varepsilon, \\
|I(cf_1) - cI(f_1)| &\leq |I(cf_1) - I_{\phi}(cf_1)| + |I_{\phi}(cf_1) - cI_{\phi}(f_1)| + |cI_{\phi}(f_1) - cI(f_1)| < \varepsilon(1 + c),
\end{align*}
以及
\begin{multline*}
|I(f_1 + f_2) - I(f_1) - I(f_2)| \leq |I_{\phi}(f_1 + f_2) - I_{\phi}(f_1) - I_{\phi}(f_2)| \\
+ |I(f_1 + f_2) - I_{\phi}(f_1 + f_2)| + |I(f_1) - I_{\phi}(f_1)| + |I(f_2) - I_{\phi}(f_2)| < 4\varepsilon.
\end{multline*}
由于 $\varepsilon$ 任意, 这蕴含 $I(L_xf_1) = I(f_1)$, $I(cf_1) = cI(f_1)$ 且 $I(f_1 + f_2) = I(f_1) + I(f_2)$.

现在我们把它延拓成线性泛函 $\mathcal{C}_c(G) \to \CC$, 仍记作 $I$. 设 $f \in \mathcal{C}_c(G)$. 则我们可以把 $f$ 写成 $f = (f_1 - f_2) + i(g_1 - g_2)$, 其中 $f_1, f_2, g_1, g_2 \in \mathcal{C}_c^+(G) \cup \{0\}$ (例如取 $f_1 = \max(0, \operatorname{Re}(f))$, $f_2 = \max(0, -\operatorname{Re}(f))$, $g_1 = \max(0, \operatorname{Im}(f))$, $g_2 = \max(0, \operatorname{Im}(f))$). 我们令 $I(f) = I(f_1) - I(f_2) + i(I(g_1) - I(g_2))$ (约定 $I(0) = 0$). 若 $f = (F_1 - F_2) + i(G_1 - G_2)$, 其中 $F_1, F_2, G_1, G_2 \in \mathcal{C}_c^+(G) \cup \{0\}$, 则 $F_1 + f_2 = F_2 + f_1$ 且 $G_1 + g_2 = G_2 + g_1$, 所以 $I(f)$ 得到相同的结果. 容易验证 $I$ 是从 $\mathcal{C}_c(G)$ 到 $\CC$ 的线性泛函, 且由构造它是正的. 由 Riesz 表示定理, 存在 $G$ 上的正则 Borel 测度 $\mu$ 使得 $I(f) = \int_G f \,\dd\mu$. 由命题 \ref{prop:I.2.6}, 这个测度是左 Haar 测度.

我们现在证明定理的第二部分 (左 Haar 测度在相差一个常数意义下唯一). 设 $\mu_1$, $\mu_2$ 是 $G$ 上的两个左 Haar 测度. 由 Riesz 表示定理中的唯一性 (以及 $\mathcal{C}_c^+(G)$ 生成 $\mathcal{C}_c(G)$ 这一事实), 只需找到一个正实数 $c$ 使得对每个 $f \in \mathcal{C}_c^+(G)$ 有 $\int_G f \,\dd\mu_1 = c\int_G f \,\dd\mu_2$. 由命题 \ref{prop:I.2.6}, 对每个 $f \in \mathcal{C}_c^+(G)$ 有 $\int_G f \,\dd\mu_2 > 0$. 所以只需证明: 若 $f, g \in \mathcal{C}_c^+(G)$, 则
\begin{equation}
\label{eq:I.2.star}
\frac{\int_G f \,\dd\mu_1}{\int_G f \,\dd\mu_2} = \frac{\int_G g \,\dd\mu_1}{\int_G g \,\dd\mu_2} \quad (*).
\end{equation}

设 $f, g \in \mathcal{C}_c^+(G)$. 设 $V_0$ 是 $1$ 的对称紧邻域, 并令
\[
A = (\operatorname{supp}(f))V_0 \cup V_0(\operatorname{supp}(f))
\]
及
\[
B = (\operatorname{supp}(g))V_0 \cup V_0(\operatorname{supp}(g)).
\]
由命题 \ref{prop:I.1.7}, $A$ 和 $B$ 紧. 若 $y \in V_0$, 则函数 $x \longmapsto f(xy) - f(yx)$ 与 $x \longmapsto g(xy) - g(yx)$ 分别支撑在 $A$ 和 $B$ 上.

设 $\varepsilon > 0$. 由命题 \ref{prop:I.1.12}, 存在 $1$ 的对称邻域 $V \subset V_0$ 使得对每个 $x \in G$ 和每个 $y \in V$, 有 $|f(xy) - f(yx)| < \varepsilon$ 且 $|g(xy) - g(yx)| < \varepsilon$. 设 $h \in \mathcal{C}_c^+(G)$ 满足 $\operatorname{supp}(h) \subset V$ 且对每个 $x \in G$ 有 $h(x) = h(x^{-1})$. 则
\begin{align*}
\Bigl(\int_G h \,\dd\mu_2\Bigr)\Bigl(\int_G f \,\dd\mu_1\Bigr) &= \int_{G \times G} h(y)f(x) \,\dd\mu_1(x)\,\dd\mu_2(y) \\
&= \int_{G \times G} h(y)f(yx) \,\dd\mu_1(x)\,\dd\mu_2(y).
\end{align*}
(我们用了 $\mu_1$ 的左不变性. 又, 我们可以应用 Fubini 定理, 因为所有函数都支撑在紧集上, 而紧集有有限测度.) 类似地,
\begin{align*}
\Bigl(\int_G h \,\dd\mu_1\Bigr)\Bigl(\int_G f \,\dd\mu_2\Bigr) &= \int_{G \times G} h(x)f(y) \,\dd\mu_1(x)\,\dd\mu_2(y) \\
&= \int_{G \times G} h(y^{-1}x)f(y) \,\dd\mu_1(x)\,\dd\mu_2(y) \\
&= \int_{G \times G} h(x^{-1}y)f(y) \,\dd\mu_1(x)\,\dd\mu_2(y) \\
&= \int_{G \times G} h(y)f(xy) \,\dd\mu_1(x)\,\dd\mu_2(y).
\end{align*}

因此
\begin{align*}
\Bigl(\int_G h \,\dd\mu_1\Bigr)\Bigl(\int_G f \,\dd\mu_2\Bigr) - \Bigl(\int_G h \,\dd\mu_2\Bigr)\Bigl(\int_G f \,\dd\mu_1\Bigr) &= \int_{G \times G} h(y)(f(xy) - f(yx)) \,\dd\mu_1(x)\,\dd\mu_2(y) \\
&\leq \varepsilon\mu_1(A)\int_G h \,\dd\mu_2,
\end{align*}
因为 $\operatorname{supp}(h) \subset V$. 除以 $\bigl(\int_G f \,\dd\mu_2\bigr)\bigl(\int_G h \,\dd\mu_2\bigr)$, 得到
\[
\Bigl(\int_G h \,\dd\mu_1\Bigr)\Bigl(\int_G h \,\dd\mu_2\Bigr)^{-1} - \Bigl(\int_G f \,\dd\mu_1\Bigr)\Bigl(\int_G f \,\dd\mu_2\Bigr)^{-1} \leq \varepsilon\mu_1(A)\Bigl(\int_G f \,\dd\mu_2\Bigr)^{-1}.
\]
类似地,
\[
\Bigl(\int_G h \,\dd\mu_1\Bigr)\Bigl(\int_G h \,\dd\mu_2\Bigr)^{-1} - \Bigl(\int_G g \,\dd\mu_1\Bigr)\Bigl(\int_G g \,\dd\mu_2\Bigr)^{-1} \leq \varepsilon\mu_1(B)\Bigl(\int_G g \,\dd\mu_2\Bigr)^{-1}.
\]
两式相加给出
\[
\left|\frac{\int_G f \,\dd\mu_1}{\int_G f \,\dd\mu_2} - \frac{\int_G g \,\dd\mu_1}{\int_G g \,\dd\mu_2}\right| \leq \varepsilon\left(\frac{\mu_1(A)}{\int_G f \,\dd\mu_2} + \frac{\mu_1(B)}{\int_G g \,\dd\mu_2}\right).
\]
由于 $\varepsilon$ 任意, 这给出所要的等式 $(*)$.
\end{proof}

我们现在来比较左、右 Haar 测度.

\begin{proposition}[\index{模函数}模函数 (modular function)]
\label{prop:I.2.8}
设 $G$ 是局部紧群, $x \in G$. 则存在 $\Delta(x) \in \R_{>0}$ 使得对 $G$ 的每个左 Haar 测度 $\mu$, 有 $\mu(Ex) = \Delta(x)\mu(E)$. 而且, $\Delta \colon G \to \R_{>0}$ 是连续群同态 (其中 $\R_{>0}$ 上的群结构由乘法给出), 且对 $G$ 的每个左 Haar 测度 $\mu$, 每个 $x \in G$ 和每个 $f \in L^1(G, \mu)$, 有
\[
\int_G R_xf \,\dd\mu = \Delta(x^{-1})\int_G f \,\dd\mu.
\]
\end{proposition}

\begin{proof}
设 $x \in G$, $\mu$ 是 $G$ 上的左 Haar 测度. 则由 $\mu_x(E) = \mu(Ex)$ 定义的测度 $\mu_x$ 也是 $G$ 上的左 Haar 测度, 所以由定理 \ref{thm:I.2.7} 中的唯一性, 存在 $\Delta(x) \in \R_{>0}$ 使得 $\mu_x = \Delta(x)\mu$, 即对 $G$ 的每个 Borel 子集 $E$ 有 $\mu(Ex) = \Delta(x)\mu(E)$. 假设 $\lambda$ 是 $G$ 上的另一个左 Haar 测度. 则再次由定理 \ref{thm:I.2.7}, 存在 $c > 0$ 使得 $\lambda = c\mu$, 于是对 $G$ 的每个 Borel 子集 $E$,
\[
\lambda(Ex) = c\mu(Ex) = c\Delta(x)\mu(E) = \Delta(x)\lambda(E).
\]
这就证明了第一个结论.

我们证明 $\Delta$ 是群同态. 设 $x, y \in G$, $E$ 是 $G$ 的满足 $\mu(E) \neq 0$ 的 Borel 子集. 则
\[
\Delta(xy)\mu(E) = \mu(Exy) = \Delta(y)\mu(Ex) = \Delta(y)\Delta(x)\mu(E),
\]
因此 $\Delta(xy) = \Delta(x)\Delta(y)$.

我们证明最后一个结论. 若 $E$ 是 $G$ 的 Borel 子集且 $x \in G$, 则 $R_x\mathbb{1}_E = \mathbb{1}_{Ex^{-1}}$, 所以由 $\Delta$ 的定义,
\[
\int_G R_x\mathbb{1}_E \,\dd\mu = \mu(Ex^{-1}) = \Delta(x^{-1})\mu(E) = \Delta(x)^{-1}\int_G \chi_E \,\dd\mu.
\]
这就对 $f = \chi_E$ 证明了结论. 一般情形通过用函数 $\mathbb{1}_E$ 的线性组合逼近 $f$ 得到.

最后, 我们证明 $\Delta$ 连续. 设 $f \in \mathcal{C}_c^+(G)$. 我们知道函数 $G \to \mathcal{C}_c(G)$, $x \longmapsto R_{x^{-1}}f$ 连续 (见注记 \ref{rem:I.1.13}), 所以函数 $G \to \CC$, $x \longmapsto \int_G R_{x^{-1}}f \,\dd\mu$ 也连续. 但刚才已见 $\int_G R_{x^{-1}}f \,\dd\mu = \Delta(x)\int_G f \,\dd\mu$, 且由命题 \ref{prop:I.2.6} 知 $\int_G f \,\dd\mu > 0$. 因此 $\Delta$ 连续.
\end{proof}

\begin{definition}[模函数与\index{幺模群}幺模群 (unimodular group)]
\label{def:I.2.9}
前一命题中的函数 $\Delta$ 称为 $G$ 的模函数 (modular function). 若 $\Delta = 1$ (即 $G$ 的某个 (或任何) 左 Haar 测度也是右 Haar 测度), 则称群 $G$ 是幺模 (unimodular) 的.
\end{definition}

\begin{remark}[推广的模因子]
\label{rem:I.2.10}
假设 $\alpha \colon G \to G$ 是同胚, 使得对每个 $x \in G$, 存在 $\beta(x) \in G$ 满足: 对每个 $y \in G$, $\alpha(xy) = \beta(x)\alpha(y)$. (例如, $\alpha$ 可以是 $G$ 中固定元素的右平移, 或具有连续逆的连续群同构.) 则我们可以推广命题 \ref{prop:I.2.8} 的构造, 得到 $\Delta(\alpha) \in \R_{>0}$ 满足: 对每个 $f \in \mathcal{C}_c(G)$, 对 $G$ 的每个左 Haar 测度 $\mu$,
\[
\Delta(\alpha)\int_G f(\alpha(x)) \,\dd\mu(x) = \int_G f(x) \,\dd\mu(x)
\]
(或等价地, 对 $G$ 的每个 Borel 子集 $E$, $\mu(\alpha(E)) = \Delta(\alpha)\mu(E)$). 而且, 若 $\beta \colon G \to G$ 满足与 $\alpha$ 相同的条件, 则 $\alpha \circ \beta$ 也满足, 且 $\Delta(\alpha \circ \beta) = \Delta(\alpha)\Delta(\beta)$.
\end{remark}

\begin{example}[幺模群的例子]
\label{exa:I.2.11}
\begin{enumerate}
\item[(1)] 任何紧群都是幺模的. 事实上, 若 $G$ 紧, 则 $\Delta(G)$ 是 $\R_{>0}$ 的紧子群, 但 $\R_{>0}$ 唯一的紧子群是 $\{1\}$. 特别地, 紧群 $G$ 有唯一的满足 $\mu(G) = 1$ 的左、右 Haar 测度 $\mu$; 我们称这个测度为 $G$ 的\index{规范化 Haar 测度}规范化 Haar 测度 (normalized Haar measure).
\item[(2)] 任何离散群都是幺模的. 事实上, 我们有 $G$ 上的一个同时是右 Haar 测度的左 Haar 测度: 计数测度.
\item[(3)] 若 $G$ 交换, 则 $G$ 上的左平移与右平移相等, 所以 $G$ 是幺模的.
\item[(4)] 群 $\GL_n(\R)$ 和 $\GL_n(\CC)$ 是幺模的. (对 $\GL_n(\R)$ 在习题 \ref{ex:I.5.3.2}(c) 中证明, 同样的证明对 $\GL_n(\CC)$ 也成立.)
\item[(5)] $M_2(\R)$ 中可逆上三角矩阵的群不是幺模的 (见习题 \ref{ex:I.5.3.2}(d)). 事实上, 其模函数为
  \[
  \Delta \colon \begin{pmatrix} a & b \\ 0 & c \end{pmatrix} \longmapsto |ac^{-1}|.
  \]
\item[(6)] 回忆换位子群 (commutator subgroup) $[G, G]$ 是由所有 $xyx^{-1}y^{-1}$ ($x, y \in G$) 生成的子群. 它是 $G$ 的正规子群, 且从 $G$ 到交换群的任何群同态在 $[G, G]$ 上都平凡. 特别地, 模函数 $\Delta$ 在 $[G, G]$ 上平凡, 所以若 $G = [G, G]$, 则 $G$ 是幺模的. 更一般地, 利用第一个例子, 我们看到若商群 $G/[G, G]$ 紧, 则 $G$ 是幺模的.
\end{enumerate}
\end{example}

\begin{proposition}[$f(x^{-1})$ 的积分公式]
\label{prop:I.2.12}
设 $G$ 是局部紧群, $\mu$ 是 $G$ 上的左 Haar 测度. 我们用 $\nu(E) = \mu(E^{-1})$ 定义 $G$ 上的右 Haar 测度 $\nu$ (见命题 \ref{prop:I.2.6}).
则对每个 $f \in \mathcal{C}_c(G)$, 有
\[
\int_G f(x^{-1}) \,\dd\mu_G(x) = \int_G f(x) \,\dd\nu(x) = \int_G \Delta(x^{-1})f(x) \,\dd\mu(x).
\]
\end{proposition}

我们也将此性质记为 $d\nu(x) = \Delta(x^{-1})d\mu(x)$, 或 $d\mu(x^{-1}) = \Delta(x^{-1})d\mu(x)$.

% ============ I.3 表示 (原书 pp.19-30) ============
\section{表示}
在本节中, $G$ 是一个拓扑群.

\subsection{连续表示}

\begin{definition}[$\Hom(V,W)$ 与 $\End(V)$ 上的拓扑]
\label{def:I.3.1.1}
若 $V$ 和 $W$ 是赋范 $\CC$-向量空间, 我们记 $\Hom(V, W)$ 为从 $V$ 到 $W$ 的有界线性算子组成的 $\CC$-向量空间, 并在其上赋予由算子范数 (operator norm) $\|\cdot\|_{\mathrm{op}}$ 给出的拓扑. 又记 $\End(V) = \Hom(V, V)$, $\GL(V) = \End(V)^{\times}$, 赋予 $\End(V)$ 的拓扑所诱导的拓扑.
\end{definition}

\begin{definition}[连续表示 (continuous representation)]
\label{def:I.3.1.2}
设 $V$ 是赋范 $\CC$-向量空间. 则 $G$ 在 $V$ 上的 (连续) 表示是群同态 $\rho \colon G \to \Aut_{\CC}(V)$ 使得作用映射 $G \times V \to V$, $(g, v) \longmapsto \rho(g)(v)$ 连续.
\end{definition}

我们用 $(\rho, V)$, $\rho$ 或常常简单地用 $V$ 来指代表示. 有时我们不显式给出映射 $\rho$ 的名字, 而把 $G$ 在 $V$ 上的作用写成 $(g, v) \longmapsto gv$.

\begin{remark}[定义的一般性]
\label{rem:I.3.1.3}
\begin{itemize}
\item 若 $V$ 是任意拓扑向量空间 (在拓扑域上), 此定义仍有意义.
\item 若 $(\rho, V)$ 是 $G$ 的连续表示, 则每个 $g \in G$ 在 $V$ 上的作用是 $V$ 的连续自同态, 所以我们得到一个群同态 $\rho \colon G \to \GL(V)$. 但除非 $V$ 有限维, 这个同态不一定连续 (见命题 \ref{prop:I.3.5.1}). 例 \ref{exa:I.3.1.11} 中定义的 $G$ 在 $L^p(G)$ 上的正则表示就是这样的例子.
\item 若 $\rho \colon G \to \GL(V)$ 是关于 $\End(V)$ 上的弱*拓扑连续的群同态, 则它不一定是连续表示. (例如, 取 $G = \GL(V)$, 赋予 $\End(V)$ 上弱*拓扑诱导的拓扑, $\rho = \mathrm{id}$. 若 $V$ 无穷维, 则这不是 $G$ 在 $V$ 上的连续表示.)
\end{itemize}
\end{remark}

\begin{example}[连续表示的例子]
\label{exa:I.3.1.4}
\begin{itemize}
\item $G$ 在 $V$ 上的\index{平凡表示}平凡表示 (trivial representation) 是由 $\rho(x) = \mathrm{id}_V$ (对每个 $x \in G$) 给出的表示. (它是连续表示.)
\item 若 $V$ 有限维, 则 $\GL(V)$ 的恒等映射是 $\GL(V)$ 在 $V$ 上的连续表示.
\item 若 $G = S^1$ 且 $n \in \Z$, 则映射 $G \to \CC$, $z \longmapsto z^n$ 是 $G$ 在 $\CC$ 上的连续表示.
\item 映射 $\rho \colon \R \to \GL_2(\CC)$, $x \longmapsto \begin{pmatrix} 1 & x \\ 0 & 1 \end{pmatrix}$ 是 $\R$ 在 $\CC^2$ 上的连续表示.
\item 关于 $G$ 在其函数空间上的表示, 见例 \ref{exa:I.3.1.11}.
\end{itemize}
\end{example}

\begin{definition}[\index{交错算子}交错算子 (intertwining operator)]
\label{def:I.3.1.5}
设 $(\rho_1, V_1)$ 和 $(\rho_2, V_2)$ 是 $G$ 的两个表示. 从 $V_1$ 到 $V_2$ 的交错算子 (或 $G$-\index{等变映射}等变映射 (G-equivariant map)) 是从 $V_1$ 到 $V_2$ 的有界 $\CC$-线性映射 $T$ 使得对每个 $g \in G$ 和每个 $v \in V_1$, 有 $T(\rho_1(g)v) = \rho_2(g)T(v)$.

我们记 $\Hom_G(V_1, V_2)$ 为从 $V_1$ 到 $V_2$ 的交错算子空间, $\End_G(V_1)$ 为从 $V_1$ 到自身的交错算子空间.

若存在交错算子 $T \colon V_1 \to V_2$ 与 $T' \colon V_2 \to V_1$ 使得 $T' \circ T = \mathrm{id}_{V_1}$ 且 $T \circ T' = \mathrm{id}_{V_2}$, 则称表示 $(\rho_1, V_1)$ 与 $(\rho_2, V_2)$ \index{同构的表示}同构 (isomorphic) (或\index{等价的表示}等价 (equivalent)).
\end{definition}

\begin{definition}[\index{子表示}子表示与不可约等概念]
\label{def:I.3.1.6}
设 $(\rho, V)$ 是 $G$ 的一个表示.
\begin{enumerate}
\item[(1)] $V$ 的子表示 (subrepresentation) (或 $G$-不变子空间 (G-invariant subspace)) 是满足对每个 $g \in G$ 都有 $\rho(g)(W) \subset W$ 的线性子空间 $W$.
\item[(2)] 若 $V \neq 0$ 且其唯一的闭 $G$-不变子空间是 $0$ 和 $V$, 则称表示 $(\rho, V)$ 不可约 (irreducible). 否则称其为可约 (reducible) 的.
\item[(3)] 若每当 $V = W_1 \oplus W_2$ (其中 $W_1$, $W_2$ 是 $V$ 的两个闭 $G$-不变子空间) 就有 $W_1 = 0$ 或 $W_2 = 0$, 则称表示 $(\rho, V)$ \index{不可分解表示}不可分解 (indecomposable).
\item[(4)] 若存在一族 $V$ 的闭 $G$-不变子空间 $(W_i)_{i \in I}$ 处于直和中且 $\bigoplus_{i \in I} W_i$ 在 $V$ 中稠密, 则称表示 $(\rho, V)$ \index{半单表示}半单 (semisimple). (若 $I$ 有限, 则直和也在 $V$ 中闭, 所以这蕴含 $V = \bigoplus_{i \in I} W_i$.)
\end{enumerate}
\end{definition}

\begin{remark}[$G$-稳定子空间的闭包]
\label{rem:I.3.1.7}
若 $(\rho, V)$ 是 $G$ 的表示且 $W \subset V$ 是 $G$-稳定子空间, 则其闭包 $\overline{W}$ 也被 $G$ 稳定.
\end{remark}

\begin{example}[不可分解但不不可约的表示]
\label{exa:I.3.1.8}
$\R$ 在 $\CC^2$ 上由 $\rho(x) = \begin{pmatrix} 1 & x \\ 0 & 1 \end{pmatrix}$ 给出的表示 $\rho$ 是不可分解的但不是不可约的.
\end{example}

\begin{lemma}[交错算子的核与像]
\label{lem:I.3.1.9}
设 $(\rho_1, V_1)$ 和 $(\rho_2, V_2)$ 是 $G$ 的两个表示, $T \colon V_1 \to V_2$ 是一个交错算子. 则 $\Ker(T)$ 是 $V_1$ 的子表示, $\Im(T)$ 是 $V_2$ 的子表示.
\end{lemma}

\begin{proof}
设 $v \in \Ker(T)$, $g \in G$. 则 $T(\rho_1(g)(v)) = \rho_2(g)(T(v)) = 0$, 所以 $\rho_1(g)(v) \in \Ker(T)$.

现在设 $w \in \Im(T)$, 并选取 $v \in V_1$ 使得 $w = T(v)$. 则 $\rho_2(g)(w) = T(\rho_1(g)(v)) \in \Im(T)$.
\end{proof}

\begin{proposition}[连续表示的一个判定]
\label{prop:I.3.1.10}
设 $V$ 是赋范向量空间, $\rho \colon G \to \End(V)$ 是乘法映射. 我们记 $\|\cdot\|_{\mathrm{op}}$ 为 $\End(V)$ 上的算子范数. 假设:
\begin{enumerate}
\item[(a)] 对每个 $g \in G$, 有 $\|\rho(g)\|_{\mathrm{op}} \leq 1$;
\item[(b)] 对每个 $v \in V$, 映射 $G \to V$, $g \longmapsto \rho(g)(v)$ 连续.
\end{enumerate}
则 $(\rho, V)$ 是 $G$ 的连续表示.
\end{proposition}

\begin{proof}
设 $g_0 \in G$, $v_0 \in V$, $\varepsilon > 0$. 我们要找 $g$ 在 $G$ 中的邻域 $U$ 和 $\delta > 0$ 使得: $g \in U$ 且 $\|v - v_0\| < \delta \Rightarrow \|\rho(g)(v) - \rho(g_0)(v_0)\| < \varepsilon$.

选取 $g$ 在 $G$ 中的邻域 $U$ 使得: $g \in U \Rightarrow \|\rho(g)(v_0) - \rho(g_0)(v_0)\| < \varepsilon/2$, 并取 $\delta = \varepsilon/2$. 则若 $g \in U$ 且 $\|v - v_0\| < \delta$, 有
\begin{multline*}
\|\rho(g)(v) - \rho(g_0)(v_0)\| \leq \|\rho(g)(v) - \rho(g)(v_0)\| + \|\rho(g)(v_0) - \rho(g_0)(v_0)\| \\
< \|\rho(g)\|_{\mathrm{op}}\|v - v_0\| + \varepsilon/2 < \varepsilon/2 + \varepsilon/2 = \varepsilon,
\end{multline*}
因为 $\|\rho(g)\|_{\mathrm{op}} \leq 1$.
\end{proof}

\begin{example}[正则表示]
\label{exa:I.3.1.11}
\begin{enumerate}
\item[(1)] 我们已对每个 $x \in G$ 定义了函数空间上的两个自同态 $L_x$ 和 $R_x$, 且这些自同态保持 $\|\cdot\|_{\infty}$. 所以由命题 \ref{prop:I.3.1.10} 和注记 \ref{rem:I.1.13}, 它们定义了 $G$ 在 $\mathcal{C}_c(G)$ 上的两个表示.
\item[(2)] 假设 $G$ 是局部紧 Hausdorff 的. 我们固定 $G$ 上的左 Haar 测度 $dx$, 并对 $1 \leq p \leq \infty$ 记 $L^p(G)$ 为此测度对应的 $L^p$ 空间. 测度的左不变性蕴含算子 $L_x$ 保持 $L^p$ 范数, 所以我们得到 $G$ 在 $L^p(G)$ 上的 $\CC$-线性左作用; 由命题 \ref{prop:I.3.1.10}, 要证明它是一个表示, 只需证明: 若 $f \in L^p(G)$, 则映射 $G \to L^p(G)$, $x \longmapsto L_xf$ 连续. 当 $p = \infty$ 时这不一定成立, 但对 $1 \leq p < \infty$ 成立, 见下面的命题 \ref{prop:I.3.1.13}. 所以对 $1 \leq p < \infty$ 我们得到 $G$ 在 $L^p(G)$ 上的表示.
  若我们改取 $G$ 上的右 Haar 测度, 则算子 $R_x$ 将对 $1 \leq p < \infty$ 定义 $G$ 在 $L^p(G)$ 上的表示. 所以若 $G$ 是幺模的, 我们得到 $G$ 在 $L^p(G)$ 上的两个可交换的表示.
\end{enumerate}
\end{example}

\begin{definition}[左/\index{右正则表示}右正则表示 (left/right regular representation)]
\label{def:I.3.1.12}
设 $G$ 是带有左 (相应地右) Haar 测度 $dx$ 的局部紧群, $L^2(G)$ 是对应的 $L^2$ 空间. 由算子 $L_x$ (相应地 $R_x$) 给出的 $G$ 在 $L^2(G)$ 上的表示称为 $G$ 的左 (相应地右) 正则表示.
\end{definition}

\begin{proposition}[$L^p$ 中平移的连续性]
\label{prop:I.3.1.13}
设 $G$ 是局部紧群, $\mu$ 是 $G$ 上的左 Haar 测度, $L^p(G)$ 是对应的 $L^p$ 空间. 假设 $1 \leq p < \infty$.
则对每个 $f \in L^p(G)$, 当 $x \to 1$ 时有 $\|L_xf - f\|_p \to 0$ 且 $\|R_xf - f\|_p \to 0$.
\end{proposition}

\begin{proof}
先假设 $f \in \mathcal{C}_c(G)$, 并固定 $1$ 的紧邻域 $V$. 则 $K := V(\operatorname{supp} f) \cup (\operatorname{supp} f)V$ 由命题 \ref{prop:I.1.7} 是紧的, 所以 $\mu(K) < +\infty$. 对每个 $x \in V$, 有 $\operatorname{supp}(f), \operatorname{supp}(L_xf), \operatorname{supp}(R_xf) \subset K$, 所以 $\|L_xf - f\|_p \leq \mu(K)^{1/p}\|L_xf - f\|_{\infty}$ 且 $\|R_xf - f\|_p \leq \mu(K)^{1/p}\|R_xf - f\|_{\infty}$. 结论由命题 \ref{prop:I.1.12} 得出.

现在设 $f$ 是 $L^p(G)$ 的任意元素. 我们仍固定 $1$ 的紧邻域 $V$, 并令 $C = \sup_{x \in V}\Delta(x)^{-1/p}$. 设 $\varepsilon > 0$. 存在 $g \in \mathcal{C}_c(G)$ 使得 $\|f - g\|_p < \varepsilon$. 则对 $x \in V$,
\[
\|L_xf - f\|_p \leq \|L_x(f - g)\|_p + \|L_xg - g\|_p + \|g - f\|_p \leq 2\varepsilon + \|L_xg - g\|_p
\]
(因为 $\|L_x(f - g)\|_p = \|f - g\|_p$) 且
\[
\|R_xf - f\|_p \leq \|R_x(f - g)\|_p + \|R_xg - g\|_p + \|g - f\|_p \leq (1 + C)\varepsilon + \|R_xg - g\|_p
\]
(因为 $\|R_x(f - g)\|_p = \Delta(x)^{-1/p}\|f - g\|_p$). 由证明的第一部分, 当 $x$ 趋于 $1$ 时 $\|L_xg - g\|_p$ 和 $\|R_xg - g\|_p$ 趋于 $0$, 所以我们可以找到 $1$ 的邻域 $U \subset V$ 使得对 $x \in U$ 有 $\|L_xf - f\|_p \leq 3\varepsilon$ 且 $\|R_xf - f\|_p \leq (2 + C)\varepsilon$.
\end{proof}

\subsection{\index{酉表示}酉表示}

回忆 (复) \index{希尔伯特空间}希尔伯特空间是带 Hermitian 内积\footnote{我们总假定 Hermitian 内积对第一个变量是 $\CC$-线性的.} 的 $\CC$-向量空间 $V$, 使得 $V$ 对相应范数完备. 若 $V$ 是带 Hermitian 内积的有限维 $\CC$-向量空间, 则它自动完备, 因此是希尔伯特空间. 我们通常用 $\langle \cdot, \cdot \rangle$ 表示所有 Hermitian 内积空间上的内积 (除非另有说明).

\begin{notation}[伴随算子与正交补]
\label{not:I.3.2.1}
设 $V$ 和 $W$ 是 Hermitian 内积空间. 对每个连续 $\CC$-线性映射 $T \colon V \to W$, 若 $T$ 的伴随 (adjoint) 存在, 我们记 $T^* \colon W \to V$. 回忆对每个 $v \in V$ 和 $w \in W$ 有 $\langle T(v), w \rangle = \langle v, T^*(w) \rangle$, 且若 $V$ 和 $W$ 是希尔伯特空间则 $T^*$ 总存在.
若 $V'$ 是 $V$ 的子空间, 我们记 $V'^{\perp}$ 为 $V'$ 的正交补 (orthogonal complement); 它由
\[
(V')^{\perp} = \{v \in V \mid \forall v' \in V',\, \langle v, v' \rangle = 0\}
\]
定义.

最后, 我们记 $\UU(V)$ 为 $V$ 的酉自同构 (unitary endomorphism) 群, 即保持内积的 $V$ 的自同构 $T$ (对所有 $v, w \in V$ 有 $\langle T(v), T(w) \rangle = \langle v, w \rangle$) 的群. 酉自同构 $T$ 自动有界且可逆 (逆等于 $T^*$).
\end{notation}

下面的结果是命题 \ref{prop:I.3.1.10} (以及酉算子范数为 $1$ 这一事实) 的直接推论.

\begin{corollary}[酉表示的等价条件]
\label{cor:I.3.2.2}
若 $V$ 是希尔伯特空间且 $\rho \colon G \to \UU(V)$ 是群同态, 则以下条件等价:
\begin{enumerate}
\item[(1)] 映射 $G \times V \to V$, $(g, v) \longmapsto \rho(g)(v)$ 连续.
\item[(2)] 对每个 $v \in V$, 映射 $G \to V$, $g \longmapsto \rho(g)(v)$ 连续.
\end{enumerate}
\end{corollary}

\begin{definition}[酉表示 (unitary representation)]
\label{def:I.3.2.3}
若 $V$ 是希尔伯特空间, 则 $G$ 在 $V$ 上的酉表示是满足推论 \ref{cor:I.3.2.2} 的条件的群同态 $\rho \colon G \to \UU(V)$.
\end{definition}

这些表示是我们研究的主要对象.

\begin{example}[$L^2$ 空间上的酉表示]
\label{exa:I.3.2.4}
若 $(X, \mu)$ 是任意测度空间, 则 $L^2(X)$ 是希尔伯特空间, 内积为
\[
\langle f, g \rangle = \int_X f(x)\overline{g(x)} \,\dd\mu(x).
\]
所以若 $G$ 是局部紧群, 则 $G$ 的\index{左正则表示}左正则表示和右正则表示是 $G$ 的酉表示 (若 $G$ 幺模则在同一个空间上).
\end{example}

\begin{remark}[$\rho$ 不一定连续]
\label{rem:I.3.2.5}
注意 $\rho$ 一般来说仍不一定是连续映射. (除非 $\dim_{\CC}V < +\infty$.) 例如, 对 $S^1$ 的左正则表示它不连续.
又注意, 证明中不需要 $V$ 的完备性, 所以推论 \ref{cor:I.3.2.2} 对任意 Hermitian 内积空间都成立.
\end{remark}

\begin{lemma}[不变子空间的正交补]
\label{lem:I.3.2.6}
设 $(\rho, V)$ 是 $G$ 的酉表示. 则对 $V$ 的每个 $G$-不变子空间 $W$, 子空间 $W^{\perp}$ 也被 $G$ 不变.

特别地, 若 $W$ 是 $V$ 的闭 $G$-不变子空间, 则 $V = W \oplus W^{\perp}$, 其中 $W^{\perp}$ 是闭 $G$-不变子空间.
\end{lemma}

\begin{proof}
设 $v \in W^{\perp}$, $g \in G$. 则对每个 $w \in W$,
\[
\langle \rho(g)(v), w \rangle = \langle v, \rho(g)^{-1}w \rangle = 0
\]
(最后一个等式来自 $\rho(g)^{-1}w \in W$), 因此 $\rho(g)(v) \in W^{\perp}$.
\end{proof}

\begin{lemma}[交错算子的伴随]
\label{lem:I.3.2.7}
设 $(\rho_1, V_1)$ 和 $(\rho_2, V_2)$ 是 $G$ 的两个酉表示, $T \colon V_1 \to V_2$ 是交错算子. 则 $T^* \colon V_2 \to V_1$ 也是交错算子.
\end{lemma}

\begin{proof}
设 $w \in V_2$, $g \in G$. 则对每个 $v \in V_1$,
\begin{multline*}
\langle v, T^*(\rho_2(g)(w)) \rangle = \langle T(v), \rho_2(g)(w) \rangle = \langle \rho_2(g)^{-1}T(v), w \rangle = \langle T(\rho_1(g)^{-1}(v)), w \rangle \\
= \langle \rho_1(g)^{-1}(v), T^*(w) \rangle = \langle v, \rho_1(g)T^*(w) \rangle.
\end{multline*}
所以 $T^*(\rho_2(g)(w)) = \rho_1(g)(T^*(w))$.
\end{proof}

\begin{theorem}[平均化内积]
\label{thm:I.3.2.8}
假设群 $G$ 是紧 Hausdorff 的. 设 $(V, \langle \cdot, \cdot \rangle_0)$ 是希尔伯特空间, $\rho \colon G \to \GL(V)$ 是 $G$ 在 $V$ 上的连续表示. 则存在 $V$ 上的 Hermitian 内积 $\langle \cdot, \cdot \rangle$ 满足以下性质:
\begin{enumerate}
\item[(1)] 存在实数 $c, C > 0$ 使得对每个 $v \in V$ 有 $c|hv, vi_0| \leq |\langle v, v \rangle| \leq C|\langle v, v \rangle_0|$. 换句话说, 两个内积给出的范数等价, 所以 $V$ 对内积 $\langle \cdot, \cdot \rangle$ 仍是希尔伯特空间.
\item[(2)] 表示 $\rho$ 对内积 $\langle \cdot, \cdot \rangle$ 是酉的.
\end{enumerate}
\end{theorem}

\begin{remarkstar}[I.3.1]
(a) 若 $V$ 不可约, 则由 \index{Schur 引理}Schur 引理 (见定理 \ref{thm:I.3.4.1}) 知, 这个内积在相差一个常数意义下唯一.

(b) 对非紧群这是错的. 例如, 考虑 $\R$ 在 $\CC^2$ 上由 $\rho(t) = \begin{pmatrix} 1 & t \\ 0 & 1 \end{pmatrix}$ 给出的表示 $\rho$. 不存在使该表示为酉的 $\CC^2$ 上的内积 (否则 $\rho(\R)$ 将是该内积的酉群的闭子群, 从而是紧的, 但这是不可能的, 因为 $\rho(\R) \simeq \R$).
\end{remarkstar}

\begin{proof}[定理的证明]
我们用下面的公式定义 $\langle \cdot, \cdot \rangle \colon V \times V \to \CC$: 对所有 $v, w \in V$,
\[
\langle v, w \rangle = \int_G \langle \rho(g)v, \rho(g)w \rangle_0 \,\dd g,
\]
其中 $dg$ 是 $G$ 上的规范化 Haar 测度. 这定义了 $V$ 上的 Hermitian 形式, 且由测度的左不变性, 对每个 $v, w \in V$ 和 $g \in G$ 有 $\langle \rho(g)v, \rho(g)w \rangle = \langle v, w \rangle$.

若我们证明了性质 (1), 它也将蕴含 $\langle \cdot, \cdot \rangle$ 正定 (从而是内积), 于是大功告成. 设 $v \in V$. 则把 $v$ 分别映到 $\rho(g)(v)$ 和 $\rho(g)^{-1}(v)$ 的两个映射 $G \to V$ 连续. 由于 $G$ 紧, 它们都有界. 由共鸣定理 (定理 \ref{thm:I.3.2.11}), 存在 $A, B \in \R_{>0}$ 使得对每个 $g \in G$ 有 $\|\rho(g)^{-1}\|_{\mathrm{op}} \leq A$ 且 $\|\rho(g)\|_{\mathrm{op}} \leq B$. 由算子范数的次乘性, 第一个不等式蕴含对每个 $g \in G$ 有 $\|\rho(g)\|_{\mathrm{op}} \leq A^{-1}$. 所以 $\langle \cdot, \cdot \rangle$ 的定义 (以及 $G$ 体积为 $1$ 的事实) 给出性质 (1), 其中 $c = A^{-2}$, $C = B^2$.
\end{proof}

\begin{corollary}[紧群的有限维表示半单]
\label{cor:I.3.2.9}
若 $G$ 是紧 Hausdorff 的, 则 $G$ 的每个非零有限维连续表示都是半单的.
\end{corollary}

\begin{proof}
由定理, 我们可以假设表示是酉的. 我们对 $\dim V$ 归纳证明此推论. 若 $\dim V \leq 1$ 结论显然, 所以假设 $\dim V \geq 2$ 且已知结论对所有维数更小的空间成立. 若 $V$ 不可约, 则已完成. 否则, 存在 $G$-不变子空间 $W \subsetneq V$ 且 $W \neq 0$. 该子空间因有限维而闭, 且由引理 \ref{lem:I.3.2.6} 有 $V = W \oplus W^{\perp}$, 其中 $W^{\perp}$ 不变. 由于 $\dim(W), \dim(W^{\perp}) < \dim(V)$, 可以对 $W$ 和 $W^{\perp}$ 应用归纳假设并得出它们半单. 但此时它们的直和 $V$ 也半单.
\end{proof}

\begin{remark}[半单性的进一步讨论]
\label{rem:I.3.2.10}
这对紧群的无穷维酉表示仍然成立 (但更难证明, 见定理 \ref{thm:IV.2.1}), 但对非紧群的无穷维酉表示是错的 (例如, 若 $G$ 交换且不紧, 则由推论 \ref{cor:I.3.4.4} 和习题 \ref{ex:III.6.1.2}(c), 其正则表示不半单), 对非紧群的有限维 (非酉) 表示也是错的 (见例 \ref{exa:I.3.1.8}).
\end{remark}

\begin{theorem}[共鸣定理 (Uniform Boundedness Principle)]
\label{thm:I.3.2.11}
设 $V$ 和 $W$ 是赋范向量空间, 并假设 $V$ 是巴拿赫空间 (即它对范数完备). 设 $(T_i)_{i \in I}$ 是从 $V$ 到 $W$ 的有界线性算子族.
若族 $(T_i)_{i \in I}$ 逐点有界 (即对每个 $v \in V$ 有 $\sup_{i \in I}\|T_i(v)\| < +\infty$), 则有界 (即 $\sup_{i \in I}\|T_i\|_{\mathrm{op}} < +\infty$).
\end{theorem}

\begin{proof}
假设 $\sup_{i \in I}\|T_i\|_{\mathrm{op}} = +\infty$, 并选取 $I$ 的元素序列 $(i_n)_{n \geq 0}$ 使得 $\|T_{i_n}\|_{\mathrm{op}} \geq 4^n$. 我们用如下方式定义 $V$ 的元素序列 $(v_n)_{n \geq 0}$:
\begin{itemize}
\item $v_0 = 0$;
\item 对 $n \geq 1$, 由下面的引理 \ref{lem:I.3.2.12}, 可以找到 $V$ 的元素 $v_n$ 使得 $\|v_n - v_{n-1}\| \leq 3^{-n}$ 且 $\|T_{i_n}(v_n)\| \geq \frac{2}{3} \cdot 3^{-n}\|T_{i_n}\|_{\mathrm{op}}$.
\end{itemize}
对 $m \geq n$ 有 $\|v_n - v_m\| \leq \frac{1}{2} \cdot 3^{-n}$, 所以序列 $(v_n)_{n \geq 0}$ 是 Cauchy 列; 由于 $V$ 完备, 它有极限 $v$, 且对每个 $n \geq 0$ 有 $\|v_n - v\| \leq \frac{1}{2} \cdot 3^{-n}$. 不等式 $\|T_{i_n}(v_n)\| \geq \frac{2}{3} \cdot 3^{-n}\|T_{i_n}\|_{\mathrm{op}}$ 和三角不等式现在蕴含 $\|T_{i_n}(x)\| \geq \frac{1}{6} \cdot 3^{-n}\|T_{i_n}\|_{\mathrm{op}} \geq \frac{1}{6}\bigl(\frac{4}{3}\bigr)^n$, 所以序列 $(\|T_{i_n}(x)\|)_{n \geq 0}$ 无界, 这与假设矛盾.
\end{proof}

\begin{lemma}[范数的局部下界]
\label{lem:I.3.2.12}
设 $V$ 和 $W$ 是赋范向量空间, $T \colon V \to W$ 是有界线性算子. 则对任意 $v \in V$ 和 $r > 0$, 有
\[
\sup_{v' \in B(v,r)} \|T(v')\| \geq r\|T\|_{\mathrm{op}},
\]
其中 $B(v, r) = \{v' \in V \mid \|v - v'\| < r\}$.
\end{lemma}

\begin{proof}
对每个 $x \in V$, 有
\[
\|T(x)\| \leq \frac{1}{2}(\|T(v + x)\| + \|T(v - x)\|) \leq \max(\|T(v + x)\|, \|T(v - x)\|).
\]
对 $x \in B(0, r)$ 取上确界就得到引理中的不等式.
\end{proof}

最后, 我们有下面的结果, 其证明用到 Schur 引理 (定理 \ref{thm:I.3.4.1}), 并见习题 \ref{ex:I.5.5.9}.

\begin{theorem}[紧群的不可约酉表示有限维]
\label{thm:I.3.2.13}
若 $G$ 是紧群, 则 $G$ 的每个不可约酉表示都是有限维的.\footnote{取自 Alan Sokal 的一篇论文.}\footnote{精确的参考文献.}
\end{theorem}

\subsection{\index{循环表示}循环表示}

\begin{definition}[循环表示 (cyclic representation)]
\label{def:I.3.3.1}
设 $(\rho, V)$ 是 $G$ 的连续表示, $v \in V$. 则 $\operatorname{Span}\{\rho(g)(v) \mid g \in G\}$ 的闭包 $W$ 是 $V$ 的子表示, 称为由 $v$ 生成的\index{循环子空间}循环子空间 (cyclic subspace).
若 $V = W$, 则称 $V$ 是循环表示, $v$ 是 $V$ 的\index{循环向量}循环向量 (cyclic vector).
\end{definition}

\begin{example}[循环表示的例]
\label{exa:I.3.3.2}
\index{不\index{可约表示}可约表示}不可约表示是循环的, 且其每个非零向量都是循环向量.
其逆不真. 例如, 考虑置换群 $S_n$ 在 $\CC^n$ 上由 $\rho(\sigma)(x_1, \ldots, x_n) = (x_{\sigma^{-1}(1)}, \ldots, x_{\sigma^{-1}(n)})$ 定义的表示 $\rho$, 并取 $v = (1, 0, \ldots, 0) \in \CC^n$. 则集合 $\rho(S_n)(v)$ 是 $\CC^n$ 的典范基, 因此生成 $\CC^n$, 所以 $v$ 是 $\rho$ 的循环向量. 但 $\rho$ 不可约吗? 不是, 因为 $\CC(1, 1, \ldots, 1)$ 是一个子表示.
\end{example}

\begin{proposition}[循环分解]
\label{prop:I.3.3.3}
$G$ 的每个酉表示都是循环表示的直和.
\end{proposition}

若指标集无穷, 我们理解的直和是闭直和 (即代数直和的闭包).

\begin{proof}
设 $(\pi, V)$ 是 $G$ 的酉表示. 由 Zorn 引理, 可以找到 $V$ 的一族两两正交的循环子空间 $(W_i)_{i \in I}$ 的极大族. 假设 $V$ 不是 $W_i$ 的直和, 则存在非零向量 $v \in \bigl(\bigoplus_{i \in I} W_i\bigr)^{\perp}$. 由引理 \ref{lem:I.3.2.6}, 由 $v$ 生成的循环子空间包含于 $\bigl(\bigoplus_{i \in I} W_i\bigr)^{\perp}$, 这与族 $(W_i)_{i \in I}$ 的极大性矛盾. 因此 $V = \bigoplus_{i \in I} W_i$.
\end{proof}

\subsection{Schur 引理}

下面的定理是基本的. 在我们得到希尔伯特空间正规自同态的\index{谱}谱定理 (定理 \ref{thm:II.4.1}) 之前, 我们无法完全证明它.

\begin{theorem}[Schur 引理 (Schur's lemma)]
\label{thm:I.3.4.1}
设 $(\rho_1, V_1)$ 和 $(\rho_2, V_2)$ 是 $G$ 的两个表示, $T \colon V_1 \to V_2$ 是交错算子.

\begin{enumerate}
\item[(1)] 若 $V_1$ 不可约, 则 $T$ 要么是零要么是单射.
\item[(2)] 若 $V_2$ 不可约, 则 $T$ 是零或有稠密像.
\item[(3)] 假设 $V_1$ 是酉的. 则 $V_1$ 不可约当且仅当 $\End_G(V_1) = \CC \cdot \mathrm{id}_{V_1}$.
\item[(4)] 假设 $V_1$ 和 $V_2$ 是酉且不可约的. 则 $\Hom_G(V_1, V_2)$ 的维数是 $0$ (若 $V_1$ 和 $V_2$ 不同构) 或 $1$ (若 $V_1$ 和 $V_2$ 同构).
\end{enumerate}
\end{theorem}

\begin{proof}
我们证明前两点. 由引理 \ref{lem:I.3.1.9}, $\Ker(T)$ 和 $\Im(T)$ 是 $V_1$ 和 $V_2$ 的 $G$-不变子空间. 而且 $\Ker(T)$ 是 $V_1$ 的闭子空间. 若 $V_1$ 不可约, 则其唯一的闭不变子空间是 $0$ 和 $V_1$; 这给出第一点. 若 $V_2$ 不可约, 则其唯一的闭不变子空间是 $0$ 和 $V_2$; 这给出第二点.

我们证明第三点. 先假设 $V_1$ 不可约. 则它有闭不变子空间 $W$ 满足 $0 \neq W \neq V_1$, 且到 $W$ 上的正交投影是 $G$-等变自同态 (由引理 \ref{lem:I.3.4.3}). 所以 $\End_G(V_1)$ 严格包含 $\CC \cdot \mathrm{id}_{V_1}$.

现在假设 $V_1$ 不可约, 设 $T \in \End_G(V_1)$. 我们要证明 $T \in \CC\mathrm{id}_{V_1}$. 若 $V_1$ 有限维, 则 $T$ 有特征值 $\lambda$, 于是 $\Ker(T - \lambda\mathrm{id}_{V_1})$ 是 $V_1$ 的非零 $G$-不变子空间, 从而是 $V_1$, 得到 $T = \lambda\mathrm{id}_{V_1}$. 一般情形下, 我们仍知道每个 $T \in \End(V)$ 都有非空谱 (由定理 \ref{thm:II.1.1.3}), 但若 $\lambda$ 在 $T$ 的谱中, 我们只知道 $T - \lambda\mathrm{id}_V$ 不可逆, 不知道 $\Ker(T - \lambda\mathrm{id}_V) \neq 0$. 所以不能应用同样的策略. 我们改用谱定理 (定理 \ref{thm:II.4.1}) 的一个推论. 注意 $\End(V_1)$ 的子群 $\rho_1(G)$ 满足推论 \ref{cor:II.4.4} 的假设, 因为 $V_1$ 不可约, 所以它在 $\End(V_1)$ 中的中心化子等于 $\CC\mathrm{id}_{V_1}$; 但这个中心化子正是 $\End_G(V_1)$, 于是得证.

我们证明第四点. 设 $T \colon V_1 \to V_2$ 是交错算子. 则由引理 \ref{lem:I.3.2.7}, $T^* \colon V_2 \to V_1$ 也是交错算子, 所以 $T^*T \in \End_G(V_1)$ 且 $TT^* \in \End_G(V_2)$. 由第三点, 存在 $c \in \CC$ 使得 $T^*T = c\mathrm{id}_{V_1}$. 若 $c \neq 0$, 则 $T$ 单射且 $\Im(T)$ 闭 (因为对每个 $v \in V_1$ 有 $\|T(v)\| \geq \frac{|c|}{\|T\|_{\mathrm{op}}}\|v\|$, 见引理 \ref{lem:I.3.4.2}), 所以由第二点 $T$ 是同构, 其逆为 $c^{-1}T^*$; 因此 $V_1$ 和 $V_2$ 同构, 且 $\Hom_G(V_1, V_2) \simeq \End_G(V_1)$ 是 $1$ 维的. 假设 $c = 0$. 若 $T \neq 0$, 则由第二点它有稠密像, 但此时由第一点 $T^* = 0$, 从而 $T = (T^*)^* = 0$, 这荒谬; 所以 $T = 0$. 因此我们证明了: 若 $\Hom_G(V_1, V_2) \neq 0$, 则 $V_1$ 和 $V_2$ 必同构; 这就完成了第四点的证明.
\end{proof}

\begin{lemma}[压缩映射的闭像]
\label{lem:I.3.4.2}
设 $V$, $W$ 是赋范向量空间, $T \colon V \to W$ 是有界线性算子. 假设 $V$ 完备. 若存在 $c > 0$ 使得对每个 $v \in V$ 有 $\|T(v)\| \geq c\|v\|$, 则 $\Im(T)$ 闭.
\end{lemma}

\begin{proof}
设 $(v_n)_{n \in \N}$ 是 $V$ 的元素序列, 使得序列 $(T(v_n))_{n \in \N}$ 收敛到 $w \in W$. 我们要证明 $w \in \Im(T)$. 注意对所有 $n, m \in \N$, 有 $\|v_n - v_m\| \leq c^{-1}\|T(v_n) - T(v_m)\|$. 这蕴含 $(v_n)_{n \in \N}$ 是 Cauchy 列, 所以它因 $V$ 完备而有极限 $v \in V$. 由于 $T$ 连续, $w = \lim_{n \to +\infty} T(v_n) = T(v)$, 所以 $w \in \Im(T)$.
\end{proof}

\begin{lemma}[正交投影的等变性]
\label{lem:I.3.4.3}
设 $(\rho, V)$ 是 $G$ 的酉表示, $W$ 是 $V$ 的闭子空间, $\pi$ 是到 $W$ 上的正交投影, 看作 $V$ 的线性自同态.
则 $W$ 被 $G$ 不变当且仅当 $\pi$ 是 $G$-等变的.
\end{lemma}

\begin{proof}
假设 $\pi$ 是 $G$-等变的. 设 $w \in W$, $g \in G$. 则 $\rho(g)(w) = \rho(g)(\pi(w)) = \pi(\rho(g)(w)) \in W$. 所以 $W$ 被 $G$ 不变.

反之, 假设 $W$ 被 $G$ 不变. 由引理 \ref{lem:I.3.2.6}, 其正交补 $W^{\perp}$ 也被 $G$ 不变. 设 $v \in V$, $g \in G$. 记 $w = \pi(g)$ 且 $w' = g - \pi(g)$. 则 $\rho(g)(v) = \rho(g)(w) + \rho(g)(w')$, 其中 $\rho(g)(w) \in W$ 且 $\rho(g)(w') \in W^{\perp}$, 所以 $\pi(\rho(g)(v)) = \rho(g)(w)$.
\end{proof}

\begin{corollary}[交换群的不可约酉表示]
\label{cor:I.3.4.4}
若 $G$ 交换, 则 $G$ 的每个不可约酉表示都是 $1$ 维的.
\end{corollary}

因此 $G$ 的每个不可约酉表示都等价于唯一一个连续群同态 $G \to S^1$.

\begin{proof}
设 $(\rho, V)$ 是不可约酉表示. 由于 $G$ 交换, 对所有 $x, y \in G$, 算子 $\rho(x)$ 与 $\rho(y)$ 交换, 所以对每个 $x \in G$ 有 $\rho(x) \in \End_G(V)$. 由 Schur 引理, 这蕴含对每个 $x \in G$ 有 $\rho(x) \in \CC \cdot \mathrm{id}_V$. 特别地, $V$ 的每个线性子空间都被 $G$ 不变. 由于 $V$ 不可约, 它没有非平凡的闭不变子空间, 所以它必是 $1$ 维的.
\end{proof}

\begin{example}[$\R$ 的不可约酉表示]
\label{exa:I.3.4.5}
设 $G = \R$. 则 $G$ 的每个不可约酉表示都具有 $\rho_y \colon x \longmapsto e^{ixy}$ 的形式, 其中 $y \in \R$. 表示 $\rho_y$ 通过 $S^1 \simeq \R/\Z$ 分解当且仅当 $y \in 2\pi\Z$. (见习题 \ref{ex:I.5.4.1}(c) 和 (d).)
\end{example}

\subsection{有限维表示}

回忆若 $V$ 是有限维 $\CC$-向量空间, 则 $V$ 上的所有范数等价.\footnote{参考文献 ?} 所以 $V$ 有典范拓扑, $\End(V)$ 也有 (因为它也是有限维向量空间).

\begin{proposition}[有限维表示的条件]
\label{prop:I.3.5.1}
设 $V$ 是赋范 $\CC$-向量空间, $\rho \colon G \to \GL(V)$ 是群同态. 考虑以下条件.
\begin{enumerate}
\item[(i)] 映射 $G \times V \to V$, $(g, v) \longmapsto \rho(g)(v)$ 连续 (即 $\rho$ 是 $G$ 在 $V$ 上的连续表示).
\item[(ii)] 对每个 $v \in V$, 映射 $G \to V$, $g \longmapsto \rho(g)(v)$ 连续.
\item[(iii)] 映射 $\rho \colon G \to \GL(V)$ 连续.
\end{enumerate}
则 $(iii) \Rightarrow (i) \Rightarrow (ii)$. 若此外 $V$ 有限维, 则三个条件等价.
\end{proposition}

\begin{proof}
(i)$\Rightarrow$(ii) 显然.
(ii)$\Rightarrow$(iii): 假设 $V$ 有限维, 设 $(e_1, \ldots, e_n)$ 是 $V$ 的基, 设 $\|\cdot\|$ 是由 $\Bigl\|\sum_{i=1}^n x_ie_i\Bigr\| = \sup_{1 \leq i \leq n}|x_i|$ 定义的 $V$ 上的范数. 我们在 $\End(V)$ 上使用对应的算子范数并仍记之为 $\|\cdot\|$. 设 $g_0 \in G$, $\varepsilon > 0$; 我们要找 $g_0 \in G$ 的邻域 $U$ 使得: $g \in U \Rightarrow \|\rho(g) - \rho(g_0)\| \leq \varepsilon$.
对每个 $i \in \{1, \ldots, n\}$, 由假设, 函数 $G \to V$, $g \longmapsto \rho(g)(e_i)$ 连续, 所以存在 $g_0$ 在 $G$ 中的邻域 $U_i$ 使得: $g \in U_i \Rightarrow \|\rho(g)(e_i) - \rho(g_0)(e_i)\| \leq \varepsilon/n$. 令 $U = \bigcap_{i=1}^n U_i$. 则若 $g \in U$, 对每个 $v = \sum_{i=1}^n x_ie_i \in V$,
\[
\|\rho(g)(v) - \rho(g_0)(v)\| \leq \sum_{i=1}^n |x_i|\|\rho(g)(e_i) - \rho(g_0)(e_i)\| < \sum_{i=1}^n |x_i|\varepsilon/n \leq \varepsilon\|v\|,
\]
这意味着 $\|\rho(g) - \rho(g_0)\| \leq \varepsilon$.
(iii)$\Rightarrow$(i): 设 $g_0 \in G$, $v_0 \in V$, $\varepsilon > 0$. 我们要找 $g$ 在 $G$ 中的邻域 $U$ 和 $\delta > 0$ 使得: $g \in U$ 且 $\|v - v_0\| < \delta \Rightarrow \|\rho(g)(v) - \rho(g_0)(v_0)\| < \varepsilon$.
选取 $\delta$ 使得 $0 < \delta \leq \frac{\varepsilon}{2\|\rho(g_0)\|}$, 设 $U$ 是 $g_0$ 在 $G$ 中的邻域使得: $g \in U \Rightarrow \|\rho(g) - \rho(g_0)\| < \frac{\varepsilon}{2(\|v_0\| + \delta)}$. 则若 $g \in U$ 且 $\|v - v_0\| < \delta$, 有 $\|v\| \leq \|v_0\| + \delta$, 从而
\begin{align*}
\|\rho(g)(v) - \rho(g_0)(v_0)\| &\leq \|\rho(g)(v) - \rho(g_0)(v)\| + \|\rho(g_0)(v) - \rho(g_0)(v_0)\| \\
&\leq \|\rho(g) - \rho(g_0)\|\|v\| + \|\rho(g_0)\|\|v - v_0\| \\
&< \frac{\varepsilon}{2(\|v_0\| + \delta)}(\|v_0\| + \delta) + \|\rho(g_0)\|\delta \\
&\leq \varepsilon/2 + \varepsilon/2 = \varepsilon.
\end{align*}
\end{proof}

% ============ I.4 \index{卷积}卷积与\index{群代数}群代数 (原书 pp.30-44) ============
\section{卷积与群代数}

设 $G$ 是局部紧群, $dx$ 是 $G$ 上的左 Haar 测度. 我们记 $L^p(G)$ 为此测度对应的 $L^p$ 空间. 又记 $\Delta$ 为 $G$ 的模函数.

\subsection{\texorpdfstring{$L^1(G)$ 上的卷积与 $G$ 的群代数}{L^1(G) 上的卷积与 G 的群代数}}

\begin{definition}[卷积 (convolution)]
\label{def:I.4.1.1}
设 $f$ 和 $g$ 是从 $G$ 到 $\CC$ 的函数. $f$ 与 $g$ 的卷积 (convolution), 记为 $f * g$, 是函数 $x \longmapsto \int_G f(y)g(y^{-1}x)\,\dd y$ (如果它有定义).
\end{definition}

\begin{proposition}[卷积的 $L^1$ 范数估计]
\label{prop:I.4.1.2}
设 $f, g \in L^1(G)$. 则积分 $\int_G f(y)g(y^{-1}x)\,\dd y$ 对 $G$ 中几乎每个 $x$ 绝对收敛, 所以 $f * g$ 几乎处处有定义, 且 $f * g \in L^1(G)$, 并有
\[
\|f * g\|_1 \leq \|f\|_1\|g\|_1.
\]
\end{proposition}

\begin{proof}
由 Fubini-Tonelli 定理和 $G$ 上测度的左不变性, 函数 $G \times G \to \CC$, $(x, y) \longmapsto f(y)g(y^{-1}x)$ 可积, 且
\[
\int_{G \times G}|f(y)g(y^{-1}x)|\,\dd x\,\dd y = \int_{G \times G}|f(y)||g(x)|\,\dd x\,\dd y = \|f\|_1\|g\|_1.
\]
所以第一个结论也由 Fubini 定理得出, 而第二个结论显然.
\end{proof}

注意卷积对两个变量都显然是线性的.

\begin{proposition}[卷积的性质]
\label{prop:I.4.1.3}
设 $f, g \in L^1(G)$.
\begin{enumerate}
\item[(1)] 对 $G$ 中几乎每个 $x$,
\begin{align*}
f * g(x) &= \int_G f(y)g(y^{-1}x)\,\dd y \\
&= \int_G f(xy)g(y^{-1})\,\dd y \\
&= \int_G f(y^{-1})g(yx)\Delta(y^{-1})\,\dd y \\
&= \int_G f(xy^{-1})g(y)\Delta(y^{-1})\,\dd y \\
&= \int_G f(y)L_yg(x)\,\dd y \\
&= \int_G g(y^{-1})R_yf(x)\,\dd y.
\end{align*}
\item[(2)] 对每个 $h \in L^1(G)$, 有
\[
(f * g) * h = f * (g * h).
\]
(换句话说, 卷积乘积是结合的.)
\item[(3)] 对每个 $x \in G$, 有
\[
L_x(f * g) = (L_xf) * g
\]
且
\[
R_x(f * g) = f * (R_xg).
\]
\item[(4)] 若 $G$ 交换, 则 $f * g = g * f$.
\end{enumerate}
\end{proposition}

\begin{proof}
(1). 前四行的等式由代换 $y \longmapsto xy$ 和 $y \longmapsto y^{-1}$, $dy$ 的左不变性及命题 \ref{prop:I.2.12} 得到. 最后两行只是前两行的改写.

(2). 对 $G$ 中几乎每个 $x$,
\begin{align*}
((f * g) * h)(x) &= \int_G (f * g)(y)h(y^{-1}x)\,\dd y \\
&= \int_{G \times G} f(z)g(z^{-1}y)h(y^{-1}x)\,\dd z\,\dd y \\
&= \int_G f(z)\Bigl(\int_G g(z^{-1}y)h(y^{-1}x)\,\dd y\Bigr)\,\dd z \\
&= \int_G f(z)\Bigl(\int_G g(y)h(y^{-1}z^{-1}x)\,\dd y\Bigr)\,\dd z \\
&= \int_G f(z)(g * h)(z^{-1}x)\,\dd z \\
&= (f * (g * h))(x).
\end{align*}

(3). 这由定义和 (1) 中前两行的等式立即得出.

(4). 这由 (1) 和 $\Delta = 1$ 的事实得出.
\end{proof}

\begin{definition}[\index{Banach 代数}Banach 代数 (Banach algebra)]
\label{def:I.4.1.4}
($\CC$ 上的) Banach 代数是结合的 $\CC$-代数 $A$, 带有范数 $\|\cdot\|$ 使 $A$ 成为巴拿赫空间 (即完备赋范向量空间), 且对每个 $x, y \in A$ 有 $\|xy\| \leq \|x\|\|y\|$ (即范数是次乘性的). 若 $A$ 有单位元 $e$, 我们还要求 $\|e\| = 1$.
\end{definition}

注意我们不假定 $A$ 有单位元. 若它有, 则称 $A$ 是\index{含幺}含幺 (unital) 的.

\begin{example}[Banach 代数的例子]
\label{exa:I.4.1.5}
(a) 若 $V$ 是巴拿赫空间, 则 $\End(V)$ 是含幺 Banach 代数.
(b) 由命题 \ref{prop:I.4.1.2} 和 \ref{prop:I.4.1.3}, 空间 $L^1(G)$ 带卷积乘积是 Banach 代数. 我们称它为 $G$ 的 ($L^1$) 群代数 (group algebra).
\end{example}

\begin{remark}[\index{测度代数}测度代数 $\mathcal{M}(G)$]
\label{rem:I.4.1.6}
若群 $G$ 离散且 $dx$ 是计数测度, 则 $\delta_1 := \mathbb{1}_{\{1\}}$ 是卷积乘积的单位元. 一般地, $L^1(G)$ 不一定有单位元. (当且仅当 $G$ 离散时它有.)\footnote{参考文献 ?} 实际上我们可以把它看作一个确实有单位元的更大的 Banach 代数的子代数, 即 $G$ 的测度代数 (measure algebra) $\mathcal{M}(G)$ (例如见 {[}11{]} 的第 2.5 节):
回忆 $G$ 上的 (复) Radon 测度是 $\mathcal{C}_0(G)$ (带范数 $\|\cdot\|_{\infty}$) 上的有界线性泛函. 我们记 $\mathcal{M}(G)$ 为 Radon 测度的空间, 记 $\|\cdot\|$ 为其范数 (即算子范数); 这是巴拿赫空间. 若 $\mu$ 是 Radon 测度, 我们记 $f \longmapsto \int_G f(x)\,\dd\mu(x)$ 为 $\mathcal{C}_0(G)$ 上对应的线性泛函. 我们定义两个 Radon 测度 $\mu$ 与 $\nu$ 的卷积乘积 $\mu * \nu$ 为线性泛函
\[
f \longmapsto \int_{G \times G} f(xy)\,\dd\mu(x)\,\dd\nu(y).
\]
则不难验证 $\|\mu * \nu\| \leq \|\mu\|\|\nu\|$ 且 $\mathcal{M}(G)$ 上的卷积乘积结合. 这使得 $\mathcal{M}(G)$ 成为 Banach 代数, 且 $1$ 处的 Dirac 测度是 $\mathcal{M}(G)$ 的单位元.
又注意 $\mathcal{M}(G)$ 交换当且仅当 $G$ 交换. 事实上, 由 $*$ 的定义显然 $\mathcal{M}(G)$ 在 $G$ 交换时交换. 为证其逆, 我们记 $\delta_x$ 为 $x$ 处的 Dirac 测度 (于是 $\int_G f\,\dd\delta_x = f(x)$). 则显然对每个 $x, y \in G$ 有 $\delta_x * \delta_y = \delta_{xy}$. 所以若 $\mathcal{M}(G)$ 交换, 则对每个 $x, y \in G$ 有 $\delta_{xy} = \delta_{yx}$, 这蕴含 $G$ 交换.
尽管 $L^1(G)$ 不含 $\mathcal{M}(G)$ 的单位元, 我们有一族称为``\index{近似单位元}近似单位元 (approximate identities)''的函数, 它们在实际上几乎和 $\delta_1$ 一样好. 特别地, 我们将能证明 $L^1(G)$ 交换当且仅当 $G$ 交换 (见推论 \ref{cor:I.4.1.10}).
\end{remark}

\begin{definition}[近似单位元 (approximate identity)]
\label{def:I.4.1.7}
带支撑在 $1$ 在 $G$ 中的邻域基 $\mathcal{U}$ 中的 (对称、连续的) 近似单位元是 $\mathcal{C}_c^+(G)$ 中的函数族 $(\psi_U)_{U \in \mathcal{U}}$ 使得对每个 $U \in \mathcal{U}$:
\begin{itemize}
\item $\operatorname{supp}(\psi_U) \subset U$;
\item $\psi_U(x^{-1}) = \psi_U(x)$, $\forall x \in G$;
\item $\int_G \psi_U(x)\,\dd x = 1$.
\end{itemize}
\end{definition}

对某些结果, 我们不需要 $\psi_U$ 的连续性或 $\psi_U(x^{-1}) = \psi_U(x)$ 这一事实.

\begin{proposition}[近似单位元的存在性]
\label{prop:I.4.1.8}
对 $1$ 在 $G$ 中的每个邻域基 $\mathcal{U}$, 存在支撑在 $\mathcal{U}$ 中的近似单位元.
\end{proposition}

\begin{proof}
设 $U \in \mathcal{U}$. 则 $U$ 包含 $1$ 的对称邻域 $V \subset U$ 和 $1$ 的紧邻域 $K \subset V$, 且由推论 \ref{cor:A.3.1}, 存在支撑包含于 $V$ 的连续函数 $f \colon X \to [0, 1]$ 使得 $f|_K = 1$. 特别地, $f \neq 0$, 所以 $f \in \mathcal{C}_c^+(X)$. 定义 $g \colon X \to [0, 2]$ 为 $g(x) = f(x) + f(x^{-1})$. 则 $g \in \mathcal{C}_c^+(X)$ (因为 $g|_K = 2$) 且 $\operatorname{supp}(g) \subset V \subset U$. 现在取 $\psi_U = \frac{g}{\int_G g(x)\,\dd x}$.
\end{proof}

\begin{proposition}[近似单位元的逼近性质]
\label{prop:I.4.1.9}
设 $\mathcal{U}$ 是 $1$ 在 $G$ 中的邻域基, $(\psi_U)_{U \in \mathcal{U}}$ 是支撑在 $\mathcal{U}$ 中的近似单位元.
\begin{enumerate}
\item[(1)] 对每个 $f \in L^1(G)$, 当 $U \to \{1\}$ 时 $\|\psi_U * f - f\|_1 \to 0$ 且 $\|f * \psi_U - f\|_1 \to 0$. 事实上, 我们有:
\[
\|\psi_U * f - f\|_1 \leq \sup_{y \in U}\|L_yf - f\|_1
\]
且
\[
\|f * \psi_U - f\|_1 \leq \sup_{y \in U}\|R_yf - f\|_1.
\]
\item[(2)] 若 $f \in L^{\infty}(G)$ 且 $f$ 左 (相应地右) 一致连续, 则当 $U \to \{1\}$ 时 $\|\psi_U * f - f\|_{\infty} \to 0$ (相应地 $\|f * \psi_U - f\|_{\infty} \to 0$). 事实上, 我们有:
\[
\|\psi_U * f - f\|_{\infty} \leq \sup_{y \in U}\|L_yf - f\|_{\infty}
\]
且
\[
\|f * \psi_U - f\|_{\infty} \leq \sup_{y \in U}\|R_yf - f\|_{\infty}.
\]
\end{enumerate}
\end{proposition}

在 (2) 中, 注意若 $f \colon G \to \CC$ 有界且 $g \in \mathcal{C}_c(G)$, 则对每个 $x \in G$, 定义 $(f * g)(x)$ 的积分绝对收敛.

\begin{proof}
(1). 设 $U \in \mathcal{U}$. 对每个 $x \in G$,
\[
(\psi_U * f)(x) - f(x) = \int_G \psi_U(y)(L_yf(x) - f(x))\,\dd y
\]
(因为 $\int_G \psi_U(y)\,\dd y = 1$). 所以
\begin{align*}
\|\psi_U * f - f\|_1 &= \int_G \Bigl|\int_G \psi_U(y)(L_yf(x) - f(x))\,\dd y\Bigr|\,\dd x \\
&\leq \int_{G \times G} \psi_U(y)|L_yf(x) - f(x)|\,\dd y\,\dd x \\
&\leq \int_G \psi_U(y)\|L_yf - f\|_1\,\dd y \\
&\leq \sup_{y \in U}\|L_yf - f\|_1.
\end{align*}
第一个收敛结果由当 $y \to 1$ 时 $\|L_yf - f\|_1 \to 0$ 得出, 即命题 \ref{prop:I.3.1.13}.
第二个收敛结果的证明类似 (我们得到 $\|f * \psi_U - f\|_1 \leq \sup_{y \in U}\|R_yf - f\|_1$ 并应用命题 \ref{prop:I.3.1.13}).

(2). 设 $U \in \mathcal{U}$. 则对每个 $x \in G$,
\[
|(\psi_U * f)(x) - f(x)| \leq \int_G \psi_U(y)|L_yf(x) - f(x)|\,\dd y.
\]
由于对 $y \notin U$ 有 $\psi_U(y) = 0$, 这蕴含
\[
|(\psi_U * f)(x) - f(x)| \leq \bigl(\sup_{y \in U}|L_yf(x) - f(x)|\bigr)\Bigl(\int_G \psi_U(y)\,\dd y\Bigr) = \sup_{y \in U}|L_yf(x) - f(x)|.
\]
对 $x \in G$ 取上确界给出
\[
\|\psi_U * f - f\|_{\infty} \leq \sup_{y \in U}\|L_yf - f\|_{\infty}.
\]
所以第一个结论由左一致连续的定义立即得出. 第二个结论的证明类似.
\end{proof}

\begin{corollary}[$L^1(G)$ 的交换性与\index{理想}理想]
\label{cor:I.4.1.10}
\begin{enumerate}
\item[(1)] Banach 代数 $L^1(G)$ 交换当且仅当群 $G$ 交换.
\item[(2)] 设 $I$ 是 $L^1(G)$ 的闭线性子空间. 则 $I$ 是左 (相应地右) 理想当且仅当它在算子 $L_x$ (相应地 $R_x$) ($x \in G$) 下稳定.
\end{enumerate}
\end{corollary}

\begin{proof}
(1). 若 $G$ 交换, 则我们已见 $L^1(G)$ 交换. 反之, 假设 $L^1(G)$ 交换. 设 $x, y \in G$. 设 $f \in \mathcal{C}_c(G)$, 并选取近似单位元 $(\psi_U)_{U \in \mathcal{U}}$. 由命题 \ref{prop:I.4.1.3}, 对每个 $U \in \mathcal{U}$,
\[
(R_xf) * (R_y\psi_U) = R_y((R_xf) * \psi_U) = R_y(\psi_U * (R_xf)) = R_yR_x(\psi_U * f) = R_{yx}(f * \psi_U)
\]
且
\[
(R_xf) * (R_y\psi_U) = (R_y\psi_U) * (R_xf) = R_x(f * (R_y\psi_U)) = R_xR_y(f * \psi_U) = R_{xy}(f * \psi_U).
\]
在 $1$ 处取值给出 $(f * \psi_U)(xy) = (f * \psi_U)(yx)$. 但命题 \ref{prop:I.4.1.9} (以及命题 \ref{prop:I.1.12}) 蕴含当 $U \to \{1\}$ 时 $\|f * \psi_U - f\|_{\infty} \to 0$, 所以我们得到
\[
f(xy) = \lim_{U \to \{1\}}(f * \psi_U)(xy) = \lim_{U \to \{1\}}(f * \psi_U)(yx) = f(yx).
\]
由于这对每个 $f \in \mathcal{C}_c(G)$ 成立, 必有 $xy = yx$ (这由局部紧性和 Urysohn 引理得出).\footnote{参考文献 ?}

(2). 我们证明左理想的情形 (右理想的证明类似). 假设 $I$ 是左理想, 设 $x \in G$. 选取近似单位元 $(\psi_U)_{U \in \mathcal{U}}$. 我们知道当 $U \to \{1\}$ 时 $\psi_U * f \to f$ ($L^1(G)$ 中), 所以当 $U \to \{1\}$ 时 $L_x(\psi_U * f) \to L_xf$ (因为 $L_x$ 保持 $L^1$ 范数). 但由命题 \ref{prop:I.4.1.3}, $L_x(\psi_U * f) = (L_x\psi_U) * f$; 由于 $I$ 是左理想, 对每个 $U \in \mathcal{U}$ 有 $(L_x\psi_U) * f \in I$, 又因 $I$ 闭, 这最终蕴含 $L_xf \in I$.
反之, 假设 $I$ 被所有算子 $L_x$ ($x \in G$) 稳定. 设 $f \in L^1(G)$, $g \in I$. 由命题 \ref{prop:I.4.1.3}, $f * g = \int_G f(y)L_yg\,\dd y$. 由积分的定义, 函数 $f * g$ 在 $L_yg$ ($y \in G$) 的线性张成的空间的闭包中, 所以由假设 (以及 $I$ 闭) 它在 $I$ 中.
\end{proof}

\subsection{\texorpdfstring{$G$ 的表示与 $L^1(G)$ 的表示}{G 的表示与 L^1(G) 的表示}}

\begin{definition}[\index{Banach $*$-代数}Banach $*$-代数与\index{对合}对合]
\label{def:I.4.2.1}
Banach $*$-代数 (Banach $\ast$-algebra) 是带对合反自同构 $*$ 的 Banach 代数 $A$. (即对每个 $x, y \in A$ 和 $\lambda \in \CC$, 有 $(x + y)^* = x^* + y^*$, $(\lambda x)^* = \bar{\lambda}x^*$, $(xy)^* = y^*x^*$ 且 $(x^*)^* = x$.)
反自同构 $*$ 称为 Banach 代数 $A$ 上的对合 (involution).
\end{definition}

\begin{example}[Banach $*$-代数的例子]
\label{exa:I.4.2.2}
(a) $\CC$, 带对合 $z^* = \bar{z}$.
(b) 若 $G$ 是带左 Haar 测度的局部紧群, 则 $L^1(G)$ 带卷积乘积和由 $f^*(x) = \Delta(x)^{-1}\overline{f(x^{-1})}$ 定义的对合 $*$ 是 Banach $*$-代数 (注意 $f^*$ 属于 $L^1(G)$, 且由命题 \ref{prop:I.2.12} 有 $\int_G f^*(x)\,\dd x = \int_G \overline{f(x)}\,\dd x$ 和 $\int_G |f^*(x)|\,\dd x = \int_G |f(x)|\,\dd x$; 所以 $\|f^*\|_1 = \|f\|_1$). 它交换当且仅当 $G$ 交换, 它有单位元当且仅当 $G$ 离散.
(c) 若 $X$ 是局部紧 Hausdorff 空间, 空间 $\mathcal{C}_0(X)$ 带范数 $\|\cdot\|_{\infty}$, 通常的 (逐点) 乘法和由 $f^*(x) = \overline{f(x)}$ 定义的对合 $*$ 是交换 Banach $*$-代数. 它有单位元当且仅当 $X$ 紧 (单位元是常值函数 $1$).
(d) 设 $H$ 是希尔伯特空间. 则 $\End(H)$ 带算子范数和对合 $T \longmapsto T^*$ (其中 $T^*$ 是如上 $T$ 的伴随) 是含幺 Banach $*$-代数. 它交换当且仅当 $\dim_{\CC}(H) = 1$.
\end{example}

\begin{definition}[\index{$*$-同态}$*$-同态与表示]
\label{def:I.4.2.3}
(i) 若 $A$ 和 $B$ 是两个 Banach $*$-代数, 从 $A$ 到 $B$ 的 $*$-同态 ($*$-homomorphism) 是 $\CC$-代数同态 $u \colon A \to B$, 它作为线性算子有界, 且对每个 $x \in A$ 满足 $u(x^*) = u(x)^*$.
(ii) Banach $*$-代数 $A$ 在希尔伯特空间 $H$ 上的表示是从 $A$ 到 $\End(H)$ 的 $*$-同态 $\pi$. 若对每个 $v \in H - \{0\}$, 存在 $x \in A$ 使得 $\pi(x)(v) \neq 0$, 则称该表示非退化 (nondegenerate).
\end{definition}

我们将需要下面的结果, 它将在下一节中证明. (见推论 \ref{cor:II.3.9}.)

\begin{proposition}[正规算子的\index{谱半径}谱半径公式]
\label{prop:I.4.2.4}
设 $V$ 是希尔伯特空间. 则对每个满足 $TT^* = T^*T$ 的 $T \in \End(H)$, 有
\[
\|T\|_{\mathrm{op}} = \lim_{n \to \infty}\|T^n\|_{\mathrm{op}}^{1/n}.
\]
\end{proposition}

\begin{corollary}[表示自动收缩]
\label{cor:I.4.2.5}
设 $A$ 是对每个 $x \in A$ 满足 $\|x^*\| = \|x\|$ 的 Banach $*$-代数, $\pi$ 是 $A$ 在希尔伯特空间 $V$ 上的表示. 则 $\|\pi\|_{\mathrm{op}} \leq 1$.
\end{corollary}

\begin{proof}
由定义, 算子 $\pi$ 有界; 设 $C = \|\pi\|_{\mathrm{op}}$. 设 $x \in A$, $T = \pi(x^*x) \in \End(H)$. 注意 $T = T^*$. 对每个 $n \geq 1$,
\[
\|T^n\| \leq C\|(x^*x)^n\|_1 \leq C\|x\|^{2n}
\]
(因为 $\|x^*\| = \|x\|$). 另一方面, 由命题 \ref{prop:I.4.2.4},
\[
\|T\|_{\mathrm{op}} = \lim_{n \to +\infty}\|T^n\|_{\mathrm{op}}^{1/n},
\]
从而
\[
\|\pi(x)\|_{\mathrm{op}} = \|\pi(x)^*\pi(x)\|_{\mathrm{op}}^{1/2} = \|T\|_{\mathrm{op}}^{1/2} \leq \bigl(\lim_{n \to +\infty}C^{1/n}\|x\|^{2n/n}\bigr)^{1/2} = \|x\|.
\]
换句话说, $\|\pi\|_{\mathrm{op}} \leq 1$.
\end{proof}

我们现在如前固定一个局部紧群 $G$. 我们将使用习题 \ref{ex:I.5.6.1} 中定义的向量值积分, 以及习题 \ref{ex:I.5.6.2} 和 \ref{ex:I.5.6.3} 中证明的性质.

\begin{theorem}[$G$ 的酉表示与 $L^1(G)$ 的表示]
\label{thm:I.4.2.6}
\begin{enumerate}
\item[(1)] 设 $(\pi, V)$ 是 $G$ 的酉表示. 我们按下述方式定义从 $L^1(G)$ 到 $V$ 的线性自同态空间的映射, 仍记为 $\pi$: 若 $f \in L^1(G)$, 我们令
\[
\pi(f) = \int_G f(x)\pi(x)\,\dd x,
\]
其含义是对每个 $v \in V$,
\[
\pi(f)(v) = \int_G f(x)\pi(x)(v)\,\dd x
\]
(积分由习题 \ref{ex:I.5.6.3} 收敛). 则这是 Banach $*$-代数 $L^1(G)$ 在 $V$ 上的\index{非退化表示}非退化表示, 而且对每个 $x \in G$ 和每个 $f \in L^1(G)$,
\[
\pi(L_xf) = \pi(x)\pi(f) \quad \text{且} \quad \pi(R_xf) = \Delta(x)^{-1}\pi(f)\pi(x)^{-1}.
\]
\item[(2)] Banach $*$-代数 $L^1(G)$ 在希尔伯特空间 $V$ 上的每个非退化表示 $\pi$ 都由 (1) 中那样的群 $G$ 的酉表示 $\pi$ 得到.
而且, 若 $(\psi_U)_{U \in \mathcal{U}}$ 是近似单位元, 则对每个 $x \in G$ 和每个 $v \in V$,
\[
\pi(x)(v) = \lim_{U \to \{1\}}\pi(L_x\psi_U)(v).
\]
\item[(3)] 设 $(\pi, V)$ 是 $G$ 的酉表示, $\pi \colon L^1(G) \to \End(V)$ 是伴随的 $*$-同态. 则 $V$ 的闭子空间 $W$ 被 $G$ 不变当且仅当对每个 $f \in L^1(G)$ 有 $\pi(f)(W) \subset W$.
\item[(4)] 设 $(\pi_1, V_1)$ 和 $(\pi_2, V_2)$ 是 $G$ 的酉表示, $\pi_i \colon L^1(G) \to \End(V_i)$, $i = 1, 2$, 是伴随的 $*$-同态. 则有界线性映射 $T \colon V_1 \to V_2$ 是 $G$-等变的当且仅当对每个 $f \in L^1(G)$ 有 $T \circ \pi_1(f) = \pi_2(f) \circ T$.
\end{enumerate}
\end{theorem}

\begin{proof}
(1). 若 $f \in L^1(G)$, 则映射 $\pi(f) \colon V \to V$ 显然是 $\CC$-线性的, 且对每个 $v \in V$:
\[
\|\pi(f)(v)\| = \Bigl\|\int_G f(x)\pi(x)(v)\,\dd x\Bigr\| \leq \int_G |f(x)|\|v\|\,\dd x \leq \|v\|\|f\|_1,
\]
所以 $V$ 的自同态 $\pi(f)$ 有界且 $\|\pi(f)\|_{\mathrm{op}} \leq \|f\|_1$. 又容易看出把 $f$ 映到 $\pi(f)$ 的映射 $\pi \colon L^1(G) \to \End(H)$ 是 $\CC$-线性的, 而不等式 $\|\pi(f)\|_{\mathrm{op}} \leq \|f\|_1$ 蕴含它也有界 (我们也看到 $\|\pi\|_{\mathrm{op}}$ 被 $1$ 控制, 正如推论 \ref{cor:I.4.2.5} 所应有).
设 $f, g \in L^1(G)$. 则对每个 $v \in V$,
\begin{align*}
\pi(f * g)(v) &= \int_{G \times G} f(y)g(y^{-1}x)\pi(x)(v)\,\dd x\,\dd y \\
&= \int_G f(y)\Bigl(\int_G g(y^{-1}x)\pi(x)(v)\,\dd x\Bigr)\,\dd y \\
&= \int_G f(y)\Bigl(\int_G g(x)\pi(yx)(v)\,\dd x\Bigr)\,\dd y \\
&= \int_G f(y)\pi(y)(\pi(g)(v))\,\dd y \\
&= \pi(f)(\pi(g)(v)).
\end{align*}
所以 $\pi(f * g) = \pi(f) \circ \pi(g)$. 又,
\begin{align*}
\pi(f^*)(v) &= \int_G \Delta(x)^{-1}\overline{f(x^{-1})}\pi(x)(v)\,\dd x \\
&= \int_G \overline{f(x)}\pi(x^{-1})(v)\,\dd x \quad \text{(由命题 \ref{prop:I.2.12})} \\
&= \int_G \overline{f(x)}\pi(x)^*(v)\,\dd x,
\end{align*}
所以对 $w \in V$,
\[
\langle \pi(f^*)(v), w \rangle = \int_G \langle \overline{f(x)}\pi(x)^*(v), w \rangle\,\dd x = \int_G \langle v, \overline{f(x)}\pi(x)(w) \rangle\,\dd x = \langle v, \pi(f)(w) \rangle.
\]
这意味着 $\pi(f^*) = \pi(f)^*$. 所以我们已经证明了 $\pi$ 是 $*$-同态.
设 $f \in L^1(G)$, $x \in G$. 则对每个 $v \in V$,
\[
\pi(x)(\pi(f)(v)) = \int_G f(y)\pi(x)(\pi(y)(v))\,\dd y = \int_G f(x^{-1}y)\pi(y)(v)\,\dd y = \pi(L_xf)(v)
\]
且
\[
\pi(f)(\pi(x)^{-1}(v)) = \int_G f(y)\pi(y)(\pi(x^{-1})(v))\,\dd y = \Delta(x)\int_G f(yx)\pi(y)(v)\,\dd y = \Delta(x)\pi(R_xf)(v).
\]

最后, 我们证明表示 $\pi \colon L^1(G) \to \End(V)$ 非退化. 设 $v \in V - \{0\}$, 选取 $1$ 在 $G$ 中的紧邻域 $K$ 使得对每个 $x \in K$ 有 $\|\pi(x)(v) - v\| \leq \frac{1}{2}\|v\|$. 令 $f = \operatorname{vol}(K)^{-1}\mathbb{1}_K$. 则
\[
\|\pi(f)(v) - v\| = \frac{1}{\operatorname{vol}(K)}\Bigl\|\int_K (\pi(x)(v) - v)\,\dd x\Bigr\| \leq \frac{1}{2}\|v\|,
\]
特别地 $\pi(f)(v) \neq 0$.

(2). 设 $\pi$ 是 Banach $*$-代数 $L^1(G)$ 在希尔伯特空间 $V$ 上的非退化表示. 选取 $G$ 的近似单位元 $(\psi_U)_{U \in \mathcal{U}}$. 证明的想法是 $\pi(x)$ 应当是当 $U$ 趋于 $\{1\}$ 时 $\pi(L_x\psi_U)$ 的极限.

我们现在把上述证明想法严格化. 注意由推论 \ref{cor:I.4.2.5}, $\|\pi\|_{\mathrm{op}} \leq 1$. 设 $W$ 是 $\pi(f)(v)$ ($f \in L^1(G)$, $v \in V$) 的线性张成. 我断言 $W$ 在 $V$ 中稠密. 事实上, 设 $v \in W^{\perp}$. 则对每个 $f \in L^1(G)$, 对所有 $v' \in V$ 有 $\langle \pi(f)(v), v' \rangle = \langle v, \pi(f^*)(v') \rangle = 0$, 因此 $\pi(f)(v) = 0$. 由于 $\pi$ 非退化, 只有 $v = 0$ 才可能. 因此 $W^{\perp} = 0$, 这意味着 $W$ 在 $V$ 中稠密.
设 $x \in G$. 我们要定义元素 $\tilde{\pi}(x) \in \End(V)$ 使得对每个 $f \in L^1(G)$ 有 $\tilde{\pi}(x)\pi(f) = \pi(L_xf)$. 这迫使我们把 $W$ 的元素 $w = \sum_{j=1}^n \pi(f_j)(v_j)$ ($f_j \in L^1(G)$, $v_j \in V$) 上的 $\tilde{\pi}(x)$ 定义为
\[
\tilde{\pi}(x)(w) = \sum_{j=1}^n \pi(L_xf_j)(v_j).
\]
这是良定义的, 因为对每个 $n \geq 1$ 和所有 $f_1, \ldots, f_n \in L^1(G)$, $v_1, \ldots, v_n \in V$,
\begin{align*}
\sum_{j=1}^n \pi(L_xf_j)(v_j) &= \lim_{U \to \{1\}}\sum_{j=1}^n \pi(L_x(\psi_U * f_j))(v_j) \\
&= \lim_{U \to \{1\}}\sum_{j=1}^n \pi((L_x\psi_U) * f_j)(v_j) \\
&= \lim_{U \to \{1\}}\pi(L_x\psi_U)\Bigl(\sum_{j=1}^n \pi(f_j)(v_j)\Bigr),
\end{align*}
所以若 $\sum_{j=1}^n \pi(f_j)(v_j) = 0$, 则 $\sum_{j=1}^n \pi(L_xf_j)(v_j) = 0$.

而且, 由于对每个 $U \in \mathcal{U}$ 有 $\|\pi(L_x\psi_U)\|_{\mathrm{op}} \leq \|\pi\|_{\mathrm{op}}\|\psi_U\|_1 \leq 1$, 对每个 $w \in W$ 有 $\|\tilde{\pi}(x)(w)\| \leq \|w\|$, 所以 $\tilde{\pi}(x)$ 是 $W$ 上范数 $\leq 1$ 的有界线性算子, 因此由连续性延拓成范数 $\leq 1$ 的有界线性算子 $\tilde{\pi}(x) \in \End(V)$.
其次, 利用 $L_{xy} = L_x \circ L_y$, 我们看到对所有 $x, y \in G$, 在 $W$ 上有 $\tilde{\pi}(xy) = \tilde{\pi}(x)\tilde{\pi}(y)$, 从而在 $V$ 上成立. 类似地, $L_1 = \mathrm{id}_{L^1(G)}$ 这一事实蕴含 $\tilde{\pi}(1) = \mathrm{id}_V$. 又对每个 $x \in G$, 若 $v \in V$, 则
\[
\|v\| = \|\tilde{\pi}(x^{-1})\tilde{\pi}(x)(v)\| \leq \|\tilde{\pi}(x^{-1})\|_{\mathrm{op}}\|\tilde{\pi}(x)(v)\| \leq \|\tilde{\pi}(x)(v)\| \leq \|v\|,
\]
所以 $\|\tilde{\pi}(x)(v)\| = \|v\|$, 即 $\tilde{\pi}(x)$ 是酉算子.
设 $v \in V$. 我们要证明映射 $G \to V$, $x \longmapsto \tilde{\pi}(x)(v)$ 连续. 由命题 \ref{prop:I.3.1.10}, 这将蕴含 $\tilde{\pi} \colon G \to \End(V)$ 是 $G$ 在 $V$ 上的酉表示. 我们先假设 $v = \pi(f)(v')$, 其中 $f \in L^1(G)$, $v' \in V$. 则 $\tilde{\pi}(x)(v) = \pi(L_xf)(v')$, 所以结论由映射 $G \to L^1(G)$, $x \longmapsto L_xf$ 的连续性 (见命题 \ref{prop:I.3.1.13})、$\pi$ 的连续性及求值映射 $\End(V) \to V$, $T \longmapsto T(v')$ 的连续性得出. 由于连续函数 $G \to V$ 的有限和仍连续, 我们对每个 $v \in W$ 得到结论. 现在处理一般情形. 设 $x \in G$, $\varepsilon > 0$. 我们必须找到 $x$ 在 $G$ 中的邻域 $U$ 使得对每个 $y \in U$ 有 $\|\tilde{\pi}(y)(v) - \tilde{\pi}(x)(v)\| < \varepsilon$. 选取 $w \in W$ 使得 $\|v - w\| < \varepsilon/3$, 以及 $x$ 在 $G$ 中的邻域 $U$ 使得对每个 $y \in U$ 有 $\|\tilde{\pi}(y)(w) - \tilde{\pi}(x)(w)\| < \varepsilon/3$ (由本段的第一部分这是可能的). 则对每个 $y \in U$,
\begin{multline*}
\|\tilde{\pi}(y)(v) - \tilde{\pi}(x)(v)\| \leq \|\tilde{\pi}(y)(v) - \tilde{\pi}(y)(w)\| + \|\tilde{\pi}(y)(w) - \tilde{\pi}(x)(w)\| + \|\tilde{\pi}(x)(w) - \tilde{\pi}(x)(v)\| \\
< \|v - w\| + \varepsilon/3 + \|v - w\| \quad \text{(因为 $\tilde{\pi}(x)$ 和 $\tilde{\pi}(y)$ 酉)} < \varepsilon,
\end{multline*}
如所愿.

我们证明由 $\tilde{\pi}$ 诱导的 $L^1(G)$ 的表示 $\tilde{\pi}$ 就是我们一开始出发的表示 $\pi$. 设 $f, g \in L^1(G)$. 则对每个 $v \in V$,
\begin{align*}
\tilde{\pi}(f)\pi(g)(v) &= \int_G f(x)\tilde{\pi}(x)\pi(g)(v)\,\dd x \\
&= \int_G f(x)\pi(L_xg)(v)\,\dd x \\
&= \int_G \pi(f(x)L_xg)(v)\,\dd x \\
&= \pi\Bigl(\int_G f(x)L_xg\,\dd x\Bigr)(v) \\
&= \pi(f * g)(v) \\
&= \pi(f)\pi(g)(v).
\end{align*}
所以, 若 $f \in L^1(G)$, 则 $\tilde{\pi}(f)$ 和 $\pi(f)$ 在 $W$ 上相等. 由于 $W$ 在 $V$ 中稠密, 这蕴含 $\tilde{\pi}(f) = \pi(f)$.

最后, 我们证明最后一条结论. 设 $(\psi_U)_{U \in \mathcal{U}}$ 是如上近似单位元. 设 $x \in G$. 我们已见对每个 $v \in W$,
\[
\tilde{\pi}(x)(v) = \lim_{U \to \{1\}}\pi(L_x\psi_U)(v).
\]
由于两边都是 $v \in V$ 的连续函数 (对右端, 我们用 $\|\pi(L_x\psi_U)\|_{\mathrm{op}} = 1$ 这一事实), 此恒等式延拓到所有 $v \in V$.

(3). 假设 $W$ 被 $G$ 不变. 设 $f \in L^1(G)$, $w \in W$. 由于 $\pi(f)(w) = \int_G f(x)\pi(x)(w)\,\dd x$ 是形如 $\pi(x)(w)$ ($x \in G$) 的元素的线性组合的极限, 它仍在 $W$ 中.
反之, 假设对每个 $f \in L^1(G)$ 有 $\pi(f)(W) \subset W$. 设 $x \in G$, $(\psi_U)_{U \in \mathcal{U}}$ 是近似单位元. 则由 (2) 的最后一条结论, 对每个 $w \in W$,
\[
\pi(x)(w) = \lim_{U \to \{1\}}\pi(L_x\psi_U)(w) \in W.
\]
所以 $W$ 被 $G$ 不变.

(4). 设 $T \colon V_1 \to V_2$ 是有界线性映射, $W \subset V_1 \times V_2$ 是 $T$ 的图 (graph); 这是 $V_1 \times V_2$ 的闭线性子空间. 则 $T$ 是 $G$-等变的当且仅当 $W$ 被 $G$ 不变, 而 $T$ 是 $L^1(G)$-等变的当且仅当 $W$ 被所有 $\pi_1(f) \times \pi_2(f)$ ($f \in L^1(G)$) 稳定. 所以结论由 (3) 得出.
\end{proof}

\begin{example}[左正则表示与卷积]
\label{exa:I.4.2.7}
设 $\pi$ 是 $G$ 的由 $\pi(x)(f) = L_xf$ 给出的表示 (见例 \ref{exa:I.3.1.11}). 则对每个 $f, g \in L^1(G)$, 有 $\pi(f)(g) = f * g$. 事实上, 由 $\pi(f)$ 的定义有 $\pi(f)(g) = \int_G f(x)L_xg\,\dd x$, 所以结论由习题 \ref{ex:I.5.6.4} 得出.
\end{example}

\subsection{\texorpdfstring{其他 $L^p$ 空间上的卷积}{其他 L^p 空间上的卷积}}

我们只介绍以后证明紧群 Peter-Weyl 定理所需的一些结果. 最重要的情形是 $L^2(G)$.
大多数结果基于 Minkowski 不等式, 它在习题 \ref{ex:I.5.6.7} 中证明. 这里我们只对 $G$ 上的函数陈述它.

\begin{proposition}[Minkowski 不等式 (Minkowski's inequality)]
\label{prop:I.4.3.1}
设 $p \in [1, +\infty)$, $\phi$ 是从 $G \times G$ 到 $\CC$ 的函数. 则
\[
\Bigl(\int_G \Bigl(\int_G \phi(x, y)\,\dd\mu(y)\Bigr)^p \,\dd\mu(x)\Bigr)^{1/p} \leq \int_G \Bigl(\int_G |\phi(x, y)|^p \,\dd\mu(x)\Bigr)^{1/p} \,\dd\mu(y),
\]
其含义是: 若右端有限, 则 $\int_G \phi(x, y)\,\dd\mu(y)$ 对几乎所有 $x \in G$ 绝对收敛, 左端有限且不等式成立.
\end{proposition}

\begin{corollary}[$L^1 * L^p \subset L^p$]
\label{cor:I.4.3.2}
设 $p \in [1, +\infty)$, $f \in L^1(G)$, $g \in L^p(G)$.
\begin{enumerate}
\item[(1)] 定义 $f * g(x)$ 的积分对几乎所有 $x \in G$ 绝对收敛, 且 $f * g \in L^p(G)$, $\|f * g\|_p \leq \|f\|_1\|g\|_p$.
\item[(2)] 若 $G$ 幺模, 则把 $f * g$ 换成 $g * f$ 后同样的结论成立.
\end{enumerate}
\end{corollary}

\begin{proof}
(1). 我们对函数 $\phi(x, y) = f(y)g(y^{-1}x)$ 应用 Minkowski 不等式. 对每个 $y \in G$,
\[
\int_G |\phi(x, y)|^p \,\dd\mu(x) = |f(y)|^p \int_G |g(x)|^p \,\dd\mu(x) = |f(y)|^p\|g\|_p^p
\]
由 $\mu$ 的左不变性, 所以
\[
\int_G \Bigl(\int_G |\phi(x, y)|^p \,\dd\mu(x)\Bigr)^{1/p} \,\dd\mu(y) = \|g\|_p \int_G |f(y)|\,\dd\mu(y) = \|f\|_1\|g\|_p.
\]
Minkowski 不等式首先说明 $\int_G \phi(x, y)\,\dd\mu(y) = f * g(x)$ 对几乎所有 $x \in G$ 绝对收敛, 这是第一个结论. Minkowski 不等式的其余部分正是 $\|f * g\|_p \leq \|f\|_1\|g\|_p$.

(2). 假设 $G$ 幺模. 则
\[
g * f(x) = \int_G g(y)f(y^{-1}x)\,\dd\mu(x) = \int_G g(xy^{-1})f(y)\,\dd\mu(y).
\]
所以证明与 (1) 相同, 对函数 $\phi(x, y) = g(xy^{-1})f(y)$ 应用 Minkowski 不等式.
\end{proof}

现在我们把命题 \ref{prop:I.4.1.9} 推广到其他 $L^p$ 空间.

\begin{corollary}[$L^p$ 中的近似单位元]
\label{cor:I.4.3.3}
设 $\mathcal{U}$ 是 $1$ 在 $G$ 中的邻域基, $(\psi_U)_{U \in \mathcal{U}}$ 是支撑在 $\mathcal{U}$ 中的近似单位元. 则对每个 $1 \leq p < +\infty$, 若 $f \in L^p(G)$, 则当 $U \to \{1\}$ 时 $\|\psi_U * f - f\|_p \to 0$ 且 $\|f * \psi_U - f\|_p \to 0$.
\end{corollary}

\begin{proof}
设 $U \in \mathcal{U}$, $f \in L^p(G)$. 则对每个 $x \in G$,
\[
\psi_U * f(x) - f(x) = \int_G \psi_U(y)(L_yf(x) - f(x))\,\dd\mu(y)
\]
(因为 $\int_G \psi_U\,\dd\mu = 1$). 对函数 $\phi(x, y) = \psi_U(y)(L_yf(x) - f(x))$ 应用 Minkowski 不等式, 得到
\[
\|\psi_U * f - f\|_p \leq \int_G \|L_yf - f\|_p\psi_U(y)\,\dd\mu(y) \leq \sup_{y \in U}\|L_yf - f\|_p.
\]
类似地,
\begin{align*}
f * \psi_U(x) - f(x) &= \int_G f(xy)\psi_U(y^{-1})\,\dd\mu(y) - f(x)\int_G \psi_U(y)\,\dd\mu(y) \\
&= \int_G (R_yf(x) - f(x))\psi_U(y)\,\dd\mu(y).
\end{align*}
所以对函数 $\phi(x, y) = (R_yf(x) - f(x))\psi_U(y)$ 应用 Minkowski 不等式给出
\[
\|f * \psi_U - f\|_p \leq \int_G \|R_yf - f\|_p\psi_U(y)\,\dd\mu(y) \leq \sup_{y \in U}\|R_yf - f\|_p.
\]
因此两个结论都由命题 \ref{prop:I.3.1.13} 得出.
\end{proof}

最后, 我们证明在某些情形下卷积乘积使函数更光滑. 对我们最重要的情形是 $G$ 紧且 $p = q = 2$.

\begin{proposition}[卷积使函数光滑]
\label{prop:I.4.3.4}
假设 $G$ 幺模. 设 $p, q \in (1, +\infty)$ 满足 $p^{-1} + q^{-1} = 1$, $f \in L^p(G)$, $g \in L^q(G)$.
则 $f * g$ 存在, $f * g \in \mathcal{C}_0(G)$ 且 $\|f * g\|_{\infty} \leq \|f\|_p\|g\|_q$.
\end{proposition}

\begin{proof}
设 $x \in G$. 我们有
\[
f * g(x) = \int_G f(y)R_xg(y^{-1})\,\dd\mu(y).
\]
由于 $G$ 幺模, 函数 $y \longmapsto R_xg(y^{-1})$ 仍在 $L^q(G)$ 中且与 $g$ 有相同的 $L^q$ 范数. 所以由 H\"older 不等式,\footnote{当 $p = q = 2$ 时它化为 Cauchy-Schwarz 不等式.} 上面的积分绝对收敛, 且 $|f * g(x)| \leq \|f\|_p\|g\|_q$. 这证明了 $f * g$ 的存在性和关于其范数的结论. 它也说明双线性映射 $L^p(G) \times L^q(G) \to L^{\infty}(G)$, $(f, g) \longmapsto f * g$ 连续. 由于 $\mathcal{C}_0(G)$ 在 $L^{\infty}(G)$ 中闭且 $\mathcal{C}_c(G)$ 在 $L^p(G)$ 和 $L^q(G)$ 中都稠密, 只需证明 $f, g \in \mathcal{C}_c(G)$ 时 $f * g \in \mathcal{C}_0(G)$.
所以设 $f, g \in \mathcal{C}_c(G)$. 设 $x \in G$, $\varepsilon > 0$, 选取 $x$ 的邻域 $U$ 使得对每个 $y \in G$ 和 $x' \in U$, 有 $|g(yx) - g(yx')| \leq \varepsilon$. 则若 $x' \in U$,
\[
|f * g(x) - f * g(x')| \leq \int_G |f(y)||g(y^{-1}x) - g(y^{-1}x')|\,\dd\mu(y) \leq \varepsilon \int_G |f(y)|\,\dd\mu(y).
\]
这证明 $f * g$ 连续. 设 $K = (\operatorname{supp} g)(\operatorname{supp} f)$; 这是 $G$ 的紧子集. 我们要证明 $\operatorname{supp}(f * g) \subset K$, 这将完成证明. 设 $x \in G$, 假设 $f * g(x) \neq 0$. 则存在 $y \in G$ 使得 $f(y)g(y^{-1}x) \neq 0$. 必有 $y \in \operatorname{supp} f$ 且 $y^{-1}x \in \operatorname{supp} g$, 所以 $x \in y(\operatorname{supp} g) \subset K$.
\end{proof}

% ============ I.5 习题 (原书 pp.44-118) ============
\section{习题}

\subsection{拓扑群的例子}

\begin{exercise}[$GL(V)$ 的拓扑]
\label{ex:I.5.1.1}
设 $V$ 是 $\CC$ 上的巴拿赫空间. (即 $V$ 是赋范 $\CC$-向量空间, 对由范数给出的度量完备.) 我们记 $\End(V)$ 为从 $V$ 到自身的有界线性算子空间, 赋予算子范数. 回忆若 $\|\cdot\|$ 是 $V$ 上的范数, 则算子范数 $\|\cdot\|_{\mathrm{op}}$ 定义为: 对每个 $f \in \End(V)$,
\[
\|f\|_{\mathrm{op}} = \inf\{c \in \R_{\geq 0} \mid \forall v \in V,\, \|f(v)\| \leq c\|v\|\} = \sup_{v \in V, \|v\| = 1}\|f(v)\|.
\]
设 $\GL(V)$ 为 $\End(V)$ 中可逆元素的群, 赋予 $\End(V)$ 的拓扑所诱导的拓扑.
\begin{enumerate}
\item[(a)] 证明 $\GL(V)$ 是 $\End(V)$ 的开子集.
\item[(b)] 证明 $\GL(V)$ 是拓扑群.
\item[(c)] 证明 $\GL(V)$ 局部紧当且仅当 $V$ 有限维.
\end{enumerate}
\end{exercise}

\begin{solution}
(a). 注意由算子范数的定义, 对所有 $x, y \in \End(V)$ 有 $\|xy\|_{\mathrm{op}} \leq \|x\|_{\mathrm{op}}\|y\|_{\mathrm{op}}$. (此性质称为``次乘性 (submultiplicativity)''.) 所以若 $x \in \End(V)$ 满足 $\|x\|_{\mathrm{op}} < 1$, 则级数 $\sum_{n \geq 0}x^n$ 收敛 (按约定取 $x^0 = \mathrm{id}_V$), 且 $(\mathrm{id}_V - x)\bigl(\sum_{n \geq 0}x^n\bigr) = \bigl(\sum_{n \geq 0}x^n\bigr)(\mathrm{id}_V - x) = \mathrm{id}_V$. 因此若 $\|x\|_{\mathrm{op}} < 1$, 则 $\mathrm{id}_V - x \in \GL(V)$.
现在设 $x \in \GL(V)$. 我们要证明 $\GL(V)$ 包含 $x$ 在 $\End(V)$ 中的邻域. 设 $y \in \End(V)$ 满足 $\|y\|_{\mathrm{op}} < \|x^{-1}\|_{\mathrm{op}}$. 则 $\|x^{-1}y\|_{\mathrm{op}} < 1$, 所以 $\mathrm{id}_V - x^{-1}y$ 可逆, 从而 $x - y = x(\mathrm{id}_V - x^{-1}y)$ 可逆. 所以 $\End(V)$ 中满足 $\|x - x'\|_{\mathrm{op}} < \|x^{-1}\|_{\mathrm{op}}$ 的每个元素 $x'$ 都在 $\GL(V)$ 中, 这证明了结论.
若 $V$ 有限维, 我们也可以用行列式是连续映射 $\det \colon \End(V) \to \CC$ 且 $\GL(V)$ 是 $\CC$ 的开子集 $\CC^{\times}$ 的原像这一事实.

(b). 我们来证明乘法是从 $\End(V) \times \End(V)$ 到 $\End(V)$ 的连续映射. (这立即蕴含乘法在 $\GL(V)$ 上连续.) 这由算子范数的次乘性立即得出. 事实上, 若 $x, x', y, y' \in \End(V)$, 则
\[
\|xy - x'y'\|_{\mathrm{op}} = \|x(y - y') + (x - x')y'\|_{\mathrm{op}} \leq \|x\|_{\mathrm{op}}\|y - y'\|_{\mathrm{op}} + \|x - x'\|_{\mathrm{op}}\|y'\|_{\mathrm{op}}.
\]
利用
\[
\|y'\|_{\mathrm{op}} = \|y + (y' - y)\|_{\mathrm{op}} \leq \|y\|_{\mathrm{op}} + \|y - y'\|_{\mathrm{op}},
\]
我们看到, 若固定 $x$ 和 $y$, 则当 $(\|x - x'\|_{\mathrm{op}}, \|y - y'\|_{\mathrm{op}})$ 趋于 $(0,0)$ 时 $\|xy - x'y'\|_{\mathrm{op}}$ 趋于 $0$.

我们来证明取逆在 $\GL(V)$ 上连续. 设 $x \in \GL(V)$. 设 $y \in \End(V)$, 记 $h = x - y$, $c = \|h\|_{\mathrm{op}}\|x^{-1}\|_{\mathrm{op}}$. 则 $y = x - h = x(\mathrm{id}_V - x^{-1}h)$. 我们在 (a) 的解答中已见, 若 $c < 1$, 则 $y$ 可逆且 $y^{-1} = \bigl(\sum_{n \geq 0}(x^{-1}h)^n\bigr)x^{-1} = x^{-1} + \sum_{n \geq 1}(x^{-1}h)^nx^{-1}$; 特别地,
\[
\|y^{-1} - x^{-1}\|_{\mathrm{op}} \leq \sum_{n \geq 1}\|(x^{-1}h)^nx^{-1}\|_{\mathrm{op}} = \|x^{-1}\|_{\mathrm{op}}\sum_{n \geq 1}c^n = \frac{c}{1-c}\|x^{-1}\|_{\mathrm{op}}.
\]
这证明: 若 $x$ 固定, 则当 $\|x - y\|_{\mathrm{op}}$ 趋于 $0$ 时 $\|x^{-1} - y^{-1}\|_{\mathrm{op}}$ 趋于 $0$, 这蕴含结论.
若 $V$ 有限维, 还有另一种证明第二点的方法. 事实上, 此时我们可以假设 $V = \CC^n$ (对某个 $n \in \N$), 于是 $\GL(V) = \GL_n(\CC)$. 然后利用: 若 $x \in \GL_n(\CC)$, 则 $x$ 的逆等于 $(\det x)^{-1}y^T$, 其中 $y$ 是 $x$ 的余因子矩阵. 由于 $y$ 的系数是 $x$ 的连续函数 (因为它们是 $x$ 的子矩阵的行列式的 $\pm 1$ 倍), 这证明 $x^{-1}$ 的系数是 $x$ 的连续函数.

(c). 由 I.5.1.1(a) 的解答, 拓扑群局部紧当且仅当其单位元有紧邻域. 由于由 (a) $\GL(V)$ 在 $\End(V)$ 中开, 这蕴含 $\GL(V)$ 局部紧当且仅当 $e$ 在 $\End(V)$ 中有紧邻域. 由于 $\End(V)$ 的拓扑由范数定义, 这等价于 $\End(V)$ 中的闭球紧. 由 Riesz 引理, 这等价于 $\End(V)$ 有限维. 若 $V$ 有限维, 则 $\End(V)$ 也有限维. 若 $V$ 无穷维, 则由 Hahn-Banach 定理可知 $\End(V)$ 也无穷维.
\end{solution}

\begin{exercise}[拓扑群的乘积]
\label{ex:I.5.1.2}
设 $(G_i)_{i \in I}$ 是一族拓扑群.
\begin{enumerate}
\item[(a)] 证明 $\prod_{i \in I} G_i$ (赋予乘积拓扑) 是拓扑群.
\item[(b)] 若所有 $G_i$ 都局部紧, $\prod_{i \in I} G_i$ 是否总局部紧? (给出证明或反例.)
\end{enumerate}
\end{exercise}

\begin{solution}
(a). 我们来证明乘法连续. 设 $(x_i), (y_i) \in \prod_{i \in I} G_i$. 设 $U$ 是 $(x_iy_i)$ 在 $\prod_{i \in I} G_i$ 中的邻域. 由乘积拓扑的定义, 存在 $I$ 的有限子集 $J$ 和 $x_iy_i$ 在 $G_i$ 中的开邻域 $U_i$ ($i \in J$) 使得 $U \supset \bigl(\prod_{i \in J} U_i\bigr) \times \bigl(\prod_{i \in I-J} G_i\bigr)$. 由 $G_i$ ($i \in J$) 上乘法的连续性, 对每个 $i \in J$ 可以找到 $x_i$ 和 $y_i$ 的开邻域 $V_i$ 和 $W_i$ 使得 $V_iW_i \subset U_i$. 设 $V = \bigl(\prod_{i \in J} V_i\bigr) \times \bigl(\prod_{i \in I-J} G_i\bigr)$, $W = \bigl(\prod_{i \in J} W_i\bigr) \times \bigl(\prod_{i \in I-J} G_i\bigr)$. 则 $V$ 和 $W$ 是 $(x_i)$ 和 $(y_i)$ 在 $\prod_{i \in I} G_i$ 中的开邻域, 且有 $VW \subset U$.

我们来证明取逆连续. (证明类似.) 设 $(x_i) \in \prod_{i \in I} G_i$. 设 $U$ 是 $(x_i^{-1})$ 在 $\prod_{i \in I} G_i$ 中的邻域. 由乘积拓扑的定义, 存在 $I$ 的有限子集 $J$ 和 $x_i$ 在 $G_i$ 中的开邻域 $U_i$ ($i \in J$) 使得 $U \supset \bigl(\prod_{i \in J} U_i\bigr) \times \bigl(\prod_{i \in I-J} G_i\bigr)$. 由 $G_i$ ($i \in J$) 上取逆的连续性, 对每个 $i \in J$ 可以找到 $x_i$ 的开邻域 $V_i$ 使得 $V_i^{-1} \subset U_i$. 设 $V = \bigl(\prod_{i \in J} V_i\bigr) \times \bigl(\prod_{i \in I-J} G_i\bigr)$. 则 $V$ 是 $(x_i)$ 在 $\prod_{i \in I} G_i$ 中的开邻域, 且有 $V^{-1} \subset U$.

(b). 只要有无穷多个 $G_i$ 不紧, 答案就是``否''. 事实上, 记 $p_j \colon \prod_{i \in I} G_i \to G_j$ 为投影映射. 这些映射连续, 所以把紧集映到紧集. 现在假设使得 $G_i$ 不紧的 $i \in I$ 的集合无穷. 若 $\prod_{i \in I} G_i$ 局部紧, 则其单位元有紧邻域 $K$. 由乘积拓扑的定义, $K$ 必包含形如 $\bigl(\prod_{i \in J} U_i\bigr) \times \bigl(\prod_{i \in I-J} G_i\bigr)$ 的集合 $U$, 其中 $J$ 是 $I$ 的有限子集, 且对每个 $i \in J$, $U_i$ 是 $e$ 在 $G_i$ 中的邻域. 由假设, 存在 $i \in I - J$ 使得 $G_i$ 不紧. 但 $G_i \supset p_i(K) \supset p_i(U) = G_i$, 所以 $G_i = p_i(K)$ 紧, 这荒谬.
反之, 假设存在 $I$ 的有限子集 $J$ 使得对每个 $i \in I - J$, $G_i$ 紧. 则 $\prod_{i \in I} G_i$ 局部紧. 事实上, 有 $\prod_{i \in I} G_i = \bigl(\prod_{i \in J} G_i\bigr) \times \bigl(\prod_{i \in I-J} G_i\bigr)$, 且由 Tychonoff 定理 $\prod_{i \in I-J} G_i$ 紧, 所以只需证明 $\prod_{i \in J} G_i$ 局部紧. 换句话说, 我们可以假设 $I$ 有限. 但此时, 若 $(x_i) \in \prod_{i \in I} G_i$, 且对每个 $i \in I$, $K_i$ 是 $x_i$ 的紧邻域, 则乘积 $\prod_{i \in I} K_i$ 是 $(x_i)$ 的紧邻域.
\end{solution}

\begin{exercise}[投射极限 (projective limit)]
\label{ex:I.5.1.3}
设 $(I, \leq)$ 是一个有序集. 考虑一族集合 $(X_i)_{i \in I}$ 和一族映射 $(u_{ij} \colon X_i \to X_j)_{i \geq j}$ 满足:
\begin{itemize}
\item 对每个 $i \in I$, $u_{ii} = \mathrm{id}_{X_i}$;
\item 对所有 $i \geq j \geq k$, $u_{ik} = u_{ij} \circ u_{jk}$.
\end{itemize}
这称为由有序集 $I$ 指标的集合的投射系统 (projective system). 此投射系统的投射极限 (projective limit) 是 $\prod_{i \in I} X_i$ 的如下子集:
\[
\varprojlim_{i \in I} X_i = \Bigl\{(x_i)_{i \in I} \in \prod_{i \in I} X_i \;\Big|\; \forall i, j \in I \text{ 且 } i \geq j,\, u_{ij}(x_i) = x_j\Bigr\}.
\]
\begin{enumerate}
\item[(a)] 若所有 $X_i$ 都是 Hausdorff 拓扑空间且所有 $u_{ij}$ 都连续, 证明 $\varprojlim_{i \in I} X_i$ 是 $\prod_{i \in I} X_i$ 的闭子集. 从现在起, 我们总在 $\varprojlim_{i \in I} X_i$ 上赋予诱导拓扑.
\item[(b)] 若所有 $X_i$ 都是紧 Hausdorff 拓扑空间且所有 $u_{ij}$ 都连续, 证明 $\varprojlim_{i \in I} X_i$ 也紧 Hausdorff. (提示: Tychonoff 定理.)
\item[(c)] 若所有 $X_i$ 都是群 (相应地环) 且所有 $u_{ij}$ 都是群 (相应地环) 同态, 证明 $\varprojlim_{i \in I} X_i$ 是 $\prod_{i \in I} X_i$ 的子群 (相应地子环).
\item[(d)] 若所有 $X_i$ 都是拓扑群且所有 $u_{ij}$ 都是连续群同态, 证明 $\varprojlim_{i \in I} X_i$ 是拓扑群.
\item[(e)] 设 $p$ 是一个素数. 取 $I = \N$ (通常序), $X_n = \Z/p^n\Z$, $u_{nm} \colon \Z/p^n\Z \to \Z/p^m\Z$ 是模 $p^m$ 的约化映射. 证明 $\Z_p := \varprojlim_{i \in I} X_i$ 是环, 且对加法是紧拓扑群.
\end{enumerate}
\end{exercise}

\begin{solution}
(a). 记 $X = \prod_{i \in I} X_i$, $X' = \varprojlim_{i \in I} X_i$. 对每个 $i \in I$, 设 $p_i \colon X \to X_i$ 是投影; 这是连续映射. 因此, 若 $i, j \in I$ 满足 $i \geq j$, 则 $X$ 的子集 $\{x \in X \mid u_{ij} \circ p_i(x) = p_j(x)\}$ 闭 (因为它是连续映射 $(u_{ij} \circ p_i, p_j) \colon X \to X_j \times X_j$ 下对角线的原像, 而对角线因 $X_j$ Hausdorff 而闭). 但由投射极限的定义,
\[
X' = \bigcap_{i,j \in I,\, i \geq j}\{x \in X \mid u_{ij} \circ p_i(x) = p_j(x)\}.
\]
所以 $X'$ 也闭.

(b). 若所有 $X_i$ 都是紧 Hausdorff 拓扑空间, 则由 Tychonoff 定理 $X$ 紧 Hausdorff. 由 (a), 投射极限是 $X$ 的闭子空间, 所以它也紧 Hausdorff.

(c). 我们沿用 (a) 的记号. 则所有投影 $p_i$ 都是群 (相应地环) 同态, 所以对所有满足 $i \geq j$ 的 $i, j \in I$, $X$ 的子集 $\{x \in X \mid u_{ij} \circ p_i(x) = p_j(x)\}$ 是子群 (相应地子环). 由投射极限的定义,
\[
X' = \bigcap_{i,j \in I,\, i \geq j}\{x \in X \mid u_{ij} \circ p_i(x) = p_j(x)\}.
\]
所以 $X'$ 也是子群 (相应地子环).

(d). 由习题 \ref{ex:I.5.1.2}(a), 直积是拓扑群. 由 (c), 投射极限 $X'$ 是 $X$ 的子群. 因此 $X'$ 是拓扑群.

(e). 集合 $\Z_p$ 由 (c) 是环, 由 (d) 是拓扑群. 由 (b) 它紧 (注意有限集赋予离散拓扑是紧 Hausdorff 的).
\end{solution}

\begin{exercise}[$p$-进数域 $\Q_p$]
\label{ex:I.5.1.4}
设 $p$ 是一个素数. 我们在 $\Q$ 上如下定义 $p$-进范数 (p-adic norm) $|\cdot|_p$:
\begin{itemize}
\item $|0|_p = 0$;
\item 若 $x$ 是非零有理数, 把 $x$ 写成 $x = p^ny$, 其中 $y$ 的分子和分母都与 $p$ 互素的有理数, 并令 $|x|_p = p^{-n}$.
\end{itemize}
\begin{enumerate}
\item[(a)] 证明对每个 $x, y \in \Q$ 有:
  \begin{itemize}
  \item $|x + y|_p \leq \max(|x|_p, |y|_p)$, 且当 $|x|_p \neq |y|_p$ 时等号成立;
  \item $|xy|_p = |x|_p|y|_p$.
  \end{itemize}
特别地, $p$-进距离函数 $d(x, y) = |x - y|_p$ 是 $\Q$ 上的度量. 我们记 $\Q_p$ 为 $\Q$ 对此度量的完备化.
\item[(a)] 证明 $p$-进范数 $|\cdot|_p$、$\Q$ 的加法和乘法通过连续性延拓到 $\Q_p$, 且 $\Q_p$ 是域 (称为 $p$-进数域 (field of p-adic numbers)), (a) 的结论延拓到 $\Q_p$.
\item[(b)] 证明 $\Q_p$ 的加法群是拓扑群.
\item[(c)] 计算 $\R$ 的子集 $|\Q_p|_p$.
\item[(d)] 证明 $\Q_p$ 中的每个开球也是闭球, 且 $\Q_p$ 中每个正半径的闭球也是开球.
\item[(e)] 证明 $\Q_p$ 完全不连通 (totally disconnected) (即其唯一的非空连通子集是单点集) 但不离散.
\item[(f)] 证明级数 $\sum_{n \geq 0}x_n$ 收敛当且仅当 $\lim_{n \to +\infty}|x_n|_p = 0$.
\item[(g)] 若 $m \in \Z$ 且 $(c_n)_{n \geq m}$ 是一族整数, 证明级数 $\sum_{n \geq m}c_np^n$ 在 $\Q_p$ 中收敛, 且其 $p$-进绝对值 $\leq p^{-m}$, 当 $c_m$ 与 $p$ 互素时等号成立.
\item[(h)] 设 $x \in \Q_p - \{0\}$. 证明存在唯一的 $m \in \Z$ 和唯一的 $\{0, 1, \ldots, p-1\}$ 的元素族 $(c_n)_{n \geq m}$ 使得 $x_m \neq 0$ 且 $x = \sum_{n \geq m}c_np^n$, 且 $|x|_p = p^{-m}$.
\item[(i)] 设 $B = \{x \in \Q_p \mid \|x\|_p \leq 1\}$. 证明这是 $\Q_p$ 的子环, 且是 $\Z$ 在 $\Q_p$ 中的闭包.
\item[(j)] 我们如下定义从 $B$ 到 $\prod_{n \geq 0}\Z/p^n\Z$ 的映射 $u$: 若 $x \in B$, 则由 (e) 可以找到收敛到 $x$ 的 $\Z$ 的元素 Cauchy 序列 $(x_n)_{n \geq 0}$. 取子序列后, 可以假设对每个 $n$ 有 $|x - x_n|_p \leq p^{-n}$. 我们令 $u(x) = (x_n \bmod p^n\Z)_{n \geq 0}$.
  证明 $u$ 是良定义的, 是从 $B$ 到 $\Z_p$ 的同胚, 且也是环同态. 我们将用此来等同 $B$ 与 $\Z_p$.
\item[(k)] 我们把 $M_n(\Q_p)$ 等同于 $\Q_p^{n^2}$, 在其上赋予乘积拓扑, 并在 $\GL_n(\Q_p)$ 上使用诱导拓扑. 证明 $\GL_n(\Q_p)$ 是局部紧拓扑群.
\item[(l)] 证明 $\GL_n(\Z_p)$ 是 $\GL_n(\Q_p)$ 的开紧子群. (提示: 证明 $\Z_p^{\times}$ 在 $\Z_p$ 中闭.)
\end{enumerate}
\end{exercise}

\begin{solution}
(a). 我们首先注意, 若 $x \in \Z - \{0\}$, 则可以把 $x$ 写成 $x = p^mx'$, 其中 $m \geq 0$ 且 $x' \in \Z$ 与 $p$ 互素, 所以 $|x|_p = p^{-m} \leq 1$. 当然, 若 $x = 0$, 也有 $|x|_p \leq 1$.
又注意由 $|\cdot|_p$ 的定义立即得出: 若 $x \in \Q^{\times}$, 则 $|x^{-1}|_p = |x|_p^{-1}$.

设 $x, y \in \Q_p$. 若 $x = 0$, 则 $x + y = y$ 且 $xy = 0$, 所以两点都显然; $y = 0$ 的情形类似. 所以我们假设 $x$ 和 $y$ 都非零, 并把 $x$ 写成 $x = p^nx'$, $y = p^my'$, 其中 $x'$ 和 $y'$ 是分子分母都与 $p$ 互素的有理数. 则 $x'y'$ 的分子分母也都与 $p$ 互素, 且 $xy = p^{n+m}x'y'$, 所以
\[
|xy|_p = p^{-n-m} = p^{-n}p^{-m} = |x|_p|y|_p.
\]
为证明第一个恒等式, 注意由于该恒等式对 $x$ 和 $y$ 对称, 可以假设 $n \leq m$. (注意此时 $p^{-n} = |x|_p = \max(|x|_p, |y|_p)$.) 把 $x'$ 和 $y'$ 写成 $x' = \frac{a}{b}$, $y' = \frac{c}{d}$, 其中 $a, b, c, d \in \Z$ 都与 $p$ 互素. 则
\[
x + y = p^n(x' + p^{m-n}y) = p^n\frac{ad + p^{m-n}cb}{bd},
\]
因此, 由我们已经证明的结论,
\[
|x + y|_p = |p^n|_p|ad + p^{m-n}bc|_p|bd|_p^{-1}.
\]
由于 $bd$ 与 $p$ 互素, 有 $|bd|_p = 1$ (由 $|\cdot|_p$ 的定义). 由于 $ad + p^{m-n}bc \in \Z$, 由开头的注记有 $|ad + p^{m-n}cb|_p \leq 1$. 最后得到
\[
|x + y|_p \leq |p^n|_p = p^{-n} = \max(|x|_p, |y|_p).
\]
最后, 若 $|x|_p \neq |y|_p$, 则 $n < m$, 所以 $ad + p^{m-n}cb$ 与 $p$ 互素, 由 $|\cdot|_p$ 的定义得 $|x + y|_p = p^{-n}$.

(b). 由定义, $\Q_p$ 是 $\Q$ 的元素 Cauchy 序列 $(x_n)_{n \geq 0}$ (对 $p$-进距离给出的度量) 商去等价关系 $\sim$ 的商集, 其中 $\sim$ 定义为: $(x_n)_{n \geq 0} \sim (y_n)_{n \geq 0}$ 当且仅当 $n \to +\infty$ 时 $|x_n - y_n|_p \to 0$.
注意 (a) 的第二个恒等式蕴含三角不等式: 对所有 $x, y \in \Q$, 有 $|x + y|_p \leq |x|_p + |y|_p$.
设 $x \in \Q_p$, $(x_n)_{n \geq 0}$ 是表示 $x$ 的 Cauchy 序列. 由三角不等式, 对所有 $n, m \in \N$ 有 $||x_n|_p - |x_m|_p| \leq |x_n - x_m|_p$. 所以 $(|x_n|_p)_{n \geq 0}$ 是 $\R$ 中的 Cauchy 序列, 且因 $\R$ 完备它有极限. 设 $(y_n)_{n \geq 0}$ 是表示 $x$ 的另一个 Cauchy 序列. 由三角不等式, 对每个 $n \geq 0$ 有 $||x_n|_p - |y_n|_p| \leq |x_n - y_n|_p$, 所以 $(|x_n|_p)_{n \geq 0}$ 和 $(|y_n|_p)_{n \geq 0}$ 的极限相等. 因此我们可以定义 $|x|_p = \lim_{n \to +\infty}|x_n|_p$.
现在设 $x, y \in \Q_p$, 并选取表示 $x$ 和 $y$ 的 Cauchy 序列 $(x_n)_{n \geq 0}$ 和 $(y_n)_{n \geq 0}$. 首先注意序列 $(|x_n|_p)_{n \geq 0}$ 和 $(|y_n|_p)_{n \geq 0}$ 有界 (例如因为它们如上所述收敛). 现在利用 (a), 对所有 $n, m \in \N$:
\[
|(x_n + y_n) - (x_m + y_m)|_p \leq \max(|x_n - x_m|_p, |y_n - y_m|_p)
\]
且
\[
|x_ny_n - x_my_m|_p = |x_n(y_n - y_m) + (x_n - x_m)y_m|_p \leq \max(|x_n|_p|y_n - y_m|_p, |y_m|_p|x_n - x_m|_p).
\]
所以序列 $(x_n + y_n)_{n \geq 0}$ 和 $(x_ny_n)_{n \geq 0}$ 是 Cauchy 序列 (对第二个, 用 $(|x_n|_p)_{n \geq 0}$ 和 $(|y_n|_p)_{n \geq 0}$ 有界的事实), 因此它们表示 $\Q_p$ 的元素. 我们想把这些元素称为 $x + y$ 和 $xy$, 但首先必须验证它们与表示 $x$ 和 $y$ 的 Cauchy 序列的选取无关. 所以设 $(x'_n)_{n \geq 0}$ 和 $(y'_n)_{n \geq 0}$ 是分别表示 $x$ 和 $y$ 的另外两个 Cauchy 序列. 则对每个 $n \geq 0$,
\[
|(x_n + y_n) - (x'_n + y'_n)|_p \leq \max(|x_n - x'_n|_p, |y_n - y'_n|_p)
\]
且
\[
|x_ny_n - x'_ny'_n|_p = |x_n(y_n - y'_n) + (x_n - x'_n)y'_n|_p \leq \max(|x_n|_p|y_n - y'_n|_p, |y'_n|_p|x_n - x'_n|_p).
\]
所以两个序列 $((x_n + y_n) - (x'_n + y'_n))_{n \geq 0}$ 和 $(x_ny_n - x'_ny'_n)_{n \geq 0}$ 都收敛到 $0$, 这意味着序列 $(x_n + y_n)_{n \geq 0}$ 与 $(x'_n + y'_n)_{n \geq 0}$ (相应地 $(x_ny_n)_{n \geq 0}$ 与 $(x'_ny'_n)_{n \geq 0}$) 有相同的极限, 所以 $x + y$ 和 $xy$ 的定义有意义.
$\Q_p$ 的环公理由运算的定义立即得出. 我们来验证 $\Q_p$ 是域. 设 $x \in \Q_p - \{0\}$, 选取表示 $x$ 的 Cauchy 序列 $(x_n)_{n \geq 0}$. 由于 $x \neq 0$, 序列 $(|x_n|_p)_{n \geq 0}$ 不能收敛到 $0$, 所以其极限 (即 $|x|_p$) 非零. 所以对足够大的 $n, m$ 有 $|x_n - x_m|_p \leq |x|_p/2$, 且经过用等价的 Cauchy 序列替换 $(x_n)_{n \geq 0}$ 后, 可以假设它对所有 $n, m \geq 0$ 成立. 设 $n, m \geq 0$. 由 (a), 有 $|x_n|_p \leq \max(|x_m|_p, |x_n - x_m|_p)$. 令 $m \to +\infty$ 取极限, 得 $|x_n|_p \leq |x|_p$. 类似地, 令 $n \to +\infty$ 取极限, 并利用 $|x|_p > |x|_p/2 \geq \lim_{n \to +\infty}|x_n - x_m|_p$, 得 $|x|_p \leq |x_m|_p$. 这蕴含对每个 $n \geq 0$ 有 $|x|_p = |x_n|_p$. 现在, 若我们能证明序列 $(x_n^{-1})_{n \geq 0}$ 是 Cauchy 序列, 则它表示的 $\Q_p$ 的元素显然将是 $x$ 的逆. 但对所有 $n, m \geq 0$,
\[
|x_n^{-1} - x_m^{-1}|_p = |x_n^{-1}x_m^{-1}|_p|x_m - x_n|_p = |x|_p^{-2}|x_m - x_n|_p,
\]
所以 $|x_n^{-1} - x_m^{-1}|_p$ 在 $n, m \to +\infty$ 时确实收敛到 $0$.

我们最后证明 (a) 的恒等式在 $\Q_p$ 中仍成立. 若 $x, y \in \Q_p - \{0\}$, 则刚才已见可以找到收敛到 $x$ 和 $y$ 的 Cauchy 序列 $(x_n)_{n \geq 0}$ 和 $(y_n)_{n \geq 0}$ 使得对每个 $n \geq 0$, $|x|_p = |x_n|_p$ 且 $|y|_p = |y_n|_p$. 当然, 若 $x$ 或 $y$ 为 $0$ 这也成立. 则恒等式由 (a) 和 $\Q_p$ 上运算的定义立即得出.

(c). $\Q_p$ 上的加法由定义连续 (它定义为映射的连续性延拓). 取逆映射 $x \longmapsto -x$ 连续因为 $|x|_p = |-x|_p$.

(d). 我们已见: 若 $x \in \Q_p$, 则存在收敛到 $x$ 的 $\Q$ 的 Cauchy 序列且对每个 $n \geq 0$ 有 $|x|_p = |x_n|_p$. 所以 $|\Q_p|_p = |\Q|_p = \{0\} \cup p^{\Z}$.

(e). 若 $x \in \Q_p$, $r \in \R$, 记 $B(x, r)$ 和 $\overline{B}(x, r)$ 为以 $x$ 为心、$r$ 为半径的开球和闭球.
设 $x \in \Q_p$, $r \in \R$. 若 $r \leq 0$, 则 $B(x, r) = \emptyset = \overline{B}(x, -1)$. 假设 $r > 0$, 设 $n$ 是满足 $p^n < r \leq p^{n+1}$ 的唯一整数. 由前一问, 若 $a \in \Q_p$ 满足 $|a|_p < r$, 则 $|a|_p \leq p^n$, 且显然其逆也成立. 所以 $B(x, r) = \overline{B}(x, p^n)$. 现在设 $m$ 是满足 $p^m \leq r < p^{m+1}$ 的唯一整数. 则类似地有 $B(x, r) = \overline{B}(x, p^{m+1})$.

(f). 设 $x, y \in \Q_p$ 且 $x \neq y$. 则 $|x - y|_p > 0$, 所以可以选取 $r > 0$ 使得 $r < |x - y|_p$. 则 $B(x, r)$ 是 $\Q_p$ 的包含 $x$ 而不包含 $y$ 的开闭子集, 所以 $x$ 和 $y$ 不能在 $\Q_p$ 的同一连通子集中. 这证明 $\Q_p$ 完全不连通.
为证明 $\Q_p$ 不离散, 只需证明其子集 $\{0\}$ 不开. 这由序列 $(p^n)_{n \geq 0}$ 在 $\Q_p$ 中收敛到 $0$ 且对每个 $n \in \Z$ 有 $p^n \neq 0$ 的事实得出.

(g). 定义有理数序列 $(S_n)_{n \geq 0}$ 为 $S_n = \sum_{i=0}^n x_i$. 则级数 $\sum_{n \geq 0}x_n$ 收敛当且仅当序列 $(S_n)_{n \geq 0}$ 收敛. 特别地, 若级数收敛, 则 $|x_n|_p = |S_{n+1} - S_n|_p$ 在 $n \to +\infty$ 时趋于 $0$.
反之, 假设 $\lim_{n \to +\infty}|x_n|_p = 0$. 对所有 $n, n' \in \N$, 若 $n \leq n'$, 则 (利用 (b)):
\[
|S_{n'} - S_n|_p = \Bigl|\sum_{i=n+1}^{n'}x_i\Bigr|_p \leq \max_{n+1 \leq i \leq n'}|x_i|_p \leq \sup_{i \geq n+1}|x_i|_p.
\]
当 $n \to +\infty$ 时这趋于 $0$, 所以 $(S_n)_{n \geq 0}$ 是 Cauchy 序列, 因此它在 $\Q_p$ 中收敛, 级数 $\sum_{n \geq 0}x_n$ 也随之收敛.

(h). 收敛性由前一问以及 (b) 的证明中注意到的 $|c|_p \leq 1$ (对每个 $c \in \Z$) 得出. 设 $x = \sum_{n \geq m}c_np^n$. 由定义,
\[
x = \lim_{n \to +\infty}\sum_{i=m}^n c_ip^i.
\]
对每个 $n \geq m$,
\[
\Bigl|\sum_{i=m}^n c_ip^i\Bigr|_p \leq \max_{m \leq i \leq n}|c_i|_p|p^i|_p = p^{-m},
\]
所以 $|x|_p \leq p^{-m}$. 假设 $c_m$ 与 $p$ 互素; 则 $|c_m|_p = 1$. 因此 $|c_mp^m|_p = p^{-m} > |c_ip^i|_p$ (对所有 $i > m$), 所以再次利用 (b), 对每个 $n \geq m$,
\[
\Bigl|\sum_{i=m}^n c_ip^i\Bigr|_p = p^{-m}.
\]
这给出 $|x|_p = p^{-m}$.

(i). 让我们证明存在性. 可以假设 $x \neq 0$ (否则结论平凡). 我们知道对某个 $m \in \Z$ 有 $|x|_p = p^{-m}$. 选取收敛到 $x$ 的 Cauchy 序列 $(x_n)_{n \geq 0}$. 取子序列后, 可以假设对每个 $n \geq 0$ 有 $|x - x_n|_p < p^{-n}$.
设 $n \geq 0$. 把 $x_n$ 写成 $p$ 进制形式 $x_n = \sum_{i=a_n}^{b_n}c_{i,n}p^i$, 其中 $a_n, b_n \in \Z$, $c_{i,n} \in \{0, 1, \ldots, p-1\}$. 可以假设 $c_{a_n,n} \neq 0$. 则 $c_{a_n,n}$ 与 $p$ 互素, 所以 $|c_{a_n,n}p^{a_n}|_p = p^{-a_n} > |c_{i,n}p^i|_p$ (对所有 $i > a_n$), 于是由 (b),
\[
p^{-m} = |x_n|_p = p^{-a_n},
\]
最终得到 $m = a_n$.
又, 通过在 $i > b_n$ 时令 $c_{i,n} = 0$, 可以把 $b_n$ 换成 $+\infty$.
设 $n, n' \in \N$, 并假设 $n \geq n'$. 则 $|x_n - x_{n'}|_p \leq \max(|x_n - x|_p, |x - x_{n'}|_p) < p^{-n'}$. 另一方面, $x_n - x_{n'} = \sum_{i \geq m}(c_{i,n} - c_{i,n'})p^i$. 注意 $c_{i,n} - c_{i,n'}$ 属于 $\{1 - p, \ldots, p - 1\}$, 所以它们要么是 $0$ 要么与 $p$ 互素. 由 (h), 这蕴含 $|x_n - x_{n'}|_p = p^{-r}$, 其中 $r$ 是使得 $c_{r,n} - c_{r,n'} \neq 0$ 的最小整数. 这又蕴含 $n' < r$, 即对 $m \leq i \leq n'$ 有 $c_{i,n} = c_{i,n'}$.
我们现在对 $i \geq m$ 定义整数 $c_i$: 若 $i \leq 0$ 取 $c_i = c_{0,i}$, 若 $i \geq 0$ 取 $c_i = c_{i,i}$. 由上一段, 若 $0 \leq i \leq n$ 则 $c_i = c_{i,n}$. 对每个 $n \geq 0$, 令 $y_n = \sum_{i=m}^n c_ip^i$. 则
\[
x_n - y_n = \sum_{i \geq m}(c_{i,n} - c_i)p^i = \sum_{i \geq n+1}(c_{i,n} - c_i)p^i,
\]
所以由 (h) 有 $|x_n - y_n|_p \leq p^{-n-1}$. 因此序列 $(y_n)_{n \geq 0}$ 也收敛到 $x$, 这证明 $x = \sum_{i \geq m}c_ip^i$.
让我们证明唯一性. 假设有两个整数序列 $(c_n)_{n \geq m}$, $(d_n)_{n \geq m}$ 使得 $x = \sum_{n \geq m}c_np^n = \sum_{n \geq m}d_np^n$ 且对每个 $n$ 有 $c_n, d_n \in \{0, \ldots, p-1\}$. 则 $0 = \sum_{n \geq m}(c_n - d_n)p^n$. 又, 对每个 $n$, $c_n - d_n$ 属于 $\{1 - p, \ldots, p - 1\}$, 所以它是 $0$ 或与 $p$ 互素. 若存在 $n$ 使得 $c_n - d_n \neq 0$, 则由 (h) 这将蕴含 $|0|_p \neq 0$, 这不可能. 所以对每个 $n$ 有 $c_n = d_n$.

(j). $B$ 是子环这一事实由 (b) (以及 $|-x|_p = |x|_p$, 这在定义上显然) 得出, 且已在 (b) 的证明中看到 $\Z \subset B$. 又, $B$ 是闭球, 所以它在 $\Q_p$ 中闭, 因此包含 $\Z$ 的闭包.
设 $x \in B$. 由 (i), 可以把 $x$ 写成 $x = \sum_{n \geq 0}c_np^n$, 其中 $c_n \in \{0, \ldots, p-1\}$. 这意味着 $x$ 是整数序列 $\bigl(\sum_{i=0}^n c_ip^i\bigr)_{n \geq 0}$ 的极限, 所以 $x$ 在 $\Z$ 于 $\Q_p$ 中的闭包里.

(k). 我们证明 $u$ 良定义. 设 $x \in B$, $(x_n)_{n \geq 0}$, $(x'_n)_{n \geq 0}$ 是如题设的两个序列. 设 $n \geq 0$. 则 $|x_n - x'_n|_p \leq \max(|x_n - x|_p, |x - x'_n|_p) \leq p^{-n}$, 这意味着 $p^n$ 整除 $x_n - x'_n$, 所以 $x_n$ 和 $x'_n$ 在 $\Z/p^n\Z$ 中有相同的像. 这证明 $u(x)$ 良定义.
$u$ 是环同态这一事实由 $\Q_p$ 上环运算的定义以及对每个 $n$ 从 $\Z$ 到 $\Z/p^n\Z$ 的约化是环同态这一事实立即得出.
我们证明 $u$ 单射. 设 $x, y \in B$ 满足 $u(x) = u(y)$, 并选取满足题设条件的收敛到 $x$, $y$ 的序列 $(x_n)_{n \geq 0}$, $(y_n)_{n \geq 0}$. 则对每个 $n \geq 0$ 有 $x_n \equiv y_n \pmod{p^n}$, 所以 $p^n$ 整除 $x_n - y_n$, 即 $|x_n - y_n|_p \leq p^{-n}$. 令 $n \to +\infty$ 取极限, 得 $|x - y|_p = 0$. 但已在 (b) 中看到 $\Q_p$ 中 $p$-进范数为 $0$ 的唯一元素是 $0$, 所以 $x = y$.
我们证明 $u$ 满射. 设 $(x_n + p^n\Z)_{n \geq 0}$ 是 $\Z_p$ 的元素. 对每个 $n \geq 0$, 在 $\{0, \ldots, p^n - 1\}$ 中选取 $x_n + p^n\Z$ 的代表元, 记为 $x_n$. 又把 $x_n$ 写成 $p$ 进制形式 $x_n = \sum_{i=0}^{n-1}c_{i,n}p^i$, 其中 $0 \leq c_{i,n} \leq p - 1$. 设 $m \geq n$. 我们知道 $x_m \equiv x_n \pmod{p^n}$, 所以对 $0 \leq i \leq n$ 有 $c_{i,m} = c_{i,n}$. 我们定义序列 $(c_i)_{i \geq 0}$ 为 $c_i = c_{i,0} = c_{i,1} = \cdots = c_{i,i}$, 并令 $x = \sum_{i \geq 0}c_ip^i \in \Q_p$. 则由 (c), $x \in B$, 且容易验证 $(x_n)_{n \geq 0}$ 是收敛到 $x$ 且满足题设条件的整数 Cauchy 序列. 因此 $u(x) = (x_n + p^n\Z)_{n \geq 0}$.
我们证明 $u$ 连续. $\Z_p$ 中的每个开集都是形如 $\Z_p \cap \bigl((\prod_{n \geq m+1}\Z/p^n\Z) \times \{(x_m, \ldots, x_0)\}\bigr)$ 的开集的并, 其中 $m \geq 0$, 且对 $0 \leq i \leq m$ 有 $x_i \in \Z/p^i\Z$. 所以只需证明这种形式的集合的原像在 $B$ 中开. 记 $A = \bigl(\prod_{n \geq m+1}\Z/p^n\Z\bigr) \times \{(x_m, \ldots, x_0)\}$. 选取 $x \in \Z$ 使得 $x \equiv x_m \pmod{p^m}$. 则对 $0 \leq n \leq m$ 有 $x \equiv x_n \pmod{p^n}$ (因为 $x_n \equiv x_m \pmod{p^n}$). 我们将证明 $y \in B$ 属于 $u^{-1}(A)$ 当且仅当 $|x - y|_p < p^{-m+1}$, 这证明 $u^{-1}(A)$ 开. 首先注意, 由于 $|\cdot|_p$ 的值总是 $0$ 或 $p$ 的整数次幂, 条件 $|x - y|_p < p^{-m+1}$ 等价于 $|x - y|_p \leq p^{-m}$.\footnote{所以在 $\Q_p$ 中, 每个开球都是闭球, 反之亦然.} 设 $y \in B$, $(y_n)_{n \geq 0}$ 是如 $u$ 的定义中那样收敛到 $y$ 的 Cauchy 序列. 假设 $|x - y|_p \leq p^{-m}$. 则对 $n \in \{0, \ldots, m\}$, 有 $|y_n - x|_p \leq \max(|y_n - y|_p, |y - x|_p) \leq p^{-n}$, 所以 $y_n \equiv x \equiv x_n \pmod{p^n}$. 所以 $u(y) \in A$. 反之, 假设 $u(y) \in A$. 则 $y_m \equiv x_m \equiv x \pmod{p^m}$, 所以 $|y_m - x|_p \leq p^{-m}$. 由于 $|y - y_m|_p \leq p^{-m}$, 这蕴含 $|x - y|_p \leq p^{-m}$.
最后, 我们证明 $u$ 是开映射. 由于 $u$ 双射, 只需证明开球的像是开的. 我们几乎已经做了这件事: 设 $x \in B$, $r \in \R_{>0}$, $A'$ 是以 $x$ 为心、$r$ 为半径的开球. 若 $m$ 是满足 $p^{-m} < r$ 的最小整数, 则 $A'$ 也是以 $x$ 为心、$p^{-m}$ 为半径的闭球 (因为 $|\cdot|_p$ 取值于 $\{0\} \cup p^{\Z}$). 设 $y$ 是满足 $|x - y|_p < p^{-m}$ 的整数. 则 (a) 的第二个恒等式蕴含对 $z \in \Q_p$, $|x - z|_p \leq p^{-m}$ 当且仅当 $|y - z|_p \leq p^{-m}$, 这意味着在 $A'$ 的定义中可以用 $y$ 代替 $x$. 则上面已见 $u(A') = \bigl(\prod_{n \geq m+1}\Z/p^n\Z\bigr) \times \{(x_m, \ldots, x_0)\}$, 其中对 $0 \leq n \leq m$ 有 $x_n = y + p^n\Z$.

(l). $\GL_n(\Q_p)$ 是拓扑群的证明与习题 \ref{ex:I.5.1.2}(b) 相同 (有限维情形). 它也在 $M_n(\Q_p)$ 中开, 因为它是连续函数 $\det$ 下 $\Q_p$ 的开子集 $\Q_p^{\times}$ 的原像. 所以要证明 $\GL_n(\Q_p)$ 局部紧, 只需证明 $M_n(\Q_p)$ 局部紧, 而这将由 $\Q_p$ 局部紧得出. 但对每个 $x \in \Q_p$, 以 $x$ 为心、半径为 $1$ 的闭球 $x + \Z_p$ 是 $x$ 的紧邻域: 它紧因为 $\Z_p$ 紧且平移 $x$ 连续 (由度量的定义), 它开因为它等于以 $x$ 为心、$p$ 为半径的开球.

(m). 回忆 $\GL_n(\Z_p)$ 是 $M_n(\Z_p)$ 的可逆元群, 所以
\[
\GL_n(\Z_p) = \{A \in \GL_n(\Q_p) \mid A \text{ 与 } A^{-1} \in M_n(\Z_p)\}.
\]
换句话说, $\GL_n(\Z_p)$ 是连续映射 $\GL_n(\Q_p) \to M_n(\Q_p) \times M_n(\Q_p)$, $A \longmapsto (A, A^{-1})$ 下 $M_n(\Q_p) \times M_n(\Q_p)$ 的开子集 $M_n(\Z_p) \times M_n(\Z_p)$ 的原像, 所以它在 $\GL_n(\Q_p)$ 中开. 由于 $M_n(\Z_p)$ 紧 (因为 $\Z_p$ 紧), 为证明 $\GL_n(\Z_p)$ 紧, 只需证明它在 $M_n(\Z_p)$ 中闭. 由于它是连续映射 $\det \colon M_n(\Z_p) \to \Z_p$ 下 $\Z_p^{\times}$ 的原像, 只需证明 $\Z_p^{\times}$ 在 $\Z_p$ 中闭. 但 $\Z_p^{\times} = \{x \in \Z_p \mid |x|_p = 1\}$ (这蕴含它闭). 事实上, 设 $x \in \Z_p$. 若 $x$ 在 $\Z_p$ 中有逆, 则 $|x^{-1}|_p = |x|_p^{-1} \leq 1$, 所以 $|x|_p \geq 1$, 从而 $|x|_p = 1$. 反之, 若 $|x|_p = 1$, 则 $|x^{-1}|_p = 1$, 所以 $x^{-1} \in \Z_p$.
\end{solution}

\subsection{Van Dantzig 定理}
\label{sec:I.5.2}
本习题来自 Terry Tao 的博客.

\begin{exercise}[紧 Hausdorff 完全不连通空间]
\label{ex:I.5.2.1}
在此题中, $X$ 是紧 Hausdorff 完全不连通拓扑空间. (回忆``完全不连通 (totally disconnected)''意味着 $X$ 的唯一的非空连通子集是单点集.)
设 $x \in X$, 设 $A$ 是 $X$ 的所有包含 $x$ 的开闭子集的交集. 证明 $A = \{x\}$. (提示: 这等价于证明 $A$ 连通. 还要回忆在任何紧 Hausdorff 空间中, 不相交的闭集可被开集分离.)
\end{exercise}

\begin{solution}
注意 $A$ 在 $X$ 中闭. 假设 $A = A_1 \cup A_2$, 其中 $A_1$ 和 $A_2$ 在 $X$ 中闭且不相交, 且 $x \in A_1$. 由于任何紧 Hausdorff 空间正规, 可以找到 $X$ 的开子集 $U_1 \supset A_1$ 和 $U_2 \supset A_2$ 使得 $U_1 \cap U_2 = \emptyset$. 让我们找一个包含 $x$ 的开闭邻域 $V$ 使得 $V \cap \partial U_2 = \emptyset$. 对每个 $y \in \partial U_2$, 由于 $y \notin A$, 可以找到 $x$ 的开闭邻域 $V_y$ 使得 $y \notin V_y$. 注意 $X - V_y$ ($y \in \partial U_2$) 构成覆盖 $\partial U_2$ 的 $X$ 的开子集族; 由于 $\partial U_2$ 紧, 此族有仍然覆盖 $\partial U_2$ 的有限子族, 设为 $(X - V_{y_1}, \ldots, X - V_{y_n})$. 设 $V = V_{y_1} \cap \cdots \cap V_{y_n}$; 则 $V$ 仍开且闭, $x \in V$, 且 $V \cap \partial U_2 = \emptyset$. 最后这条性质蕴含 $B := V - U_2$ 也等于 $V - \overline{U_2}$, 所以它仍开且闭. 又, 有 $x \in B$ (因为 $x \notin U_2$), 且 $A_2 \cap B = \emptyset$. 但由定义 $A$ 必须包含于 $B$, 所以 $A_2 = \emptyset$. 这证明 $A$ 连通, 从而是单点集, 所以等于 $\{x\}$.
\end{solution}

\begin{exercise}[局部紧完全不连通群的紧开子群]
\label{ex:I.5.2.2}
在此题中, $G$ 是局部紧完全不连通拓扑群.
\begin{enumerate}
\item[(a)] 证明 $G$ 的单位元有紧开邻域 $K$.
\item[(b)] 证明存在包含于 $K$ 的开子群 $G_0$.
\item[(c)] 证明 $G$ 的紧开子群构成 $1$ 在 $G$ 中的邻域基.
\item[(d)] 设 $G$ 是习题 \ref{ex:I.5.1.4} 中的群 $\GL_n(\Q_p)$. 找出 $1$ 在 $G$ 中由紧开子群组成的邻域基.
\end{enumerate}
\end{exercise}

\begin{solution}
(a). 设 $V$ 是 $1$ 的紧邻域. 则 $\partial V$ 也紧且不含 $1$. 由习题 \ref{ex:I.5.2.1}, 对每个 $y \in \partial V$, 存在 $1$ 的开闭邻域 $B_y$ 使得 $B_y \cap \partial V = \emptyset$. 由于 $\partial V \subset \bigcup_{y \in \partial V}(X - B_y)$ 且 $X - B_y$ 开, 存在 $y_1, \ldots, y_n \in \partial V$ 使得 $\partial V \cap B = \emptyset$, 其中 $B = \bigcap_{i=1}^n B_{y_i}$. 注意 $B$ 仍开且闭, 且 $1 \in B$. 又, 由于 $\partial V \cap B = \emptyset$, 有 $B \cap V = B \cap \mathring{V}$, 所以 $K := B \cap V$ 开且紧 (因为它在 $V$ 中闭) 且包含 $1$.

(b). 设 $U$ 是 $1$ 的开对称邻域使得 $UK \subset K$, 设 $G_0$ 是由 $U$ 生成的 $G$ 的子群. 让我们证明 $G_0$ 是 $G$ 的开紧子群且 $G_0 \subset K$. 首先证明 $G_0$ 开. 设 $g \in G_0$; 则 $gU \subset G_0$ 且 $gU$ 在 $G$ 中开, 所以 $G_0$ 包含 $g$ 的邻域. 由于拓扑群的每个开子群也闭, 我们又得 $G_0$ 闭. 所以要证明它紧, 只需证明它包含于 $K$. 注意, 由于 $U$ 对称且含 $1$, 有 $G_0 = \bigcup_{n \geq 1}U^n$. 由于 $U \subset K$ (因为 $1 \in K$) 且 $UK \subset K$, 简单归纳证明对每个 $n \geq 1$ 有 $U^n \subset K$. 所以 $G_0 \subset K$.

(c). (a) 的解答中的论证实际上证明 $1$ 的每个紧邻域都包含 $1$ 的开紧邻域, 然后 (b) 蕴含它也包含 $G$ 的紧开子群. 因此, 由于 $G$ 局部紧, $1$ 在 $G$ 中的每个邻域都包含 $G$ 的紧开子群.

(d). 让我们在 $M_n(\Q_p)$ 上选取诱导乘积拓扑的范数. 例如, 由
\[
\|(a_{ij})_{1 \leq i,j \leq n} - (b_{ij})_{1 \leq i,j \leq n}\| = \sup_{1 \leq i,j \leq n}|a_{ij} - b_{ij}|_p
\]
定义的范数 $\|\cdot\|$ 即可. 对每个整数 $m \geq 1$, 设 $K_m = I_n + p^mM_n(\Z_p)$. 按我们选的范数, 这正是 $M_n(\Q_p)$ 中以 $I_n$ 为心、$p^{-m+1}$ 为半径的开球 (也是以 $I_n$ 为心、$p^{-m}$ 为半径的闭球). 特别地, 集合 $K_m$ ($m \geq 1$) 构成 $I_n$ 在 $M_n(\Q_p)$ 中的开邻域族, 从而 $K_m$ ($m \gg 0$) 构成 $I_n$ 在 $\GL_n(\Q_p)$ 中的开邻域族 (因为 $\GL_n(\Q_p)$ 在 $M_n(\Q_p)$ 中开, 它是连续映射 $\det$ 下 $\Q_p$ 的开子集 $\Q_p^{\times}$ 的原像).
又注意 $K_m$ 与 $M_n(\Z_p) \simeq \Z_p^{n^2}$ 同胚 (通过映射 $I_n + p^mX \longmapsto X$), 所以它也紧.
至此, 我们已经有了由紧开子群组成的邻域基. 实际上我们可以更精确, 证明对每个 $m \geq 1$ 有 $K_m \subset \GL_n(\Q_p)$ (而不只是对足够大的 $m$), 这只需 $K_1 \subset \GL_n(\Q_p)$. 事实上, 甚至有 $K_1 \subset \GL_n(\Z_p)$. 确实, 显然 $K_1 \subset M_n(\Z_p)$. 而且, 若 $X \in K_1$, 则容易看出 $\det(X) \in 1 + p\Z_p \subset \Q_p$, 由习题 \ref{ex:I.5.1.4}(a) 这蕴含 $|\det(X)|_p = 1$, 从而 $\det(X)^{-1}$ 也在 $\Z_p$ 中, 即 $\det(X) \in \Z_p^{\times}$.
\end{solution}

\subsection{Haar 测度的例子}

\begin{exercise}[仿射平移的群上的左 Haar 测度]
\label{ex:I.5.3.1}
设 $G$ 是拓扑群. 假设存在 $G$ 与某个 $\R^N$ 的开子集的同胚 (不要求与任何群结构相容), 使得 $G$ 上的左平移由仿射映射给出. 也就是说, 若把 $G$ 等同于它在 $\R^N$ 中的像 (只是作为拓扑空间!), 则对每个 $x \in G$, 存在 $N \times N$ 矩阵 $A(x) \in M_N(\R)$ 和元素 $b(x) \in \R^N$ 使得对每个 $y \in G$, $xy = A(x)y + b(x)$.
证明 $|\det A(x)|^{-1}dx$ 是 $G$ 上的左 Haar 测度, 其中 $dx$ 表示 $\R^N$ 上的 Lebesgue 测度. (提示: 换元公式. 另外, 首先证明 $x$ 唯一确定 $A(x)$ 和 $b(x)$, 且 $x \longmapsto A(x)$ 是从 $G$ 到 $\GL_N(\R)$ 的群同态.)
\end{exercise}

\begin{solution}
设 $x \in G$. 假设有 $A, A' \in M_N(\R)$ 和 $b, b' \in \R^N$ 使得对每个 $y \in G$ 有 $xy = Ay + b = A'y + b'$. 则对每个 $y \in G$ 有 $(A - A')y = b' - b$. 但方程 $(A - A')y = b' - b$ 的解集是 $\R^N$ 的仿射子空间 (即线性子空间的平移), 所以除非它等于 $\R^N$, 否则它有空的内部. 由于 $G$ 在 $\R^N$ 中开, 这意味着必有 $(A - A')y = b' - b$ 对所有 $y \in \R^N$ 成立. 这只有在 $\Ker(A - A') = \R^N$ 时才可能, 故 $A = A'$, 然后也有 $b = b'$. 所以 $x$ 确定 $A(x)$ 和 $b(x)$.
我们证明对每个 $x \in G$, $A(x)$ 可逆. 事实上, 若 $A(x)$ 不可逆, 则映射 $G \to G$, $y \longmapsto xy$ 的像包含于 $b(x) + \Im(A(x))$, 这是 $\R^N$ 的真线性子空间经 $b(x)$ 平移后的集, 所以它有空的内部. 但此像必须等于 $G$, 从而是 $\R^N$ 的开子集, 所以我们得到矛盾, 因此 $A(x)$ 必可逆.
我们证明 $x \longmapsto A(x)$ 是群同态. 设 $x_1, x_2 \in G$. 对每个 $y \in G$,
\[
A(x_1x_2)y + b(x_1x_2) = x_1x_2y = A(x_1)A(x_2)y + A(x_1)b(x_2) + b(x_1).
\]
由第一段, 这蕴含 $A(x_1x_2) = A(x_1)A(x_2)$ 且 $b(x_1x_2) = A(x_1)b(x_2) + b(x_1)$.

现在回忆换元公式蕴含: 若 $U$ 是 $\R^N$ 的可测子集, $A \in \GL_N(\R)$, $b \in \R^N$, 且 $V$ 是 $U$ 在变换 $y \longmapsto Ay + b$ 下的像, 则对每个 $f \in L^1(V)$,
\[
\int_V f(v)\,\dd v = |\det A|\int_U f(Ay + b)\,\dd y.
\]
将其应用于 $U = V = G$, $A = A(y)$, $b = b(y)$ (对某个 $y \in G$), 对每个 $f \in L^1(G)$ 得到
\[
\int_G f(x)\,\dd x = |\det A(y)|\int_G f(yx)\,\dd x.
\]
设 $f \in \mathcal{C}_c(G)$. 则函数 $x \longmapsto |\det A(x)|^{-1}f(x)$ 也在 $\mathcal{C}_c(G)$ 中, 所以可以把上一段应用于它. 对每个 $y \in G$ 得到
\[
\int_G f(yx)|\det A(yx)|^{-1}\,\dd x = |\det A(y)|\int_G f(x)|\det A(x)|^{-1}\,\dd x.
\]
利用 $\det(A(xy)) = \det(A(x))\det(A(y))$, 两边除以 $|\det(A(y))|^{-1}$ 得
\[
\int_G f(yx)|\det A(x)|^{-1}\,\dd x = \int_G f(x)|\det A(x)|^{-1}\,\dd x,
\]
这就是所要的结果.
\end{solution}

\begin{exercise}[经典李群的 Haar 测度]
\label{ex:I.5.3.2}
在此题中, $dx$ 总表示 $\R$ 上的 Lebesgue 测度.
\begin{enumerate}
\item[(a)] 证明 $\frac{dx}{|x|}$ 是乘法群 $\R^{\times}$ 上的 Haar 测度.
\item[(b)] 证明 $\frac{dxdy}{x^2 + y^2}$ 是乘法群 $\CC^{\times}$ (坐标为 $z = x + iy$) 上的 Haar 测度.
\item[(c)] 设 $dT$ 是 $M_n(\R)$ 上的 Lebesgue 测度. 证明 $|\det T|^{-n}dT$ 是 $\GL_n(\R)$ 上的左、右 Haar 测度.
\item[(d)] 设 $G = \Bigl\{\begin{pmatrix} x & y \\ 0 & z \end{pmatrix} \;\Big|\; x, z \in \R^{\times},\, y \in \R\Bigr\}$. 证明 $\frac{dxdydz}{x^2|z|}$ 是 $G$ 上的左 Haar 测度. 它是右 Haar 测度吗?
\end{enumerate}
\end{exercise}

\begin{solution}
当然, 我们将用前一题来解每一问.

(a). 显然的包含 $R^{\times} \subset \R$ 使 $\R^{\times}$ 成为 $\R$ 的开子集. 设 $x \in \R^{\times}$. 则对每个 $y \in \R^{\times}$, $xy = A(x)y + b(x)$, 其中 $A(x) = x \in \GL_1(\R)$, $b(x) = 0$. 所以结论由 $\det(A(x)) = x$ 得出.

(b). 通过映射 $z \longmapsto (\Re(z), \Im(z))$ 把 $\CC^{\times}$ 嵌入 $\R^2$. 这使 $\CC^{\times}$ 成为 $\R^2$ 的开子集. 设 $z = x + iy \in \CC^{\times}$, 其中 $x, y \in \R$. 则 $\CC^{\times}$ 上按 $z$ 的左平移由左乘 $2 \times 2$ 矩阵 $A(z) = \begin{pmatrix} x & -y \\ y & x \end{pmatrix}$ 给出 (这只是公式 $(x + iy)(a + ib) = (xa - yb) + i(ya + xb)$). 所以结论由 $\det(A(z)) = x^2 + y^2$ 得出.

(c). 群 $\GL_n(\R)$ 是 $M_n(\R) \simeq \R^{n^2}$ 的开子集 (因为它由方程 $\det(x) \neq 0$ 给出). 设 $x \in \GL_n(\R)$. 则 $M_n(\R)$ 上按 $x$ 的左平移是线性变换, 我们需要计算它的行列式. 注意 $M_n(\R) = \R^n \oplus \cdots \oplus \R^n$, 其中我们有 $n$ 个和项, 对应 $n \times n$ 矩阵的 $n$ 个列. 按 $x$ 的左乘保持这个分解, 且它在每个和项上的作用就是 $x$ 在 $\R^n$ 上的通常作用, 其行列式为 $\det(x)$. 所以 $M_n(\R)$ 上按 $x$ 的左平移的行列式为 $\det(x)^n$.
为看出 $|\det(T)|^{-n}dT$ 也是右 Haar 测度, 我们使用前一题中把左平移换成右平移的显然类似, 并如上看到 $x \in \GL_n(\R)$ 在 $M_n(\R)$ 上按右平移的作用的行列式为 $\det(x)^n$.

(d). 我们通过把 $\begin{pmatrix} x & y \\ 0 & z \end{pmatrix}$ 送到 $(x, y, z)$ 把 $G$ 嵌入为 $\R^3$ 的开子集. 设 $g = \begin{pmatrix} x & y \\ 0 & z \end{pmatrix} \in G$. 利用
\[
\begin{pmatrix} x & y \\ 0 & z \end{pmatrix}\begin{pmatrix} a & b \\ 0 & c \end{pmatrix} = \begin{pmatrix} xa & xb + yc \\ 0 & zc \end{pmatrix},
\]
我们看到我们处于前一题的情形, 其中
\[
A(g) = \begin{pmatrix} x & 0 & 0 \\ 0 & x & y \\ 0 & 0 & z \end{pmatrix}
\]
且 $b(g) = 0$. 所以 $\det(A(g)) = x^2z$, 我们得到结论.
同 (c), 使用前一题对右平移的类似, 我们看到 $g$ 在 $G$ 上按右平移的作用是线性的, 由矩阵
\[
\begin{pmatrix} x & 0 & 0 \\ y & z & 0 \\ 0 & 0 & z \end{pmatrix}
\]
给出, 其行列式为 $xz^2$. 所以 $\frac{dxdydz}{|x|z^2}$ 是 $G$ 上的右 Haar 测度. 它不是形如 $c\frac{dxdydz}{x^2|z|}$ 的 (其中 $c$ 是常数), 所以 $\frac{dxdydz}{x^2|z|}$ 不能是右 Haar 测度.
\end{solution}

\begin{exercise}[$(\Z/2\Z)^{\N}$ 的 Haar 测度]
\label{ex:I.5.3.3}
考虑群 $G = (\Z/2\Z)^{\N}$.
\begin{enumerate}
\item[(a)] 证明存在 $G$ 上的 Haar 测度 $\mu$ 使得 $\mu(G) = 1$.
\item[(b)] 证明 $G$ 的每个开子集都是形如 $U = V \times (\Z/2\Z)^{\N_{\geq n+1}}$ 的集合的可数并, 其中 $n \in \N$, $V \subset (\Z/2\Z)^{\{0,\ldots,n\}}$, 且 $\mu(U) = \frac{|V|}{2^{n+1}}$.
\item[(c)] 考虑映射 $u \colon G \to [0,1]$, 把 $(x_n)_{n \in \N} \in G$ 送到 $\sum_{n \geq 0}x_n2^{-n-1}$. (在 $u$ 的定义中把 $\Z/2\Z$ 等同于 $\{0,1\}$.) 证明 $u$ 可测, 并把 $\mu$ 送到 $[0,1]$ 上的 Lebesgue 测度 $\lambda$. 即证明: 若 $B \subset [0,1]$ 是 Borel 集, 则 $u^{-1}(B)$ 是 Borel 集且 $\lambda(B) = \mu(u^{-1}(B))$. (提示: 证明形如 $[j2^{-k}, (j+1)2^{-k}]$ 的半开区间生成 $[0,1]$ 上的 Borel $\sigma$-代数, 并计算它们在 $u$ 下的原像.)
\end{enumerate}
\end{exercise}

\begin{solution}
(a). 设 $\mu$ 是 $G$ 上的左 Haar 测度. 由于 $G$ 交换, $\mu$ 也是右 Haar 测度. 又, 由习题 \ref{ex:I.5.1.3}(b), 群 $G$ 紧, 所以 $\mu(G) < +\infty$, 且把 $\mu$ 乘以 $\mu(G)^{-1}$ 后可以假设 $\mu(G) = 1$.

(b). 由乘积拓扑的定义, $G$ 的每个开子集都是形如 $U = V \times (\Z/2\Z)^{\N - I}$ 的集合的并, 其中 $I \subset \N$ 有限, $V \subset (\Z/2\Z)^I$. 由于 $\N$ 的每个有限子集都包含于形如 $\{0,1,\ldots,n\}$ 的集合, 可以假设 $I = \{0,1,\ldots,n\}$ (对某个 $n \in \N$). 我们还需证明并可以是可数的. 假设有 $G$ 的开集 $\Omega$ 具有形式 $\bigcup_{i \in I}U_i$, 其中 $U_i = V_i \times (\Z/2\Z)^{\N_{\geq n_i+1}}$, $V_i \subset (\Z/2\Z)^{\{0,\ldots,n_i\}}$. 对每个 $n \in \N$, 设 $I_n = \{i \in I \mid n_i = n\}$. 则
\[
\bigcup_{i \in I_n}U_i = V_n \times (\Z/2\Z)^{\N_{\geq n+1}},
\]
其中
\[
V_n = \bigcup_{i \in I_n}V_i \subset (\Z/2\Z)^{\{0,\ldots,n\}}.
\]
所以 $\Omega = \bigcup_{n \in \N}V_n \times (\Z/2\Z)^{\N_{\geq n+1}}$.

现在计算 $\mu(U)$, 其中 $U = V \times (\Z/2\Z)^{\N_{\geq n+1}}$, $V \subset (\Z/2\Z)^{\{0,\ldots,n\}}$. 记 $W = (\Z/2\Z)^{\N_{\geq n+1}}$. 若 $v, v' \in (\Z/2\Z)^{\{0,\ldots,n\}}$, 则 $\{v\} \times W = (v' - v) + \{v'\} \times W$, 所以 $\mu(\{v\} \times W) = \mu(\{v'\} \times W)$. 由于 $G = \coprod_{v \in (\Z/2\Z)^{\{0,\ldots,n\}}}\{v\} \times W$, 这蕴含对每个 $v \in (\Z/2\Z)^{\{0,\ldots,n\}}$,
\[
1 = \mu(G) = |(\Z/2\Z)^{\{0,\ldots,n\}}|\mu(\{v\} \times W),
\]
所以 $\mu(\{v\} \times W) = \frac{1}{2^{n+1}}$. 另一方面, $U = \coprod_{v \in V}\{v\} \times W$, 所以我们得到 $\mu(U) = \frac{|V|}{2^{n+1}}$.

(c). 记 $I_{j,k} = [j2^{-k}, (j+1)2^{-k})$. 我们首先证明 $I_{j,k}$ ($k \geq 0$, $0 \leq j \leq 2^k - 1$) 生成 $[0,1]$ 上的 Borel $\sigma$-代数. $[0,1]$ 的每个开子集都是开区间 $(a, b)$ ($0 \leq a < b \leq 1$) 以及可能的一个或两个半开区间 $[0, b)$ ($0 < b \leq 1$) 和 $(a, 1]$ ($0 \leq a < 1$) 的可数并. 所以我们只需验证其中任何一个都可以写成 $I_{j,k}$ 的可数并.
假设 $0 \leq a < b \leq 1$. 若 $b - a > 2^{-k}$, 且 $i = 1 + \lfloor b2^k a \rfloor$, $i' = -1 + \lceil b2^k \rceil$ (其中 $\lfloor\cdot\rfloor$ 和 $\lceil\cdot\rceil$ 是取整函数), 则 $0 < i2^{-k} - a \leq 2^{-k}$ 且 $0 < b - i'2^{-k} \leq 2^{-k}$. 这蕴含
\[
(a, b) = \bigcup_{k > -\log_2(b-a)}\ \bigcup_{j=1+\lfloor b2^k a \rfloor}^{-2+\lceil b2^k \rceil} I_{j,k}
\]
(其中 $\log_2$ 是以 $2$ 为底的对数). 类似地, 若 $0 < b \leq 1$ 且 $0 \leq a < 1$, 则
\[
[0, b) = \bigcup_{k > -\log_2(b)}\ \bigcup_{j=0}^{-2+\lceil b2^k \rceil} I_{j,k}
\]
且
\[
(a, 1] = \bigcup_{k > -\log_2(1-a)}\ \bigcup_{j=1+\lfloor b2^k a \rfloor}^{2^k - 1} I_{j,k}.
\]
这证明了关于 $[0,1]$ 的 Borel $\sigma$-代数的结论. 为完成此题, 我们只需证明对所有 $k \geq 0$ 和所有 $j \in \{0,\ldots,2^k-1\}$, 原像 $u^{-1}(I_{j,k})$ 是 $G$ 中的 Borel 集且 $\mu(u^{-1}(I_{j,k})) = 2^{-k}$. 所以我们计算这些原像.
首先注意, 若 $x \in [0,1]$, 则 $u^{-1}(x)$ 是单点集, 除非 $x$ 具有 $j2^{-k}$ 的形式 ($0 < j < 2^k$); 在后一情形, $x$ 有第二个二进制展开, 其中所有系数在某一秩之后都是 $1$.
现在设 $k \geq 0$, $j \in \{0,1,\ldots,2^k-1\}$. 若 $k = 0$, 则 $j = 0$ 且 $I_{j,k} = [0,1]$, 所以 $u^{-1}(I_{j,k}) = G$ 且由 $\mu$ 的选取有 $\mu(G) = \lambda([0,1]) = 1$. 假设 $k \geq 1$. 由于 $0 \leq j \leq 2^k - 1$, 可以把 $j$ 写成二进制形式 $j = \sum_{i=0}^{k-1}a_{k-1-i}2^i$, 其中 $a_i \in \{0,1\}$. 若 $0 < j$, 还把 $j - 1$ 写成 $j - 1 = \sum_{i=0}^{k-1}b_{k-1-i}2^i$, 其中 $b_i \in \{0,1\}$. 若 $j + 1 < 2^k$, 还把 $j + 1$ 写成 $j + 1 = \sum_{i=0}^{k-1}c_{k-1-i}2^i$, 其中 $c_i \in \{0,1\}$. 则
\[
j2^{-k} = \sum_{i=0}^{k-1}a_i2^{-(i+1)},
\]
\[
j2^{-k} = \sum_{i=0}^{k-1}b_i2^{-(i+1)} + \sum_{i=k}^{+\infty}2^{-i} \quad \text{若 } j > 0
\]
且
\[
(j+1)2^{-k} = \sum_{i=0}^{k-1}c_i2^{-(i+1)} \quad \text{若 } j + 1 < 2^k,
\]
所以 $u^{-1}(I_{j,k}) = X \cup Y$, 其中
\[
Y = \{(a_0,\ldots,a_{k-1})\} \times (\Z/2\Z)^{\N_{\geq k}}
\]
且
\[
X = \begin{cases} \{(b_0,\ldots,b_{k-1},1,1,\ldots),(c_0,\ldots,c_{k-1})\} & \text{若 } 0 < j < 2^k - 1, \\ \{(b_0,\ldots,b_{k-1},1,1,\ldots)\} & \text{若 } j = 2^k - 1, \\ \{(c_0,\ldots,c_{k-1})\} & \text{若 } j = 0. \end{cases}
\]
由于 $X$ 闭且 $Y$ 开, 这是 $G$ 的 Borel 子集. 我们又由 (b) 知道 $\mu(Y) = 2^{-k} = \lambda(I_{j,k})$, 所以只需证明 $\mu(X) = 0$. 即要证明 $G$ 中的单点集体积为 $0$. 由于所有单点集互相是平移, 只需处理 $\{0\}$ 的情形. 这由对每个 $n \geq 0$,
\[
\{0\} \subset \{0\}^{\{0,1,\ldots,n\}} \times (\Z/2\Z)^{\N_{\geq n+1}}
\]
的事实得出, 因为由 (b), 右端体积为 $2^{-(n+1)}$.
\end{solution}

\begin{exercise}[$\Q_p$ 的 Haar 测度]
\label{ex:I.5.3.4}
对 $x \in \Q_p$ 和 $r \in \R$, 记 $B(x, r) = \{y \in \Q_p \mid |x - y|_p \leq r\}$ (以 $x$ 为心、$r$ 为半径的闭球). 设 $\lambda$ 是 $\Q_p$ 上满足 $\lambda(\Z_p) = 1$ 的 Haar 测度.
\begin{enumerate}
\item[(a)] 若 $x \in \Q_p$, $m \in \Z$, 证明 $\lambda(B(x, p^m)) = p^m$.
\item[(b)] 对每个 Borel 集 $X \subset \Q_p$, 证明
\[
\lambda(X) = \inf\Bigl\{\sum_{i \geq 0}p^{m_i} \;\Big|\; \exists x_0, x_1, \ldots \in \Q_p \text{ 使得 } X \subset \bigcup_{i \geq 0}B(x_i, p^{m_i})\Bigr\}.
\]
\end{enumerate}
\end{exercise}

\begin{solution}
(a). 首先注意, 若 $x, y \in \Q_p$, 则 $B(x, p^m) = B(y, p^m) + x - y$, 所以 $\lambda(B(x, p^m)) = \lambda(B(y, p^m))$. 又注意
\[
B(0, p^m) = \{x \in \Q_p \mid |x|_p \leq p^m\} = \{x \in \Q_p \mid |p^mx|_p \leq 1\} = p^{-m}\Z_p
\]
(对每个 $m \in \Z$). 所以对每个 $x \in \Q_p$ 和每个 $m \in \Z$,
\[
B(x, p^m) = x + p^{-m}\Z_p.
\]
又, 由习题 \ref{ex:I.5.1.4}(i),
\[
\Z_p = \coprod_{a=0}^{p-1}(a + p\Z_p).
\]
乘以 $p^{-m}$ 给出
\[
B(0, p^m) = \coprod_{i=0}^{p-1}B(p^{-m}i, p^{m-1}),
\]
所以 $\lambda(B(0, p^m)) = p\lambda(B(0, p^{m-1}))$. 由于由假设 $\lambda(B(0, 1)) = \lambda(\Z_p) = 1$, 结论对 $|m|$ 归纳得出.

(b). 首先, 由习题 \ref{ex:I.5.1.4}(f), 球 $B(x, r)$ 构成 $\Q_p$ 的拓扑的一组 (开!) 集基, 所以 $\Q_p$ 的每个开子集都是球 $B(x, r)$ 的并. 由于 $\Q$ 在 $\Q_p$ 中稠密且可数, $\Q_p$ 的每个开子集都是球 $B(x, r)$ 的可数并 (且可以把 $x$ 取在 $\Q$ 中, 但这无关紧要). 又注意, 由习题 \ref{ex:I.5.1.4}(b), 若 $y \in B(x, r)$, 则 $B(y, r) = B(x, r)$. 所以若 $\Q_p$ 的两个闭球相交, 则其中一个必包含另一个. 这蕴含 $\Q_p$ 的每个开子集都是球 $B(x, r)$ 的不相交可数并. 结论现在由 (a) 和 $\lambda$ 的外正则性立即得出.
\end{solution}

\begin{exercise}[商群上的 Haar 测度]
\label{ex:I.5.3.5}
设 $G$ 是局部紧群, $H$ 是 $G$ 的闭子群. 记 $\pi$ 为从 $G$ 到 $G/H$ 的商映射. 记 $\Delta_G$ (相应地 $\Delta_H$) 为 $G$ (相应地 $H$) 的模函数, 并假设 $\Delta_G|_H = \Delta_H$. 固定 $G$ 和 $H$ 上的左 Haar 测度 $\mu_G$ 和 $\mu_H$.
\begin{enumerate}
\item[(a)] 证明对 $G/H$ 的每个紧子集 $K'$, 存在 $G$ 的紧子集 $K$ 使得 $\pi(K) = K'$.
\item[(b)] 设 $f \in L^1(G)$. 证明函数 $G \to \CC$, $x \longmapsto \int_H f(xh)\,\dd\mu_H(h)$ 被 $H$ 的元素的右平移不变. 因此它定义了一个函数 $G/H \to \CC$, 记为 $f^H$.
\item[(c)] 若 $f \in \mathcal{C}_c(G)$, 证明 $f^H \in \mathcal{C}_c(G/H)$.
\item[(d)] 证明映射 $\mathcal{C}_c(G) \to \mathcal{C}_c(G/H)$, $f \longmapsto f^H$ 满射. (提示: 可以利用对 $G$ 的每个紧子集 $K$, 存在 $\phi \in \mathcal{C}_c^+(G)$ 使得对每个 $x \in K$ 有 $\phi(x) > 0$ 这一事实.)
\item[(e)] 若 $f \in \mathcal{C}_c(G)$ 满足 $f^H = 0$, 证明 $\int_G f(x)\,\dd\mu_G(x) = 0$. (提示: 用 $\mathcal{C}_c(G/H)$ 中在 $\pi(\operatorname{supp}(f))$ 上等于 $1$ 的函数, 以及命题 \ref{prop:I.2.12}.)
\item[(f)] 证明存在 $G/H$ 上唯一的正则 Borel 测度 $\mu_{G/H}$, 它被 $G$ 的元素的左平移不变, 且对每个 $f \in \mathcal{C}_c(G)$ 有 $\int_G f(x)\,\dd\mu_G(x) = \int_{G/H}f^H(y)\,\dd\mu_{G/H}(y)$.
\item[(g)] 若 $P$ 是 $G$ 的闭子群且 $\pi$ 诱导同胚 $P \xrightarrow{\sim} G/H$, 证明 $\mu_{G/H}$ 在此同胚下的原像是 $P$ 上的左 Haar 测度.
\item[(h)] 若 $P$ 是 $G$ 的闭子群, 映射 $P \times H \to G$, $(p, h) \longmapsto ph$ 是同胚, 且 $d\mu_P$ 是 $P$ 上的左 Haar 测度, 证明线性泛函 $\mathcal{C}_c(G) \to \CC$, $f \longmapsto \int_H\int_P f(ph)\,\dd\mu_P(p)\,\dd\mu_H(h)$ 定义 $G$ 上的左 Haar 测度.
\end{enumerate}
\end{exercise}

\begin{solution}
(a). 设 $V$ 是 $1$ 在 $G$ 中的紧邻域. 则 $\pi(V)$ 是 $\pi(1)$ 在 $G/H$ 中的紧邻域. 我们有 $K' \subset \bigcup_{x \in \pi^{-1}(K')}\pi(xV)$. 由于 $K'$ 紧, 可以找到 $x_1, \ldots, x_n$ 使得 $K' \subset \bigcup_{i=1}^n\pi(x_iV)$. 设 $K = \pi^{-1}(K') \cap \bigl(\bigcup_{i=1}^n x_iV\bigr)$. 则 $K$ 是紧集 $\bigcup_{i=1}^n x_iV$ 的闭子集, 所以它紧, 且有 $\pi(K) = K'$.

(b). 设 $x \in H$. 则对每个 $g \in G$,
\[
\int_H f(gxh)\,\dd\mu_H(h) = \int_H f(gh)\,\dd\mu_H(h)
\]
由 $\mu_H$ 的左不变性.

(c). 我们需要证明 $f^H$ 连续且有紧支撑.
固定 $1$ 的对称紧邻域 $V_0$, 注意 $A := \operatorname{supp} f \cup V_0(\operatorname{supp} f)$ 紧. 设 $\varepsilon > 0$. 由于 $f$ 左一致连续, 存在 $1$ 的邻域 $V \subset V_0$ 使得对每个 $x \in G$ 和每个 $y \in V$, 有 $|f(yx) - f(x)| \leq \varepsilon$. 则对每个 $x \in G$ 和每个 $y \in V$,
\[
|f^H(\pi(yx)) - f^H(x)| = \Bigl|\int_H (f(yxh) - f(xh))\,\dd\mu_H(h)\Bigr| \leq \varepsilon\mu_H(x^{-1}A \cap H),
\]
因为除非 $y \in (x^{-1}\operatorname{supp} f) \cup (x^{-1}y^{-1}\operatorname{supp} f) \subset x^{-1}A$, 有 $f(yxh) = f(xh) = 0$. 由于 $x^{-1}A \cap H$ 紧, 它有有限测度, 且上面的计算蕴含 $f^H$ 在点 $\pi(x)$ 处连续.
现在我们证明 $f^H$ 有紧支撑. 由 $f^H$ 的定义, 若 $x \notin KH$ 则 $f^H(\pi(x)) = 0$. 所以 $f^H$ 的支撑包含于 $\pi(KH) = \pi(K)$, 从而是紧的.

(d). 设 $g \in \mathcal{C}_c(G/H)$, $K'$ 是它的支撑. 由 (a), 存在 $G$ 的紧子集 $K$ 使得 $\pi(K) = K'$. 设 $\phi \in \mathcal{C}_c^+(G)$ 使得对每个 $x \in K$ 有 $\phi(x) > 0$. 我们证明对每个 $y \in K'$ 有 $\phi^H(y) > 0$. 设 $y \in K'$, 把 $y$ 写成 $y = \pi(x)$ ($x \in K$). 由于 $\phi(x) > 0$ 且 $\phi$ 连续, 可以找到 $1$ 在 $G$ 中的开邻域 $V$ 和 $c \in \R_{>0}$ 使得对每个 $x' \in xV$ 有 $\phi(x') \geq c$. 特别地,
\[
\phi^H(y) = \int_H \phi(xh)\,\dd\mu_H(h) \geq \int_{H \cap V}\phi(xh)\,\dd\mu_H(h) \geq c \cdot \mu_H(U \cap H) > 0
\]
(因为 $U \cap H$ 是 $H$ 的非空开子集, 有 $\mu_H(U \cap H) > 0$).
我们如下定义函数 $F \colon G \to \CC$:
\[
F(x) = \begin{cases} \frac{g(\pi(x))}{\phi^H(\pi(x))} & \text{若 } \phi^H(\pi(x)) > 0, \\ 0 & \text{否则}. \end{cases}
\]
注意 $F$ 在开子集 $U_1 = \{x \in G \mid \phi^H(\pi(x)) > 0\}$ 和 $U_2 = G - \operatorname{supp}(g \circ \pi)$ 上连续 (在第二个子集上它恒为零). 由于 $U_1 \supset \pi^{-1}(K')$ 且 $\pi^{-1}(K') = \operatorname{supp}(g \circ \pi)$, 有 $U_1 \cup U_2 = G$, 函数 $F$ 在 $G$ 上连续. 最后, 取 $f = F\phi$. 则 $f \in \mathcal{C}_c(G)$, 我们只需证明 $f^H = g$.
设 $x \in G$. 若 $\phi^H(\pi(x)) = 0$, 则对每个 $h \in H$ 有 $f(xh) = 0$, 所以 $f^H(\pi(x)) = 0$. 我们已见 $\phi^H$ 在 $K' = \operatorname{supp}(g)$ 上取正值, 所以我们也有 $x \notin \operatorname{supp}(g)$, 即 $g(x) = 0 = f^H(x)$. 现在假设 $\phi^H(\pi(x)) > 0$. 注意函数 $H \to \CC$, $h \longmapsto F(xh)$ 是常数. 所以
\[
f^H(\pi(x)) = F(x)\int_H \phi(xh)\,\dd\mu_H(h) = \frac{g(\pi(x))}{\phi^H(\pi(x))}\phi^H(\pi(x)) = g(\pi(x)).
\]
最后, 注意若 $g \in \mathcal{C}_c^+(G/H)$ 则 $f \in \mathcal{C}_c^+(G)$, 且我们顺带证明了若 $f \in \mathcal{C}_c^+(G)$ 则 $f^H \in \mathcal{C}_c^+(G/H)$ (我们是对 $\phi$ 证明的).

(e). 设 $\psi \in \mathcal{C}_c(G/H)$ 使得对每个 $y \in \pi(\operatorname{supp} f)$ 有 $\psi(y) = 1$. 由 (d), 存在 $\phi \in \mathcal{C}_c(G)$ 使得 $\phi^H = \psi$. 我们有
\begin{align*}
\int_G f(x)\,\dd\mu_G(x) &= \int_G f(x)\phi^H(\pi(x))\,\dd\mu_G(x) \\
&= \int_{G \times H} f(x)\phi(xh)\,\dd\mu_G(x)\,\dd\mu_H(h) \\
&= \int_H \Bigl(\int_G f(x)\phi(xh)\,\dd\mu_G(x)\Bigr)\,\dd\mu_H(h) \\
&= \int_H \Bigl(\Delta_G(h)^{-1}\int_G f(xh^{-1})\phi(x)\,\dd\mu_G(x)\Bigr)\,\dd\mu_H(h) \\
&= \int_G \phi(x)\Bigl(\int_H \Delta_H(h)^{-1}f(xh^{-1})\,\dd\mu_H(h)\Bigr)\,\dd\mu_G(x) \\
&= \int_G \phi(x)\Bigl(\int_H f(xh)\,\dd\mu_H(h)\Bigr)\,\dd\mu_G(x) \quad \text{(由命题 \ref{prop:I.2.12})} \\
&= 0 \quad \text{(因为 } f^H = 0\text{)}.
\end{align*}

(f). 由 (e), 正线性函数 $\mathcal{C}_c(G) \to \CC$, $f \longmapsto \int_G f\,\dd\mu_G$ 穿过线性映射 $\mathcal{C}_c(G) \to \mathcal{C}_c(G/H)$, $f \longmapsto f^H$. 由 (d) (及其解答末尾的注记), 它定义了 $\mathcal{C}_c(G/H)$ 上的正线性泛函. 由 Riesz 表示定理, 这来自 $G/H$ 上的正则 Borel 测度 $\mu_{G/H}$. 展开定义, 对每个 $f \in \mathcal{C}_c(G)$ 得到
\[
\int_G f\,\dd\mu_G = \int_{G/H}f^H\,\dd\mu_{G/H}.
\]
由 $\mu_G$ 的左不变性和 (d), 若 $f \in \mathcal{C}_c(G/H)$, $x \in G$, 则
\[
\int_{G/H}f(xy)\,\dd\mu_{G/H}(y) = \int_{G/H}f(y)\,\dd\mu_{G/H}.
\]
使用 Riesz 表示定理的唯一性部分 (如课堂上那样), 看到对 $G/H$ 的每个 Borel 子集 $E$ 有 $\mu_{G/H}(xE) = \mu_{G/H}(E)$.

(g). 设 $\nu$ 是同胚 $\alpha \colon P \xrightarrow{\sim} G/H$ 下 $\mu_{G/H}$ 的原像. 它是正则 Borel 测度, 因为 $\alpha$ 是同胚. 又注意对每个 $x \in P$ 有 $\alpha(xy) = x\alpha(y)$ (这在 $\alpha$ 的定义上显然). 由于 $\mu_{G/H}$ 被 $P$ 的元素的左平移不变, $\nu$ 也一样.

(h). 假设蕴含 $\pi$ 诱导同胚 $P \xrightarrow{\sim} G/H$, 因此如 (g) 得到 $P$ 上的左 Haar 测度 $\nu$. 由左 Haar 测度的唯一性, 有 $\mu_P = c\nu$ (对某个 $c \in \R_{>0}$). 所以对每个 $f \in \mathcal{C}_c(G)$,
\begin{align*}
\int_H\int_P f(ph)\,\dd\mu_P(p)\,\dd\mu_H(h) &= c\int_P\Bigl(\int_H f(ph)\,\dd\mu_H(h)\Bigr)\,\dd\nu(p) \\
&= c\int_{G/H}f^H(y)\,\dd\mu_{G/H}(y) \\
&= c\int_G f(x)\,\dd\mu_G(x).
\end{align*}
所以泛函 $f \longmapsto \int_H\int_P f(ph)\,\dd\mu_P(p)\,\dd\mu_H(h)$ 正, 且对应于左 Haar 测度 $c\mu_G$.
\end{solution}

\begin{exercise}[半直积的模函数]
\label{ex:I.5.3.6}
设 $G$ 是局部紧群. 设 $A$ 和 $N$ 是 $G$ 的两个闭子群, 使得 $A \times N \to G$, $(a, n) \longmapsto an$ 是同胚, 且 $A$ 正规化 $N$ (即对每个 $a \in A$ 和 $n \in N$, 有 $ana^{-1} \in N$).
\begin{enumerate}
\item[(a)] 若 $\mu_A$ 和 $\mu_N$ 是 $A$ 和 $N$ 上的左 Haar 测度, 证明线性泛函 $\mathcal{C}_c(G) \to \CC$, $f \longmapsto \int_A\int_N f(an)\,\dd\mu_A\,\dd\mu_N$ 定义 $G$ 上的左 Haar 测度.
\item[(b)] 设 $a \in A$. 证明存在 $\alpha(a) \in \R_{>0}$ 使得对每个 $f \in \mathcal{C}_c(N)$,
\[
\int_N f(ana^{-1})\,\dd\mu_N(n) = \alpha(a)\int_N f(n)\,\dd\mu_N(n).
\]
\item[(c)] 若 $\Delta_G$, $\Delta_A$ 和 $\Delta_N$ 分别是 $G$, $A$ 和 $N$ 的模函数, 证明当 $a \in A$, $n \in N$ 时 $\Delta_G(an) = \alpha(a)\Delta_A(a)\Delta_N(n)$.
\end{enumerate}
\end{exercise}

\begin{solution}
(a). 情形与习题 \ref{ex:I.5.3.5} 非常相似 (例如让 $N$ 扮演 $H$ 的角色), 区别是我们对模函数不做任何假设. 尽管如此, 习题 \ref{ex:I.5.3.5} 的 (a)-(d) 的结果仍然成立, 因为它们的证明不用模函数的假设. 特别地, 我们得到从 $\mathcal{C}_c(G)$ 到 $\mathcal{C}_c(G/N) \simeq \mathcal{C}_c(A)$ 的满线性变换 $f \longmapsto f^N$, 且它把 $\mathcal{C}_c^+(G)$ 送到 $\mathcal{C}_c^+(A)$. 题设的线性函数把 $f \in \mathcal{C}_c(G)$ 送到 $\int_A f^N(a)\,\dd\mu_A(a)$, 所以它正, 且 Riesz 表示定理说存在 $G$ 上唯一的正则 Borel 测度 $\mu_G$ 使得对每个 $f \in \mathcal{C}_c(G)$,
\[
\int_G f\,\dd\mu_G = \int_A\int_N f(an)\,\dd\mu_A(a)\,\dd\mu_N(n).
\]
由于 $\mu_A$ 是 $A$ 上的左 Haar 测度, 上面的公式蕴含对每个 $f \in \mathcal{C}_c(G)$ 和每个 $a \in A$ 有 $\int_G L_af\,\dd\mu_G = \int_G f\,\dd\mu_G$. 我们证明 $\mu_G$ 被 $N$ 左不变. 设 $x \in N$, $f \in \mathcal{C}_c(G)$. 则
\begin{align*}
\int_G L_xf\,\dd\mu_G &= \int_A\int_N f(xan)\,\dd\mu_A(a)\,\dd\mu_N(n) \\
&= \int_A\Bigl(\int_N f(a(a^{-1}xa)n)\,\dd\mu_N(n)\Bigr)\,\dd\mu_A(a) \\
&= \int_A\Bigl(\int_N f(an)\,\dd\mu_N(n)\Bigr)\,\dd\mu_A(a) \quad \text{因为 } a^{-1}xa \in N \text{ 且 } \mu_N \text{ 左不变} \\
&= \int_G f\,\dd\mu_G.
\end{align*}
由于 $G = AN$, 这蕴含对每个 $x \in G$ 和每个 $f \in \mathcal{C}_c(G)$ 有 $\int_G L_xg\,\dd\mu_G = \int_G f\,\dd\mu_G$. 由命题 \ref{prop:I.2.6}, $\mu_G$ 是 $G$ 上的左 Haar 测度.

(b). 注意映射 $N \to N$, $n \longmapsto a^{-1}na$ 是同胚. 所以公式 $E \longmapsto \mu_N(a^{-1}Ea)$ 定义 $N$ 上的正则 Borel 测度, 记为 $\nu$. 若 $E$ 是 Borel 子集且 $n \in N$, 则
\[
\nu(nE) = \mu(a^{-1}nEa) = \mu((a^{-1}na)a^{-1}Ea) = \mu(a^{-1}Ea) = \nu(E).
\]
所以 $\nu$ 是 $N$ 上的左 Haar 测度, 因此存在 $\alpha(a) \in \R_{>0}$ 使得 $\nu = \alpha(a)\mu_N$. 现在, 若 $E$ 是 $N$ 的 Borel 子集且 $f = \mathbb{1}_E$, 则函数 $n \longmapsto f(ana^{-1})$ 是 $a^{-1}Ea$ 的特征函数, 所以
\[
\int_N f(ana^{-1})\,\dd\mu_N(n) = \mu(a^{-1}Ea) = \alpha(a)\mu(E) = \alpha(a)\int_N f\,\dd\mu_N.
\]
这按通常方式延拓到所有函数 $f \in L^1(N)$, 特别地延拓到 $f \in \mathcal{C}_c(N)$.

(c). 设 $a \in A$, $n \in N$, 固定 $f \in \mathcal{C}_c^+(G)$. 则
\begin{align*}
\Delta_G(an)\int_G f\,\dd\mu_G &= \int_G R_{an}(f)\,\dd\mu_G \\
&= \int_A\int_N f(bman)\,\dd\mu_A(b)\,\dd\mu_N(m) \\
&= \int_A\Bigl(\int_N f(ba(a^{-1}ma)n)\,\dd\mu_N(m)\Bigr)\,\dd\mu_A(b) \\
&= \alpha(a)^{-1}\int_A\Bigl(\int_N f(bamn)\,\dd\mu_N(m)\Bigr)\,\dd\mu_A(b) \quad \text{由 (b)} \\
&= \alpha(a)^{-1}\Delta_N(n)^{-1}\int_A\Bigl(\int_N f(bam)\,\dd\mu_N(m)\Bigr)\,\dd\mu_A(b) \quad \text{由 } \Delta_N \text{ 的定义} \\
&= \alpha(a)^{-1}\Delta_N(n)^{-1}\int_N\Bigl(\int_A f(bam)\,\dd\mu_A(b)\Bigr)\,\dd\mu_N(m) \\
&= \alpha(a)^{-1}\Delta_N(n)^{-1}\Delta_A(a)^{-1}\int_N\Bigl(\int_A f(bm)\,\dd\mu_A(b)\Bigr)\,\dd\mu_N(m) \quad \text{由 } \Delta_A \text{ 的定义} \\
&= \alpha(a)^{-1}\Delta_N(n)^{-1}\Delta_A(a)^{-1}\int_G f\,\dd\mu_G.
\end{align*}
由于 $\int_G f\,\dd\mu_G > 0$, 这蕴含 $\Delta_G(an) = \alpha(a)\Delta_A(a)\Delta_N(n)$.
\end{solution}

\begin{exercise}[$\SL_n(\R)$ 的 Iwasawa 分解与模函数]
\label{ex:I.5.3.7}
设 $G = \SL_n(\R)$, $H = \SO(n)$, $P \subset G$ 是对角线元素为正的上三角矩阵 (行列式为 $1$) 的子群.
\begin{enumerate}
\item[(a)] 证明映射 $P \times H \to G$, $(p, h) \longmapsto ph$ 是同胚. (提示: Gram-Schmidt.)
\item[(b)] 给出与习题 \ref{ex:I.5.3.2}(d) 中公式类似的 $P$ 上左 Haar 测度的公式.
\item[(c)] 计算 $P$ 的模函数.
\item[(d)] 证明 $G$ 幺模. (有几种方法.)
\item[(e)] 若 $n = 2$, 证明 $\SO(n) \simeq S^1$ (圆群), 并给出 $G$ 上的左 Haar 测度.
\end{enumerate}
\end{exercise}

\begin{solution}
(a). 在此题中, 记 $\langle \cdot, \cdot \rangle$ 为 $\R^n$ 上的通常 Euclid 内积, $\|\cdot\|$ 为对应的范数.
记题设的映射 $P \times H \to G$ 为 $\alpha$. 此映射连续, 因为 $\SL_n(\R)$ 是拓扑群. 我们首先证明它单射. 假设有 $p, p' \in P$ 和 $h, h' \in H$ 使得 $ph = p'h'$. 则 $p^{-1}p' = h(h')^{-1} \in P \cap H$ 是对角线元素为正的上三角特殊正交矩阵. 这样的矩阵必为单位阵. 事实上, 设 $(v_1, \ldots, v_n)$ 是它的列, $(e_1, \ldots, e_n)$ 是 $\R^n$ 的典范基. 我们要证明 $(v_1, \ldots, v_n) = (e_1, \ldots, e_n)$. 由于 $v_1$ 是范数 $1$ 的向量且是 $e_1$ 的正倍数, 必有 $v_1 = e_1$. 由于向量 $v_2, \ldots, v_n$ 与 $v_1$ 正交, 它们的第一个分量全为 $0$. 所以 $v_2$ 是 $e_2$ 的正倍数; 由于 $v_2$ 范数 $1$, 必有 $v_2 = e_2$. 现在向量 $v_3, \ldots, v_n$ 与 $v_2$ 正交, 所以它们的第二个分量为零, 于是 $v_3$ 是 $e_3$ 的正倍数, 等等.
现在回忆 Gram-Schmidt 正交化过程. 若 $(v_1, \ldots, v_n)$ 是 $\R^n$ 的基, 它按下述方式产生正交基 $(w_1, \ldots, w_n)$ 和标准正交基 $(u_1, \ldots, u_n)$:
\begin{itemize}
\item $w_1 = v_1$ 且 $u_1 = \frac{1}{\|w_1\|}w_1$;
\item 对 $1 \leq k \leq n - 1$, $u_{k+1} = \frac{1}{\|w_{k+1}\|}w_{k+1}$, 其中 $w_{k+1} = v_{k+1} - \sum_{i=1}^k\frac{\langle w_i, v_{k+1} \rangle}{\langle w_i, w_i \rangle}w_i$.
\end{itemize}
特别地, 若 $A$ (相应地 $B$, 相应地 $C$) 是列向量为 $(v_1, \ldots, v_n)$ (相应地 $(w_1, \ldots, w_n)$, 相应地 $(u_1, \ldots, u_n)$) 的矩阵, 则 $B = AN$ 且 $C = AND$, 其中 $N$ 是对角线为 $1$ 的上三角矩阵, $D$ 是对角线为 $(\|w_1\|^{-1}, \ldots, \|w_n\|^{-1})$ 的对角矩阵. 又注意 $N$ 和 $D$ 的元素是 $v_1, \ldots, v_n$ 的连续函数, 所以 $B$ 和 $C$ 的元素也是, 且 $C$ 是正交矩阵. 若 $x \in \SL_n(\R)$, 把此过程应用于 $x$ 的列, 得到一个正交矩阵 $h$ 和一个矩阵 $p \in P$, 两者都连续依赖于 $x$, 使得 $h = xp$, 即 $x = hp^{-1}$. 又, $\det(h) = \det(xp) = \det(p) > 0$, 所以 $h$ 实际上在 $\SO(n)$ 中. 由于 $g \longmapsto g^{-1}$ 是 $\GL_n(\R)$ (从而是其子群 $P$) 上的连续函数, 我们构造了连续映射 $\beta \colon G \to P \times H$ 使得 $\alpha \circ \beta = \mathrm{id}_G$. 特别地, 映射 $\alpha$ 满射, 所以它双射. 则它的逆必为 $\beta$, 且我们知道 $\beta$ 连续. 所以 $\alpha$ 是同胚.

(b). 注意 $P$ 是 $M_n(\R)$ 中上三角矩阵构成的 $\R$-向量空间 $V$ 的开子集. 而且, 对每个 $p \in P$, $P$ 上按 $p$ 的左平移是线性自同态 $T_p \colon V \to V$, $x \longmapsto px$ 的限制. 所以我们可以应用习题 \ref{ex:I.5.3.1} 把 $P$ 上的 Haar 测度定义为 $|\det(T_p)|^{-1}dV(p)$, 其中 $dV$ 是 $V$ 上的 Lebesgue 测度.
我们还需计算 $p \in P$ 时的 $\det(T_p)$. 设 $p \in P$, $a_1, \ldots, a_n$ 是它的对角线元素. 设 $(e_1, \ldots, e_n)$ 如前为 $\R^n$ 的典范基, $V_i = \operatorname{Span}(e_1, \ldots, e_i) \subset \R^n$ ($1 \leq i \leq n$). 注意 $p \in \GL_n(\R)$ 的作用保持子空间 $V_1, \ldots, V_n$, 且 $p$ 诱导的 $V_i$ 的自同态的行列式为 $a_1 \cdots a_i$. 用矩阵的列分解 $V$ (如习题 \ref{ex:I.5.3.2}(c) 的解答), 得到同构 $V \simeq V_1 \oplus V_2 \oplus \cdots \oplus V_n$ 使得自同态 $T_p$ 对应于 $p$ 在每个 $V_i$ 上的作用. 所以我们得到
\[
\det(T_p) = \prod_{i=1}^n\prod_{r=1}^i a_r = a_1^na_2^{n-2}\cdots a_{n-1}^2a_n = \prod_{i=1}^n a_i^{n+1-i}.
\]

(c). 我们将使用习题 \ref{ex:I.5.3.6}, 取 $G = P$, $N$ 为幂幺上三角矩阵群 (即对角线为 $1$ 的上三角矩阵群), $A$ 为对角线元素为正的对角矩阵群. 设 $\alpha \colon A \times N \to P$ 是由 $\alpha(a, n) = an$ 定义的映射. 让我们证明 $\alpha$ 是同胚. 映射 $\alpha$ 显然连续, 且它单射, 因为 $N \cap A = \{1\}$. 设 $x \in P$, 设 $a \in A$ 是与 $x$ 有相同对角线元素的矩阵. 则 $n := a^{-1}x \in N$, 且 $\alpha(a, n) = x$. 所以 $\alpha$ 双射. 而且, 矩阵 $a$ 连续依赖于 $x$, 所以 $n$ 也如此, 故 $\alpha$ 的逆连续, 最后 $\alpha$ 是同胚.
我们想应用习题 \ref{ex:I.5.3.6}(c). 为此需要计算 $A$ 和 $N$ 的模函数以及函数 $\alpha \colon A \to \R_{>0}$.
首先, 由于 $A$ 交换, 有 $\Delta_A = 1$.
对 $N$, 有几种做法. 例如, 你可以注意到 $N$ (仅作为拓扑空间) 显然与 $M_n(\R)$ 中对角线为零的上三角矩阵构成的 $\R$-向量空间 $W$ 同胚. (只需忘掉矩阵的对角线项.) 而且, 对 $n \in N$, $N$ 上按 $n$ 的左平移对应于 $W$ 的线性自同态 $U_n$, 由 $U_n(X) = nX$ ($X \in W$) 给出. 注意 $W$ 是前一问中空间 $V$ 的子空间, 且 $U_n$ 是 $T_n$ 的限制. 所以可以用与前一问相同的方法计算 $\det(U_n)$, 得到 $\det(U_n) = 1$.
所以 $W$ 上的 Lebesgue 测度是 $N$ 上的左 Haar 测度. 我们可以用右平移代替左平移重做一切, 得到 $W$ 上的 Lebesgue 测度也是 $N$ 上的右 Haar 测度. 这意味着 $N$ 幺模, 所以 $\Delta_N = 1$.
最后, 我们需要计算函数 $\alpha$. 回忆它由
\[
\int_N f(ana^{-1})\,\dd n = \alpha(a)\int_N f(n)\,\dd n
\]
定义 (对每个 $f \in \mathcal{C}_c(N)$), 其中 $dn$ 是 $W$ 上的 Lebesgue 测度 (我们刚看到它是 $N$ 上的 Haar 测度). 注意 $c_a \colon X \longmapsto aXa^{-1}$ 是 $W$ 的线性自同态, 所以一旦知道 $\det(c_a)$ 就能用换元公式计算 $\int_N f(ana^{-1})\,\dd n$. 我们得到 $\det(c_a)\int_N f(ana^{-1})\,\dd n = \int_N f(n)\,\dd n$, 所以 $\alpha(a) = \det(c_a)^{-1}$. 但容易看出, 若 $a$ 的对角线元素为 $(a_1, \ldots, a_n)$, 则
\[
\det(c_a) = a_1^{n-1}a_2^{n-3}\cdots a_n^{1-n} = \prod_{i=1}^n a_i^{n-2i+1}.
\]
所以最终对 $p \in P$,
\[
\Delta_P(p) = a_1^{1-n}a_2^{3-n}\cdots a_n^{n-1} = \prod_{i=1}^n a_i^{2i-n-1},
\]
其中 $a_1, \ldots, a_n$ 是 $p$ 的对角线元素.

(d). 若你知道 (或知道如何证明) $\SL_n(\R)$ 等于它的换位子群, 这就很容易. 这里给出另一种方法: 设 $\GL_n(\R)^+$ 是行列式为正的 $n \times n$ 矩阵群. 这是 $\GL_n(\R)$ 的开子群 (它是连续群同态 $\det \colon \GL_n(\R) \to \R^{\times}$ 下 $\R_{>0}$ 的原像), 所以若 $\mu$ 是 $\GL_n(\R)$ 上的 Haar 测度 (回忆由习题 \ref{ex:I.5.3.2}(c), $\GL_n(\R)$ 幺模), 则它到 $\GL_n(\R)^+$ 的限制是非零正则 Borel 测度, 且显然是 $\GL_n(\R)^+$ 上的左、右 Haar 测度. 现在注意有拓扑群的同构 $\R_{>0} \times \SL_n(\R) \to \GL_n(\R)^+$, $(\lambda, x) \longmapsto \lambda x$ (其逆由 $x \longmapsto (\det(x)^{1/n}, \det(x)^{-1/n}x)$ 给出), 所以可以应用习题 \ref{ex:I.5.3.6}, 取 $G = \GL_n(\R)^+$, $A = \R_{>0}I_n$, $N = \SL_n(\R)$. 由于 $A$ 和 $N$ 交换, 有 $\alpha = 1$. 我们知道 $A$ 幺模 (因为它交换), 且刚才已见 $\GL_n(\R)^+$ 幺模, 所以由习题 \ref{ex:I.5.3.6}(c), $\SL_n(\R)$ 也幺模.

(e). 众所周知 $\R^2$ 中的旋转群 (即 $\SO(2)$) 与圆群 $S^1$ 同构. 同构把 $e^{2i\pi\theta} \in S^1$ 送到矩阵 $\begin{pmatrix} \cos\theta & \sin\theta \\ -\sin\theta & \cos\theta \end{pmatrix}$. 又, 我们在课堂上见过可以通过把 $f \in \mathcal{C}_c(S^1)$ 送到 $\int_0^1 f(e^{2i\pi\theta})\,\dd\theta$ 的线性泛函来定义 $S^1$ 上的 Haar 测度, 其中 $d\theta$ 是 $\R$ 上的 Lebesgue 测度.
这当然的重点是, 习题 \ref{ex:I.5.3.5} 现在允许你定义 $\SL_2(\R)$ 上的 Haar 测度. 为处理 $\SL_n(\R)$ 的情形, 我们需要 $\SO(n)$ 上的 Haar 测度. 这种测度的一个例子在习题 \ref{ex:I.5.3.9} 中给出.
\end{solution}

\begin{exercise}[$\GL_n(\Q_p)$ 上的 Haar 测度]
\label{ex:I.5.3.8}
(回忆习题 \ref{ex:I.5.1.4}, \ref{ex:I.5.3.1}, \ref{ex:I.5.3.2} 和 \ref{ex:I.5.3.4}.) 记 $dx$ 为加法群 $\Q_p$ 上的 Haar 测度. 又记 $dx$ (相应地 $dA$) 为 $\Q_p^n$ (相应地 $M_n(\Q_p) \simeq \Q_p^{n^2}$) 上的乘积测度; 注意它是相应加法群的 Haar 测度.
\begin{enumerate}
\item[(a)] 证明对每个 $f \in L^1(\Q_p)$ 和每个 $a \in \Q_p^{\times}$, $b \in \Q_p$,
\[
\int_{\Q_p}f(x)\,\dd x = |a|_p\int_{\Q_p}f(ax + b)\,\dd x.
\]
\item[(b)] 设 $n \geq 1$. 证明若 $f \in L^1(\Q_p^n)$, $A \in \GL_n(\Q_p)$, $b \in \Q_p^n$, 则
\[
\int_{\Q_p^n}f(x)\,\dd x = |\det(A)|_p\int_{\Q_p^n}f(Ax + b)\,\dd x.
\]
\item[(c)] 证明 $|\det(A)|_p^{-n}dA$ 是 $\GL_n(\Q_p)$ 上的左、右 Haar 测度.
\item[(d)] 设 $B$ 是 $\GL_n(\Q_p)$ 中上三角矩阵群. 找出 $B$ 上的左 Haar 测度并计算 $B$ 的模函数.
\end{enumerate}
\end{exercise}

\begin{solution}
(a). 首先, 利用 $dx$ 的平移不变性, 看到对每个 $f \in L^1(\Q_p)$ 和 $a, b \in \Q_p$,
\[
\int_{\Q_p}f(ax + b)\,\dd x = \int_{\Q_p}f(ax)\,\dd x.
\]
设 $a \in \Q_p^{\times}$. 我们使用习题 \ref{ex:I.5.3.4} 的记号. 若 $x \in \Q_p$, $m \in \Z$, 则
\[
aB(x, p^m) = \{ay \mid |x - y|_p \leq p^m\} = \{y \in \Q_p \mid |ax - y|_p \leq |a|_pp^m\} = B(ax, |a|_pp^m),
\]
所以由 \ref{ex:I.5.3.4}(a), $\operatorname{vol}(aB(x, p^m)) = |a|_p\operatorname{vol}(B(x, p^m))$. 使用同一题 (b), 对 $\Q_p$ 的每个 Borel 子集 $E$ 得到 $\operatorname{vol}(aE) = |a|_p\operatorname{vol}(E)$. 假设 $f = \mathbb{1}_E$, 其中 $E$ 是 $\Q_p$ 的 Borel 子集. 则
\[
\int_{\Q_p}f(ax)\,\dd x = \operatorname{vol}(a^{-1}E) = |a|_p^{-1}\int_{\Q_p}f(x)\,\dd x,
\]
所以对这个函数 $f$ 得到所要的结果. 结论现在对每个 $f \in L^1(\Q_p)$ 由积分的线性和连续性得出.

(b). 像 (a) 那样使用 $dx$ 的平移不变性, 看到只需在 $b = 0$ 的情形证明结论. 设 $A \in \GL_n(\Q_p)$. 首先注意 $A = A_1A_2$, 且若我们对 $A_1$ 和 $A_2$ 知道结论, 就对 $A$ 知道结论; 事实上, 对每个 $f \in L^1(\Q_p^n)$,
\begin{multline*}
\int_{\Q_p^n}f(x)\,\dd x = |\det(A_1)|_p\int_{\Q_p^n}f(A_1x)\,\dd x \\
= |\det(A_1)|_p|\det(A_2)|_p\int_{\Q_p^n}f(A_1(A_2x))\,\dd x = |\det(A)|_p\int_{\Q_p^n}f(Ax)\,\dd x.
\end{multline*}
Gauss 算法 (解线性方程组) 说可以通过初等行变换 (对应左乘一个下三角矩阵) 和行置换 (对应左乘一个置换矩阵) 把 $A$ 变成上三角. 所以由上面的观察, 只需对上三角矩阵、下三角矩阵和置换矩阵证明结论.
先假设 $A$ 是置换矩阵. 则存在置换 $\sigma \in S_n$ 使得对每个 $x = (x_1 \ldots, x_n) \in \Q_p^n$, $Ax = (x_{\sigma(1)}, \ldots, x_{\sigma(n)})$. 由于 $dx$ 是 $\Q_p^n$ 的 $n$ 个因子 $\Q_p$ 上的相同测度的乘积, 对每个 $f \in L^1(\Q_p)$ 有 $\int_{\Q_p^n}f(Ax)\,\dd x = \int_{\Q_p^n}f(x)\,\dd x$. 结论由 $\det(A) = \pm 1$ 得出.
假设 $A$ 是上三角的, 记 $A = (a_{ij})_{1 \leq i,j \leq n}$. 设 $f \in L^1(\Q_p)$. 则
\begin{multline*}
\int_{\Q_p^n}f(A(x_1, \ldots, x_n)) = \int_{\Q_p}\cdots\int_{\Q_p}f(a_{11}x_1 + \cdots + a_{1n}x_n, \ldots, \\
a_{n-1,n-1}x_{n-1} + a_{n-1,n}x_n, a_{nn}x_n)\,\dd x_n\,\dd x_{n-1}\cdots\dd x_1.
\end{multline*}
使用 (a), 看到最后的积分等于
\[
|a_{11}|_p^{-1}\cdots|a_{n-1,n-1}|_p^{-1}|a_{nn}|_p^{-1}\int_{\Q_p^n}f(x)\,\dd x = |\det(A)|_p^{-1}\int_{\Q_p^n}f(x)\,\dd x.
\]
下三角矩阵的情形类似 (只需把 $dx_i$ 反序).

(c). 一旦有了 (b) 的换元公式, 就可以在习题 \ref{ex:I.5.3.1} 和 \ref{ex:I.5.3.2} 中用 $\Q_p$ 代替 $\R$, 所有结果用完全相同的证明仍然成立. (习题 \ref{ex:I.5.3.2}(b) 除外, 它对 $\Q_p$ 没有意义.) 特别地, 我们得到 $|\det(A)|_p^{-n}dA$ 是 $\GL_n(\Q_p)$ 上的左、右 Haar 测度.

(d). 同样地, 我们只需应用习题 \ref{ex:I.5.3.7} 的 (b) 和 (c) 的证明 (以及习题 \ref{ex:I.5.3.1} 的 $\Q_p$ 版本) 就能得到结果. 假设习题 \ref{ex:I.5.3.7} 中没有符号错误, 则 $B$ 上的左 Haar 测度是 $\prod_{i=1}^n|a_{ii}|_p^{i-n-1}dA$, 其中 $dA$ 是上三角矩阵构成的 $\Q_p$-向量空间 $V \simeq \Q_p^{n(n+1)/2}$ 上的乘积测度, $a_{ij}$ 是矩阵的元素. 而 $B$ 的模函数由
\[
\Delta(A) = \prod_{i=1}^n|a_{ii}|_p^{2i-n-1}
\]
给出.
\end{solution}

\begin{exercise}[$\SO(n)$ 上的 Haar 测度]
\label{ex:I.5.3.9}
本题的目标是给出 $\SO(n)$ 上 Haar 测度的公式. (对酉群 $\UU(n)$ 可以做类似的事.)
\begin{enumerate}
\item[(a)] 对 $X \in M_n(\R)$, 令 $\Phi(X) = (I_n - X)(I_n + X)^{-1}$. 证明当 $-1$ 不是 $X$ 的特征值时这有意义, 且只要有意义就有 $\Phi(\Phi(X)) = X$.
\item[(b)] 记 $A_n$ 为 $n \times n$ 反对称矩阵 (即满足 $X^T = -X$ 的 $X \in M_n(\R)$) 构成的 $\R$-向量空间, $U$ 为 $\SO(n)$ 中不以 $-1$ 为特征值的元素的集合. 证明 $U$ 是 $\SO(n)$ 的开稠密子集, 且 $\Phi$ 诱导同胚 $A_n \xrightarrow{\sim} U$.
\item[(c)] 设 $X \in A_n$. 证明存在 $A_n$ 的开稠密子集 $V$ 和 $W$ 使得公式 $\Phi(L_XY) = \Phi(X)\Phi(Y)$ 定义微分同胚 $L_X \colon V \xrightarrow{\sim} W$, 且 $0 \in V$.
\item[(d)] 设 $dX$ 是 $A_n$ 上的 Lebesgue 测度. 对每个 $X \in A_n$ 和每个 $L_X$ 在其上有定义的 $Y \in A_n$, 记 $L'_X(Y)$ 为 $L_X$ 在 $Y$ 处的微分. 它是从 $A_n$ 到 $A_n$ 的线性变换, 使得对每个 $H \in A_n$,
\[
L_X(Y + tH) = L_X(Y) + tL'_X(Y)(H) + o(t).
\]
固定 $X \in A_n$. 我们要计算 $\det(L'_X(0))$. 回忆 $L'_X(0)$ 是 $A_n$ 的线性自同态, 且注意 $A_n \otimes_{\R}\CC$ 是 $M_n(\CC)$ 中反对称矩阵的空间.
\begin{enumerate}
\item[(i)] 证明 $\det(L'_X(0))$ 良定义且非零.
\item[(ii)] 证明对每个 $H \in A_n$, $L'_X(0)(H) = (I_n - X)H(I_n + X)$.
\item[(iii)] 证明 $X$ 有 (复) 特征向量基 $(v_1, \ldots, v_n)$, 使得对应特征值具有 $i\lambda_1, \ldots, i\lambda_n$ 的形式, 其中 $\lambda_1, \ldots, \lambda_n \in \R$.
\item[(iv)] 对 $j, k \in \{1, \ldots, n\}$, 令 $Y_{jk} = v_jv_k^T - v_kv_j^T$. 证明 $Y_{jk} \in A_n \otimes_{\R}\CC$, 且它是 $L'_X(0)$ 的特征向量, 对应特征值为 $(1 - i\lambda_j)(1 - i\lambda_k)$.
\item[(v)] 证明 $(Y_{jk})_{1 \leq j < k \leq n}$ 是 $A_n \otimes_{\R}\CC$ 的基.
\item[(vi)] 证明 $\det(L'_X(0)) = \det(I_n - iX)^{n-1}$.
\end{enumerate}
\item[(e)] 证明把 $f \in \mathcal{C}_c(\SO(n))$ 送到
\[
\int_{A_n}f(\Phi(X))\frac{1}{|\det L'_X(0)|}\,\dd X
\]
的线性泛函定义 $\SO(n)$ 上的左 Haar 测度. (提示: 注意 $(L_X \circ L_Y)(0) = L_X(Y)$, 并使用链式法则.)
\end{enumerate}
\end{exercise}

\begin{solution}
(a). 若 $X \in M_n(\R)$, 则 $-1$ 不是 $X$ 的特征值当且仅当 $I_n + X$ 可逆, 即当且仅当定义 $\Phi(X)$ 的公式有意义. 所以 $\Phi$ 的定义域是由方程 $\det(I_n + X) \neq 0$ 定义的开集. 又注意 $I_n - X$ 与 $I_n + X$ 交换, 所以 $I_n - X$ 与 $(I_n + X)^{-1}$ 交换 (若后者有定义), 因此也有 $\Phi(X) = (I_n + X)^{-1}(I_n - X)$.
设 $X \in M_n(\R)$ 使得 $\Phi(X)$ 有定义. 则
\[
I_n + \Phi(X) = ((I_n + X) + (I_n - X))(I_n + X)^{-1} = 2(I_n + X)^{-1}
\]
且
\[
I_n - \Phi(X) = ((I_n + X) - (I_n - X))(I_n + X)^{-1} = 2X(I_n + X)^{-1}.
\]
特别地, $I_n + \Phi(X)$ 可逆, 所以 $\Phi(\Phi(X))$ 有意义, 且
\[
\Phi(\Phi(X)) = (I_n - \Phi(X))(I_n + \Phi(X))^{-1} = 2X(I_n + X)^{-1}(2(I_n + X)^{-1})^{-1} = X.
\]

(b). 设 $g \in \SO(n)$. 则可以找到 $P \in \GL_n(\R)$ 使得
\[
PgP^{-1} = \begin{pmatrix} r_1 & 0 & \cdots & 0 \\ 0 & r_2 & 0 & 0 \\ \vdots & 0 & \ddots & 0 \\ 0 & \cdots & 0 & r_m \end{pmatrix},
\]
其中:
\begin{itemize}
\item 若 $n$ 偶, 则 $m = n/2$ 且 $r_1, \ldots, r_m$ 是形如 $\begin{pmatrix} \cos\theta_i & \sin\theta_i \\ -\sin\theta_i & \cos\theta_i \end{pmatrix}$ ($\theta_i \in [0, 2\pi)$) 的 $2 \times 2$ 矩阵;
\item 若 $n$ 奇, 则 $m = (n+1)/2$, 矩阵 $r_m$ 是 $1 \times 1$ 矩阵 $1$, 且 $r_1, \ldots, r_{m-1}$ 是形如 $\begin{pmatrix} \cos\theta_i & \sin\theta_i \\ -\sin\theta_i & \cos\theta_i \end{pmatrix}$ ($\theta_i \in [0, 2\pi)$) 的 $2 \times 2$ 矩阵.
\end{itemize}
两种情形下, $-1$ 是 $g$ 的特征值当且仅当至少有一个 $\theta_i$ 等于 $\pi$. 所以, 通过变化 $\theta_i$, 可以找到不以 $-1$ 为特征值的 $\SO(n)$ 的元素序列收敛到 $g$. 这证明 $U$ 在 $\SO(n)$ 中稠密.
其次, 由于反对称矩阵只有纯虚特征值, 函数 $\Phi$ 在 $A_n$ 上有定义. 又注意在 $\Phi$ 的定义上显然 $\Phi$ 在其定义开集上连续. 由 (a) 的第二部分, $\Phi$ 单射, 且为证明 $\Phi$ 是从 $A_n$ 到 $U$ 的同胚, 只需证明它是从 $A_n$ 到 $U$ 的双射 (因为此时它的逆将是 $\Phi$). 所以我们只需证明 $\Phi(A_n) = U$. 再次利用只要有意义就有 $\Phi(\Phi(X)) = X$, 看到只需证明 $\Phi(A_n) \subset U$ 和 $\Phi(U) \subset A_n$.
设 $X \in A_n$. 则 $X^T = -X$, 所以 $\Phi(X)^T = (I_n + X^T)^{-1}(I_n - X^T) = (I_n - X)^{-1}(I_n + X)$, 从而 $\Phi(X)^T\Phi(X) = I_n$, 这意味着 $\Phi(X) \in \Or(n)$. 由于 $\Phi$ 连续且 $A_n$ 连通, $\Phi(A_n)$ 连通. 但 $I_n = \Phi(0) \in \Phi(A_n)$, 所以 $\Phi(A_n)$ 包含于 $\SO(n)$.
设 $X \in \SO(n)$ 不以 $-1$ 为特征值. 则 $X^T = X^{-1}$, 所以
\[
\Phi(X)^T = (I_n - X^T)(I_n + X^T)^{-1} = (I_n - X^{-1})(I_n + X^{-1})^{-1} = (X - I_n)(X + I_n)^{-1} = -\Phi(X).
\]
所以我们有 $\Phi(U) \subset A_n$.

(c). 固定 $X \in A_n$. 注意公式 $\Phi(L_XY) = \Phi(X)\Phi(Y)$ 也可以写成 $L_XY = \Phi(\Phi(X)\Phi(Y))$ (由 (a)).
对 $Y \in A_n$, $\Phi(X)\Phi(Y)$ 在 $\Phi$ 下有像 (由 (b) 将自动在 $A_n$ 中) 当且仅当 $\Phi(Y) \in \Phi(X)^{-1}U$. 所以可以取 $V = \Phi(U \cap (\Phi(X)^{-1}U))$; 由于由 (b) $U \cap (\Phi(X)^{-1}U)$ 在 $\SO(n)$ 中稠密, 它在 $A_n$ 中稠密. 则 $V$ 在映射 $L_X \colon Y \longmapsto \Phi(\Phi(X)\Phi(Y))$ 下的像是 $W := \Phi((\Phi(X)U) \cap U)$.
映射 $L_X \colon V \to W$ 连续且满射. 事实上, 由于 $\Phi$ 无穷次可微 (由公式``对每个 $A \in \GL_n(\R)$, $A^{-1}$ 等于 $\det(A)^{-1}$ 乘以其余因子矩阵的转置'', 它由参数的元素的合理函数给出), 映射 $L_X$ 也无穷次可微.
设 $X' = \Phi(\Phi(X)^{-1}) \in A_n$. 则如上得到连续满射 $L_{X'} \colon W \to V$, 由公式 $L_{X'}(Y) = \Phi(\Phi(X)^{-1}\Phi(Y))$ 定义. 映射 $L_X$ 与 $L_{X'}$ 互为逆, 特别地它们都是微分同胚.
最后, 若 $Y = 0$, 则 $\Phi(Y) = I_n$. 所以 $\Phi(Y) \in U$, 且也有 $\Phi(Y) \in \Phi(X)^{-1}U$, 因为 $\Phi(X)\Phi(Y) = \Phi(X) \in U$. 这证明 $0 \in V$.

(d). (i) 设 $X \in A_n$. 由于 $L_X$ 在点 $0$ 处有定义, 微分 $L'_X(0)$ 有意义; 又, 由于 $L_X$ 是微分同胚, $\det(L'_X(0)) \neq 0$.

(ii) 注意对 $Y \in A_n$,
\[
(I_n + X + Y)^{-1} = (I_n + X)^{-1}(I_n + (I_n + X)^{-1}Y) = (I_n + X)^{-1}(I_n - Y(I_n + X)^{-1} + o(Y)),
\]
所以
\begin{align*}
\Phi(X + Y) &= (I_n - X - Y)(I_n + X + Y)^{-1} \\
&= ((I_n - X) - Y)(I_n + X)^{-1}(I_n - Y(I_n + X)^{-1} + o(Y)) \\
&= \Phi(X) - \Phi(X)Y(I_n + X)^{-1} - Y(I_n + X)^{-1} + o(Y).
\end{align*}
特别地 (取 $X = 0$),
\[
\Phi(Y) = I_n - 2Y + o(Y).
\]
所以
\[
\Phi(X)\Phi(Y) = \Phi(X) - 2\Phi(X)Y + o(Y),
\]
且
\begin{align*}
L_X(Y) &= \Phi(\Phi(X)\Phi(Y)) = \Phi(\Phi(X) - 2\Phi(X)Y + o(Y)) \\
&= \Phi(\Phi(X)) - \Phi(\Phi(X))(-2\Phi(X)Y)(I_n + \Phi(X))^{-1} - (-2\Phi(X)Y)(I_n + \Phi(X))^{-1} + o(Y).
\end{align*}
利用 $\Phi(\Phi(X)) = X$ 和 $I_n + \Phi(X) = 2(I_n + X)^{-1}$ (见 (a)), 可以把最后这个表达式化简为
\begin{multline*}
X + X\Phi(X)Y(I_n + X)^{-1} + \Phi(X)Y(I_n + X)^{-1} + o(Y) \\
= X + (I_n + X)\Phi(X)Y(I_n + X)^{-1} + o(Y) = X + (I_n - X)Y(I_n + X) + o(Y).
\end{multline*}
但然后 $L'_X(0)(Y) = (I_n - X)Y(I_n + X)$ 的结论由微分的定义立即得出.

(iii) 由于 $X$ 反对称且元素为实数, 它正规, 所以谱定理说 $X$ 在 $\CC^n$ 的标准正交基下可对角化; 换句话说, 存在酉矩阵 $P$ 使得 $PXP^{-1}$ 对角. 我们已经用过 $X$ 的特征值是纯虚的的事实, 但容易快速重验: $X^* = -X$, $P^* = P^{-1}$, 且 $(PXP^{-1})^* = (P^*)^{-1}X^*P^* = -PXP^{-1}$. 由于 $PXP^{-1}$ 对角, 这意味着它的对角线元素 (即 $X$ 的特征值) 全是纯虚的.

(iv) 由 $Y_{jk}$ 的定义直接得出 $Y_{jk}^T = -Y_{jk}$, 所以 $Y_{jk} \in A_n \otimes_{\R}\CC$. 而且, 由 (ii),
\begin{align*}
L'_X(0)(Y_{ij}) &= (I_n - X)Y_{ij}(I_n + X) \\
&= (I_n - X)(v_jv_k^T)(I_n - X^T) - (I_n - X)(v_kv_j^T)(I_n - X^T) \\
&= (1 - i\lambda_j)(v_jv_k^T)(1 - i\lambda_k) - (1 - i\lambda_k)(v_kv_j^T)(1 - i\lambda_j) \\
&= (1 - i\lambda_j)(1 - i\lambda_k)Y_{ij}.
\end{align*}

(v) 由于 $(v_1, \ldots, v_n)$ 是 $\CC^n$ 的基, 矩阵 $v_jv_k^T$ ($1 \leq j, k \leq n$) 构成 $M_n(\CC)$ 的基. 所以矩阵 $Y_{jk} = (v_jv_k^T) - (v_jv_k^T)^T$ ($1 \leq j, k \leq n$) 生成 $A_n \otimes_{\R}\CC$. 注意 $Y_{jj} = 0$ 且 $Y_{kj} = -Y_{jk}$, 所以 $A_n \otimes_{\R}\CC$ 实际上由矩阵 $Y_{jk}$ ($1 \leq j < k \leq n$) 张成. 由于这样的矩阵有 $n(n-1)/2$ 个且 $\dim_{\CC}(A_n \otimes_{\R}\CC) = \dim_{\R}(A_n) = n(n-1)/2$, 它们构成 $A_n \otimes_{\R}\CC$ 的基.

(vi) 由 (iv) 和 (v),
\[
\det(L'_X(0)) = \prod_{1 \leq j < k \leq n}(1 - i\lambda_j)(1 - i\lambda_k) = \prod_{r=1}^n(1 - i\lambda_r)^{n-1}
\]
(因为每个 $1 - i\lambda_r$ 在第一个大乘积中出现 $n-1$ 次: 作为第一个因子 $(1 - i\lambda_j)$ 出现 $n - r$ 次, 作为第二个因子 $(1 - i\lambda_k)$ 出现 $r - 1$ 次). 为得到结果, 只需注意 $I_n - iX$ 的特征值是 $1 - i\lambda_1, \ldots, 1 - i\lambda_n$, 所以
\[
\det(I_n - X) = \prod_{r=1}^n(1 - i\lambda_r).
\]

(e). 记此泛函为 $\Lambda$. 首先, 由 (d), 函数 $X \longmapsto |\det(L'_X(0))|^{-1}$ 在 $A_n$ 上处处有定义且连续, 所以定义 $\Lambda$ 的积分有意义.
我们需要验证 $\Lambda$ 正且被左平移不变. 首先验证正性. 设 $f \in \mathcal{C}_c^+(\SO(n))$. 则可以找到 $\varepsilon > 0$ 和 $\SO(n)$ 的非空开子集 $\Omega$ 使得 $f|_{\Omega} \geq \varepsilon$. 由于 $U$ 在 $\SO(n)$ 中开稠, 它与 $\Omega$ 的交开且非空, 所以 $\Phi(U \cap \Omega)$ 在 $A_n$ 中开且非空, 且
\[
\Lambda(f) \geq \varepsilon\int_{\Phi(U \cap \Omega)}\frac{1}{|\det(L'_X(0))|}\,\dd X > 0
\]
(因为函数 $X \longmapsto |\det(L'_X(0))|^{-1}$ 在 $\Phi(U \cap \Omega)$ 上连续且正).
现在验证左不变性. 固定 $f \in \mathcal{C}_c(\SO(n))$. 设 $g \in U$. 则 $\Lambda(L_gf) = \int_{A_n}f(g^{-1}\Phi(Y))\frac{1}{|\det(L'_Y(0))|}\,\dd Y$. 选取 $X, X' \in A_n$ 使得 $\Phi(X) = g^{-1}$, $\Phi(X') = g$. 则
\begin{align*}
\Lambda(L_gf) &= \int_{A_n}f(\Phi(X)\Phi(Y))\frac{1}{|\det(L'_Y(0))|}\,\dd Y \\
&= \int_V f(\Phi(X)\Phi(Y))\frac{1}{|\det(L'_Y(0))|}\,\dd Y \quad \text{(因为 } \operatorname{vol}(A_n - V) = 0\text{)} \\
&= \int_V f(\Phi(L_XY))\frac{1}{|\det(L'_Y(0))|}\,\dd Y.
\end{align*}
现在注意, 若 $Y \in V$, 则 $L_Y(0) = Y$ 也在 $V$ 中, 所以 $L_X(Y) = L_X \circ L_Y(0) = L_{L_XY}(0)$ 有意义, 且由链式法则
\[
L'_{L_XY}(0) = L'_X(Y) \circ L'_Y(0),
\]
从而特别地
\[
\frac{1}{|\det(L'_Y(0))|} = \frac{|\det(L'_X(Y))|}{|\det(L'_{L_XY}(0))|}.
\]
这蕴含
\[
\Lambda(L_gf) = \int_V f(\Phi(L_XY))\frac{|\det(L'_X(Y))|}{|\det(L'_{L_XY}(0))|}\,\dd Y.
\]
使用代换 $Z = L_XY$, 看到这等于
\[
\int_W f(\Phi(Z))\frac{1}{|\det(L'_Z(0))|}\,\dd Z.
\]
由于 $\operatorname{vol}(A_n - W) = 0$, 最后的积分等于 $\int_{A_n}f(\Phi(Z))\frac{1}{|\det(L'_Z(0))|}\,\dd Z$, 即 $\Lambda(f)$.
所以我们已证明函数 $\SO(n) \to \CC$, $g \longmapsto \Lambda(L_gf)$ 在开稠子集 $U$ 上为常数. 由于此函数连续 (它是连续函数 $\SO(n) \to \mathcal{C}_c(\SO(n))$, $g \longmapsto L_gf$ 与连续线性泛函 $\Lambda \colon \mathcal{C}_c(\SO(n)) \to \CC$ 的复合), 它在整个 $\SO(n)$ 上为常数, 这意味着对每个 $g \in \SO(n)$ 有 $\Lambda(L_gf) = \Lambda(f)$.
\end{solution}

\begin{exercise}[$\SU(2)$ 上的 Haar 测度]
\label{ex:I.5.3.10}
设 $G = \SU(2)$.
\begin{enumerate}
\item[(a)] 证明 $G$ 的每个元素都具有 $\begin{pmatrix} a & -\bar{b} \\ b & \bar{a} \end{pmatrix}$ 的形式, 其中 $a, b \in \CC$ 且 $|a|^2 + |b|^2 = 1$.
若按通常方式把 $\CC$ 等同于 $\R^2$, 前一问给出 $\SU(2)$ 与 $S^3$ ($\R^4$ 中的单位球面) 之间的同胚 $\alpha$.
\item[(a)] 若 $g \in \SU(2)$, 证明 $\SU(2)$ 上按 $g$ 的左平移通过 $\alpha$ 对应于 $\SO(4)$ 的元素在 $\R^4$ 上的作用到 $S^3$ 的限制 (即存在 $A \in \SO(4)$ 使得对每个 $h \in \SU(2)$, $gh = A\alpha(h)$).
\item[(b)] 设 $\mu$ 是 $S^3$ 上的通常球面测度; 即若 $\lambda$ 是 $\R^4$ 上的 Lebesgue 测度, 则由定义对 $S^3$ 的每个 Borel 子集 $E$,
\[
\mu(E) = \frac{2}{\pi^2}\lambda(\{tx, t \in [0,1], x \in E\})
\]
(注意 $\R^4$ 中单位球的体积是 $\frac{\pi^2}{2}$). 证明 $\mu$ 在 $\alpha$ 下的原像是 $\SU(2)$ 上的左、右 Haar 测度.
\item[(c)] 我们在 $S^3$ 上使用下列 (超球面) 坐标: 若 $(x_1, x_2, x_3, x_4) \in S^3$, 写
\[
\begin{cases} x_1 = \cos\theta, \\ x_2 = \sin\theta\cos\psi, \\ x_3 = \sin\theta\sin\psi\cos\phi, \\ x_4 = \sin\theta\sin\psi\sin\phi, \end{cases}
\]
其中 $0 \leq \theta \leq \pi$, $0 \leq \psi \leq \pi$, $0 \leq \phi \leq 2\pi$. 证明对每个 $f \in \mathcal{C}_c(S^3)$,
\[
\int_{S^3}f\,\dd\mu = \frac{1}{2\pi^2}\int_0^{\pi}\int_0^{\pi}\int_0^{2\pi}f(\cos\theta, \sin\theta\cos\psi, \sin\theta\sin\psi\cos\phi, \sin\theta\sin\psi\sin\phi)\sin^2\theta\sin\psi\,\dd\theta\,\dd\psi\,\dd\phi.
\]
(你可以自由使用计算机计算任何大行列式.)
\end{enumerate}
\end{exercise}

\begin{solution}
(a). 显然题设形式的每个矩阵都在 $\SU(2)$ 中. 让我们证明其逆. 设 $A = \begin{pmatrix} a & c \\ b & d \end{pmatrix} \in M_2(\CC)$. 则 $A \in \UU(2)$ 当且仅当 $A^*A = I_2$, 这意味着 $A$ 的两个列向量对 $\CC^2$ 上的通常 Hermitian 内积正交且范数为 $1$. 由于 $\CC^2$ 中一条线的正交补是 $1$ 维的, 这蕴含存在 $\lambda \in \CC^{\times}$ 使得 $\begin{pmatrix} c \\ d \end{pmatrix} = \lambda\begin{pmatrix} -\bar{b} \\ \bar{a} \end{pmatrix}$. 列的范数条件给出 $\bar{a}a + \bar{b}b = \lambda\bar{\lambda}(\bar{a}a + \bar{b}b) = 1$, 而 $\det(A) = 1$ 的条件给出 $\lambda(\bar{a}a + \bar{b}b) = 1$. 所以我们得到 $\lambda = 1$, 如所愿.

(b). 设 $V$ 是形如 $\begin{pmatrix} a & -\bar{b} \\ b & \bar{a} \end{pmatrix}$ ($a, b \in \CC$) 的矩阵的空间. 则 $\alpha$ 延拓成从 $V$ 到 $\CC^2$ 的 $\CC$-线性同构, 从而到从 $V$ 到 $\R^4$ 的 $\R$-线性同构, 把 $\begin{pmatrix} a & -\bar{b} \\ b & \bar{a} \end{pmatrix}$ 送到 $(\Re(a), \Im(a), \Re(b), \Im(b))$. 若 $A \in \SU(2)$, 则 $A$ 在 $V$ 上的左乘作用就是 $A$ 在 $\CC^2$ 上的通常作用, 所以它对应于 $\R^4$ 的保持通常 Euclid 范数的线性自同构, 即属于 $\Or(4)$. 又, 此作用的行列式就是 $\det(A) = 1$, 所以 $\R^4$ 的对应自同构在 $\SO(4)$ 中.

(c). 首先注意 $\mu$ 是 $S^3$ 上的正则 Borel 测度 ($S^3$ 的子集 $E$ 是 Borel 子集当且仅当 $\{tx, t \in [0,1], x \in E\}$ 是 $\R^4$ 的 Borel 子集, 它紧当且仅当 $\{tx, t \in [0,1], x \in E\}$ 紧, 开当且仅当 $\{tx, t \in (0,1], x \in E\}$ (与 $\{tx, t \in [0,1], x \in E\}$ 有相同测度) 开).
由 $\R^4$ 中的换元公式, 测度 $\mu$ 被 $\SO(4)$ 在 $S^3$ 上的作用不变. 由 (b), 它在 $\alpha$ 下的原像被 $\SU(2)$ 上的左平移不变, 所以是左 Haar 测度. 但群 $\SU(2)$ 紧, 所以每个左 Haar 测度也是右 Haar 测度.

(d). 设 $B^4$ 是 $\R^4$ 中的闭单位球. 设 $f \in \mathcal{C}_c(S^3)$. 我们定义函数 $g \in L^1(B^4)$:
\begin{multline*}
g(r\cos\theta, r\sin\theta\cos\psi, r\sin\theta\sin\psi\cos\phi, r\sin\theta\sin\psi\sin\phi) \\
= f(\cos\theta, \sin\theta\cos\psi, \sin\theta\sin\psi\cos\phi, \sin\theta\sin\psi\sin\phi)
\end{multline*}
对 $0 \leq r \leq 1$. (注意: $g$ 在 $0$ 处可能没有良好定义, 但这无所谓, 因为 $\{0\}$ 体积为 $0$.) 则由 $\mu$ 的定义, $\int_{S^3}f\,\dd\mu = \frac{2}{\pi^2}\int_{B^4}g\,\dd\lambda$. 我们可以用换元公式计算这最后一个积分 (并避开此换元不双射的集合, 反正它的体积为 $0$). 若 $\beta$ 是把 $(r, \theta, \phi, \psi)$ 送到 $(r\cos\theta, r\sin\theta\cos\psi, r\sin\theta\sin\psi\cos\phi, r\sin\theta\sin\psi\sin\phi)$ 的映射, 则
\[
|D\beta(r, \theta, \phi, \psi)| = r^3(\sin\theta)^2\sin\psi,
\]
所以 $\int_{B^4}g\,\dd\lambda$ 等于
\begin{multline*}
\int_0^1\int_0^{\pi}\int_0^{\pi}\int_0^{2\pi}f(\cos\theta, \sin\theta\cos\psi, \sin\theta\sin\psi\cos\phi, \sin\theta\sin\psi\sin\phi)r^3\sin^2\theta\sin\psi\,\dd r\,\dd\theta\,\dd\psi\,\dd\phi \\
= \frac{1}{4}\int_0^{\pi}\int_0^{\pi}\int_0^{2\pi}f(\cos\theta, \sin\theta\cos\psi, \sin\theta\sin\psi\cos\phi, \sin\theta\sin\psi\sin\phi)\sin^2\theta\sin\psi\,\dd\theta\,\dd\psi\,\dd\phi.
\end{multline*}
乘以 $\frac{2}{\pi^2}$ 就得到结果.
\end{solution}

\subsection{局部紧交换群的对偶}

\begin{exercise}[对偶群 (dual group)]
\label{ex:I.5.4.1}
设 $G$ 是交换拓扑群. 记 $\widehat{G}$ 为从 $G$ 到 $S^1$ 的连续群同态的集合.
由于从 $G$ 到 $S^1$ 的两个连续同态的乘积也是从 $G$ 到 $S^1$ 的连续同态 (因为 $S^1$ 交换), 集合 $\widehat{G}$ 有自然的群结构. 我们在 $\widehat{G}$ 上赋予紧收敛拓扑 (topology of compact convergence); 即若 $\chi \in \widehat{G}$, 则 $\chi$ 的邻域基由集合 $\{\psi \in \widehat{G} \mid \sup_{x \in K}|\chi(x) - \psi(x)| < c\}$ 给出, 其中 $K$ 遍历 $G$ 的紧子集, $c$ 遍历正实数.
\begin{enumerate}
\item[(a)] 证明 $\widehat{G}$ 是拓扑群. 它称为 $G$ 的对偶群 (dual group).
\item[(b)] 假设 $G = \R$.
  \begin{enumerate}
  \item[(i)] 设 $\rho \colon G \to \GL_n(\CC)$ 是连续群同态. 证明存在唯一的 $A \in M_n(\CC)$ 使得对每个 $t \in \R$, $\rho(t) = \exp(tA)$. (有几种方法. 一种方法是注意到若结论成立, 则 $c_0(0)$ 必须存在且等于 $A$, 然后由此反推.)
  \item[(ii)] 证明 $\rho$ 的像包含于 $\UU(n)$ 当且仅当 $A^* = -A$.
  \item[(iii)] 证明映射 $\R \to \widehat{G}$, $x \longmapsto (t \longmapsto e^{ixt})$ 是拓扑群的同构 (即是群同构也是同胚).
  \end{enumerate}
\item[(c)] 证明存在拓扑群的同构 $\widehat{S^1} \simeq \Z$, 它把 $\mathrm{id}_{S^1}$ 送到 $1$.
\item[(d)] 拓扑群 $\widehat{\Z}$ 是什么?
\item[(e)] 假设 $G = \Q_p$ (参看习题 \ref{ex:I.5.1.4}). 我们如下定义映射 $\chi_1 \colon \Q_p \to S^1$: 若 $x \in \Q_p$, 可以把 $x$ 写成 $x = \sum_{n=-\infty}^{+\infty}c_np^n$, 其中 $0 \leq c_n \leq p - 1$ 且对足够小的 $n$ 有 $c_n = 0$, 且这唯一确定 $c_n$ (见习题 \ref{ex:I.5.1.4}(h)). 我们令
  \[
  \chi_1(x) = \exp\Bigl(2\pi i\sum_{n=-\infty}^{-1}c_np^n\Bigr).
  \]
  \begin{enumerate}
  \item[(i)] 证明 $\chi_1 \colon \Q_p \to S^1$ 是连续群同态且 $\Ker(\chi_1) = \Z_p$.
  \item[(ii)] 对每个 $y \in \Q_p$, 定义 $\chi_y \colon \Q_p \to S^1$ 为 $\chi_y(x) = \chi_1(xy)$. 证明这也是连续群同态, 并求其核.
  \item[(iii)] 设 $\chi \in \widehat{\Q_p}$. 证明存在 $k \in \Z$ 使得 $\chi$ 在 $\{x \in \Q_p \mid |x|_p \leq p^{-k}\}$ 上等于 $1$.
  \item[(iv)] 设 $\chi \in \widehat{\Q_p}$ 满足 $\chi(1) = 1$ 且 $\chi(p^{-1}) \neq 1$. 证明存在整数序列 $(c_r)_{r \geq 0}$ 使得 $1 \leq c_0 \leq p - 1$, 对 $r \geq 1$ 有 $0 \leq c_r \leq p - 1$, 且对每个 $k \in \Z_{\geq 1}$,
  \[
  \chi(p^{-k}) = \exp\Bigl(2\pi i\sum_{r=1}^k c_{k-r}p^{-r}\Bigr).
  \]
  \item[(v)] 设 $\chi \in \widehat{\Q_p}$ 满足 $\chi(1) = 1$ 且 $\chi(p^{-1}) \neq 1$. 证明存在 $y \in \Q_p$ 使得 $|y|_p = 1$ 且 $\chi = \chi_y$.
  \item[(vi)] 证明映射 $\Q_p \to \widehat{\Q_p}$, $y \longmapsto \chi_y$ 是拓扑群的同构.
  \item[(vii)] 证明 $\chi_y|_{\Z_p} = \chi_{y'}|_{\Z_p}$ 当且仅当 $y - y' \in \Z_p$, 且映射 $y \longmapsto \chi_y$ 诱导拓扑群的同构 $\Q_p/\Z_p \xrightarrow{\sim} \widehat{\Z_p}$, 其中 $\Q_p/\Z_p$ 赋予离散拓扑.
  \end{enumerate}
\end{enumerate}
\end{exercise}

\begin{solution}
(a). 我们需要验证 $\widehat{G}$ 的群运算连续. 从乘法开始. 设 $\chi_1, \chi_2 \in \widehat{G}$, 选取 $\chi_1\chi_2$ 的形如 $\{\psi \in \widehat{G} \mid \sup_{x \in K}|\chi(x) - \psi(x)| < c\}$ ($K \subset G$ 紧, $c > 0$) 的邻域 $U$. 我们需要找到 $\chi_1$ 的邻域 $U_1$ 和 $\chi_2$ 的邻域 $U_2$ 使得 $U_1U_2 \subset U$. 取
\[
U_i = \{\psi \in \widehat{G} \mid \sup_{x \in K}|\chi_i(x) - \psi(x)| < c/2\}
\]
($i = 1, 2$). 设 $\psi_1 \in U_1$, $\psi_2 \in U_2$. 则若 $x \in K$,
\begin{multline*}
|(\psi_1\psi_2)(x) - (\chi_1\chi_2)(x)| = |\psi_1(x)(\psi_2(x) - \chi_2(x)) + \chi_2(x)(\psi_1(x) - \chi_1(x))| \\
\leq |\psi_1(x)||\psi_2(x) - \chi_2(x)| + |\chi_2(x)||\psi_1(x) - \chi_1(x)| = |\psi_2(x) - \chi_2(x)| + |\psi_1(x) - \chi_1(x)| < c,
\end{multline*}
因为 $\psi_1$ 和 $\chi_2$ 酉. 所以 $\psi_1\psi_2 \in U$.
取逆的证明类似. 设 $\chi \in \widehat{G}$, 选取 $\chi^{-1}$ 的形如 $\{\psi \in \widehat{G} \mid \sup_{x \in K}|\chi^{-1}(x) - \psi(x)| < c\}$ ($K \subset G$ 紧, $c > 0$) 的邻域 $U$. 我们需要找到 $\chi$ 的邻域 $V$ 使得 $V^{-1} \subset U$. 取 $V = \{\psi \in \widehat{G} \mid \sup_{x \in K}|\chi(x) - \psi(x)| < c\}$. 设 $\psi \in V$. 则对每个 $x \in K$,
\[
|\psi^{-1}(x) - \chi^{-1}(x)| = |\psi^{-1}(x)||\chi^{-1}(x)||\chi(x) - \psi(x)| = |\chi(x) - \psi(x)| < c.
\]
所以 $\psi^{-1} \in U$.

(b). (i) 在 $M_n(\CC)$ 上选取范数 $\|\cdot\|$. 由于 $\GL_n(\CC)$ 在 $M_n(\CC)$ 中开, 可以选取以 $I_n$ 为心、包含于 $\GL_n(\CC)$ 的非空开球 $B$. 我们只关心 $B$ 是 $M_n(\CC)$ 的凸子集这一事实. 由于 $\rho$ 连续且 $\rho(0) = I_n$, 可以找到 $c > 0$ 使得 $\rho([0, c]) \subset B$. 则
\[
\int_0^1\rho(cx)\,\dd x = c\int_0^c\rho(x)\,\dd x \in B,
\]
所以 $X := \int_0^c\rho(x)\,\dd x \in \GL_n(\CC)$. 对每个 $t \in \R$,
\[
X\rho(t) = \int_0^c\rho(x + t)\,\dd x = \int_t^{t+c}\rho(x)\,\dd x.
\]
特别地, $\rho$ 连续可微, 且
\[
\rho'(t) = X^{-1}(\rho(t + c) - \rho(t)) = X^{-1}(\rho(c) - I_n)\rho(t).
\]
满足初值条件 $\rho(0) = I_n$ 的这个微分方程的唯一解是 $\rho(t) = \exp(tA)$, 其中 $A = X^{-1}(\rho(c) - I_n)$. 最后, 矩阵由 $\rho$ 唯一确定, 因为必须有 $A = \rho'(0)$.

(ii) 若 $A^* = -A$, 则对每个 $t \in \R$,
\[
\rho(t)\rho(t)^* = \exp(tA)\exp(tA^*) = \exp(t(A + A^*)) = \exp(0) = I_n
\]
(我们用 $tA$ 与 $tA^*$ 交换这一事实得到等式 $\exp(tA)\exp(tA^*) = \exp(tA + tA^*)$), 所以 $\rho(t) \in \UU(n)$.
反之, 假设 $\rho(\R) \subset \UU(n)$. 注意 $A = \lim_{t \to 0}\frac{1}{t}(\rho(t) - I_n)$, 所以 $A^* = \lim_{t \to 0}\frac{1}{t}(\rho(t)^* - I_n)$. 由于
\[
\rho(t)^* - I_n = \rho(t)^{-1} - I_n = -\rho(t)^{-1}(\rho(t) - I_n)
\]
且当 $t \to 0$ 时 $\rho(t)^{-1} \to I_n$, 这蕴含 $A^* = -A$.

(iii) 记题设的映射 $\R \to \widehat{G}$ 为 $\alpha$.
我们在 (i) 和 (ii) 中已见每个连续群同态 $\rho \colon \R \to S^1$ 都具有 $\rho(t) = e^{zt}$ 的形式, 其中 $z \in \CC$ 唯一且满足 $\bar{z} = -z$; 最后一个条件意味着 $z = ix$ (对某个 $x \in \R$). 这意味着 $\alpha$ 双射. 又容易看出 $\alpha$ 是群同态, 所以我们只需证明 $\alpha$ 是同胚.
首先证明 $\alpha$ 连续. 设 $x \in \R$, 考虑 $\alpha(x)$ 的形如 $\{\rho \in \widehat{G} \mid \forall t \in K,\, |\alpha(x)(t) - \rho(t)| < c\}$ ($K \subset \R$ 紧, $c > 0$) 的邻域 $U$. 则对每个 $y, t \in \R$,
\[
|\alpha(x)(t) - \alpha(y)(t)|^2 = |e^{ixt} - e^{iyt}|^2 = |1 - e^{it(x-y)}|^2 = (1 - \cos(t(x-y)))^2 + (\sin(t(x-y)))^2.
\]
选取 $\varepsilon > 0$ 使得对每个 $t \in K$ 和 $z \in (-\varepsilon, \varepsilon)$, 有 $(1 - \cos(tz))^2 + (\sin(tz))^2 < c^2$. 则若 $|x - y| < \varepsilon$, 有 $\alpha(y) \in U$.
现在我们证明 $\alpha$ 开. 设 $x \in \R$, 选取 $x$ 的形如 $(x - \varepsilon, x + \varepsilon)$ 的邻域 $V$. 我们要证明 $\alpha(V)$ 包含 $\alpha(x)$ 的邻域. 选取 $\delta > 0$ 使得函数 $t \longmapsto \sin(t)$ 和 $t \longmapsto 1 - \cos(t)$ 在 $[0, 2\delta\varepsilon]$ 上都递增, 并设
\[
U = \{\rho \in \widehat{G} \mid \forall t \in K,\, |\alpha(x)(t) - \rho(t)| < c\},
\]
其中 $K = [-\delta, \delta]$, $c = \sup_{t \in [0, \varepsilon\delta/2]}((1 - \cos(t))^2 + (\sin(t))^2)^{1/2}$ (注意这也是 $[-\varepsilon\delta/2, \varepsilon\delta/2]$ 上的上确界, 因为我们估计的函数是偶的). 设 $y \in \R$ 满足 $|x - y| \geq \varepsilon$. 我们要证明 $\alpha(y) \notin U$. 可以找到 $t \in K$ 使得 $\varepsilon\delta \leq t(x - y) \leq 2\varepsilon\delta$. 则
\[
|\alpha(x)(t) - \alpha(y)(t)| = ((1 - \cos(t(x - y)))^2 + (\sin(t(x - y)))^2)^{1/2} > c,
\]
由 $\delta$ 和 $c$ 的选取. 所以 $\alpha(y) \notin U$.

(c). 注意有拓扑群的同构 $\R/2\pi\Z \xrightarrow{\sim} S^1$, 由 $t \longmapsto e^{it}$ 给出. 所以得到群同构
\[
\widehat{S^1} \simeq \{\rho \in \widehat{\R} \mid \rho(2\pi\Z) = \{1\}\} \simeq \{x \in \R \mid \forall t \in 2\pi\Z,\, e^{ixt} = 1\} = \Z
\]
(其中第二个同构来自 (b)). 还需证明这是拓扑群的同构, 即 $\widehat{S^1}$ 离散. 若你已提前读过, 你知道这是习题 \ref{ex:I.5.4.2}(e) 的特殊情形 (在 $S^1$ 的情形我不知道更简单的证明).

(d). 由于 $\Z$ 离散, 从 $\Z$ 到 $S^1$ 的连续群同态就是从 $\Z$ 到 $S^1$ 的群同态. 由于 $\Z$ 是由 $1 \in \Z$ 生成的自由交换群, 映射 $\rho \longmapsto \rho(1)$ 是从 $\Z \to S^1$ 的群同态集合到 $S^1$ 的同构. 所以作为群, $\widehat{\Z}$ 与 $S^1$ 同构. 记此同构为 $\beta \colon S^1 \to \widehat{\Z}$ (于是 $\beta$ 把 $z \in S^1$ 送到同态 $\Z \to S^1$, $n \longmapsto z^n$). 若我们证明 $\beta$ 连续, 则它将自动是同胚, 因为 $S^1$ 紧. 但 $\Z$ 的紧子集恰是它的有限子集, 所以 $\beta$ 的连续性由映射 $S^1 \to S^1$, $z \longmapsto z^n$ 的连续性立即得出.

(e). (i) 设 $x, x' \in \Q_p$, 写 $x = \sum_{n=-\infty}^{+\infty}c_np^n$, $x' = \sum_{n=-\infty}^{+\infty}c'_np^n$ (对 $c_n$ 和 $c'_n$ 的条件与题设相同). 则由习题 \ref{ex:I.5.1.4}(h), 对每个 $N \in \Z$, 若对每个 $n \leq N - 1$ 有 $c_n = c'_n$, 则 $|x - x'|_p \leq p^{-N}$. 特别地, 若 $|x - x'|_p \leq 1$ 则 $\chi_1(x) = \chi_1(x')$, 所以 $\chi_1$ 连续且把每个 $x \in \Z_p$ 送到 $1 = \chi_1(0)$.
我们还需证明 $\chi_1$ 是群同态. 设 $G_0$ 是 $\Q_p$ 的由满足 $x = \sum_{n=-\infty}^{+\infty}a_np^n$ (其中 $a_n \in \Z$ 且对足够大的 $|n|$ 有 $a_n = 0$) 的元素组成的子群. 这是稠密子群 (因为 $\sum_{n=-\infty}^{+\infty}c_np^n$ 是 $N \to +\infty$ 时 $\sum_{n=-N}^{N}c_np^n$ 的极限), 且包含于 $\Q$ 中. 由于我们知道 $\chi_1$ 连续, 只需对 $x, y \in G_0$ 证明 $\chi_1(x + y) = \chi_1(x)\chi_1(y)$. 但注意若 $x \in G_0$, 则 $\chi_1(x) = \exp(2\pi ix)$, 其中把 $x$ 看作 $\Q$ 的元素. 这蕴含结果.
最后, 我们必须证明 $\Ker(\chi_1) = \Z_p$. 我们已经看到 $\Z_p \subset \Ker(\chi_1)$. 反之, 设 $x \in \Q_p$, 如上写 $x = \sum_{n=-\infty}^{+\infty}c_np^n$. 假设 $x \notin \Z_p$; 则存在 $m < 0$ 使得 $c_m \neq 0$. 选取这样的 $m$. 我们有
\[
p^{-m} \leq c_mp^{-m} < \sum_{n=-\infty}^{-1}c_np^n \leq (p - 1)\sum_{r \geq 1}p^{-r} = 1
\]
(第二个不等式严格, 因为对足够小的 $n$ 有 $c_n = 0$). 所以 $\sum_{n=-\infty}^{-1}c_np^n \in (0, 1)$, 且 $\chi_1(x) = \exp\bigl(2\pi i\sum_{n=-\infty}^{-1}c_np^n\bigr) \neq 1$.

(ii) 映射 $\chi_y$ 是连续群同态, 因为它是连续群同态 $\chi_1$ 和 $m_y \colon \Q_p \to \Q_p$, $x \longmapsto xy$ 的复合. $\Q_p$ 的元素 $x$ 在 $\chi_y$ 的核中当且仅当 $xy \in \Ker(\chi_1) = \Z_p$. 所以若 $y = 0$, 则 $\Ker(\chi_y) = \Q_p$; 若 $y \neq 0$, 则
\[
\Ker(\chi_y) = y^{-1}\Z_p = |y|_p\Z_p = \{x \in \Q_p \mid |x|_p \leq |y|_p^{-1}\}.
\]

(iii) 选取 $\CC^{\times}$ 的邻域 $U$ 使得 $\CC^{\times}$ 包含于 $U$ 的唯一子群是平凡群. (见 (b).) 则 $\chi^{-1}(U \cap S^1)$ 是 $1$ 在 $\Q_p$ 中的邻域, 所以存在 $k \in \Z$ 使得 $\chi^{-1}(U \cap S^1) \supset \{x \in \Q_p \mid |x|_p \leq p^k\}$. 但由于 $\{x \in \Q_p \mid |x|_p \leq p^k\}$ 是 $\Q_p$ 的子群, 它在 $\chi$ 下的像是 $S^1$ 的包含于 $U$ 的子群, 所以等于 $\{1\}$.

(iv) 对每个整数 $r \geq 0$ 写 $z_r = \chi(p^{-r})$. 则 $z_r \in S^1$, 且对每个 $r \geq 0$,
\[
z_{r+1}^p = \chi(p^{-r-1})^p = \chi(p^{-r}) = z_r.
\]
我们将对 $r \geq 0$ 归纳构造整数 $c_r$. 首先注意由假设 $z_1 \neq 1 = z_0$, 所以可以找到 $c_0 \in \{1, \ldots, p-1\}$ 使得 $z_1 = \exp(2\pi ic_0p^{-1})$. 假设已经找到 $c_0, \ldots, c_{r-1}$ ($r \geq 1$) 使得对 $1 \leq s \leq r$ 有
\[
\chi(p^{-s}) = z_s = \exp\Bigl(2\pi i\sum_{k=1}^s c_{s-k}p^{-k}\Bigr).
\]
我们必须找到 $c_r \in \{0, \ldots, p-1\}$ 使得
\[
z_{r+1} = \exp\Bigl(2\pi i\sum_{k=1}^{r+1}c_{r+1-k}p^{-k}\Bigr) = \exp\Bigl(2\pi ip^{-(r+1)}\sum_{s=0}^r c_sp^s\Bigr).
\]
由于 $z_{r+1}^p = z_r$, 有
\[
\Bigl(z_{r+1}\exp\Bigl(2\pi ip^{-r-1}\sum_{s=0}^{r-1}c_sp^s\Bigr)\Bigr)^p = 1,
\]
所以存在 $c_r \in \{0, \ldots, p-1\}$ 使得
\[
z_{r+1}\exp\Bigl(2\pi ip^{-r-1}\sum_{s=0}^{r-1}c_sp^s\Bigr) = \exp(2\pi ip^{-1}c_r),
\]
即
\[
z_{r+1} = \exp\Bigl(2\pi ip^{-(r+1)}\sum_{s=0}^r c_sp^s\Bigr).
\]

(v) 设 $(c_r)_{r \geq 0}$ 如 (iv), 令 $y = \sum_{r=0}^{+\infty}c_rp^r$. 由于 $c_0 \in \{1, \ldots, p-1\}$, 有 $|y|_p = 1$. 又, 对每个 $r \geq 1$,
\[
\chi(p^{-r}) = \exp\Bigl(2\pi ip^{-(r+1)}\sum_{k=-r}^{-1}c_{r+k}p^k\Bigr) = \chi_1(p^{-r}y) = \chi_y(p^{-r}),
\]
因为
\[
p^{-r}y = \sum_{s \geq 0}c_sp^{s-r} = \sum_{n=-r}^{+\infty}c_{r+n}p^n.
\]
另一方面, 若 $r \geq 0$, 则
\[
\chi(p^r) = \chi(1)^{p^r} = 1 = \chi_y(p^r).
\]
由于 $\chi$ 和 $\chi_y$ 是连续群同态, 且族 $(p^r)_{r \in \Z}$ 生成 $\Q_p$ 的稠密子群, 这蕴含 $\chi = \chi_y$.

(vi) 记映射 $\Q_p \to \widehat{\Q_p}$, $y \longmapsto \chi_y$ 为 $\alpha$. 容易看出 $\alpha$ 是群同态 (这由 $\chi_1$ 是群同态和 $\Q_p$ 上乘法的分配律立即得出).
首先证明 $\Ker(\alpha) = \{0\}$. 设 $y \in \Q_p - \{0\}$. 则 $y = \sum_{n=m}^{+\infty}c_np^n$, 其中 $m \in \Z$, $0 \leq c_n \leq p - 1$ 且 $c_m \geq 1$. 所以
\[
p^{-m-1}y = c_mp^{-1} + \sum_{n \geq 0}c_{n+m+1}p^n,
\]
且 $\chi_y(p^{-m-1}) = \exp(2\pi ip^{-1}c_m) \neq 1$. 这证明 $y \notin \Ker(\alpha)$.
现在我们证明 $\alpha$ 满射. 设 $\chi \in \widehat{\Q_p}$. 若 $\chi = 1$, 则 $\chi = \chi_0$, 所以假设 $\chi \neq 1$. 由 (iii), 存在 $k \in \Z$ 使得 $\chi$ 在 $\{x \in \Q_p \mid |x|_p \leq p^{-k}\}$ 上等于 $1$. 选取最小的这种 $k$ (这是可能的, 否则 $\chi$ 在整个 $\Q_p$ 上等于 $1$, 与我们 $\chi \neq 1$ 的假设矛盾). 则存在 $a \in \Q_p$ 使得 $|a|_p = p^{-k+1}$ 且 $\chi(a) \neq 1$. 定义 $\psi \in \widehat{\Q_p}$ 为 $\psi(x) = \chi(pax)$. 则 $\psi(p^{-1}) = \chi(a) \neq 1$ 且 $\psi(1) = \chi(pa) = 1$ (因为 $|pa|_p = p^{-k}$). 由 (v), 存在 $y \in \Z_p$ 使得 $\psi = \chi_y$. 换句话说, 对每个 $x \in \Q_p$,
\[
\chi(x) = \psi(p^{-1}a^{-1}x) = \chi_1(p^{-1}a^{-1}yx),
\]
即 $\chi = \alpha(p^{-1}a^{-1}y)$.
我们证明 $\alpha$ 连续. 设 $y \in \Q_p$, 选取 $\alpha(y)$ 的形如
\[
U = \{\chi \in \widehat{\Q_p} \mid \forall x \in K,\, |\chi(x) - \chi_y(x)| < c\}
\]
的邻域, 其中 $K$ 是 $\Q_p$ 的紧子集, $c > 0$. 我们要找 $y$ 在 $\Q_p$ 中的邻域 $V$ 使得 $\alpha(V) \subset U$.
由于 $\Q_p = \bigcup_{k \in \Z}p^k\Z_p$, 可以假设 $K = p^N\Z_p$ (对某个 $N \in \Z$). 我们知道 $\chi_1$ 在 $\Q_p$ 中 $\Z_p$ 的陪集上为常数, 所以若 $x \in p^N\Z_p$, 则 $\chi_x$ 在 $\Q_p$ 中 $p^{-N}\Z_p$ 的陪集上为常数. 所以若 $y' \in y + p^{-N}\Z_p$, 则对每个 $x \in K = p^N\Z_p$,
\[
|\chi_{y'}(x) - \chi_y(x)| = |\chi_x(y') - \chi_x(y)| = 0 < c.
\]
换句话说, $\alpha(y + p^{-N}\Z_p) \subset U$.
最后, 我们证明 $\alpha$ 开. 设 $y \in \Q_p$, $V$ 是 $y$ 的邻域. 可以假设 $V$ 具有 $y + p^N\Z_p = \{y' \in \Q_p \mid |y' - y|_p \leq p^{-N}\}$ 的形式 (对某个 $N \in \Z$). 我们要证明 $\alpha(V)$ 包含 $\alpha(y)$ 的邻域. 由于 $\alpha$ 是群同态, 可以假设 $y = 0$. 设
\[
U = \{\chi \in \widehat{\Q_p} \mid \forall x \in p^{-N}\Z_p,\, |\chi(x) - 1| < c\},
\]
其中 $c = \min_{1 \leq r \leq p-1}|1 - e^{2\pi irp^{-1}}|$, 让我们证明 $\alpha(p^N\Z_p) \supset U$. 设 $y' \notin p^N\Z_p$, 我们要证明 $\chi_{y'} \notin U$. 把 $y'$ 写成 $y' = \sum_{n=m}^{+\infty}c_np^n$, 其中对每个 $n \geq m$ 有 $c_n \in \{0, \ldots, p-1\}$ 且 $c_m \geq 1$. 则对 $y'$ 的假设说 $m < N$. 设 $x = p^{-m-1}$. 则 $x \in p^{-N}\Z_p$, 且
\[
\chi_{y'}(x) = \chi_1(xy') = \exp(2\pi ip^{-1}c_m),
\]
所以 $|\chi_{y'}(x) - 1| \geq c$ 且 $\chi_{y'} \notin U$.

(vii) 由于映射 $y \longmapsto \chi_y$ 是群同态, 第一个结论等价于 $\chi_y|_{\Z_p} = 1$ 当且仅当 $y \in \Z_p$ 这一事实. 我们知道 $\Ker(\chi_1) = \Z_p$, 所以对每个 $y \in \Z_p$ 有 $\Ker(\chi_y) \supset \Z_p$. 反之, 设 $y \in \Q_p - \Z_p$. 则 $|y|_p > 1$, 所以 $|y|_p \geq p$, 从而 $|py|_p \geq 1$, 且 $p^{-1}y^{-1} \in \Z_p$. 由于 $\chi_y(p^{-1}y^{-1}) = \chi_1(p^{-1}) = \exp(2\pi ip^{-1}) \neq 1$, 有 $\Ker(\chi_y) \not\supset \Z_p$.
所以映射 $y \longmapsto \chi_y$ 诱导从 $\Q_p/\Z_p$ 到 $\widehat{\Z_p}$ 的单射群同态. 我们知道 (或很快会知道) 由习题 \ref{ex:I.5.4.2}(e), $\widehat{\Z_p}$ 离散, 所以只需证明 $\widehat{\Z_p}$ 的每个元素都具有 $\chi_y|_{\Z_p}$ 的形式 (对某个 $y \in \Q_p$).
设 $\chi \in \widehat{\Z_p}$. 同 (iii), 可以找到 $k \in \N$ 使得 $\Ker(\chi) \supset p^k\Z_p$. 设 $z = \chi(1)$. 则 $z^{p^k} = \chi(p^k) = 1$, 所以可以找到 $c \in \{0, \ldots, p^k - 1\}$ 使得 $z = e^{2\pi icp^{-k}}$. 把 $c$ 写成 $c = \sum_{r=0}^{k-1}c_rp^r$, 其中 $c_r \in \{0, \ldots, p-1\}$. 则
\[
\chi(1) = \exp\Bigl(2\pi i\sum_{r=0}^{k-1}c_rp^{r-k}\Bigr) = \exp\Bigl(2\pi i\sum_{n=-k}^{-1}c_{k+n}p^n\Bigr).
\]
设 $y = \sum_{n=-k}^{-1}c_{k+n}p^n$. 则 $\chi(1) = \chi_y(1)$. 由于 $\chi_y|_{\Z_p}$ 和 $\chi$ 是 $\Z_p$ 上的连续群同态, 且 $1$ 生成 $\Z_p$ 的稠密子群, 这蕴含 $\chi = \chi_y|_{\Z_p}$.
\end{solution}

\begin{exercise}[对偶群的局部紧性]
\label{ex:I.5.4.2}
我们使用前一题的记号, 并假设 $G$ 是交换局部紧群, 固定 $G$ 上的 Haar 测度 $\mu$.
回忆有同构 $L^{\infty}(G) \to L^1(G)^{\vee} := \Hom(L^1(G), \CC)$, 把 $f \in L^{\infty}(G)$ 送到 $L^1(G)$ 上的有界算子 $g \longmapsto \int_G fg\,\dd\mu$. (这不使用 $G$ 是交换群这一事实.) 所以我们可以在 $L^{\infty}(G)$ 上考虑弱*拓扑 (或逐点收敛拓扑): 对 $f \in L^{\infty}(G)$, $f$ 的邻域基由集合 $U_{g_1,\ldots,g_n,c} = \{f' \in L^{\infty}(G) \mid |\int_G (f - f')g_i\,\dd\mu| < c,\, 1 \leq i \leq n\}$ 给出, 其中 $n \in \Z_{\geq 1}$, $g_1, \ldots, g_n \in L^1(G)$, $c > 0$.
\begin{enumerate}
\item[(a)] 证明 $\widehat{G} \subset L^{\infty}(G)$, 且 $\widehat{G}$ 的拓扑由 $L^{\infty}(G)$ 的弱*拓扑诱导.\footnote{难题.}
\item[(b)] 证明 $L^{\infty}(G)$ 的子集 $\widehat{G} \cup \{0\}$ 对弱*拓扑闭. (提示: 把它等同于 Banach $*$-代数 $L^1(G)$ 在 $\CC$ 上的表示集合.)
\item[(c)] 证明 $\widehat{G}$ 是局部紧拓扑群. (提示: Alaoglu 定理.)
\item[(d)] 若 $G$ 离散, 证明 $\widehat{G}$ 紧.
\item[(e)] 若 $G$ 紧, 证明 $\widehat{G}$ 离散.
\end{enumerate}
\end{exercise}

\begin{solution}
(a). $\widehat{G}$ 的元素是从 $G$ 到 $S^1$ 的连续函数, 从而是 $G$ 上的连续有界函数, 因此是 $L^{\infty}(G)$ 的元素. 现在我们证明 $\widehat{G}$ 上的两个拓扑一致.
首先设 $\chi \in \widehat{G}$, 设 $f_1, \ldots, f_n \in L^1(G)$, $c > 0$. 这定义 $\chi$ 的弱*开邻域
\[
U = \Bigl\{\psi \in \widehat{G} \mid \forall i \in \{1, \ldots, n\},\, \Bigl|\int_G \chi f_i\,\dd\mu - \int_G \psi f_i\,\dd\mu\Bigr| < c\Bigr\}.
\]
我们要找紧收敛拓扑下 $\chi$ 的开邻域 $V$ 使得 $V \subset U$. 设 $\varepsilon > 0$. 选取 $G$ 的紧子集 $K$ 使得对每个 $i \in \{1, \ldots, n\}$ 有 $\int_{G-K}|f_i|\,\dd\mu < \varepsilon$ (由 $\mu$ 的内正则性这是可能的). 设
\[
V = \{\psi \in \widehat{G} \mid \forall x \in K,\, |\chi(x) - \psi(x)| < \varepsilon\}.
\]
则若 $\psi \in V$, $i \in \{1, \ldots, n\}$, 有
\begin{align*}
\Bigl|\int_G \chi f_i\,\dd\mu - \int_G \psi f_i\,\dd\mu\Bigr| &\leq \int_K |f_i(x)||\chi(x) - \psi(x)|\,\dd\mu(x) + \int_{G-K}|f_i(x)||\chi(x) - \psi(x)|\,\dd x \\
&\leq \varepsilon\int_K |f_i(x)|\,\dd\mu(x) + 2\int_{G-K}|f_i(x)|\,\dd\mu(x) \\
&\leq \varepsilon(\|f_i\|_1 + 2).
\end{align*}
所以若取 $\varepsilon$ 足够小, 将得到 $V \subset U$.
现在我们证明其逆. 设 $\chi \in \widehat{G}$, $K$ 是 $G$ 的紧子集, $c > 0$. 考虑 $\chi$ 在紧收敛拓扑下的邻域
\[
V = \{\psi \in \widehat{G} \mid \forall x \in K,\, |\chi(x) - \psi(x)| < c\}.
\]
必须找到包含于它的弱*邻域. 设 $\eta > 0$ (稍后调整), 选取 $1$ 的紧邻域 $A$ 使得对每个 $y \in A$ 有 $|\chi(y) - 1| < \eta$. 设 $f = \mathbb{1}_A$; 由于 $A$ 紧, $f \in L^1(G)$. 注意对每个 $x \in G$,
\[
|\mu(A)\chi(x) - f * \chi(x)| = \Bigl|\int_A (\chi(x) - \chi(y^{-1}x))\,\dd y\Bigr| \leq \int_A |1 - \chi(y)|\,\dd y \leq \eta\mu(A).
\]
现在我们试着找到 $\widehat{G}$ 中 $\chi$ 的弱*邻域, 使其元素 $\psi$ 对 $x \in K$ 满足类似的不等式. 注意若 $\psi \in \widehat{G}$, $x \in G$, 则
\[
f * \psi(x) = \int_A \psi(y^{-1}x)\,\dd y = \psi(x)\int_A \psi(y)\,\dd y = \int_G \psi(y^{-1})f(xy)\,\dd y = \int_G \psi(y)L_{x^{-1}}f(y)\,\dd y
\]
(用 $G$ 交换且 $\psi$ 是从 $G$ 到 $S^1$ 的群同态这一事实). 现在回忆映射 $G \to L^1(G)$, $x \longmapsto L_{x^{-1}}f$ 连续 (命题 \ref{prop:I.3.1.13}). 由于 $K$ 紧, 可以找到 $x_1, \ldots, x_n$ 使得对每个 $x \in K$, 存在 $i \in \{1, \ldots, n\}$ 满足 $\|L_{x^{-1}}f - L_{x_i^{-1}}f\|_1 < \eta\mu(A)$. 考虑 $\chi$ 的如下弱*邻域:
\[
U = \Bigl\{\psi \in \widehat{G} \mid \forall x \in \{1, x_1, \ldots, x_n\},\, \Bigl|\int_G \chi(y)L_{x^{-1}}f(y)\,\dd y - \int_G \psi(y)L_{x^{-1}}f(y)\,\dd y\Bigr| < \eta\mu(A)\Bigr\}.
\]
设 $\psi \in U$. 首先, 对每个 $x \in G$,
\begin{align*}
|\mu(A)\psi(x) - f * \psi(x)| &= \Bigl|\int_A (\psi(x) - \psi(y^{-1}x))\,\dd y\Bigr| \\
&= \Bigl|\int_A (1 - \psi(y))\,\dd y\Bigr| \\
&\leq \int_A |1 - \chi(y)|\,\dd y + \Bigl|\int_A (\chi(y) - \psi(y))\,\dd y\Bigr| \\
&\leq 2\eta\mu(A).
\end{align*}
其次, 要对 $x \in K$ 估计 $|f * \chi(x) - f * \psi(x)|$. 固定 $x \in K$. 设 $i \in \{1, \ldots, n\}$ 满足 $\|L_{x^{-1}}f - L_{x_i^{-1}}f\|_1 \leq \eta\mu(A)$. 则:
\begin{align*}
|f * \chi(x) - f * \psi(x)| &= \Bigl|\int_G (\chi(y) - \psi(y))L_{x^{-1}}f(y)\,\dd y\Bigr| \\
&\leq \Bigl|\int_G (\chi(y) - \psi(y))L_{x_i^{-1}}f(y)\,\dd y\Bigr| + \Bigl|\int_G (\chi(y) - \psi(y))(L_{x_i^{-1}}f(y) - L_{x^{-1}}f(y))\,\dd y\Bigr| \\
&< \eta\mu(A) + 2\int_G |L_{x_i^{-1}}f(y) - L_{x^{-1}}f(y)|\,\dd y \\
&\leq 3\eta\mu(A).
\end{align*}
把所有估计放在一起, 对 $x \in K$ 得到 $|\mu(A)\chi(x) - \mu(A)\psi(x)| < 6\eta\mu(A)$, 即 $|\chi(x) - \psi(x)| < 6\eta$. 在一开始选取 $\eta$ 使得 $6\eta \leq c$, 得到 $U \subset V$, 如所愿.
(b). 我们在课堂上见过 $\widehat{G} \subset L^{\infty}(G) \simeq L^1(G)^{\vee}$ 是 Banach $*$-代数 $L^1(G)$ 的非退化表示集合. 设 $\pi \colon L^1(G) \to \CC$ 是 $L^1(G)$ 在 $\CC$ 上的表示, 假设它不是非退化的. 则存在 $v \in \CC - \{0\}$ 使得对每个 $f \in L^1(G)$ 有 $\pi(f)v = 0$. 但这蕴含 $\pi = 0$. 所以我们看到 $\widehat{G} \cup \{0\} \subset L^{\infty}(G)$ 确实是 $L^1(G)$ 在 $\CC$ 上的表示集合. 但说一个有界线性泛函 $\Lambda \colon L^1(G) \to \CC$ 是表示的条件在弱*拓扑下都是闭条件 (因为它们都断言 $\Lambda$ 在 $L^1(G)$ 的一些点处的取值相等), 所以 $L^1(G)$ 的表示集合是 $L^{\infty}(G)$ 的弱*闭子集.
(c). Alaoglu 定理\footnote{参考文献 ?} 说 $L^{\infty}(G)$ 的闭单位球 (对来自 $\|\cdot\|_1$ 的算子范数, 即 $\|\cdot\|_{\infty}$) 对弱*拓扑紧 Hausdorff. 但 $\widehat{G} \cup \{0\}$ 显然包含于这个闭单位球, 所以它对弱*拓扑紧 Hausdorff. 因此它的开子集 $\widehat{G}$ 对弱*拓扑局部紧, 且已在 (a) 中看到 $\widehat{G}$ 上的弱*拓扑等于紧收敛拓扑, 所以我们完成了.
(d). 考虑映射 $\alpha \colon \widehat{G} \to (S^1)^G$, 把 $\chi$ 送到族 $(\chi(x))_{x \in G}$. 这显然单射. 由于 $G$ 离散, 它的紧子集恰是有限子集, 所以紧收敛拓扑恰是 $(S^1)^G$ 上乘积拓扑诱导的拓扑. 又, 由 Tychonoff 定理, $(S^1)^G$ 紧 Hausdorff. 所以要得到结果, 只需证明 $\alpha$ 的像在 $(S^1)^G$ 中闭. 但 $\alpha$ 的像是子集族
\[
\{(a_x)_{x \in G} \in (S^1)^G \mid a_{x_0}a_{y_0} = a_{x_0y_0}\}
\]
(对所有 $x_0, y_0 \in G$) 的交, 且每个这样的子集都闭, 所以 $\Im(\alpha)$ 闭.
(e). 假设 $G$ 紧. 则 $\widehat{G}$ 的拓扑是一致收敛拓扑 (由范数 $\|\cdot\|_{\infty}$ 诱导). 要证明 $\widehat{G}$ 离散, 只需证明其子集 $\{1\}$ 开 (因为 $\widehat{G}$ 是拓扑群). 设 $c > 0$ 使得 $\CC^{\times}$ 包含于 $\{z \in \CC^{\times} \mid |1 - z| < c\}$ 的唯一子群是平凡群. 设 $U = \{\chi \in \widehat{G} \mid \|\chi - 1\|_{\infty} < c\}$. 这是 $1$ 在 $\widehat{G}$ 中的开邻域. 另一方面, 若 $\chi \in U$, 则 $\chi(G) \subset \{z \in \CC^{\times} \mid |1 - z| < c\}$; 由于 $\chi(G)$ 是 $\CC^{\times}$ 的子群, 这意味着 $\chi(G) = \{1\}$, 即 $\chi = 1$, 于是 $U = \{1\}$.
\end{solution}

\subsection{表示}

若 $G$ 是一个群, 我们说 $G$ 的表示 $(\pi, V)$ 是忠实的 (faithful), 如果 $\pi \colon G \to \GL(V)$ 是单射.

\begin{exercise}[$\SU(2)$ 的不可约表示]
\label{ex:I.5.5.1}
设 $G = \SU(2)$. 群 $G$ 通过包含 $G \subset \GL_2(\CC)$ 作用于 $\CC^2$, 我们把这个作用记为 $(g, (z_1, z_2)) \longmapsto g(z_1, z_2)$. (这称为 $G$ 的标准表示 (standard representation).)
对每个整数 $n \geq 0$, 设 $V_n$ 为 $n$ 次齐次 (homogeneous) 多项式 $P \in \CC[t_1, t_2]$ 的空间 (即 $P(t_1, t_2) = \sum_{r=0}^n a_rt_1^rt_2^{n-r}$, 其中 $a_0, \ldots, a_n \in \CC$).
\begin{enumerate}
\item[(a)] 若 $P \in V_n$, $g \in G$, 证明函数 $\CC^2 \to \CC$, $(z_1, z_2) \longmapsto P(g^{-1}(z_1, z_2))$ 仍由 $V_n$ 中的多项式给出, 且这定义 $G$ 在 $V_n$ 上的连续表示.
\item[(b)] 证明对每个 $n \geq 0$, $G$ 的表示 $V_n$ 不可约.
\item[(c)] 对哪些 $n$, 表示 $V_n$ 忠实?
\end{enumerate}
\end{exercise}

\begin{remarkstar}
我们将在后面 (见习题 \ref{ex:IV.9.1}) 看到 $\SU(2)$ 的每个不可约酉表示都同构于某个 $V_n$.
\end{remarkstar}

\begin{solution}
(a). 先取 $P = t_1^rt_2^{n-r}$, 其中 $0 \leq r \leq n$. 设 $x = \begin{pmatrix} a & b \\ c & d \end{pmatrix} \in \SU(2)$. 由于 $\det(x) = 1$, 有 $x^{-1} = \begin{pmatrix} d & -b \\ -c & a \end{pmatrix}$. 所以
\[
P(x^{-1}(z_1, z_2)) = (dz_1 - bz_2)^r(-cz_1 + az_2)^{n-r}.
\]
这仍是 $z_1$ 和 $z_2$ 的 $n$ 次齐次多项式, 记之为 $P \circ x^{-1}$. 又, 由上式显然映射 $G \to V_n$, $x \longmapsto P \circ x^{-1}$ 连续 (这意味着 $P \circ x^{-1}$ 的系数是矩阵 $x$ 的元素的连续函数).
由于单项式 $t_1^rt_2^{n-r}$ ($0 \leq r \leq n$) 生成 $V_n$, 上一段蕴含对每个 $P \in V_n$ 和每个 $x \in \SU(2)$, 函数 $\CC^2 \to \CC$, $(z_1, z_2) \longmapsto P(x^{-1}(z_1, z_2))$ 仍由 $V_n$ 的元素给出 (记之为 $P \circ x^{-1}$); 它也蕴含映射 $G \to V_n$, $x \longmapsto P \circ x^{-1}$ 连续.
对每个 $x \in G$, 映射 $V_n \to V_n$, $P \longmapsto P \circ x^{-1}$ 显然是 $\CC$-线性的 (事实上我们已经用过这一事实). 又对每个 $P \in V_n$ 和所有 $x, y \in G$, 有 $P \circ (xy)^{-1} = (P \circ y^{-1}) \circ x^{-1}$. 所以由命题 \ref{prop:I.3.5.1}, 映射 $G \times V_n \to V_n$, $(x, P) \longmapsto P \circ x^{-1}$ 连续, 即定义 $G$ 在 $V_n$ 上的连续表示.

(b). 设 $W$ 是 $V$ 的 $G$-不变子空间. 设 $P = \sum_{r=0}^n c_rt_1^rt_2^{n-r} \in W$. 我们证明对每个使得 $c_r \neq 0$ 的 $r \in \{0, \ldots, n\}$, 有 $t_1^rt_2^{n-r} \in W$. 我们对 $P$ 的非零系数的个数归纳证明. 若 $P$ 有 $0$ 或 $1$ 个非零系数, 则已完成. 假设 $P$ 至少有 $2$ 个非零系数. 固定使得 $c_r \neq 0$ 的 $r \in \{0, \ldots, n\}$. 只需找到 $W$ 的另一个元素 $Q$, 使得 $t_1^rt_2^{n-r}$ 的系数非零, 且 $Q$ 的非零系数个数比 $P$ 少; 然后可以对 $Q$ 应用归纳假设. 选取 $s \in \{0, \ldots, n\} - \{r\}$ 使得 $c_s \neq 0$. 考虑 $x_a = \begin{pmatrix} a & 0 \\ 0 & a \end{pmatrix}$, 其中 $a \in S^1$. 则 $x_a \in \SU(2)$, 且
\[
P \circ x_a^{-1} = \sum_{i=0}^n a^ia^{n-i}c_it_1^it_2^{n-i} = \sum_{i=0}^n a^{n-2i}c_it_1^it_2^{n-i}.
\]
选取 $a, a' \in S^1$ 使得 $a^{n-2s}c_s - (a')^{n-2s}c_s = 0$ 且 $a^{n-2r}c_r - (a')^{n-2r}c_r \neq 0$. 则 $Q := P \circ x_a^{-1} - P \circ x_{a'}^{-1} \in W - \{0\}$ 具有所要的性质.

现在假设 $W \neq 0$. 由上一段, 可以找到 $r \in \{0, \ldots, n\}$ 使得 $P := t_1^rt_2^{n-r} \in W$. 设 $x = \begin{pmatrix} a & b \\ -\bar{b} & \bar{a} \end{pmatrix}$, 其中 $a, b \in \CC$ 且 $|a|^2 + |b|^2 = 1$. 则 $x \in \SU(2)$ 且 $P \circ x^{-1} = (at_1 - bt_2)^r(\bar{b}t_1 + \bar{a}t_2)^{n-r} \in W$. 若把 $P \circ x^{-1}$ 写成 $P \circ x^{-1} = \sum_{i=0}^n c_it_1^it_2^{n-i}$, 则
\[
c_i = \sum_{j=\max(0,i-n+r)}^{\min(i,r)}(-1)^{r-j}\binom{r}{j}\binom{n-r}{i-j}a^{r-j}\bar{a}^{j}\bar{b}^{n-r+j-i}b^{i-j}.
\]
若取 $a = \frac{1}{\sqrt{1+t^2}}$ 且 $b = \frac{t}{\sqrt{1+t^2}}$ (其中 $t \in [-1, 1]$), 则每个 $c_i$ 是非零多项式在 $(1+t^2)^{n/2}$ 下的商, 所以只有有限多个 $t$ 使得 $c_i = 0$. 所以可以选择 $x \in \SU(2)$ 使得 $P \circ x^{-1}$ 的所有系数非零. 由第一段, 这蕴含每个单项式 $t_1^it_2^{n-i}$ ($0 \leq i \leq n$) 都在 $W$ 中. 所以 $W = V_n$.

(c). 记 $\pi_n$ 为映射 $\SU(2) \to \GL(V_n)$. 若 $n = 0$, 则 $V_n$ 是 $\SU(2)$ 的平凡表示, 所以 $\Ker(\pi_n) = \SU(2)$. 假设 $n \geq 1$, 设 $x = \begin{pmatrix} a & b \\ -\bar{b} & \bar{a} \end{pmatrix} \in \Ker(\pi_n)$. 特别地, 若 $P = t_1^n$, 则必须有 $P \circ x^{-1} = P$. 由于 $P \circ x^{-1} = \sum_{i=0}^n(-1)^{n-i}\binom{n}{i}a^i\bar{b}^{n-i}t_1^it_2^{n-i}$, 这蕴含 $a^n = 1$ 且对 $0 \leq i \leq n-1$ 有 $a^i\bar{b}^{n-i} = 0$. 特别地, $a \neq 0$, 所以必有 $b = 0$. 则 $a \in S^1$, 且对每个 $r \in \{0, \ldots, n\}$, $(t_1^rt_2^{n-r}) \circ x^{-1} = a^{n-2r}t_1^rt_2^{n-r}$, 所以 $a^{n-2r} = 1$. 若 $n$ 奇, 这蕴含 $a = 1$, 所以 $x = I_2$ 是 $\Ker(\pi_n)$ 的唯一元素. 若 $n$ 偶, 这只能蕴含 $a = \pm 1$, 所以 $x = \pm I_2$. 事实上, 若 $n$ 偶且非零, 容易验证 $-I_2$ 在 $V_n$ 上平凡作用, 所以 $\Ker(\pi_n) = \{\pm I_2\}$.
所以回答这个问题: 表示 $V_n$ 忠实当且仅当 $n$ 奇.
\end{solution}

\begin{exercise}[$\SL_2(\R)$ 的有限维酉表示平凡]
\label{ex:I.5.5.2}
设 $(\pi, V)$ 是 $G := \SL_2(\R)$ 的有限维酉表示. 我们要证明 $V$ 平凡 (即对每个 $x \in G$ 有 $\pi(x) = \mathrm{id}_V$).
\begin{enumerate}
\item[(a)] 考虑群同态 $\alpha \colon \R \to G$, $t \longmapsto \begin{pmatrix} 1 & t \\ 0 & 1 \end{pmatrix}$. 证明存在 $V$ 的基 $\mathcal{B}$ 和 $y_1, \ldots, y_n \in \R$ (其中 $n = \dim V$) 使得对每个 $t \in \R$, 自同构 $\pi(\alpha(t))$ 在 $\mathcal{B}$ 中对角, 对角线元素为 $e^{ity_1}, \ldots, e^{ity_n}$.
\item[(b)] 证明对每个 $t \in \R$ 有 $\pi(\alpha(t)) = \mathrm{id}_V$. (提示: 若 $u \in \R^{\times}$, $x = \begin{pmatrix} u & 0 \\ 0 & u^{-1} \end{pmatrix}$, 考虑 $x\alpha(t)x^{-1}$ 在 $V$ 上的作用.)
\item[(c)] 证明对每个 $x \in G$ 有 $\pi(x) = \mathrm{id}_V$.
\item[(d)] 若 $n \geq 3$, 证明 $\SL_n(\R)$ 的每个有限维酉表示都平凡.
\end{enumerate}
\end{exercise}

\begin{solution}
(a). $\GL(V)$ 的子群 $\pi(\alpha(\R))$ 交换, 且它的所有元素可对角化 (因为它们都酉), 所以可以找到 $V$ 的基 $\mathcal{B} = (v_1, \ldots, v_n)$ 使得 $\pi(\alpha(\R))$ 的所有元素在其中对角 (若需要还可以是标准正交基). (若不喜欢同时对角化 $\GL(V)$ 的无穷子集, 只需选取生成 $\operatorname{Span}(\pi(\alpha(\R)))$ 的 $A_1, \ldots, A_m \in \pi(\alpha(\R))$ 并同时对角化它们.)
对每个 $j \in \{1, \ldots, n\}$, 子空间 $\CC v_j$ 被 $\alpha(\R) \subset G$ 的作用稳定 (由基的选取), 所以我们得到 $\R$ 在 $\CC v_j$ 上的 $1$ 维表示, 且由习题 \ref{ex:I.5.4.1}(b) 知道这种表示具有 $t \longmapsto e^{ity_j}v_j$ 的形式 (对某个 $y_j \in \CC$).

(b). 我们有 $x\alpha(t)x^{-1} = \alpha(u^2t)$, 所以 $\pi(x\alpha(t)x^{-1})$ 在基 $\mathcal{B}$ 中对角, 对角线元素为 $e^{iu^2ty_1}, \ldots, e^{iu^2ty_n}$. 另一方面,
\[
\operatorname{Tr}(\pi(x\alpha(t)x^{-1})) = \operatorname{Tr}(\pi(x)\pi(\alpha(t))\pi(x)^{-1}) = \operatorname{Tr}(\pi(\alpha(t))),
\]
所以对每个 $t \in \R$ 和每个 $u \in \R^{\times}$,
\[
\sum_{j=1}^n e^{ity_j} = \sum_{j=1}^n e^{iu^2ty_j}.
\]
假设我们知道 $L^{\infty}(\R)$ 的子集 $\widehat{\R}$ 线性无关. 则此等式告诉我们: 对每个 $u \in \R^{\times}$, 集合 $\{y_1, \ldots, y_n\}$ 和 $\{u^2y_1, \ldots, u^2y_n\}$ 相等. 这只有在 $y_1 = \cdots = y_n = 0$ 时才可能, 而这又蕴含 $\alpha(t)$ 对每个 $t \in \R$ 在 $V$ 上平凡作用.
现在证明关于 $\widehat{\R}$ 的结论. 设 $y_1, \ldots, y_m \in \R$ 两两不同, $c_1, \ldots, c_m \in \CC$ 满足对每个 $t \in \R$ 有 $\sum_{j=1}^m c_je^{ity_j} = 0$. 我们要证明 $c_1 = \cdots = c_m = 0$. 设 $r \in \R$. 取 $t = 0, r, \ldots, r(m-1)$, 利用 Vandermonde 行列式的计算, 看到必须有 $e^{iry_1} = \cdots = e^{iry_m}$. 由于这对每个 $r \in \R$ 成立, 它蕴含 $y_1 = \cdots = y_m$ (例如对前面的等式关于 $r$ 求导然后在 $r = 0$ 处取值). 所以 $m = 1$, 然后 $c_1e^{ity_1} = 0$ (对每个 $t \in \R$) 这一事实蕴含 $c_1 = 0$.

(c). 若 $x \in \SL_2(\R)$ 是剪切 (transvection) 矩阵 (即形如 $\begin{pmatrix} 1 & t \\ 0 & 1 \end{pmatrix}$ 的矩阵的共轭), 则 $x = y\alpha(t)y^{-1}$ (对某个 $t \in \R$ 和 $y \in \SL_2(\R)$), 所以由 (b), $\pi(x) = \pi(y)\pi(\alpha(t))\pi(y)^{-1} = \pi(y)\pi(y)^{-1} = \mathrm{id}_V$. 由于 $\SL_2(\R)$ 由剪切矩阵生成, 这蕴含对每个 $x \in \SL_2(\R)$ 有 $\pi(x) = \mathrm{id}_V$.

(d). 设 $\pi \colon \SL_n(\R) \to \GL(V)$ 是有限维酉表示. 设 $x \in \SL_n(\R)$ 是剪切矩阵. 我们可以模仿 (a) 和 (b) 来证明 $\pi(x) = \mathrm{id}_V$, 但也可以这样做: 选取 $\R^n$ 的基 $(v_1, \ldots, v_n)$ 使得对应于 $x$ 的 $\R^n$ 的线性自同态的矩阵为
\[
\begin{pmatrix} 1 & 1 & 0 & \cdots & 0 \\ 0 & 1 & 0 & \cdots & 0 \\ 0 & 0 & 1 & \cdots & 0 \\ \vdots & \vdots & \vdots & \ddots & \vdots \\ 0 & \cdots & 0 & \cdots & 1 \end{pmatrix}.
\]
考虑 $\SL_n(\R)$ 的由在基 $(v_1, \ldots, v_n)$ 中矩阵形如
\[
\begin{pmatrix} a & b & 0 & \cdots & 0 \\ c & d & 0 & \cdots & 0 \\ 0 & 0 & 1 & \cdots & 0 \\ \vdots & \vdots & \vdots & \ddots & \vdots \\ 0 & \cdots & 0 & \cdots & 1 \end{pmatrix}
\]
($a, b, c, d \in \R$, $ad - bc = 1$) 的元素组成的子集 $G'$. 则 $G'$ 是子群, 且与 $\SL_2(\R)$ 同构. 由于 $\pi|_{G'}$ 是 $G'$ 在 $V$ 上的酉表示, 由 (c) 有 $G' \subset \Ker(\pi)$. 特别地, $\pi(x) = \mathrm{id}_V$.
现在用 $\SL_n(\R)$ 由剪切矩阵生成的事实得出结论: 对每个 $x \in \SL_n(\R)$ 有 $\pi(x) = \mathrm{id}_V$.
\end{solution}

\begin{exercise}[酉化、小球与忠实表示]
\label{ex:I.5.5.3}
\begin{enumerate}
\item[(a)] 设 $G$ 是 $\GL_n(\CC)$ 的紧子群. 证明存在 $x \in \GL_n(\CC)$ 使得 $xGx^{-1} \subset \UU(n)$.
\item[(b)] 在 $M_n(\CC)$ 上放一个你喜欢的范数 (反正它们都等价). 证明存在 $c > 0$ 使得 $\GL_n(\CC)$ 包含于球 $\{x \in \GL_n(\CC) \mid \|x - I_n\| < c\}$ 的唯一子群是平凡群.
\item[(c)] 证明对 $\GL_n(\Q_p)$ 在有限维 $\CC$-向量空间上的每个连续表示, 存在整数 $m \geq 0$ 使得 $\GL_n(\Q_p)$ 的子群 $I_n + p^mM_n(\Z_p)$ 平凡作用.
\item[(d)] 证明若 $(\pi, V)$ 是 $\GL_n(\Z_p)$ 的不可约酉表示, 则存在 $m \geq 1$ 使得 $\pi(I_n + p^mM_n(\Z_p)) = \{1\}$.
\item[(e)] 更一般地, 证明若 $G$ 是射影有限 (profinite) 群 (即有限离散群的投射极限, 见习题 \ref{ex:I.5.1.3}), 则只有当 $G$ 有限时 $G$ 才有忠实的不可约酉表示.
\end{enumerate}
\end{exercise}

\begin{solution}
(a). 考虑 $G$ 在 $\CC^n$ 上由包含 $G \subset \GL_n(\CC)$ 给出的表示 $\rho$. 由定理 \ref{thm:I.3.2.8}, 存在 $\CC^n$ 上的 Hermitian 内积使得该表示酉. 设 $A$ 是此 Hermitian 内积在 $\CC^n$ 的典范基下的矩阵. 则 $A$ 是 Hermitian 正定矩阵, 所以可以把它写成 $A = B^*B$, 其中 $B \in \GL_n(\CC)$. (这是谱定理的简单推论. 由于 $A$ Hermitian, 有酉矩阵 $P$ 和对角矩阵 $D$ 使得 $A = P^*DP$. 由于 $A$ 正定, $D$ 的对角线元素是正实数, 所以可以写成 $D = C^2$, 其中 $C$ 是另一个对角线元素为正的对角矩阵. 取 $B = P^*CP$, 则 $B$ 是 Hermitian 正定且 $A = B^2 = B^*B$.)
$\rho$ 对 $A$ 酉意味着对每个 $X \in G$ 有 $X^*AX = A$. 由于 $A = B^*B$, 这等价于 $(BXB^{-1})^*(BXB^{-1}) = I_n$. 所以 $BGB^{-1} \subset \UU(n)$.

(b). 设 $G$ 是 $\GL_n(\CC)$ 的包含于形如 $\{x \in \GL_n(\CC) \mid \|x - I_n\| < c\}$ 的球的子群. 则闭子群 $\overline{G}$ 包含于闭球 $\{x \in \GL_n(\CC) \mid \|x - I_n\| \leq c\}$, 所以它紧, 从而由 (a) 包含于形如 $P\UU(n)P^{-1}$ 的子群. 特别地, $G$ 的每个元素可对角化且所有特征值的模为 $1$.
在 $\CC^n$ 上固定任何范数, 考虑 $M_n(\CC)$ 上对应的算子范数 $\|\cdot\|$. 我们将使用此范数. 注意若 $X \in M_n(\CC)$ 且 $\lambda$ 是 $X$ 的特征值, 则存在范数 $1$ 的向量 $v \in \CC^n$ 使得 $Xv = \lambda v$, 所以 $\|X\| \geq |\lambda|$. 现在让我们证明 $\GL_n(\CC)$ 的每个包含于开球 $B := \{x \in \GL_n(\CC) \mid \|x - I_n\| < \sqrt{2}\}$ 的子群都平凡. 设 $G$ 是这种子群, $X \in G$. 设 $\lambda_1, \ldots, \lambda_n$ 是 $X$ 的特征值. 刚才已见 $|\lambda_1| = \cdots = |\lambda_n| = 1$. 假设存在 $r \in \{1, \ldots, n\}$ 使得 $\lambda_r \neq 1$, 则可以把 $\lambda_r$ 写成 $\lambda_r = e^{i\theta}$, 其中 $-\pi/2 < \theta < \pi/2$, 因为 $|\lambda_r - 1| \leq \|X - I_n\| < \sqrt{2}$; 但此时, 若选取整数 $m \geq 1$ 使得 $\pi/2 \leq m|\theta| \leq \pi$, 将有 $\|X^m - I_n\| \geq |\lambda_r^m - 1| \geq \sqrt{2}$, 这与 $X^m \in G$ 矛盾. 所以必有 $\lambda_1 = \cdots = \lambda_n = 1$, 这意味着 $X = I_n$.

(c). 我们已在习题 \ref{ex:I.5.2.2}(d) 的解答中看到 $K_m := I_n + p^mM_n(\Z_p)$ 确实是 $\GL_n(\Q_p)$ 的子群. 我们还在 $M_n(\Q_p)$ 上放了范数使得 $K_m$ 是以 $I_n$ 为心、$p^{-m+1}$ 为半径的开球.
设 $\rho \colon \GL_n(\Q_p) \to \GL(V)$ 是 $\GL_n(\Q_p)$ 在有限维向量空间 $V$ 上的连续表示. 由命题 \ref{prop:I.3.5.1}, 同态 $\rho$ 连续. 设 $U$ 是 $\mathrm{id}_V$ 在 $\GL(V)$ 中的开邻域, 使得 $\GL(V)$ 包含于 $U$ 的唯一子群是 $\{1\}$ (由 (b) 这存在). 则 $\rho^{-1}(U)$ 是 $I_n$ 在 $\GL_n(\Q_p)$ 中的开邻域, 所以对 $m \gg 0$ 它包含 $K_m$. 但 $K_m$ 是 $\GL_n(\Q_p)$ 的子群, 所以 $\rho(K_m)$ 是 $\GL(V)$ 的子群, 故一旦 $\rho(K_m) \subset U$ 就有 $\rho(K_m) = \{1\}$.

(d). 由习题 \ref{ex:I.5.1.4}(l), 群 $\GL_n(\Z_p)$ 紧. 所以由习题 \ref{ex:I.5.5.9}, 空间 $V$ 有限维. 现在结论的证明与习题 \ref{ex:I.5.5.3}(c) 完全相同.

(e). 由习题 \ref{ex:I.5.1.3}, 我们知道 $G$ 紧 Hausdorff (注意有限离散群紧 Hausdorff). 所以由习题 \ref{ex:I.5.5.9}, $G$ 的每个不可约酉表示有限维.
假设我们知道 $G$ 完全不连通. 设 $(\pi, V)$ 是 $G$ 的连续有限维表示. 由习题 \ref{ex:I.5.2.2}(c), $G$ 的紧开子群构成 $1$ 的邻域基. 由习题 \ref{ex:I.5.5.3}(b), 可以找到 $\mathrm{id}_V$ 在 $\GL(V)$ 中的邻域 $U$ 使得 $\GL(V)$ 包含于 $U$ 的唯一子群是 $\{\mathrm{id}_V\}$. 所以若选取紧开子群 $K$ 使得 $\pi(K) \subset U$, 则必有 $K \subset \Ker(\pi)$. 所以 $\Ker(\pi) = \bigcup_{x \in \Ker(\pi)}xK$ 是 $G$ 的开子群, 从而群 $G/\Ker(\pi)$ 离散. 由于它也紧, 它是有限群. 这证明除非 $G$ 有限, $G$ 不能有忠实的不可约酉表示.
所以还需证明 $G$ 完全不连通, 即 $G$ 的唯一的非空连通子集是单点集. 取有限群的投射系统 $((G_i)_{i \in I}, (u_{ij} \colon G_i \to G_j)_{i \geq j})$ 使得 $G = \varprojlim_{i \in I}G_i$. 设 $C \subset G$ 是非空连通子集. 则 $C$ 在每个 $G_i$ 中的像是非空连通的, 从而是单点集 $\{g_i\}$. 这蕴含 $C$ 的唯一元素是族 $(g_i)_{i \in I} \in \prod_{i \in I}G_i$ (此族自动在投射极限中).
\end{solution}

\begin{remarkstar}
有限群并不总有忠实的不可约酉表示. 例如, 若 $G$ 是有限交换群, 则由 Schur 引理, $G$ 的每个不可约酉表示都是 $1$ 维的. 取 $G = \Z/2\Z \times \Z/2\Z$, 设 $\pi \colon G \to \CC$ 是 $1$ 维表示. 由于 $(1, 0)$ 和 $(0, 1)$ 在 $G$ 中阶为 $2$, 必须有 $\pi(1, 0), \pi(0, 1) \in \{\pm 1\}$. 若想要 $\pi$ 忠实, 需要 $\pi(1, 0) = \pi(0, 1) = -1$, 但此时 $\pi(1, 1) = (-1)^2 = 1$, 所以 $\pi$ 不能忠实. (更一般地, 有限交换群有忠实的 $1$ 维表示当且仅当它循环.)
\end{remarkstar}

\begin{exercise}[闭子群的 Lie 代数]
\label{ex:I.5.5.4}
本题的目标是定义 $\GL_n(\CC)$ 的闭子群 $G$ 的 Lie 代数, 并证明矩阵指数在它们恒等元附近诱导 Lie 代数与群之间的同胚.
考虑函数 $\exp \colon M_n(\CC) \to M_n(\CC)$, 定义为 $\exp(X) = \sum_{n \geq 0}\frac{1}{n!}X^n$; 也记 $e^X = \exp(X)$. 你可以假定此函数的如下基本性质: 定义它的级数绝对收敛, 它无穷次可微, 且可以逐项求导. 你也可以假定对任何满足 $AB = BA$ 的 $A, B \in M_n(\CC)$, 有 $\exp(A + B) = \exp(A)\exp(B)$; 特别地, 对每个 $X \in M_n(\CC)$ 有 $\exp(X) \in \GL_n(\CC)$ 且 $\exp(X)^{-1} = \exp(-X)$.
我们还固定 $\GL_n(\CC)$ 的闭子群 $G$.
\begin{enumerate}
\item[(a)] 计算 $\exp$ 在点 $0 \in M_n(\CC)$ 处的微分. (回忆这是从 $M_n(\CC)$ 到自身的线性算子.)
\item[(b)] 证明 $\exp$ 诱导从 $M_n(\CC)$ 中 $0$ 的邻域到 $\GL_n(\CC)$ 中 $1$ 的邻域的微分同胚.
\item[(c)] 设 $L = \{X \in M_n(\CC) \mid \forall t \in \R,\, \exp(tX) \in G\}$. 证明 $L$ 是 $M_n(\CC)$ 的 $\R$-线性子空间. (提示: 对所有 $X, Y \in M_n(\CC)$, 证明 $\exp(X + Y) = \lim_{k \to +\infty}\bigl(\exp(\frac{1}{k}X)\exp(\frac{1}{k}Y)\bigr)^k$.)\footnote{这个 $L$ 称为 $G$ 的 Lie 代数. 也容易证明它对换子括号 $[X, Y] = XY - YX$ 稳定.}
\item[(d)] 若 $G = \UU(n)$, 证明 $L = \{X \in M_n(\CC) \mid X^* = -X\}$.
\item[(e)] 若 $G = \SO(n)$, 证明 $L = \{X \in M_n(\CC) \mid X^T = -X\}$.
\item[(f)] 再次假设 $G$ 是 $\GL_n(\CC)$ 的任意闭子群. 本题的目标是证明如下结论 $(*)$: 存在 $L$ 中 $0$ 的邻域 $U$ 使得 $\exp(U)$ 是 $1$ 在 $G$ 中的邻域, 且 $\exp$ 诱导同胚 $U \xrightarrow{\sim} \exp(U)$.
  \begin{enumerate}
  \item[(i)] 设 $L'$ 是 $M_n(\CC)$ 的 $\R$-线性子空间使得 $M_n(\CC) = L \oplus L'$, 考虑函数 $\phi \colon M_n(\CC) \to \GL_n(\CC)$, 定义为 $\phi(A + B) = e^Ae^B$ (对每个 $A \in L$, $B \in L'$). 证明存在 $0$ 在 $L$ 中的邻域 $U_0$、$0$ 在 $L'$ 中的邻域 $V$ 和 $1$ 在 $\GL_n(\CC)$ 中的邻域 $W$ 使得 $\phi$ 诱导微分同胚 $U_0 \times V \xrightarrow{\sim} W$.
  \item[(ii)] 假设 $(*)$ 不成立. 证明存在 $L$ 中 $0$ 的递减邻域序列 $U_0 \supset U_1 \supset \cdots$, $L$ 的元素序列 $(A_k)_{k \geq 0}$ 和 $L'$ 的元素序列 $(B_k)_{k \geq 0}$ 使得:
  \begin{itemize}
  \item 对每个 $k \geq 0$, $A_k \in U_k$;
  \item 对每个 $k \geq 0$, $B_k \neq 0$;
  \item 对每个 $k \geq 0$, $\phi(A_k + B_k) \in G$;
  \item 序列 $(B_k)_{k \geq 0}$ 的极限是 $0$;
  \item 对 $L$ 中 $0$ 的每个邻域 $U$, 对足够大的 $k$ 有 $U_k \subset U$.
  \end{itemize}
  \item[(iii)] 证明序列 $(\frac{1}{\|B_k\|}B_k)_{k \geq 0}$ 有收敛子序列, 且该子序列的极限 $B$ 非零.
  \item[(iv)] 对每个 $t \in \R$, 证明 $\lfloor\frac{t}{\|B_k\|}\rfloor\|B_k\| \to t$ ($k \to +\infty$). (其中对每个 $c \in \R$, 记 $\lfloor c \rfloor$ 为不超过 $c$ 的最大整数.)
  \item[(v)] 证明 $B \in L$.
  \end{enumerate}
\item[(g)] 设 $(\rho, V)$ 是 $G$ 的连续有限维表示. 对每个 $X \in L$, 证明存在唯一的元素 $u(X) \in \End(V)$ 使得对每个 $t \in \R$ 有 $\rho(e^{tX}) = e^{tu(X)}$. 又证明函数 $u \colon L \to \End(V)$ 是 $\R$-线性的.\footnote{函数 $u \colon L \to \End(V)$ 称为 $\rho$ 在 $0$ 处的微分. 稍加努力可以证明它保持换子括号.}
\end{enumerate}
\end{exercise}

\begin{solution}
(a). 设 $d\exp_0$ 是 $\exp$ 在点 $0$ 处的微分. 由微分的定义, 对每个 $H \in M_n(\CC)$,
\[
d\exp_0(H) = \lim_{t \to 0}\frac{1}{t}(e^{tH} - e^0) = \lim_{t \to 0}\frac{1}{t}(e^{tH} - I_n) = \frac{d}{dt}e^{tH}\Big|_{t=0}.
\]
但
\[
\frac{d}{dt}e^{tH} = \sum_{n \geq 0}\frac{1}{n!}\frac{d}{dt}(tH)^n = H\exp(tH) = \exp(tH)H,
\]
所以 $d\exp_0(H) = H$. 最后得到 $d\exp_0 = \mathrm{id}_{M_n(\CC)}$.

(b). 由于 $\GL_n(\CC)$ 在 $M_n(\CC)$ 中开, $1$ 在 $\GL_n(\CC)$ 中的邻域与足够小的 $1$ 在 $M_n(\CC)$ 中的邻域相同. 所以结论由 $d\exp_0$ 可逆和反函数定理得出.

(c). 先证明提示. 设 $U$ 是 $0$ 在 $M_n(\CC)$ 中的邻域, $V$ 是 $1$ 在 $\GL_n(\CC)$ 中的邻域, 使得 $\exp$ 是从 $U$ 到 $V$ 的微分同胚. 记 $\log \colon V \to U$ 为它的逆. 由于当 $H \to 0$ 时 $\exp(H) = 1 + H + o(H)$, 当 $H \to 0$ 时有 $\log(1 + H) = H + o(H)$.
设 $X, Y \in M_n(\CC)$. 则 $\exp(\frac{1}{k}X) = 1 + \frac{1}{k}X + O(\frac{1}{k^2})$, $\exp(\frac{1}{k}Y) = 1 + \frac{1}{k}Y + O(\frac{1}{k^2})$, 所以 $\exp(\frac{1}{k}X)\exp(\frac{1}{k}Y) = 1 + \frac{1}{k}(X + Y) + O(\frac{1}{k^2})$. 若 $k$ 足够大, 则 $\exp(\frac{1}{k}X)\exp(\frac{1}{k}Y) \in V$, 且 $\log(\exp(\frac{1}{k}X)\exp(\frac{1}{k}Y)) = \frac{1}{k}(X + Y) + o(\frac{1}{k})$. 所以最终
\[
(\exp(\frac{1}{k}X)\exp(\frac{1}{k}Y))^k = \exp(k\log(\exp(\frac{1}{k}X)\exp(\frac{1}{k}Y))) = \exp(X + Y + o(1)) \to \exp(X + Y)
\]
当 $k \to +\infty$ 时.
集合 $L$ 由定义对纯量乘法稳定. 设 $X, Y \in L$. 则对每个 $t \in \R$ 有 $c(t) := e^{tX}e^{tY} \in G$. 由于 $G$ 在 $\GL_n(\CC)$ 中闭, 这蕴含对每个 $t \in \R$,
\[
\exp(t(X + Y)) = \lim_{k \to +\infty}(\exp(\frac{t}{k}X)\exp(\frac{t}{k}Y))^k \in G.
\]
所以 $X + Y \in L$.

(d). 设 $X \in M_n(\CC)$ 满足 $X = -X^*$. 则对每个 $t \in \R$ 有 $tX = -(tX)^*$ (特别地 $tX$ 与 $tX^*$ 交换), 所以
\[
e^{tX}(e^{tX})^* = e^{tX}e^{tX^*} = e^{tX + tX^*} = e^0 = I_n,
\]
即 $e^{tX} \in \UU(n)$. 所以 $X \in L$.
反之, 设 $X \in L$. 则对每个 $t \in \R$ 有 $e^{tX}e^{tX^*} = I_n$. 对此表达式求导 (并使用 (a) 的证明中的导数表达式) 给出
\[
0 = Xe^{tX}e^{tX^*} + e^{tX}e^{tX^*}X^* = X + X^*.
\]

(e). 这与 (d) 的证明完全相同, 只需把 ``$*$'' 换成 ``$T$''.

(f). (i) 由反函数定理, 只需证明 $\phi$ 在 $0 \in M_n(\CC)$ 处的微分可逆. 若 $A \in L$, $B \in L'$, 由微分的定义
\[
d\phi_0(A + B) = \lim_{t \to 0}\frac{1}{t}(\phi(tA + tB) - \phi(0)) = \lim_{t \to 0}\frac{1}{t}(e^{tA}e^{tB} - 1) = A + B
\]
(由 (a) 中的计算), 所以 $d\phi_0 = \mathrm{id}_{M_n(\CC)}$, 它当然可逆.

(ii) 选取 $0$ 在 $L$ (相应地 $L'$) 中的邻域序列 $U_0 \supset U_1 \supset \cdots$ (相应地 $V = V_0 \supset V_1 \supset \cdots$) 使得 $L$ (相应地 $L'$) 中 $0$ 的每个邻域 $U$ (相应地 $V'$) 对足够大的 $k$ 包含 $U_k$ (相应地 $V_k$). (例如可以取 $L$ 和 $L'$ 中半径趋于 $0$ 的球.) 对每个 $k \geq 0$, 函数 $\phi$ 是从 $U_k \times V_k$ 到 $\phi(U_k \times V_k)$ 的微分同胚, 特别地 $\phi(U_k \times V_k) \cap G$ 是 $1$ 在 $G$ 中的邻域, 包含 $\exp(U_k)$. 若 $(*)$ 不成立, 则对每个 $k$, $\phi(U_k \times V_k) \cap G$ 严格包含 $\exp(U_k)$, 所以可以找到 $A_k \in U_k$ 和 $B_k \in V_k$ 使得 $\phi(A_k + B_k) \in G$ 且 $\phi(A_k + B_k) \notin \exp(U_k)$, 即 $B_k \neq 0$. 又, 由对邻域 $V_k$ 的条件, 当 $k \to +\infty$ 时 $B_k \to 0$.

(iii) 序列 $(\frac{1}{\|B_k\|}B_k)_{k \geq 0}$ 是 $L'$ 的单位球的元素序列, 而单位球紧, 所以它有收敛子序列, 其极限仍在单位球中 (且特别地非零).

(iv) 对每个 $k \geq 0$,
\[
0 \leq \frac{t}{\|B_k\|} - \Bigl\lfloor\frac{t}{\|B_k\|}\Bigr\rfloor < 1,
\]
所以
\[
0 \leq t - \Bigl\lfloor\frac{t}{\|B_k\|}\Bigr\rfloor\|B_k\| < \|B_k\|.
\]
由于 $B_k \to 0$, 有 $\|B_k\| \to 0$, 这蕴含
\[
\Bigl\lfloor\frac{t}{\|B_k\|}\Bigr\rfloor\|B_k\| \to t.
\]

(v) 取子序列后, 可以假设 $B = \lim_{k \to +\infty}\frac{1}{\|B_k\|}B_k$. 我们必须证明对每个 $t \in \R$ 有 $e^{tB} \in G$. 设 $t \in \R$. 由 (iv) 和 $B$ 的定义,
\[
e^{tB} = \lim_{k \to +\infty}\exp\Bigl(\Bigl\lfloor\frac{t}{\|B_k\|}\Bigr\rfloor\|B_k\|\frac{1}{\|B_k\|}B_k\Bigr) = \exp\Bigl(\Bigl\lfloor\frac{t}{\|B_k\|}\Bigr\rfloor B_k\Bigr).
\]
但由 (ii), 对每个 $k \geq 0$ 有 $\phi(A_k + B_k) = e^{A_k}e^{B_k} \in G$ 且 $e^{A_k} \in G$ (因为 $A_k \in L$), 所以 $e^{B_k} \in G$; 由于 $N := \lfloor\frac{t}{\|B_k\|}\rfloor$ 是整数, 这蕴含 $e^{NB_k} = (e^{B_k})^N \in G$. 最后, 由于 $G$ 在 $\GL_n(\CC)$ 中闭, 我们推出 $e^{tB} \in G$.

(g). 设 $X \in L$. 考虑映射 $\R \to \GL(V)$, $t \longmapsto \rho(e^{tX})$. 这是连续群同态, 所以由习题 \ref{ex:I.5.4.1}(b)(i), 存在唯一的 $u(X) \in \End(V)$ 使得对每个 $t \in \R$ 有 $\rho(e^{tX}) = \exp(tu(X))$.
设 $X, Y \in L$, $a \in \R$. 对每个 $t \in \R$,
\[
e^{tu(aX)} = \rho(e^{taX}) = e^{tau(X)}.
\]
在 $t = 0$ 处求导, 得到 $u(aX) = au(X)$. 现在考虑 $c \colon \R \to \GL(V)$, $t \longmapsto \rho(e^{tX})\rho(e^{tY})\rho(e^{-t(X+Y)})$. 有 $c(t) = e^{tu(X)}e^{tu(Y)}e^{-tu(X+Y)}$, 所以 $c$ 是 $\mathcal{C}^{\infty}$ 的且 $c'(0) = u(X) + u(Y) - u(X+Y)$. 另一方面, 利用 $c$ 是 $\mathcal{C}^{\infty}$ 的这一事实, 可以像 (c) 那样证明对每个 $t \in \R$,
\[
\lim_{k \to +\infty}c(\frac{t}{k})^k = e^{tc'(0)}.
\]
所以我们只需证明这个极限对每个 $t \in \R$ 等于 $\mathrm{id}_V$. 用无穷小量做简单计算表明 (若 $t$ 固定)
\[
e^{\frac{t}{k}X}e^{\frac{t}{k}Y}e^{-\frac{t}{k}(X+Y)} = I_n + O(\frac{1}{k^2}),
\]
所以
\[
(e^{\frac{t}{k}X}e^{\frac{t}{k}Y}e^{-\frac{t}{k}(X+Y)})^k = I_n + O(\frac{1}{k}),
\]
且
\[
c(\frac{t}{k}) = \rho((e^{\frac{t}{k}X}e^{\frac{t}{k}Y}e^{-\frac{t}{k}(X+Y)})^k) \xrightarrow[k \to +\infty]{} \rho(I_n) = \mathrm{id}_V.
\]
\end{solution}

我们现在来看如何找完全没有忠实有限维表示 (在任何域上) 的离散群.
设 $\Gamma$ 是一个 (离散) 群. 我们说 $\Gamma$ 剩余有限 (residually finite), 如果对每个 $x \in \Gamma - \{1\}$, 存在 $\Gamma$ 的正规子群 $\Delta$ 使得 $\Gamma/\Delta$ 有限 (我们说 $\Delta$ 在 $\Gamma$ 中有限指标) 且 $x$ 在 $\Gamma/\Delta$ 中的像不平凡.\footnote{事实上, 可以用类似的想法证明: 若 $\operatorname{char}(k) = 0$, 这样的 $\Gamma$ 对几乎每个素数 $p$ 必须实质上剩余 $p$-有限 (virtually residually p-finite), 即它有一个有限指标子群 $\Gamma'$ 使得对每个 $x \in \Gamma' - \{1\}$, 存在 $\Gamma'$ 的有限指标正规子群 $\Delta \not\ni x$ 使得 $\Gamma'/\Delta$ 是 $p$-群; 但我所知道的唯一证明使用 Noether 正规化定理.}\footnote{加一个参考文献!}
下面两个习题的目标是证明: 若 $k$ 是域且 $\Gamma \subset \GL_n(k)$ 是有限生成子群, 则 $\Gamma$ 剩余有限.

\begin{exercise}[$\GL_n(R)$ 的剩余有限性]
\label{ex:I.5.5.5}
设 $R$ 是有限生成的 $\Z$-代数且是整环. 固定整数 $n \geq 1$. 对 $R$ 的每个理想 $I$, 令
\[
\Gamma(I) = \Ker(\GL_n(R) \to \GL_n(R/I)).
\]
\begin{enumerate}
\item[(a)] 证明 $R$ 是域当且仅当 $R$ 有限.
\item[(b)] 若 $\mathfrak{m}$ 是 $R$ 的极大理想, 证明 $\Gamma(\mathfrak{m})$ 是 $\GL_n(R)$ 的有限指标正规子群.
\item[(c)] 证明 $R$ 的所有极大理想的交是 $0$. (提示: 可以假设 $R$ 不是域. 若 $a \in R - \{0\}$, 证明局部化 $R[1/a]$ 不是域, 在 $R[1/a]$ 中取极大理想, 并与 $R$ 相交.)
\item[(d)] 证明 $\GL_n(R)$ 剩余有限.
\end{enumerate}
\end{exercise}

\begin{solution}
(a). 有限整环必是域, 这是经典事实. 证明如下. 假设 $R$ 有限, 设 $a \in R - \{0\}$. 则乘以 $a$ 是从 $R$ 到自身的加法映射, 其核为 $\{0\}$ (因为 $R$ 是整环), 所以它单射; 由于 $R$ 有限, 它满射, 这意味着存在 $b \in R$ 使得 $ab = 1$, 即 $a \in R^{\times}$.
其逆由交换代数的两个经典结果得出 (例如见 Eisenbud 的书 {[}9{]} 的习题 4.30 和 4.32):
\begin{itemize}
\item 若 $K \subset L$ 是域扩张且 $L$ 作为 $K$-代数有限生成, 则 $L$ 是有限维 $K$-向量空间 (Zariski 引理).
\item 若 $R$ 是 Noether 环, $S$ 是有限生成的 $R$-代数, $T \subset S$ 是 $R$-子代数且 $S$ 是有限 $T$-代数 (即作为 $T$-模有限生成), 则 $T$ 是有限生成的 $R$-代数 (Artin-Tate).
\end{itemize}
事实上, 若 $R$ 是域, 考虑它的素域 $k$. 则 $R$ 是有限生成的 $k$-代数, 所以由 Zariski 引理是有限维 $k$-向量空间, 这由第二个结果蕴含 $k$ 是有限生成的 $\Z$-代数. 注意 $k$ 要么是 $\Q$ 要么是某个有限域 $\mathbb{F}_p$. 但 $\Q$ 不是有限生成的 $\Z$-代数 (若 $x_1, \ldots, x_n \in \Q - \{0\}$, $P$ 是整除某个 $x_i$ 的分母的 (有限) 素数集合, 则由 $x_1, \ldots, x_n$ 生成的 $\Z$-子代数的非零元素的分母的素因子也必在 $P$ 中, 所以这个 $\Z$-子代数不能等于 $\Q$). 所以 $k$ 是有限域; 由于 $R$ 是有限维 $k$-向量空间, 它也是有限域.
为完整起见, 让我们给出我们实际需要的交换代数结果部分的证明. 假设我们知道:
$(*)$ 设 $L/K$ 是域扩张, 满足:
\begin{itemize}
\item 存在 $u \in L$ 使得 $L = K(u)$ (即 $L$ 作为域由 $u$ 生成);
\item $L$ 是有限生成的 $\Z$-代数,
\end{itemize}
则扩张有限且 $K$ 也是有限生成的 $\Z$-代数.
则可以用同样的方法证明: 若 $R$ 是域, 它必有限. (只需选取生成 $R$ (在其素域 $k$ 上) 的元素 $x_1, \ldots, x_n \in R$, 并对扩张 $k(x_1, \ldots, x_{i-1}) \subset k(x_1, \ldots, x_i)$ 应用 $(*)$ 证明 $k$ 是有限生成的 $\Z$-代数. 证明的结尾如前.)
我们现在证明 $(*)$. 设 $x_1, \ldots, x_n \in L^{\times}$ 生成 $L$ (作为 $\Z$-代数). 假设 $u$ 在 $K$ 上超越; 则 $L$ 同构于 $K$ 的有理分式域. 把 $x_i$ 写成 $x_i = \frac{P_i}{Q_i}$, 其中 $P_i, Q_i \in K[u]$. 由于 $(1 + u\prod_{i=1}^nQ_i)^{-1} = L = \Z[a_1, \ldots, a_n]$, 可以写成
\[
(1 + u\prod_{i=1}^nQ_i)^{-1} = \frac{R}{Q_1^{d_1}\cdots Q_n^{d_n}},
\]
其中 $R \in K[u]$ 与所有 $Q_i$ 互素, $d_1, \ldots, d_n \in \Z_{\geq 0}$. 得到 $Q_1^{d_1}\cdots Q_n^{d_n} = R(1 + u\prod_{i=1}^nQ_i)$, 这与 $R$ 与所有 $Q_i$ 互素矛盾. 所以 $u$ 在 $K$ 上代数. 设 $X^d + a_1X^{d-1} + \cdots + a_d \in K[X]$ 是 $u$ 在 $K$ 上的最小多项式. 对每个 $i \in \{1, \ldots, n\}$, 把 $x_i$ 写成 $x_i = \sum_{j=0}^{d-1}b_{ij}u^j$, 其中 $b_{ij} \in K$. 设 $A$ 是由 $a_1, \ldots, a_d$ 和所有 $b_{ij}$ 生成的 $K$ 的 $\Z$-子代数, 让我们证明 $A = K$. 设 $y \in K$. 则 $y$ 可以写成系数在 $\Z$ 中的 $x_1, \ldots, x_d$ 的多项式, 所以它也是系数在 $A$ 中的 $u$ 的多项式, 且可以取次数 $\leq d-1$ (可以用关系 $u^d = -a_1u^{d-1} - \cdots - a_d$ 把任何次数 $\geq d$ 的项换成更低次数的项). 换句话说, 可以把 $y$ 写成 $y = \sum_{i=0}^{d-1}c_iu^i$, 其中 $c_0, \ldots, c_{d-1} \in A$. 由于族 $(1, u, \ldots, u^{d-1})$ 在 $K$ 上线性无关, 必有 $c_1 = \cdots = c_{d-1} = 0$ 且 $y = c_0 \in A$.

(b). 首先, 群 $\Gamma(I)$ 是 $\GL_n(R)$ 的正规子群 (对任何理想 $I$), 因为它是群同态的核. 假设 $\mathfrak{m}$ 是极大理想. 则由 (a), $R/\mathfrak{m}$ 是有限域. 由于 $\GL_n(R)/\Gamma(\mathfrak{m})$ 单射进 $\GL_n(R/\mathfrak{m})$, 这蕴含 $\Gamma(\mathfrak{m})$ 在 $\GL_n(R)$ 中有限指标.

(c). 若 $R$ 是域, 则 $(0)$ 是 $R$ 的极大理想, 我们完成了. 假设 $R$ 不是域; 特别地, 由 (a), 它不有限. 设 $a \in R - \{0\}$. 局部化 $R[1/a] := R[X]/(aX - 1)$ 是有限生成的 $\Z$-代数 (因为 $R$ 是), 所以由 (a) 它只有有限时才是域. 但显然的映射 $R \to R[1/a]$ 单射 (因为 $a$ 不是零因子, 回忆 $R$ 是整环), 且 $R$ 无穷, 所以 $R[1/a]$ 也无穷, 故它不是域. 设 $\mathfrak{m}'$ 是 $R[1/a]$ 的极大理想, $\mathfrak{m}$ 是它在 $R$ 中的原像. 则映射 $R/\mathfrak{m} \to R[1/a]/\mathfrak{m}'$ 单射 (因为 $R \to R[1/a]$ 是), 且 $R[1/a]/\mathfrak{m}'$ 有限 (因为由 (a) 它是域), 所以 $R/\mathfrak{m}$ 有限且是整环, 因此是域 (再次由 (a)!), $\mathfrak{m}$ 是 $R$ 的极大理想. 又注意, 由于 $a$ 在 $R[1/a]$ 中可逆, 它不能在 $\mathfrak{m}'$ 中, 所以也不能在 $\mathfrak{m}$ 中. 所以我们找到了不包含 $a$ 的 $R$ 的极大理想.

(d). 设 $x = (x_{ij})_{1 \leq i,j \leq n} \in \GL_n(R)$ 满足 $x \neq I_n$. 选取 $i, j \in \{1, \ldots, n\}$ 使得当 $i \neq j$ 时 $x_{ij} \neq 0$, 当 $i = j$ 时 $x_{ij} \neq 1$. 由 (c), 可以找到 $R$ 的极大理想 $\mathfrak{m}$ 使得当 $i \neq j$ 时 $x_{ij} \notin \mathfrak{m}$, 当 $i = j$ 时 $x_{ij} - 1 \notin \mathfrak{m}$. 换句话说, $x$ 在 $\GL_n(R)/\Gamma(\mathfrak{m})$ 中的像不是单位元. 由于由 (b) $\Gamma(\mathfrak{m})$ 是 $\GL_n(R)$ 的有限指标正规子群, 我们完成了.
\end{solution}

\begin{exercise}[有限生成子群剩余有限]
\label{ex:I.5.5.6}
设 $k$ 是域, $\Gamma$ 是 $\GL_n(k)$ 的有限生成子群.
\begin{enumerate}
\item[(a)] 证明存在 $k$ 的有限生成 $\Z$-子代数 $R$ 使得 $\Gamma \subset \GL_n(R)$.
\item[(b)] 证明 $\Gamma$ 剩余有限.
\end{enumerate}
\end{exercise}

\begin{solution}
(a). 设 $\gamma_1, \ldots, \gamma_n$ 是 $\Gamma$ 的生成元, 设 $R$ 是由 $\gamma_i$ 及其逆的元素生成的 $k$ 的 $\Z$-子代数; 由定义这是有限生成的 $\Z$-代数. 由于 $\Gamma$ 的每个元素都是元素 $\gamma_i^{\pm 1}$ 的乘积, 有 $\Gamma \subset \GL_n(R)$.

(b). 这由习题 \ref{ex:I.5.5.5}(d) 立即得出: 若 $\gamma \in \Gamma - \{1\}$, 选取 $\GL_n(R)$ 的有限指标正规子群 $\Delta$ 使得 $\gamma$ 在 $\GL_n(R)/\Delta$ 中的像不平凡. 则 $\Gamma \cap \Delta$ 是 $\Gamma$ 的正规子群, 且 $\Gamma/(\Gamma \cap \Delta)$ 单射进 $\GL_n(R)/\Delta$, 所以 $\Gamma \cap \Delta$ 在 $\Gamma$ 中有限指标且 $\gamma$ 在 $\Gamma/(\Gamma \cap \Delta)$ 中的像不平凡.
\end{solution}

当然, 如果给不出任何有限生成的非剩余有限群的例子, 前一题的结果就不太有意思了. 我们在下面两个习题中给出这种例子.

\begin{exercise}[Baumslag-Solitar 群 $BS(2,3)$ 不剩余有限]
\label{ex:I.5.5.7}
设 $\Gamma$ 是自由群 (生成元 $a$ 和 $b$) 关于关系 $a^{-1}b^2a = b^3$ 的商. 在此题中, 我们假设 $b_1 := a^{-1}ba$ 与 $b$ 在 $\Gamma$ 中不交换, 并推出 $\Gamma$ 不剩余有限.
设 $u \colon \Gamma \to \Gamma'$ 是群同态, 其中 $\Gamma'$ 有限.
\begin{enumerate}
\item[(a)] 设 $n$ 是 $u(a)$ 在 $\Gamma'$ 中的阶. 证明 $u(b)$ 的阶整除 $3^n - 2^n$.
\item[(b)] 证明存在整数 $N \geq 0$ 使得 $u(b_1) = u(b_1)^{2^N}$. (注意 $u(b)$ 的阶与 $2$ 和 $3$ 都互素.)
\item[(c)] 证明 $u(b_1)$ 与 $u(b)$ 交换.
\item[(d)] 证明 $\Gamma$ 不剩余有限.
\end{enumerate}
\end{exercise}

\begin{solution}
(a). 我们首先证明对每个 $r \in \Z_{\geq 0}$ 有 $b^{2^r} = a^rb^{3^r}a^{-r}$. $r = 0$ 的情形显然, $r = 1$ 的情形就是定义 $\Gamma$ 的关系. 设 $r \geq 1$, 假设对 $r$ 已知结论, 证明对 $r+1$. 有
\[
b^{2^{r+1}} = (b^{2^r})^2 = (ab^{3^r}a^{-1})^2 = (ab^{2^r}a^{-1})^3 = (a^{r+1}b^{3^r}a^{-(r+1)})^3 = a^{r+1}b^{3^{r+1}}a^{-(r+1)}.
\]
应用于 $r = n$ 给出 $b^{2^n} = a^nb^{3^n}a^{-n}$, 所以 $u(b)^{3^n - 2^n} = 1$, 故 $u(b)$ 的阶整除 $3^n - 2^n$.

(b). 注意 $u(b)$ 的阶奇, 因为它整除奇数 $3^n - 2^n$. 所以存在 $N \geq 1$ 使得 $u(b)^{2^N} = u(b)$. 由于 $b_1 = a^{-1}ba$, 对每个 $r \geq 0$ 有 $b_1^r = a^{-1}b^ra$, 所以 $u(b_1)^{2^N} = u(b_1)$, 如所愿.

(c). 由定义 $\Gamma$ 的关系有 $b_1^2 = b^3$, 所以由 (b) 有 $u(b_1) = u(b)^{3^N}$. 这蕴含 $u(b)$ 与 $u(b_1)$ 交换.

(d). 设 $c = b_1^{-1}bb_1b^{-1} \in \Gamma$. 则我们假设 $c \neq 1$, 但 (c) 证明对 $\Gamma$ 的每个有限指标正规子群 $\Delta$, $c$ 在 $\Gamma/\Delta$ 中的像平凡. 所以 $\Gamma$ 不剩余有限.
\end{solution}

\begin{exercise}[$b_1$ 与 $b$ 不交换]
\label{ex:I.5.5.8}
设 $\Gamma$ 是自由群 (生成元 $a$ 和 $b$) 关于关系 $a^{-1}b^2a = b^3$ 的商. 本题的目标是证明 $b_1 := a^{-1}ba$ 与 $b$ 在 $\Gamma$ 中不交换, 即 $b_1bb_1^{-1}b^{-1}$ 在 $\Gamma$ 中不平凡.\footnote{证明它的最简单方法是找 $\Gamma$ 在其上 $b_1bb_1^{-1}b$ 非平凡作用的有限维表示, 但我们做不到. 不过, 这种想法的某种变体会奏效.}
设 $F$ 是生成元 $a$ 和 $b$ 上的自由群. 回忆 $F$ 的元素是字母 $a, a^{-1}, b, b^{-1}$ 上的既约字 (reduced words). (既约字是不含冗余对 $aa^{-1}$, $a^{-1}a$, $bb^{-1}$ 或 $b^{-1}b$ 的字.) 我们把 $F$ 的元素写成 $a^{n_1}b^{m_1}\cdots a^{n_r}b^{m_r}$, 其中 $n_1, m_1, \ldots, n_r, m_r \in \Z$, 且 $m_1, n_2, m_2, \ldots, n_{r-1}, m_{r-1}, n_r \neq 0$.
设 $\Omega$ 是如下形式的既约字的集合 $b^{r_1}a^{s_1}\cdots b^{r_m}a^{s_m}b^r$, 其中:
\begin{enumerate}
\item[(i)] $m \in \Z_{\geq 0}$, $r_i, s_i, r \in \Z$;
\item[(ii)] 对每个 $i \in \{1, \ldots, m\}$ 有 $s_i \neq 0$;
\item[(iii)] 对每个 $i \in \{2, \ldots, m\}$ 有 $r_i \neq 0$;
\item[(iv)] 对每个 $i \in \{1, \ldots, m\}$, 若 $s_i > 0$ 则 $0 \leq r_i \leq 1$;
\item[(v)] 对每个 $i \in \{1, \ldots, m\}$, 若 $s_i < 0$ 则 $0 \leq r_i \leq 2$.
\end{enumerate}
由 $\Gamma$ 的定义, 有满射群同态 $F \to \Gamma$, 记为 $\phi$.
\begin{enumerate}
\item[(a)] 证明 $\phi(\Omega) = \Gamma$.
\item[(b)] 对每个 $w \in \Omega$ 和每个 $s \in \{a, a^{-1}, b, b^{-1}\}$, 找到 $w' \in \Omega$ 使得 $\phi(w') = \phi(ws)$. 在下面把此 $w'$ 记为 $w \cdot s$.
\item[(c)] 对每个 $w \in \Omega$ 和每个 $s \in \{a, a^{-1}, b, b^{-1}\}$, 证明 $(w \cdot s) \cdot s^{-1} = w$.
\item[(d)] 证明 $(w, s) \longmapsto w \cdot s$ 延拓成 $\Gamma$ 在 $\Omega$ 上的右作用.
\item[(e)] 证明 $\phi$ 诱导双射 $\Omega \xrightarrow{\sim} \Gamma$.
\item[(f)] 证明 $b_1bb_1^{-1}b^{-1} \neq 1$ 在 $\Gamma$ 中成立.
\end{enumerate}
\end{exercise}

\begin{solution}
(a). 由自由群的定义, 可以把 $F$ 的每个元素 $w$ 写成满足条件 (i), (ii) 和 (iii) 的既约字 $b^{r_1}a^{s_1}\cdots b^{r_m}a^{s_m}b^r$. 定义 $N(w)$ 为所有满足 $r_i \notin \{0, 1\}$ 的 $s_i > 0$ 的最大值; 所以若 $w$ 满足条件 (iv), 则 $N(w) = 0$. 定义 $M(w)$ 为所有满足 $r_i \notin \{0, 1, 2\}$ 的 $s_i < 0$ 的 $|s_i|$ 的最大值; 所以若 $w$ 满足条件 (v), 则 $M(w) = 0$. 我们通过对 $N(w) + M(w)$ 归纳证明存在 $w' \in \Omega$ 使得 $\phi(w) = \phi(w')$. 若 $N(w) + M(w) = 0$, 则 $w$ 满足条件 (iv) 和 (v), 所以它在 $\Omega$ 中且结论显然.
假设 $N(w) + M(w) > 0$. 若 $N(w) > 0$, 选取满足 $s_i > 0$ 且 $r_i \notin \{0, 1\}$ 的 $i \in \{1, \ldots, m\}$. 注意定义 $\Gamma$ 的关系说 $\phi(b^2a) = \phi(ab^3)$, 从而也有 $\phi(b^{-2}a) = \phi(ab^{-3})$, 这蕴含对每个 $k \in \Z$ 有 $\phi(b^{2k}a) = \phi(ab^{3k})$. 把 $r_i$ 写成 $r_i = 2k + l$, 其中 $k \in \Z$, $l \in \{0, 1\}$, 令
\[
w' = b^{r_1}a^{s_1}\cdots b^{r_{i-1}}a^{s_{i-1}}b^la b^{3k}a^{s_i-1}b^{r_{i+1}}a^{s_{i+1}}\cdots b^{r_m}a^{s_m}b^r.
\]
则由上面的观察, $\phi(w) = \phi(w')$, 且 $N(w') < N(w)$, $M(w') = M(w)$.
类似地, 若 $M(w) > 0$, 选取满足 $s_i < 0$ 且 $r_i \notin \{0, 1, 2\}$ 的 $i \in \{1, \ldots, m\}$. 对 $k \in \Z$, 等式 $\phi(b^{2k}a) = \phi(ab^{3k})$ 也可以写成 $\phi(a^{-1}b^{2k}) = \phi(b^{3k}a^{-1})$. 把 $r_i$ 写成 $r_i = 3k + l$, 其中 $k \in \Z$, $l \in \{0, 1, 2\}$, 令
\[
w' = b^{r_1}a^{s_1}\cdots b^{r_{i-1}}a^{s_{i-1}}b^la^{-1}b^{2k}a^{s_i+1}b^{r_{i+1}}a^{s_{i+1}}\cdots b^{r_m}a^{s_m}b^r.
\]
则由上面的观察, $\phi(w) = \phi(w')$, 且 $N(w') = N(w)$, $M(w') < M(w)$. 由于 $N(w)$ 和 $M(w)$ 中必有一个 $> 0$, 总可以找到 $w' \in F$ 使得 $\phi(w') = \phi(w)$ 且 $N(w') + M(w') < N(w) + M(w)$. 对 $w'$ 应用归纳假设得到结果.

(b). 设 $w = b^{r_1}a^{s_1}\cdots b^{r_m}a^{s_m}b^r \in \Omega$; 假设条件 (i)-(v) 满足. 若 $s = b$ (相应地 $s = b^{-1}$), 则 $w' = b^{r_1}a^{s_1}\cdots b^{r_m}a^{s_m}b^{r+1}$ (相应地 $w' = b^{r_1}a^{s_1}\cdots b^{r_m}a^{s_m}b^{r-1}$) 可行. 若 $s = a$, 把 $r$ 写成 $r = 2k + l$, 其中 $k \in \Z$, $l \in \{0, 1\}$, 取 $w' = b^{r_1}a^{s_1}\cdots b^{r_m}a^{s_m}b^lab^{3k}$. 若 $s = a^{-1}$, 把 $r$ 写成 $r = 3k + l$, 其中 $k \in \Z$, $l \in \{0, 1, 2\}$, 取 $w' = b^{r_1}a^{s_1}\cdots b^{r_m}a^{s_m}b^la^{-1}b^{2k}$.

(c). 若 $s \in \{b, b^{-1}\}$ 结论显然. 假设 $s = a$, 把 $r$ 写成 $r = 2k + l$, 其中 $k \in \Z$, $l \in \{0, 1\}$. 则 $w \cdot a = b^{r_1}a^{s_1}\cdots b^{r_m}a^{s_m}b^lab^{3k}$, 所以
\[
(w \cdot a) \cdot a^{-1} = b^{r_1}a^{s_1}\cdots b^{r_m}a^{s_m}b^laa^{-1}b^{2k} = w.
\]
若 $s = a^{-1}$, 把 $r$ 写成 $r = 3k + l$, 其中 $k \in \Z$, $l \in \{0, 1, 2\}$. 则 $w \cdot a = b^{r_1}a^{s_1}\cdots b^{r_m}a^{s_m}b^la^{-1}b^{2k}$, 所以
\[
(w \cdot a^{-1}) \cdot a = b^{r_1}a^{s_1}\cdots b^{r_m}a^{s_m}b^la^{-1}ab^{3k} = w.
\]

(d). 由 (c), $(w, s) \longmapsto w \cdot s$ 延拓成 $F$ 在 $\Omega$ 上的右作用. 为证明它穿过 $\Gamma$ 在 $\Omega$ 上的右作用, 只需证明 $b^{-3}a^{-1}b^2a$ 平凡作用. 设 $w = b^{r_1}a^{s_1}\cdots b^{r_m}a^{s_m}b^r \in \Omega$; 假设条件 (i)-(v) 满足. 把 $r$ 写成 $r = 3k + l$, 其中 $k \in \Z$, $l \in \{0, 1, 2\}$. 则
\[
w \cdot b^{-3} = b^{r_1}a^{s_1}\cdots b^{r_m}a^{s_m}b^{r-3} = b^{r_1}a^{s_1}\cdots b^{r_m}a^{s_m}b^{l+3(k-1)},
\]
所以
\[
w \cdot (b^{-3}a^{-1}) = b^{r_1}a^{s_1}\cdots b^{r_m}a^{s_m}b^la^{-1}b^{2(k-1)},
\]
从而
\[
w \cdot (b^{-3}a^{-1}b^2) = b^{r_1}a^{s_1}\cdots b^{r_m}a^{s_m}b^la^{-1}b^{2k},
\]
最后
\[
w \cdot (b^{-3}a^{-1}b^2a) = b^{r_1}a^{s_1}\cdots b^{r_m}a^{s_m}b^la^{-1}ab^{3k} = b^{r_1}a^{s_1}\cdots b^{r_m}a^{s_m}b^{l+3k} = w.
\]

(e). 我们已由 (a) 知道 $\phi(\Omega) = \Gamma$, 所以只需证明 $\phi|_{\Omega}$ 单射. 由 (b) 的证明中给出的作用的显式公式, 若 $w \in \Omega$, 则 $1 \cdot w = w$. 由于由 (d), $1 \cdot w$ 只依赖于 $\phi(w)$, 这证明 $\phi(w)$ 确定 $w$.

(f). 由 (e), 只需证明 $\phi(b_1bb_1^{-1}b^{-1})$ 在 $\Omega$ 中的唯一原像不平凡. 我们已在 (a) 的证明中看到把既约字变换成 $\phi$ 下同像的 $\Omega$ 元素的算法. 把它应用于 $b_1bb_1^{-1}b^{-1} = a^{-1}baba^{-1}b^{-1}ab^{-1}$, 得到 $a^{-1}baba^{-1}bab^{-4} \neq 1$ (除了一些容易犯的错误), 所以我们完成了.
\end{solution}

\begin{exercise}[紧群的不可约酉表示有限维]
\label{ex:I.5.5.9}
本题的目标是证明定理 \ref{thm:I.3.2.13}, 即紧群的每个不可约酉表示都有限维这一事实.
设 $G$ 是紧群, $dx$ 是 $G$ 上的规范化 Haar 测度, $(\pi, V)$ 是 $G$ 的非零酉表示. 固定 $u \in V - \{0\}$, 定义 $T \colon V \to V$ 为
\[
T(v) = \int_G \langle v, \pi(x)(u) \rangle\pi(x)(u)\,\dd x.
\]
\begin{enumerate}
\item[(a)] 证明 $T$ 良定义且 $T \in \End(V)$.
\item[(b)] 证明 $T$ 是 $G$-等变的.
\item[(c)] 证明对每个 $v \in V$ 有 $\langle T(v), v \rangle \geq 0$.
\item[(d)] 证明 $T \neq 0$.
\item[(e)] 证明 $T$ 在 $\{T' \in \End(V) \mid \dim_{\CC}(\Im(T')) < +\infty\}$ 的 (对 $\|\cdot\|_{\mathrm{op}}$ 的) 闭包中; 换句话说, $T$ 在有限秩自同态空间的闭包中. (提示: $G \to V$, $x \longmapsto \pi(x)(u)$ 一致连续.)
\item[(f)] 设 $B$ 是 $V$ 中的闭单位球. 证明 $T(B)$ 紧. (换句话说, 算子 $T$ 是紧算子 (compact operator). 习题 \ref{ex:I.5.6.5} 可以帮你缩短证明.)
\item[(g)] 若 $V$ 是 $G$ 的不可约表示, 证明 $V$ 有限维.
\end{enumerate}
\end{exercise}

\begin{solution}
(a). 我们必须证明定义 $T(v)$ 的积分对每个 $v \in V$ 收敛. 设 $v \in V$. 则函数 $G \to V$, $\langle v, \pi(x)(u) \rangle\pi(x)(u)$ 连续 (因为 $x \longmapsto \pi(x)(u)$ 连续); 由于 $G$ 紧, 由习题 \ref{ex:I.5.6.3} 积分存在, 而且
\[
\|T(v)\| \leq \int_G |\langle v, \pi(x)(u) \rangle|\|\pi(x)(u)\|\,\dd x \leq \|v\|\|u\|^2.
\]
函数 $T \colon V \to V$ 是 $\CC$-线性的 (因为加法和乘以纯量在 $V$ 上连续, 所以它们与积分交换, 由习题 \ref{ex:I.5.6.1}(b)), 且上面的不等式证明 $T$ 有界且 $\|T\|_{\mathrm{op}} \leq \|u\|^2$.

(b). 设 $v \in V$, $x \in G$.
\begin{align*}
T(\pi(x)(v)) &= \int_G \langle \pi(x)(v), \pi(y)(u) \rangle\pi(y)(u)\,\dd y \\
&= \int_G \langle v, \pi(x)^*\pi(y)(u) \rangle\pi(y)(u)\,\dd y \\
&= \int_G \langle v, \pi(x^{-1}y)(u) \rangle\pi(y)(u)\,\dd y \\
&= \int_G \langle v, \pi(y)(u) \rangle\pi(x)\pi(y)(u)\,\dd y
\end{align*}
(最后一个等式用 Haar 测度的左不变性). 由于 $\pi(x) \colon V \to V$ 连续且线性, 习题 \ref{ex:I.5.6.1}(b) 蕴含最后一行等于
\[
\pi(x)\Bigl(\int_G \langle v, \pi(y)(u) \rangle\pi(y)(u)\,\dd y\Bigr) = \pi(x)(T(v)),
\]
这正是我们想要的.

(c). 设 $v \in V$. 由于 $\langle v, \cdot \rangle$ 在 $V$ 上连续且线性,
\[
\langle T(v), v \rangle = \int_G \overline{\langle \pi(x)(u), v \rangle}\langle v, \pi(x)(u) \rangle\,\dd x = \int_G |\langle v, \pi(x)(u) \rangle|^2\,\dd x \geq 0.
\]

(d). 取 $v = u$. 由于 $\langle u, \pi(x)(u) \rangle = \|u\|^2 > 0$ 且 $x \longmapsto \langle u, \pi(x)(u) \rangle$ 是从 $G$ 到 $\CC$ 的连续函数, 存在 $\varepsilon > 0$ 和 $1$ 在 $G$ 中的开邻域 $U$ 使得对 $x \in U$ 有 $|\langle u, \pi(x)(u) \rangle|^2 \geq \varepsilon$. 则 (c) 的证明中的计算给出
\[
\langle T(u), u \rangle = \int_G |\langle u, \pi(x)(u) \rangle|^2\,\dd x \geq \varepsilon\mu(U) > 0.
\]
所以 $T \neq 0$.

(e). 设 $\varepsilon > 0$. 由于 $G$ 紧, 连续函数 $G \to \CC$, $x \longmapsto \pi(x)(u)$ 一致连续, 所以存在 $1$ 的邻域 $U$ 使得对 $x \in G$ 和 $y \in xU$ 有 $\|\pi(x)(u) - \pi(y)(u)\| \leq \varepsilon$. 由于 $G$ 紧且族 $(xU)_{x \in G}$ 覆盖 $G$, 可以找到 $x_1, \ldots, x_n \in G$ 使得 $G = \bigcup_{i=1}^n x_iU$. 选取 $X$ 的 Borel 子集 $E_1, \ldots, E_n$ 使得对每个 $i \in \{1, \ldots, n\}$ 有 $x_i \in E_i \subset x_iU$ 且 $X = E_1 \sqcup \cdots \sqcup E_n$ (作为集合). 若 $x \in E_i$, $v \in V$, 则
\begin{align*}
&\|\langle v, \pi(x)(u) \rangle\pi(x)(u) - \langle v, \pi(x_i)(u) \rangle\pi(x_i)(u)\| \\
&\leq \|\langle v, (\pi(x) - \pi(x_i))(u) \rangle\pi(x)(u)\| + \|\langle v, \pi(x_i)(u) \rangle(\pi(x) - \pi(x_i))(u)\| \\
&\leq \|v\|\varepsilon\|u\| + \|v\|\|u\|\varepsilon = 2\varepsilon\|v\|\|u\|.
\end{align*}
定义 $U \in \End(V)$ 为
\[
T(v) = \sum_{i=1}^n\mu(E_i)\langle v, \pi(x_i)(u) \rangle\pi(x_i)(u) = \sum_{i=1}^n\int_{E_i}\langle v, \pi(x_i)(u) \rangle\pi(x_i)(u)\,\dd x.
\]
这个算子 $U$ 有有限秩, 因为它的像包含于 $\operatorname{Span}(\pi(x_1)(u), \ldots, \pi(x_n)(u))$. 又, 由上面的计算 (以及习题 \ref{ex:I.5.6.2}), 对每个 $v \in V$,
\begin{align*}
\|T(v) - U(v)\| &\leq \sum_{i=1}^n\Bigl\|\int_{E_i}\langle v, \pi(x)(u) \rangle\pi(x)(u)\,\dd x - \int_{E_i}\langle v, \pi(x_i)(u) \rangle\pi(x_i)(u)\,\dd x\Bigr\| \\
&\leq \sum_{i=1}^n\mu(E_i)2\varepsilon\|v\|\|u\| \\
&= 2\varepsilon\|v\|\|u\|.
\end{align*}
所以 $\|T - U\|_{\mathrm{op}} \leq 2\varepsilon\|u\|$. 由于 $\varepsilon > 0$ 任意, 这证明 $T$ 是有限秩算子的极限.

(f). 由习题 \ref{ex:I.5.6.5}(e), 只需证明 $T(B)$ 完全有界. 设 $U$ 是 $0$ 的邻域, 可以假设它是半径 $\varepsilon > 0$ 的开球. 我们必须找到 $x_1, \ldots, x_n \in B$ 使得 $T(B)$ 的每个点与某个 $T(x_i)$ 的距离 $< \varepsilon$.
由 (e), 我们知道 $T$ 是有限秩算子的极限, 所以可以找到有限秩的 $U \in \End(V)$ 使得 $\|T - U\|_{\mathrm{op}} \leq \varepsilon/4$. 由于 $U$ 有有限秩, $U(B)$ 是有限维空间 $\Im(U)$ 的闭有界子集, 所以它紧. 特别地, 可以找到 $x_1, \ldots, x_n \in B$ 使得对每个 $y \in B$, 存在 $i \in \{1, \ldots, n\}$ 满足 $\|U(y) - U(x_i)\| < \varepsilon/2$.
现在设 $y \in B$, 选取满足 $\|U(y) - U(x_i)\| < \varepsilon/2$ 的 $i \in \{1, \ldots, n\}$. 则
\[
\|T(y) - T(x_i)\| \leq \|T(y) - U(y)\| + \|U(y) - U(x_i)\| + \|U(x_i) - T(x_i)\| < \|y\|\varepsilon/4 + \varepsilon/2 + \|x_i\|\varepsilon/4 \leq \varepsilon.
\]
(回忆 $y, x_i$ 在 $V$ 的闭单位球中.)

(g). 现在把一切放在一起. 假设 $V$ 是 $G$ 的不可约酉表示. 则我们构造的算子 $T \in \End(V)$ 是 $G$-等变的, 所以由 Schur 引理存在 $\lambda \in \CC$ 使得 $T = \lambda\mathrm{id}_V$. 由于 $T \neq 0$, 有 $\lambda \neq 0$. 所以 $T(\lambda^{-1}B)$ 是 $V$ 中的闭单位球. 部分 (f) 说它紧, 由 Riesz 引理这蕴含 $V$ 有限维.
\end{solution}

\subsection{向量值积分与 Minkowski 不等式}

注意: 允许不使用证明就用下列结果:
\begin{itemize}
\item Hahn-Banach 定理.
\item 希尔伯特空间 $V$ 上的每个连续线性泛函都具有 $\langle \cdot, v \rangle$ 的形式 (其中 $v \in V$) 这一事实.
\item H\"older 不等式.
\item 若 $(X, \mu)$ 是测度空间, $1 \leq p < +\infty$, $1 < q \leq +\infty$ 满足 $p^{-1} + q^{-1} = 1$, 则映射 $L^p(X, \mu) \to \Hom(L^q(X, \mu), \CC)$, $f \longmapsto (g \longmapsto \int_X fg\,\dd\mu)$ 是保范同构 (左边是 $L^p$ 范数, 右边是算子范数) 这一事实.\footnote{技术注记: 对 $p = 1$, $q = +\infty$, 若 $\mu$ 不是 $\sigma$-有限的, 这在一般情形不成立; 但对局部紧 Hausdorff 空间上的正则 Borel 测度, 可以通过稍微修改 $L^{\infty}$ 的定义来补救. 你可以忽略这一点.}
\end{itemize}

\begin{exercise}[弱可积函数与积分]
\label{ex:I.5.6.1}
设 $(X, \mu)$ 是测度空间, $V$ 是巴拿赫空间. 记 $V^{\vee} = \Hom(V, \CC)$. 若对每个 $T \in V^{\vee}$, 函数 $T \circ f \colon X \to \CC$ 都在 $L^1(X, \mu)$ 中, 则称函数 $f \colon X \to V$ 弱可积 (weakly integrable). 若 $f$ 弱可积且存在 $V$ 的元素 $v$ 使得对每个 $T \in V^{\vee}$ 有 $T(v) = \int_X T \circ f(x)\,\dd\mu(x)$, 则称 $v$ 是 $f$ 在 $X$ 上的积分 (integral), 记为 $v = \int_X f(x)\,\dd\mu(x) = \int_X f\,\dd\mu$.
\begin{enumerate}
\item[(a)] 证明 $f$ 的积分若存在则唯一.
\item[(b)] 设 $W$ 是另一个巴拿赫空间, $u \in \Hom(V, W)$. 若 $f \colon X \to V$ 弱可积且有积分 $v$, 证明 $u \circ f \colon X \to W$ 弱可积且有积分, 且积分等于 $u(v)$.
\item[(c)] 给出一个弱可积但没有积分的函数的例子.
\end{enumerate}
\end{exercise}

\begin{solution}
(a). 由 Hahn-Banach 定理, 对每个 $v \in V$, 存在 $T \in \Hom(V, \CC)$ 使得 $T(v) = \|v\|$ 且 $\|T\|_{\mathrm{op}} \leq 1$. 特别地, $V$ 的元素 $v$ 为零当且仅当对每个 $T \in \Hom(V, \CC)$ 有 $T(v) = 0$; 换句话说, 两个元素 $v, w \in V$ 相等当且仅当对每个 $T \in \Hom(V, \CC)$ 有 $T(v) = T(w)$. 这蕴含 $f$ 的积分若存在则唯一.

(b). 我们首先证明 $u \circ f$ 弱可积. 设 $T \in \Hom(W, \CC)$. 则 $T \circ u \in \Hom(V, \CC)$, 所以函数 $T \circ u \circ f \colon X \to \CC$ 可积.
现在假设 $f$ 有积分 $v$. 则对每个 $T \in \Hom(W, \CC)$, 有 $T \circ u \in \Hom(V, \CC)$, 所以 $\int_X T \circ u \circ f\,\dd\mu = T \circ u(v)$. 这意味着 $u(v)$ 是 $u \circ f$ 的积分.

(c). 取 $X = \N$ 带计数测度 $\mu$,
\[
V = c_0(\N) := \{(x_n)_{n \geq 0} \in \CC^{\N} \mid \lim_{n \to +\infty}x_n = 0\}.
\]
我们将用 $\ell^1(\N)$ 是 $c_0(\N)$ 的连续对偶这一事实, 对偶配对为 $\ell^1(\N) \times c_0(\N) \to \CC$, $((x_n), (y_n)) \longmapsto \sum_{n \geq 0}x_ny_n$, 且 $\ell^1(\N)$ 的连续对偶是 $\ell^{\infty}(\N)$ (通过类似的映射). 从 $c_0(\N)$ 到其双重对偶的映射是通常的嵌入 $c_0(\N) \subset \ell^{\infty}(\N)$.
我们定义 $f \colon X \to V$ 为 $f(n) = \mathbb{1}_{\{n\}}$. 则对每个 $(x_n)_{n \geq 0} \in \ell^1(\N)$, 若 $T \colon c_0(\N) \to \CC$ 是对应的线性泛函, 则
\[
\int_X T(f(x))\,\dd\mu(x) = \sum_{n \geq 0}x_n,
\]
它收敛因为 $(x_n)_{n \geq 0} \in \ell^1(\N)$. 所以 $f$ 弱可积. 但 $f$ 没有积分 (至少在 $c_0(\N)$ 中没有), 因为它在 $\ell^1(\N)$ 上定义的连续线性泛函由 $\ell^{\infty}(\N)$ 的某个不在 $c_0(\N)$ 中的元素 (常值序列 $1$) 表示. 由于在 $\ell^1(\N)$ 的点上取值分离 $\ell^{\infty}(\N)$ 的元素, 不可能有 $c_0(\N)$ 的元素在 $\ell^1(\N)$ 上给出相同的线性泛函.
\end{solution}

\begin{exercise}[紧支撑连续函数的积分]
\label{ex:I.5.6.2}
在此题中, $X$ 是局部紧 Hausdorff 空间, $\mu$ 是 $X$ 上的正则 Borel 测度. 设 $V$ 是巴拿赫空间, $f \colon X \to V$ 是具有紧支撑的连续函数.
\begin{enumerate}
\item[(a)] 证明 $f$ 弱可积.
\item[(b)] 若 $\mu(\operatorname{supp} f) = 0$, 证明 $\int_X f\,\dd\mu$ 存在且等于 $0$.
\end{enumerate}
本题的目标是证明:
\begin{enumerate}
\item[(i)] $f$ 有积分 $v$;
\item[(ii)] $\|v\| \leq \int_X \|f(x)\|\,\dd\mu(x)$;
\item[(iii)] 若 $\mu(\operatorname{supp} f) \neq 0$, 则 $\mu(\operatorname{supp} f)^{-1}v$ 在 $f(X)$ 的凸包的闭包中.
\end{enumerate}
由 (b), 可以 (而且将) 假设 $\mu(\operatorname{supp} f) \neq 0$.
\begin{enumerate}
\item[(a)] 证明可以假设 $X = \operatorname{supp} f$ (特别地 $X$ 紧) 且 $\mu(X) = 1$.
\end{enumerate}
从现在起, 我们假设 $X$ 紧且 $\mu(X) = 1$.
\begin{enumerate}
\item[(a)] 设 $T_1, \ldots, T_n \colon V \to \R$ 是有界 $\R$-线性泛函 (以显然的方式把 $V$ 看作 $\R$-向量空间), 定义 $a_1, \ldots, a_n \in \R$ 为 $a_i = \int_X T_i \circ f\,\dd\mu$. 证明 $(a_1, \ldots, a_n)$ 在 $\R^n$ 的紧子集 $((T_1, \ldots, T_n) \circ f)(X)$ 的凸包中. (提示: 如果不是会怎样?)
\end{enumerate}
设 $K$ 是 $f(X)$ 的凸包的闭包. 由习题 \ref{ex:I.5.6.7}, 这是 $V$ 的紧子集. 对 $\Hom_{\R}(V, \R)$ (从 $V$ 到 $\R$ 的有界 $\R$-线性泛函空间) 的每个有限子集 $\Omega$, 记 $I_{\Omega}$ 为 $K$ 的满足对每个 $T \in \Omega$ 有 $T(v) = \int_X T \circ f\,\dd\mu$ 的元素 $v$ 的集合.
\begin{enumerate}
\item[(a)] 证明对每个 $\Omega \subset \Hom_{\R}(V, \R)$, $I_{\Omega}$ 紧.
\item[(b)] 证明若 $\Omega$ 有限, 则 $I_{\Omega}$ 非空.
\item[(c)] 证明 $f$ 的积分存在且在 $K$ 中.
\item[(d)] 证明:
\[
\Bigl\|\int_X f\,\dd\mu\Bigr\| \leq \int_X \|f(x)\|\,\dd\mu(x).
\]
(提示: Hahn-Banach.)
\end{enumerate}
\end{exercise}

\begin{solution}
(a). 设 $T \in \Hom(V, \CC)$. 则 $T \circ f \colon X \to \CC$ 连续, 且其支撑包含于 $\operatorname{supp}(f)$, 从而是紧的. 所以 $T \circ f$ 可积.

(b). 假设 $\mu(\operatorname{supp} f) = 0$. 若 $T \in \Hom(V, \CC)$, 则 $T \circ f \colon X \to \CC$ 连续且 $\mu(\operatorname{supp}(T \circ f)) = 0$, 所以 $\int_X T \circ f\,\dd\mu = 0$. 这证明 $0$ 是 $f$ 的积分.

(c). 假设当 $X = \operatorname{supp} f$ 且 $\mu(X) = 1$ 时知道结论. 设 $f \colon X \to V$ 是具有紧支撑的连续函数. 我们已见可以假设 $\mu(\operatorname{supp} f) \neq 0$, 就这样做. 设 $X' = \operatorname{supp} f$, 考虑 $X'$ 上的测度 $\mu'$, 它是 $\mu$ 的限制乘以 $\mu(\operatorname{supp} f)^{-1}$. 由我们的假设, $\int_{X'}f|_{X'}\,\dd\mu'$ 存在, 记之为 $v$, 有 $\|v\| \leq \int_{X'}\|f(x)\|\,\dd\mu'(x)$, 且 $v$ 在 $f(X')$ 的凸包的闭包中.
让我们证明 $w := \mu(\operatorname{supp} f)v$ 是 $f$ 的积分. 注意 $\mu(\operatorname{supp} f)^{-1}w$ 在 $f(X') = f(X)$ 的凸包中, 且
\[
\|w\| \leq \mu(\operatorname{supp} f)\int_{X'}|f(x')|\,\dd\mu'(x') = \int_X |f(x)|\,\dd\mu(x),
\]
所以这证明了对 $f$ 的结论.
设 $T \in \Hom(V, \CC)$. 则 $\operatorname{supp}(T \circ f) \subset X'$, 所以
\[
\int_X T \circ f(x)\,\dd\mu(x) = \mu(\operatorname{supp} f)\int_{X'}T \circ f(x)\,\dd\mu'(x) = T(w).
\]
所以 $w = \int_X f\,\dd\mu$.

(d). 设 $L = ((T_1, \ldots, T_n) \circ f)(X)$. 假设 $(a_1, \ldots, a_n)$ 不在 $L$ 的凸包中. 则由超平面分离定理, 存在线性泛函 $\lambda \colon \R^n \to \R$ 和 $c > 0$ 使得对每个 $v \in L$ 有 $\lambda(a_1, \ldots, a_n) \geq c + \lambda(v)$. 设 $\lambda_1, \ldots, \lambda_n \in \R$ 是 $\lambda$ 在 $\R^n$ 的典范基向量下的像. 则对每个 $x \in X$,
\[
\sum_{i=1}^n\lambda_ia_i = \lambda(a_1, \ldots, a_n) \geq c + \lambda \circ (T_1, \ldots, T_n) \circ f(x) = c + \sum_{i=1}^n\lambda_iT_i(f(x)).
\]
对 $X$ 取积分 (利用 $\mu(X) = 1$) 给出
\[
\sum_{i=1}^n\lambda_ia_i \geq c + \sum_{i=1}^n\lambda_i\int_X T_i \circ f(x)\,\dd\mu(x) = c + \sum_{i=1}^n\lambda_ia_i,
\]
矛盾.

(e). 对每个 $T \in \Hom_{\R}(V, \R)$, $K$ 的满足 $T(v) = \int_X T \circ f\,\dd\mu$ 的元素 $v$ 的集合是 $K$ 的闭子集. 由于集合 $I_{\Omega}$ 是这种形式的集合的交, 它也是 $K$ 的闭子集, 从而因 $K$ 紧而紧.

(f). 若 $\Omega$ 有限, 写 $\Omega = \{T_1, \ldots, T_n\}$, $T_{\Omega} = (T_1, \ldots, T_n) \colon V \to \R^n$. 我们已在 (d) 中看到 $a := \int_X T_{\Omega}(f(x))\,\dd\mu(x)$ 在 $T_{\Omega}(f(X))$ 的凸包中, 所以存在 $v \in V$ 使得 $v$ 在 $f(X)$ 的凸包中 (从而在 $K$ 中) 且 $T_{\Omega}(v) = a = \int_X T_{\Omega}(f(x))\,\dd\mu(x)$. 第二个条件正好说 $v \in I_{\Omega}$.

(g). $K$ 的子集族 $(I_{\{T\}})_{T \in \Hom_{\R}(V, \R)}$ 由 (f) 有有限交性质. 由于 $K$ 紧, 这蕴含 $\bigcap_{T \in \Hom_{\R}(V, \R)}I_{\{T\}}$ 非空. 在其中选取向量 $v$. 设 $T \in \Hom(V, \CC)$. 由于 $\Re(T)$ 和 $\Im(T)$ 在 $\Hom_{\R}(V, \R)$ 中,
\begin{align*}
T(v) &= \Re(T(v)) + i\Im(T(v)) \\
&= \int_X \Re(T(f(x)))\,\dd\mu(x) + i\int_X \Im(T(f(x)))\,\dd\mu(x) \\
&= \int_X T(f(x))\,\dd\mu(x),
\end{align*}
所以 $v = \int_X f\,\dd\mu$. 又, $v \in K$, 因为由定义所有 $I_{\{T\}}$ 都包含于 $K$.

(h). 由 Hahn-Banach 定理, 存在 $T \in \Hom(V, \CC)$ 使得 $T(v) = \|v\|$ 且 $\|T\|_{\mathrm{op}} \leq 1$. 则
\[
\|v\| = |T(v)| = \Bigl|\int_X T(f(x))\,\dd\mu(x)\Bigr| \leq \int_X |T(f(x))|\,\dd\mu(x) \leq \int_X \|f(x)\|\,\dd\mu(x).
\]
\end{solution}

\begin{exercise}[$L^1$ 函数乘有界连续函数的积分]
\label{ex:I.5.6.3}
在此题中, $X$ 是局部紧 Hausdorff 空间, $\mu$ 是 $X$ 上的正则 Borel 测度. 设 $V$ 是巴拿赫空间, $f \colon X \to \CC$ 是 $L^1(X, \mu)$ 中的函数, $G \colon X \to V$ 是有界连续函数.
本题的目标是证明:
\begin{enumerate}
\item[(i)] 函数 $fG \colon X \to V$ 有积分 $v$;
\item[(ii)] $\|v\| \leq \bigl(\sup_{x \in X}\|G(x)\|\bigr)\int_X |f(x)|\,\dd\mu(x)$;
\item[(iii)] $v \in \overline{\operatorname{Span}(G(X))}$.
\end{enumerate}
\begin{enumerate}
\item[(a)] 证明 $fG$ 弱可积.
\item[(b)] 设 $(f_n)_{n \geq 0}$ 是 $\mathcal{C}_c(X)$ 的函数序列, 在 $L^1(X, \mu)$ 中收敛到 $f$. 证明对每个 $n \geq 0$, $\int_X f_nG\,\dd\mu$ 存在, 且 $(\int_X f_nG\,\dd\mu)_{n \geq 0}$ 是 Cauchy 序列.
\item[(c)] 证明上面的断言 (i), (ii) 和 (iii).
\end{enumerate}
\end{exercise}

\begin{solution}
(a). 设 $T \in \Hom(V, \CC)$. 则对每个 $x \in X$,
\[
|T(f(x)G(x))| \leq |f(x)||T(G(x))| \leq |f(x)|\|T\|_{\mathrm{op}}\|G(x)\| \leq |f(x)|\|T\|_{\mathrm{op}}\sup_{y \in X}\|G(y)\|.
\]
由于 $\sup_{y \in X}\|G(y)\| < +\infty$ 且 $f \in L^1(X, \mu)$, 函数 $T \circ (fG)$ 可积. 所以 $fG$ 弱可积.

(b). 对每个 $h \in \mathcal{C}_c(X)$, 函数 $hG \colon X \to V$ 连续且支撑包含于 $\operatorname{supp}(h)$, 从而是紧的. 由习题 \ref{ex:I.5.6.2}, 此函数可积, 且
\[
\Bigl\|\int_X (hG)(x)\,\dd\mu(x)\Bigr\| \leq \int_X |h(x)|\|G(x)\|\,\dd\mu(x) \leq \sup_{y \in X}\|G(y)\|\|h\|_1.
\]
把它应用于 $f_n$ 证明 $f_nG$ 可积, 把它应用于 $f_n - f_m$ 证明
\[
\Bigl\|\int_X (f_nG)\,\dd\mu - \int_X (f_mG)\,\dd\mu\Bigr\| \leq \sup_{y \in X}\|G(y)\|\|f_n - f_m\|_1 \xrightarrow[n,m \to +\infty]{} 0
\]
因为 $(f_n)_{n \geq 0}$ 在 $L^1(X, \mu)$ 中收敛.

(c). 由于 $V$ 完备, Cauchy 序列 $(\int_X (f_nG)\,\dd\mu)$ 在 $V$ 中有极限, 记之为 $v$. 对每个 $T \in \Hom(V, \CC)$,
\[
T(v) = \lim_{n \to +\infty}T\Bigl(\int_X (f_nG)\,\dd\mu\Bigr) = \lim_{n \to +\infty}\int_X T(f_n(x)G(x))\,\dd\mu(x).
\]
像 (a) 那样,
\[
\Bigl\|\int_X T(f_n(x)G(x))\,\dd\mu(x) - \int_X T(f(x)G(x))\,\dd\mu(x)\Bigr\| \leq \sup_{y \in X}\|G(y)\|\|T\|_{\mathrm{op}}\|f_n - f\|_1,
\]
所以当 $n \to +\infty$ 时这收敛到 $0$, 于是
\[
T(v) = \int_X T(f(x)G(x))\,\dd\mu(x).
\]
这证明 $v$ 是 $fG$ 的积分.
而且, 由习题 \ref{ex:I.5.6.2}, 对每个 $n \geq 0$, $\int_X (f_nG)\,\dd\mu$ 在 $G(X)$ 的线性张成的闭包中. 由于 $v$ 是这些向量的极限, 它也在 $\operatorname{Span}(G(X))$ 的闭包中.
最后, 为证明对 $\|v\|$ 的界, 可以用 Hahn-Banach 定理和 (h) 中刻画 $v$ 的性质, 或用如下事实: 对每个 $n \geq 0$,
\[
\Bigl\|\int_X (f_nG)\,\dd\mu\Bigr\| \leq \int_X |f_n(x)|\|G(x)\|\,\dd\mu(x) \leq \sup_{x \in X}\|G(x)\|\|f_n\|_1
\]
且这个积分序列收敛到 $v$.
\end{solution}

\begin{exercise}[卷积的积分公式]
\label{ex:I.5.6.4}
设 $G$ 是局部紧群, $\mu$ 是 $G$ 上的左 Haar 测度, $L^1(G) = L^1(G, \mu)$. 设 $f, g \in L^1(G)$.
\begin{enumerate}
\item[(a)] 证明函数 $G \to L^1(G)$, $y \longmapsto f(y)L_yg$ 弱可积且有积分.
\item[(b)] 证明
\[
f * g = \int_G f(y)L_yg\,\dd\mu(y).
\]
\end{enumerate}
\end{exercise}

\begin{solution}
(a). 注意函数 $G \to L^1(G)$, $y \longmapsto L_yg$ 连续, 且 $\sup_{y \in G}\|L_yg\|_1 = \|g\|_1 < +\infty$. 所以结论由习题 \ref{ex:I.5.6.3} 得出.

(b). 设 $F = \int_G f(y)L_yg\,\dd\mu(y) \in L^1(G)$. 由积分的定义, 对每个 $h \in L^{\infty}(G)$,
\[
\int_G h(x)F(x)\,\dd\mu(x) = \int_{G \times G}h(x)f(y)g(y^{-1}x)\,\dd\mu(y)\,\dd\mu(x) = \int_G h(x)(f * g)(x)\,\dd\mu(x).
\]
由于 $L^{\infty}(G)$ 是 $L^1(G)$ 的连续对偶, 由习题 \ref{ex:I.5.6.1}(a) 有 $f * g = F$.
\end{solution}

\begin{exercise}[完全有界与凸包]
\label{ex:I.5.6.5}
设 $V$ 是赋范向量空间. 若对 $V$ 中 $0$ 的每个邻域 $U$, 存在有限集 $F$ 使得 $A \subset F + U$, 则称 $V$ 的子集 $A$ 完全有界 (totally bounded).
\begin{enumerate}
\item[(a)] 证明 $V$ 的有限子集的凸包紧.
\item[(b)] 证明 $V$ 的每个紧子集完全有界.
\item[(c)] 若 $A \subset V$ 完全有界, 证明 $\overline{A}$ 完全有界.
\item[(d)] 若 $A$ 是 $V$ 的完全有界子集, 证明其凸包完全有界. (提示: 开球是凸的.)
\item[(e)] 若 $V$ 完备且 $A$ 完全有界, 证明 $\overline{A}$ 紧.
\item[(f)] 若 $V$ 完备且 $K \subset V$ 紧, 证明 $K$ 的凸包的闭包紧.
\end{enumerate}
\end{exercise}

\begin{solution}
(a). 设 $F$ 是 $V$ 的有限子集. 则其凸包包含于 $\operatorname{Span}(F)$, 它是有限维的. 在有限维向量空间中, 任何紧集的凸包紧, 所以有限集 $F$ 的凸包紧.

(b). 设 $K$ 是 $V$ 的紧子集, $U$ 是 $V$ 中 $0$ 的邻域. 可以假设 $U$ 开. 则 $K \subset \bigcup_{x \in K}(x + U)$. 由于 $K$ 紧, 存在 $K$ 的有限子集 $F$ 使得 $K \subset \bigcup_{x \in F}(x + U) = F + U$.

(c). 设 $A \subset V$ 完全有界, $U$ 是 $V$ 中 $0$ 的邻域. 可以假设 $U$ 是以 $0$ 为心、正半径 (设为 $c$) 的开球. 设 $U'$ 是以 $0$ 为心、半径 $c/2$ 的开球. 由于 $A$ 完全有界, 存在有限集 $F$ 使得 $A \subset F + U'$. 由于 $F$ 有限, 集合 $F + \overline{U'}$ 闭, 所以它包含 $\overline{A}$. 但 $U \supset U'$, 所以 $\overline{A} \subset F + U$.

(d). 设 $U$ 是 $V$ 中 $0$ 的邻域. 选取凸开邻域 $U'$ (例如球) 使得 $U' + U' \subset U$, 设 $F$ 是满足 $A \subset F + U'$ 的有限集. 设 $K$ 是 $F$ 的凸包, 则 $A \subset K + U'$. 由于 $K$ 和 $U'$ 凸, $K + U'$ 也凸, 所以 $A$ 的凸包包含于 $K + U'$. 另一方面, 集合 $K$ 由 (a) 紧, 从而由 (b) 完全有界, 所以存在有限集 $F'$ 使得 $K \subset F' + U'$, 从而 $K + U' \subset F' + U' + U' \subset F' + U$. 所以我们找到了有限集 $F'$ 使得 $A$ 的凸包包含于 $F' + U$.

(e). 记 $K = \overline{A}$. 对每个 $x \in V$ 和 $c > 0$, 设 $\overline{B}(x, c)$ 是以 $x$ 为心、$c$ 为半径的闭球.
设 $(U_i)_{i \in I}$ 是 $K$ 的开子集族, 满足 $K \subset \bigcup_{i \in I}U_i$, 并假设 $(U_i)_{i \in I}$ 没有覆盖 $K$ 的有限子族. 由 (c) 我们知道 $K$ 完全有界. 我们将对 $n$ 归纳构造 $K$ 的非空闭子集的递减序列 $(K_n)_{n \geq 1}$, 使得 $K_n$ 包含于半径为 $1/n$ 的球中且 $K_n$ 不能被 $(U_i)_{i \in I}$ 的有限子族覆盖.
首先, 由于 $K$ 完全有界, 存在有限集 $F$ 使得 $K \subset F + \overline{B}(0, 1) = \bigcup_{x \in F}\overline{B}(x, 1)$. 选取 $x \in F$ 使得 $K \cap \overline{B}(x, 1)$ 非空且不能被有限多个 $U_i$ 覆盖, 取 $K_1 = K \cap \overline{B}(x, 1)$.
现在假设已经构造了 $K_1, \ldots, K_n$ ($n \geq 1$). 则由于 $K_n$ 完全有界 (作为 $K$ 的子集), 存在有限集 $F$ 使得 $K_n \subset F + \overline{B}(0, (n+1)^{-1})$. 同样地, 由于 $K_n$ 不能被有限多个 $U_i$ 覆盖, 必存在 $x \in F$ 使得 $K_n \cap \overline{B}(x, (n+1)^{-1})$ 非空且也不能被有限多个 $U_i$ 覆盖, 取 $K_{n+1} = K_n \cap \overline{B}(x, (n+1)^{-1})$.
对每个 $n \geq 1$ 选取 $x_n \in K_n$. 由 $K_n$ 包含于半径为 $1/n$ 的球中这一条件, 序列 $(x_n)_{n \geq 0}$ 是 Cauchy 列. 由于 $V$ 完备, $(x_n)_{n \geq 1}$ 有极限, 设为 $x$. 由于 $x \in K$, 存在 $i \in I$ 使得 $x \in U_i$. 但此时对足够小的 $c > 0$ 有 $\overline{B}(x, c) \subset U_i$, 所以对足够大的 $n$ 有 $K_n \subset U_i$, 这与 $K_n$ 的性质矛盾.

(f). 由 (d), $K$ 的凸包完全有界, 所以其闭包由 (e) 紧.
\end{solution}

\begin{exercise}[$L^{\infty}$ 上的卷积]
\label{ex:I.5.6.6}
设 $G$ 是局部紧群, $dx$ 是 $G$ 上的左 Haar 测度, $\mathrm{UCB}(G)$ 是 $L^{\infty}(G)$ 的由 $G$ 上的左一致连续有界函数组成的子空间.
设 $f \in L^1(G)$, $\phi \in L^{\infty}(G)$.
\begin{enumerate}
\item[(a)] 证明 $f * \phi$ 存在且左一致连续有界.
\item[(b)] 若 $\phi \in \mathrm{UCB}(G)$, 证明积分 $\int_G f(y)L_y\phi\,\dd y$ 存在且等于 $f * \phi$.
\end{enumerate}
\end{exercise}

\begin{solution}
(a). 设 $x \in G$. 则定义 $f * \phi(x)$ 的积分是
\[
\int_G f(y)\phi(y^{-1}x)\,\dd\mu(y),
\]
它收敛, 因为对每个 $y \in G$ 有 $|f(y)\phi(y^{-1}x)| \leq \|\phi\|_{\infty}|f(y)|$. 这也证明对每个 $x \in G$ 有 $|f * \phi(x)| \leq \|\phi\|_{\infty}\|f\|_1$, 所以 $f * \phi$ 有界且
\[
\|f * \phi\|_{\infty} \leq \|\phi\|_{\infty}\|f\|_1.
\]
现在我们证明 $f * \phi$ 左一致连续. 设 $x \in G$. 由命题 \ref{prop:I.4.1.3}, $L_x(f * \phi) = (L_xf) * \phi$, 所以
\[
\|L_x(f * \phi) - f * \phi\|_{\infty} = \|(L_xf - f) * \phi\|_{\infty} \leq \|L_xf - f\|_1\|\phi\|_{\infty}.
\]
由命题 \ref{prop:I.3.1.13}, 当 $x$ 趋于 $1$ 时这趋于 $0$, 这正好意味着 $f * \phi$ 左一致连续.

(b). 假设 $\phi \in \mathrm{UCB}(G)$. 则映射 $G \to \mathrm{UCB}(G)$, $y \longmapsto L_y\phi$ 连续 (见注记 \ref{rem:I.1.13}), 所以由习题 \ref{ex:I.5.6.3}, 积分 $\int_G f(y)L_y\phi\,\dd\mu(y)$ 在 $\mathrm{UCB}(G)$ 中存在. 设 $h = \int_G f(y)L_y\phi\,\dd\mu(y)$.
对每个 $g \in L^1(G)$, 映射 $\psi \longmapsto \int_G g\psi\,\dd\mu$ 是 $\mathrm{UCB}(G)$ 上的连续线性泛函. 所以由积分的定义,
\[
\int_G gh\,\dd\mu = \int_{G \times G}h(x)f(y)\phi(y^{-1}x)\,\dd\mu(x)\,\dd\mu(y) = \int_G g(f * \phi)\,\dd\mu.
\]
由于 $L^1(G)$ 的元素定义的线性泛函在 $L^{\infty}(G)$ 上分离点, 这蕴含 $h = f * \phi$.
\end{solution}

\begin{exercise}[Minkowski 不等式]
\label{ex:I.5.6.7}
设 $(X, \mu)$ 和 $(Y, \nu)$ 是测度空间, 为简化我们取其为 $\sigma$-有限的.\footnote{有办法把结果推广到不一定是 $\sigma$-紧的带 Haar 测度的局部紧群.} 设 $p \in (1, +\infty)$.\footnote{Minkowski 不等式对 $p = 1$ 仍成立, 但该情形由 Fubini-Tonelli 定理立即得出.} 设 $\phi \colon X \times Y \to \CC$ 是可测函数. 假设
\[
\int_Y\Bigl(\int_X|\phi(x, y)|^p\,\dd\mu(x)\Bigr)^{1/p}\,\dd\nu(y) < \infty.
\]
\begin{enumerate}
\item[(a)] 证明函数 $\phi(\cdot, y)$ 对几乎所有 $y \in Y$ 在 $L^p(X, \mu)$ 中.
\item[(b)] 设 $Y'$ 是 $Y$ 的可测子集, 满足 $\nu(Y - Y') = 0$ 且对每个 $y \in Y'$ 有 $\phi(\cdot, y) \in L^p(X, \mu)$. 证明函数 $Y' \to L^p(X, \mu)$, $y \longmapsto \phi(\cdot, y)$ 弱可积.
\item[(c)] 证明 (b) 中函数的积分 $h \in L^p(X, \mu)$ 存在, 且对几乎所有 $x \in X$ 有
\[
h(x) = \int_Y \phi(x, y)\,\dd\nu(y).
\]
\item[(d)] 证明 Minkowski 不等式:
\[
\Bigl(\int_X\Bigl|\int_Y\phi(x, y)\,\dd\nu(y)\Bigr|^p\,\dd\mu(x)\Bigr)^{1/p} \leq \int_Y\Bigl(\int_X|\phi(x, y)|^p\,\dd\mu(x)\Bigr)^{1/p}\,\dd\nu(y).
\]
\end{enumerate}
\end{exercise}

\begin{solution}
(a). 函数 $y \longmapsto \bigl(\int_X|\phi(x, y)|^p\,\dd\mu(x)\bigr)^{1/p}$ 由假设可积, 所以它对几乎所有 $y \in Y$ 取有限值, 这意味着对几乎所有 $y \in Y$ 有 $\int_X|\phi(x, y)|^p\,\dd\mu(x) < +\infty$.

(b). 设 $q \in (1, +\infty)$ 满足 $\frac{1}{p} + \frac{1}{q} = 1$. 我们要验证对每个 $f \in L^q(X, \mu)$, 积分 $\int_{Y'}\int_X f\phi\,\dd\mu\,\dd\nu$ 收敛 (其中 $\phi = \phi(x, y)$). 设 $f \in L^q(X, \mu)$. 由 H\"older 不等式, 对每个 $y \in Y'$, $\int_X f(x)\phi(x, y)\,\dd\mu(x)$ 绝对收敛, 且
\[
\int_X |f(x)\phi(x, y)|\,\dd\mu(x) \leq \|f\|_q\|\phi(\cdot, y)\|_p.
\]
由于 $\int_{Y'}\|\phi(\cdot, y)\|_p\,\dd\nu(y)$ 由假设收敛, 这给出 $\int_{Y'}\int_X f(x)\phi(x, y)\,\dd\mu(x)\,\dd\nu(y)$ 的收敛性, 甚至是绝对收敛性.

(c). 我们已在 (b) 中看到对每个 $f \in L^q(X, \mu)$ 有 $\int_X\int_Y |f(x)\phi(x, y)|\,\dd\mu(x)\,\dd\nu(y) < +\infty$. 由 Fubini 定理, 这蕴含对每个 $f \in L^q(X, \mu)$, 对几乎所有 $x \in X$ 有 $\int_Y |f(x)\phi(x, y)|\,\dd\nu(y) = |f(x)|\int_Y|\phi(x, y)|\,\dd\nu(y) < +\infty$. 由于 $f$ 任意 (且 $\mu$ $\sigma$-有限), 我们得到对几乎所有 $x \in X$ 有 $\int_Y|\phi(x, y)|\,\dd\nu(y) < +\infty$, 设对 $x \in X'$ (其中 $\mu(X - X') = 0$).
我们定义函数 $h \colon X' \to \CC$ 为 $h(x) = \int_Y\phi(x, y)\,\dd\nu(y)$. 我们要证明这是 $y \longmapsto \phi(\cdot, y)$ 的积分. 若 $f \in L^q(X, \mu)$, 则
\[
\int_X f(x)h(x)\,\dd\mu(x) = \int_X f(x)\phi(x, y)\,\dd\mu(x)\,\dd\nu(y),
\]
所以, 用 (b) 中的 H\"older 不等式,
\[
\Bigl|\int_X f(x)h(x)\,\dd\mu(x)\Bigr| \leq \int_{X \times Y}|f(x)\phi(x, y)|\,\dd\mu(x) \leq C\|f\|_q,
\]
其中
\[
C = \int_Y\Bigl(\int_X|\phi(x, y)|^p\,\dd\mu(x)\Bigr)^{1/p}\,\dd\nu(y).
\]
这证明 $f \longmapsto \int_X f(x)h(x)\,\dd\mu(x)$ 是 $L^q(X, \mu)$ 上的有界线性泛函, 且其算子范数被 $C$ 控制. 由于 $L^q(X, \mu)$ 的连续对偶是 $L^p(X, \mu)$, 必有 $h \in L^p(X, \mu)$ 且 $\|h\|_p \leq C$. 第一个性质连同 $\int_X f(x)h(x)\,\dd\mu(x)$ 的公式说 $h$ 确实是 $y \longmapsto \phi(\cdot, y)$ 的积分.

(d). 上面证明的 $h$ 的第二个性质正是 Minkowski 不等式.
\end{solution}
